\documentclass[UTF8,12pt,a4paper,fontset=fandol]{ctexbook}


% 支持从项目根目录或当前笔记目录编译。
\makeatletter
\def\input@path{{../}{../../}{../../../}{../../../../}}
\makeatother
\usepackage{researchnotes}


\title{数学分析习题（裴礼文）笔记}
\author{}
\date{\today}


\begin{document}


\maketitle
\tableofcontents
\clearpage


\chapter{一元函数极限}
\label{chap:peiliwen-limits}

本章围绕三个问题展开：极限存在的依据是什么，如何计算极限，以及当普通极限不存在时怎样描述变量的渐近行为。前五节以定义、估计和典型题型为主，随后用上、下极限统一讨论振荡现象，最后回到支撑这些论证的实数完备性。微分、积分与级数方法集中放在相应的复习性内容中，使用时说明所需条件。以下以 $\mathbb{N}=\{1,2,\ldots\}$ 表示正整数集；未特别说明时，数列和函数均取实值，“收敛”指收敛到有限实数。

% 教材来源：数学分析/textbooks/数分习题（分章原文）/数分习题-1.pdf。
% 第 1.1 节：PDF 第 1--9 页；按主题整合正文与代表性单元练习。
\section{预备}
\label{sec:pei-preliminaries}

\subsection{函数的范围、对称性与周期}

\begin{definition}
\label{def:pei-inverse}
设 $f:X\to Y$。若不同自变量的像不同，则称 $f$ 为\emph{单射}（injection）；若 $f(X)=Y$，则称 $f$ 为\emph{满射}（surjection）；兼具两者时称为\emph{双射}（bijection）。双射的\emph{反函数}（inverse function）$f^{-1}:Y\to X$ 由 $f^{-1}(y)=x\iff f(x)=y$ 定义，满足 $f^{-1}\circ f=\operatorname{id}_X$、$f\circ f^{-1}=\operatorname{id}_Y$。
\end{definition}

讨论反函数必须同时明确原函数的定义域与陪域。函数 $x\mapsto x^2$ 作为 $\mathbb{R}\to\mathbb{R}$ 的映射既非单射也非满射，而限制为 $[0,+\infty)\to[0,+\infty)$ 后是双射，反函数为平方根函数。严格单调函数在其定义域到其像集之间必为双射，反函数保持相同的增减方向；若 $f:X\to Y$、$g:Y\to Z$ 均为双射，则由复合后依次还原各自变量可得 $(g\circ f)^{-1}=f^{-1}\circ g^{-1}$。

设 $D=-D$，即 $x\in D$ 时 $-x\in D$。满足 $f(-x)=f(x)$ 的函数称为\emph{偶函数}（even function），满足 $f(-x)=-f(x)$ 的函数称为\emph{奇函数}（odd function）。任意 $D$ 上的函数都有唯一分解
\begin{equation}
    f(x)=H(x)+G(x),\qquad
    H(x)=\frac{f(x)+f(-x)}2,\quad
    G(x)=\frac{f(x)-f(-x)}2,
    \label{eq:pei-parity-decomposition}
\end{equation}
其中 $H$ 为偶函数，$G$ 为奇函数。事实上，把 $x$ 换成 $-x$ 后相加、相减，便得到上述两个必要表达式；直接代入又验证了存在性。若奇函数 $f:D\to f(D)$ 为双射，则 $f(D)$ 也关于原点对称，且由 $f(-x)=-f(x)$ 得 $f^{-1}(-y)=-f^{-1}(y)$，故反函数仍为奇函数。

\begin{definition}
\label{def:pei-period}
若 $f$ 在 $\mathbb{R}$ 上定义，存在 $T>0$ 使每个 $x\in\mathbb{R}$ 都满足 $f(x+T)=f(x)$，则称 $f$ 为\emph{周期函数}（periodic function），$T$ 为它的一个正周期。正周期中若有最小者，则称为最小正周期。
\end{definition}

正周期不必有最小值。常值函数的每个正数都是周期；\emph{Dirichlet 函数} $D(x)=\mathbf{1}_{\mathbb{Q}}(x)$ 的每个正有理数都是周期，因此也没有最小正周期。若两个函数具有公共正周期，则它们的和、差、积具有该周期；分别知道两者周期性，并不足以直接对和、差下结论。若周期为 $T$ 的函数在 $[0,T)$ 上可逆，逆映射为 $\varphi$，则它在 $[kT,(k+1)T)$ 上到同一像集的逆映射为 $y\mapsto\varphi(y)+kT$，其中 $k\in\mathbb{Z}$。

\begin{example}
\label{ex:pei-bounded-second-difference}
设 $f:\mathbb{R}\to\mathbb{R}$ 有界，$h>0$，且对一切 $x$ 有 $f(x+h)-2f(x)+f(x-h)=0$。证明 $h$ 是 $f$ 的周期。
\end{example}
\begin{solve}
固定 $x$，记 $d=f(x+h)-f(x)$。给定关系说明沿等距点列的相邻差均为 $d$，故归纳得到 $f(x+nh)=f(x)+nd$，$n\in\mathbb{N}$。若 $d\ne0$，右端随 $n$ 无界，与 $f$ 有界矛盾，所以 $d=0$。因 $x$ 任意，结论成立。有界性不能删去：$f(x)=x$ 满足给定差分关系，却不具有正周期。
\end{solve}

图像的对称也能给出周期。若 $x=b$、$x=c$ 是两个不同的对称轴，则先后作两次反射得到 $f(x)=f(x+2(c-b))$，故 $2|c-b|$ 为周期；若 $f$ 既为奇函数，又关于 $x=a\ne0$ 对称，则 $f(x+2a)=-f(x)$，再平移一次得到周期 $4|a|$。

\subsection{常用不等式与估计}

对实数 $u,v$，三角不等式及其反向形式分别为
\[
    |u+v|\le |u|+|v|,\qquad
    \bigl||u|-|v|\bigr|\le |u-v|.
\]
第一式由 $-|u|\le u\le|u|$、$-|v|\le v\le|v|$ 相加得到；再把 $u=(u-v)+v$、$v=(v-u)+u$ 代入第一式，便得第二式。它们分别用于控制复合误差和由函数值之差反推大小差异。

\begin{theorem}[算术平均与几何平均不等式]
\label{thm:pei-amgm}
设 $a_1,\ldots,a_n\ge0$，则
\begin{equation}
    (a_1\cdots a_n)^{1/n}\le\frac{a_1+\cdots+a_n}{n},
    \label{eq:pei-amgm}
\end{equation}
等号成立当且仅当 $a_1=\cdots=a_n$。
\end{theorem}
\begin{proof}
当 $n=2$ 时，由 $(\sqrt{a_1}-\sqrt{a_2})^2\ge0$ 得结论。若结论对 $n$ 成立，把 $2n$ 个数分为两组，先对每组应用假设，再对两组的算术平均应用二元情形，即得 $2n$ 个数的结论；等号要求两组内部及两组平均都相等。因此结论对每个 $2^k$ 成立。

反向一步，假设结论对 $n+1$ 成立，记 $A=(a_1+\cdots+a_n)/n$。若 $A=0$，所有数均为零；若 $A>0$，对 $a_1,\ldots,a_n,A$ 应用假设，得 $a_1\cdots a_n A\le A^{n+1}$，约去 $A$ 即得 $n$ 个数的结论，且等号条件保持不变。对任意 $n$ 先取 $2^k\ge n$，再逐次反推即可。
\end{proof}

\begin{proposition}[对数不等式]
\label{prop:pei-log-bounds}
对 $t>-1$，有
\begin{equation}
    \frac{t}{1+t}\le\log(1+t)\le t,
    \label{eq:pei-log-bounds}
\end{equation}
两侧等号均当且仅当 $t=0$ 成立。
\end{proposition}
\begin{proof}
这里使用导数作复习性证明。函数 $t-\log(1+t)$ 与 $\log(1+t)-t/(1+t)$ 的导数分别为 $t/(1+t)$ 与 $t/(1+t)^2$；两者在 $(-1,0)$ 上递减，在 $(0,+\infty)$ 上递增，且在 $0$ 处取值零，所以结论成立。
\end{proof}

例如令 $t=1/k$，便有 $1/(k+1)<\log(1+1/k)<1/k$。此类估计保留了对数与倒数之间的联系，常能把求和问题转化为对数的相消。使用不等式时应根据最终目标选择上、下界；估计过粗而丢掉主项，即使不等式正确，也未必能确定极限。

\subsection{Wallis 积分与乘积}

以下公式使用定积分与分部积分，可在学习积分后回看。非负整数的\emph{双阶乘}（double factorial）定义为 $n!!=n(n-2)(n-4)\cdots$，并约定 $0!!=(-1)!!=1$。令
\[
    I_n=\int_0^{\pi/2}\sin^n x\,dx
       =\int_0^{\pi/2}\cos^n x\,dx.
\]
变量替换 $x\mapsto\pi/2-x$ 给出两积分相等。对 $n\ge2$ 分部积分，得到 $I_n=(n-1)(I_{n-2}-I_n)$，所以
\begin{equation}
    I_n=\frac{n-1}{n}I_{n-2}
    =\frac{(n-1)!!}{n!!}
    \begin{cases}
        \pi/2,&n\text{ 为偶数},\\
        1,&n\text{ 为奇数}.
    \end{cases}
    \label{eq:pei-wallis-integral}
\end{equation}
这里的初值是 $I_0=\pi/2$、$I_1=1$。由于 $0\le\sin x\le1$，有 $I_{2n+1}\le I_{2n}\le I_{2n-1}$，因而
\[
    1\le\frac{I_{2n}}{I_{2n+1}}
      \le\frac{I_{2n-1}}{I_{2n+1}}=\frac{2n+1}{2n}\longrightarrow1.
\]
代入式~\eqref{eq:pei-wallis-integral}，得到\emph{Wallis 乘积公式}
\begin{equation}
    \lim_{n\to\infty}
    \prod_{k=1}^n\frac{(2k)^2}{(2k-1)(2k+1)}=\frac\pi2.
    \label{eq:pei-wallis-product}
\end{equation}

对非负整数 $m,n$，还可写出统一形式
\begin{equation}
    \int_0^{\pi/2}\sin^m x\cos^n x\,dx
    =\frac{(m-1)!!(n-1)!!}{(m+n)!!}
    \begin{cases}
        \pi/2,&m,n\text{ 均为偶数},\\
        1,&\text{其余情形}.
    \end{cases}
    \label{eq:pei-mixed-wallis}
\end{equation}
记左端为 $J_{m,n}$，当 $n\ge2$ 时，对 $\sin^m x\cos x\cdot\cos^{n-1}x$ 分部积分得 $J_{m,n}=(n-1)J_{m+2,n-2}/(m+1)$。反复降阶至 $J_{m,0}=I_m$ 或 $J_{m,1}=1/(m+1)$，即可得到式~\eqref{eq:pei-mixed-wallis}。

\subsection{确界与单调映射的一个应用}

非空实数集 $E$ 的\emph{上确界}（supremum）$s=\sup E$ 指最小上界，其可操作的刻画为：每个 $x\in E$ 都有 $x\le s$，且每个 $\varepsilon>0$ 都能找到 $x_\varepsilon\in E$ 使 $s-\varepsilon<x_\varepsilon\le s$。\emph{下确界}（infimum）$i=\inf E$ 相应满足 $i\le x$，且任意 $\varepsilon>0$ 都有 $x_\varepsilon\in E$ 使 $i\le x_\varepsilon<i+\varepsilon$。确界不要求属于集合；属于集合时才分别是最大值与最小值。有界非空实数集的确界存在性将在第~\ref{sec:pei-real-completeness} 节集中讨论。

若 $\{x_n\}$、$\{y_n\}$ 有界，则逐项估计 $x_n+y_n\le\sup_kx_k+\sup_ky_k$ 给出 $\sup_n(x_n+y_n)\le\sup_nx_n+\sup_ny_n$；这里一般不能写等号，因为两列的较大值未必出现在同一指标上。若 $A,B$ 是非空有界的非负数集，允许 $x\in A$ 与 $y\in B$ 独立选取，则有 $\sup\{xy:x\in A,y\in B\}=(\sup A)(\sup B)$：一侧由上界性质得到，另一侧由分别逼近两个上确界得到。

\begin{example}
\label{ex:pei-monotone-fixed-point}
设 $f:[a,b]\to[a,b]$ 单调不减，证明存在 $c\in[a,b]$ 使 $f(c)=c$，不要求 $f$ 连续。
\end{example}
\begin{solve}
令 $S=\{x\in[a,b]:x\le f(x)\}$。因 $a\le f(a)$，集合 $S$ 非空且有上界，故 $c=\sup S\in[a,b]$。对任意 $x\in S$，单调性给出 $x\le f(x)\le f(c)$，所以 $c\le f(c)$。再次使用单调性，$f(c)\le f(f(c))$，故 $f(c)\in S$，从而 $f(c)\le c$，合起来得 $f(c)=c$。
\end{solve}

% 第 1.2 节：PDF 第 9--24 页。
\section{用定义证明极限的存在性}
\label{sec:pei-limit-existence}

\subsection{定义、量词与误差估计}

\begin{definition}
\label{def:pei-sequence-limit}
称数列 $\{x_n\}$ 收敛到 $A\in\mathbb{R}$，记为 $x_n\to A$，是指
\[
    \forall\varepsilon>0,\quad\exists N\in\mathbb{N},\quad
    \forall n>N,\quad |x_n-A|<\varepsilon.
\]
设 $a$ 是 $E\subseteq\mathbb{R}$ 的聚点，$f:E\to\mathbb{R}$。称 $x\to a$、$x\in E$ 时 $f(x)\to A$，是指
\[
    \forall\varepsilon>0,\quad\exists\delta>0,\quad
    x\in E,\ 0<|x-a|<\delta\ \Longrightarrow\ |f(x)-A|<\varepsilon.
\]
这里“聚点”意指任意去心邻域中都有 $E$ 的点，详见第~\ref{sec:pei-real-completeness} 节。函数在 $a$ 处的取值不影响此极限。
\end{definition}

还约定 $x_n\to+\infty$ 表示对每个 $M>0$，存在 $N$ 使 $n>N$ 时 $x_n>M$；$x_n\to-\infty$ 表示 $-x_n\to+\infty$。若 $|x_n|\to+\infty$，称它为\emph{无穷大量}，不要求各项最终同号。函数趋于 $\pm\infty$ 时，把上述尾部条件换成相应的去心邻域条件；讨论 $x\to+\infty$ 的函数极限时，则把 $0<|x-a|<\delta$ 换成 $x>R$。这些广义极限与有限收敛应加以区分。

证明的核心是把误差控制在事先给定的 $\varepsilon$ 内。若可以解出 $|x_n-A|<\varepsilon$ 对 $n$ 的充分条件，就直接选取 $N$；若直接求解困难，则先放大为简单的上界；还可先要求 $n>N_0$ 或 $|x-a|<\delta_0$，取得局部简化后再分配误差。\emph[tip]{所选阈值只需保证结论，不要求最小或最优。}例如，$N$ 可以依赖 $\varepsilon$，但一经选定，就必须同时控制所有 $n>N$。

\begin{example}
\label{ex:pei-root-n-definition}
用定义证明 $n^{1/n}\to1$。
\end{example}
\begin{solve}
令 $u_n=n^{1/n}-1\ge0$。对 $n\ge2$，二项式展开中各项非负，故
\[
    n=(1+u_n)^n\ge\frac{n(n-1)}2u_n^2,
    \qquad 0\le u_n\le\sqrt{\frac2{n-1}}.
\]
给定 $\varepsilon>0$，取整数 $N>1+2/\varepsilon^2$，则 $n>N$ 时 $|n^{1/n}-1|<\varepsilon$。这里先引入误差变量，再通过二项式中的一个正项取得可解的上界。
\end{solve}

\begin{example}
\label{ex:pei-delta-stages}
用定义证明 $\displaystyle\lim_{x\to1}\frac7{16x^2-9}=1$。
\end{example}
\begin{solve}
先限制 $|x-1|<1/8$，于是 $7/8<x<9/8$，从而 $16x^2-9>3$、$|x+1|<3$。因此
\[
    \left|\frac7{16x^2-9}-1\right|
    =\frac{16|x-1||x+1|}{|16x^2-9|}<16|x-1|.
\]
给定 $\varepsilon>0$，取 $\delta=\min\{1/8,\varepsilon/16\}$ 即可。先隔开分母的零点，再控制误差，是分步估计的典型做法。
\end{solve}

不收敛到某个给定实数 $A$ 的否定形式是
\[
    \exists\varepsilon_0>0,\quad\forall N\in\mathbb{N},\quad
    \exists n>N,\quad |x_n-A|\ge\varepsilon_0.
\]
它只排除这个 $A$；要证明数列没有有限极限，还须排除所有实数候选值，或改用 Cauchy 准则的否定形式。量词取反时，依次交换“任意”与“存在”，同时否定最后的不等式，不能只把 $<$ 换成 $\ge$。

\subsection{Cauchy 准则与单调有界原理}

\begin{theorem}[Cauchy 收敛准则]
\label{thm:pei-cauchy}
实数列 $\{x_n\}$ 收敛，当且仅当
\begin{equation}
    \forall\varepsilon>0,\quad\exists N\in\mathbb{N},\quad
    \forall n>N,\ \forall p\in\mathbb{N},\quad
    |x_{n+p}-x_n|<\varepsilon.
    \label{eq:pei-cauchy}
\end{equation}
满足该条件的数列称为\emph{Cauchy 列}（Cauchy sequence），也称基本列。
\end{theorem}

必要性由 $|x_{n+p}-x_n|\le|x_{n+p}-A|+|x_n-A|$ 得到；充分性依赖实数的完备性，证明见第~\ref{sec:pei-real-completeness} 节。其优点是不需要预先猜出 $A$。\emph[caution]{式~\eqref{eq:pei-cauchy} 中的 $N$ 必须与 $p$ 无关。}对每个固定的 $p$ 有 $x_{n+p}-x_n\to0$ 只是必要条件，例如 $x_n=\sqrt n$ 满足这一条件，却趋于 $+\infty$。

\begin{example}
\label{ex:pei-harmonic-divergence}
证明调和部分和 $H_n=\sum_{k=1}^n1/k$ 不收敛。
\end{example}
\begin{solve}
对每个 $n\ge1$，有 $H_{2n}-H_n=\sum_{k=n+1}^{2n}1/k\ge1/2$。取 $\varepsilon_0=1/2$，任意阈值之后总能选取这样的两项，故不满足 Cauchy 准则。这里必须允许步长 $p=n$ 随 $n$ 改变；固定 $p$ 时差值反而趋于零。
\end{solve}

\begin{theorem}[单调有界原理]
\label{thm:pei-monotone}
实数列单调不减且有上界时，收敛到其全体项的上确界；单调不增且有下界时，收敛到其全体项的下确界。
\end{theorem}
\begin{proof}
只证递增情形。令 $A=\sup\{x_n:n\in\mathbb{N}\}$。给定 $\varepsilon>0$，由确界刻画存在 $N$ 使 $A-\varepsilon<x_N\le A$。对所有 $n>N$，单调性给出 $A-\varepsilon<x_N\le x_n\le A$，即 $|x_n-A|<\varepsilon$。递减情形对 $-x_n$ 应用同一结论即可。
\end{proof}

\begin{example}
\label{ex:pei-euler-constant}
证明 $H_n-\log n$ 收敛。
\end{example}
\begin{solve}
由式~\eqref{eq:pei-log-bounds}，
\[
    (H_{n+1}-\log(n+1))-(H_n-\log n)
    =\frac1{n+1}-\log\left(1+\frac1n\right)<0.
\]
又有 $\log n=\sum_{k=1}^{n-1}\log(1+1/k)<H_{n-1}<H_n$，故该数列有下界零，因而收敛。记其极限为 $\gamma_{\mathrm E}$，称为\emph{Euler 常数}，于是 $H_n=\log n+\gamma_{\mathrm E}+o(1)$；记号 $o(1)$ 表示趋于零的量。此结论只给出存在性，没有把未知极限当作已知数值。
\end{solve}

\subsection{子列、重新编号与函数极限}

\begin{theorem}[子列判别与 Heine 判别]
\label{thm:pei-heine}
数列 $x_n\to A$ 当且仅当每个子列 $x_{n_k}$ 都趋于同一个 $A$，其中 $n_1<n_2<\cdots$。函数 $f:E\to\mathbb{R}$ 在聚点 $a$ 处极限为 $A$，当且仅当对每个满足 $t_n\in E\setminus\{a\}$、$t_n\to a$ 的数列，都有 $f(t_n)\to A$。
\end{theorem}
\begin{proof}
从原极限到子列极限，直接使用定义即可。函数情形的逆向证明采用反证法：若极限不是 $A$，存在 $\varepsilon_0>0$，使每个去心邻域内都有 $t$ 满足 $|f(t)-A|\ge\varepsilon_0$。逐次在半径 $1/n$ 的邻域中选取 $t_n$，便有 $t_n\to a$ 而 $f(t_n)$ 不趋于 $A$。数列情形用不断增大的指标作同样选择。
\end{proof}

有限个子列若覆盖原数列的全部指标，且它们均趋于同一个 $A$，则原数列趋于 $A$；证明时把各子列的阈值换算为原指标后取最大值。特别，$x_n\to A$ 等价于 $x_{2n}\to A$ 且 $x_{2n+1}\to A$。仅有每个子列都能再选出某个收敛子列，并不能保证原数列收敛，例如 $(-1)^n$；若所有这些再选出的子列都趋于同一个指定的 $A$，则可用反证法得到原列收敛到 $A$。

收敛数列还允许适当重新编号。若 $g:\mathbb{N}\to\mathbb{N}$ 的每个纤维 $g^{-1}(\{k\})$ 都是有限集，则 $g(m)\to+\infty$：对每个 $K$，集合 $\{m:g(m)\le K\}$ 是有限个有限集的并。因此 $a_n\to A$ 蕴含 $a_{g(m)}\to A$。重新编号后的指标不必单调，故它未必是通常定义下的子列。

\subsection{平均、拟合与有限运算}

\begin{theorem}[加权平均保持极限]
\label{thm:pei-weighted-cesaro}
设 $a_n\to a$，$p_n\ge0$，$P_n=\sum_{k=1}^np_k\to+\infty$。则
\begin{equation}
    \frac{\sum_{k=1}^np_ka_k}{P_n}\longrightarrow a.
    \label{eq:pei-weighted-cesaro}
\end{equation}
其中只需在 $P_n>0$ 的充分大指标上讨论。
\end{theorem}
\begin{proof}
给定 $\varepsilon>0$，先取 $K$ 使 $k>K$ 时 $|a_k-a|<\varepsilon/2$。分开有限首段与尾段，有
\[
    \left|\frac{\sum_{k=1}^np_ka_k}{P_n}-a\right|
    \le\frac{\sum_{k=1}^Kp_k|a_k-a|}{P_n}
       +\frac\varepsilon2\frac{\sum_{k=K+1}^np_k}{P_n}.
\]
固定 $K$ 后，第一项趋于零，第二项不超过 $\varepsilon/2$，故结论成立。
\end{proof}

取 $p_k=1$ 就是算术平均的\emph{Cesàro 定理}，取 $p_k=k$ 得到 $\sum_{k=1}^nka_k/\sum_{k=1}^nk\to a$。更一般，若三角形系数 $c_{n,k}$ 满足 $\sum_{k=1}^nc_{n,k}=1$、$\sup_n\sum_{k=1}^n|c_{n,k}|<+\infty$，且每个固定 $k$ 有 $c_{n,k}\to0$，则 $\sum_{k=1}^nc_{n,k}a_k\to a$。证明仍是上述分段估计：有限首段由逐列趋零控制，尾段由绝对值和的一致上界控制；这也说明为什么带符号权重需要额外限制。

\begin{example}
\label{ex:pei-fitting-sine}
设 $f(t)/t\to1$（$t\to0$，$t\ne0$），固定 $a\ne0$。证明
\[
    \sum_{k=1}^nf\left(\frac{(2k-1)a}{n^2}\right)\longrightarrow a.
\]
\end{example}
\begin{solve}
记 $t_{n,k}=(2k-1)a/n^2$，则 $\sum_kt_{n,k}=a$，且 $\max_k|t_{n,k}|\le2|a|/n\to0$。给定 $\eta>0$，充分大 $n$ 时所有 $k\le n$ 同时满足 $|f(t_{n,k})-t_{n,k}|\le\eta|t_{n,k}|$，于是总误差不超过 $\eta\sum_k|t_{n,k}|=\eta|a|$。因 $\eta$ 任意，结论成立。此法把目标值 $a$ 写成与原和式对应的形式；取 $f(t)=\sin t$ 就得到典型的三角和极限。
\end{solve}

有限极限具有和、差、积运算法则；商法则还要求分母极限不为零。若 $x_n>0$、$x_n\to a>0$、$y_n\to b\in\mathbb{R}$，则由 $x_n^{y_n}=\exp(y_n\log x_n)$ 得 $x_n^{y_n}\to a^b$。\emph[caution]{这些法则不能直接推广到项数随 $n$ 增加的和或积。}例如令 $u_{n,k}=1$（$k=n$）、$u_{n,k}=0$（$k\ne n$），每个固定 $k$ 都有 $u_{n,k}\to0$，但 $\sum_{k=1}^nu_{n,k}=1$；例~\ref{ex:pei-fitting-sine} 中对整行误差的控制正是不可省去的步骤。

% 第 1.3 节：PDF 第 24--46 页。
\section{求极限值的若干方法}
\label{sec:pei-computing-limits}

\subsection{等价替换、相消与初等变形}

当 $v$ 在极限过程的充分后期不为零，且 $u/v\to1$，称 $u$ 与 $v$ \emph{渐近等价}（asymptotically equivalent），记为 $u\sim v$；$r=o(v)$ 表示 $r/v\to0$，$r=O(v)$ 表示该比值在充分后期有界。下列等价式均针对 $t\to0$：
\begin{equation}
    \begin{gathered}
        \sin t\sim\tan t\sim\arcsin t\sim\arctan t\sim t,\\
        \log(1+t)\sim e^t-1\sim t,\qquad
        1-\cos t\sim\tfrac12t^2,\\
        a^t-1\sim t\log a\quad(a>0,\ a\ne1),\qquad
        (1+t)^\alpha-1\sim\alpha t\quad(\alpha\ne0).
    \end{gathered}
    \label{eq:pei-equivalents}
\end{equation}
其中三角函数式可由基本三角极限得到，其余可由相应函数在零点的导数定义得到；使用反三角函数时取零点附近的通常分支。

\begin{proposition}[等价替换的范围]
\label{prop:pei-equivalent-replacement}
若 $u\sim v$，且相关商有定义，则在乘除结构中把因子 $u$ 换成 $v$，不改变存在的极限。但在 $u-w$ 中作同样替换，需要进一步验证 $u-v=o(u-w)$，或给出足以控制相消后误差的其他估计。
\end{proposition}

例如 $\displaystyle\lim_{t\to0}\frac{\sin t-t}{t^3}=-1/6$，若仅把 $\sin t$ 换成 $t$，会错误地得到零；同样，$\displaystyle\lim_{t\to0}\frac{1-\cos t-t^2/2}{t^4}=-1/24$，主项抵消后必须保留四阶信息。这里的数值可由本节后面的 Taylor 公式核查。相反，$\log(1+\sin t)/t$ 可分解为 $[\log(1+\sin t)/\sin t](\sin t/t)$，两个因子都趋于一，所以原极限为一。

先作精确恒等变形，往往比近似更直接。对 $n\ge2$，
\[
    \prod_{k=2}^n\left(1-\frac1{k^2}\right)
    =\prod_{k=2}^n\frac{(k-1)(k+1)}{k^2}
    =\frac{n+1}{2n}\longrightarrow\frac12.
\]
由倍角公式，对固定 $x\ne0$ 和充分大 $n$，
\[
    \prod_{k=1}^n\cos\frac{x}{2^k}
    =\frac{\sin x}{2^n\sin(x/2^n)}\longrightarrow\frac{\sin x}{x};
\]
当 $x=0$ 时各乘积为一。根式差通常先乘共轭式，例如 $\sqrt{n^2+n}-n=1/(\sqrt{1+1/n}+1)\to1/2$；由此再利用三角函数的连续性可得 $\sin^2(\pi\sqrt{n^2+n})\to1$。

\subsection{已知极限、变量替换与夹逼}

\begin{proposition}[幂式极限]
\label{prop:pei-power-limit}
若 $u\to0$、$1+u>0$，且 $uv\to L\in\mathbb{R}$，则 $(1+u)^v\to e^L$。更一般，对任意正底数 $F$，都应先研究 $G\log F$；若它趋于 $L$，则 $F^G\to e^L$。
\end{proposition}
\begin{proof}
对 $u\ne0$，有 $v\log(1+u)=uv\,[\log(1+u)/u]\to L$，在 $u=0$ 的位置左端为零且与原式相容；再利用指数函数连续性即可。若 $u$ 在无穷多位置为零，则 $uv$ 在这些位置也为零，所以有限极限 $L$ 必为零，论证仍成立。
\end{proof}

例如 $a,b>0$ 时，$\left[(a^{1/n}+b^{1/n})/2\right]^n\to\sqrt{ab}$，因为底数为 $1+(\log a+\log b)/(2n)+o(1/n)$。若 $a=0<b$，原式为 $b/2^n\to0$，与 $\sqrt{ab}=0$ 相符；不能在这一边界情形使用 $\log a$。这说明参数边界需要单独检查。

夹逼定理要求下界与上界具有相同的极限。利用 $s-1<\lfloor s\rfloor\le s$ 可得 $|t\lfloor1/t\rfloor-1|\le|t|$，故 $t\lfloor1/t\rfloor\to1$（$t\to0$）。若 $a_1,\ldots,a_m\ge0$、$M=\max_ka_k$，则
\[
    M\le\left(\sum_{k=1}^ma_k^p\right)^{1/p}\le m^{1/p}M
    \quad(p>0),
\]
所以 $p\to+\infty$ 时该表达式趋于 $M$。如果无法一次找到同极限的两端，也可对任意 $\eta>0$ 建立相差至多 $\eta$ 的渐近上、下界，最后再令 $\eta\downarrow0$。

\begin{example}
\label{ex:pei-convolution-means}
设 $x_n\to a$、$y_n\to b$，证明 $\displaystyle\frac1n\sum_{k=1}^nx_ky_{n+1-k}\to ab$。
\end{example}
\begin{solve}
令 $x_k=a+\alpha_k$、$y_k=b+\beta_k$，展开后除常数 $ab$ 外，出现两个趋零数列的算术平均，以及 $n^{-1}\sum_k\alpha_k\beta_{n+1-k}$。前两者由定理~\ref{thm:pei-weighted-cesaro} 趋于零。又因 $\alpha_k\to0$，存在 $M$ 使 $|\alpha_k|\le M$，故最后一项的绝对值不超过 $Mn^{-1}\sum_{k=1}^n|\beta_k|\to0$。这种替换把一般极限归结到零极限，并未交换随 $n$ 变化的求和与极限。
\end{solve}

由例~\ref{ex:pei-euler-constant}，有 $H_{2n}-H_n=\log2+o(1)$，从而
\begin{equation}
    \lim_{n\to\infty}\sum_{k=n+1}^{2n}\frac1k=\log2.
    \label{eq:pei-harmonic-block}
\end{equation}
另一种证明是逐项使用式~\eqref{eq:pei-log-bounds} 并求和，所得对数相消后同样夹逼到 $\log2$。同一个问题采用不同方法，关键在于识别它究竟具有平均结构、相消结构，还是积分和结构。

\subsection{微分、积分与级数方法的复习}

本小节引用后续微积分中的工具，适合在掌握对应定理后回看；这里只说明使用条件与关键计算，不把这些工具的完整证明作为本章的前提。\emph{L'Hospital 法则}要求 $f,g$ 在有关去心单侧邻域内可微，$g'$ 不为零，且 $f,g$ 同趋零，或者 $|g|\to+\infty$；若 $f'/g'\to L$（允许 $L=\pm\infty$），则 $f/g\to L$。分母趋于无穷大的版本不要求分子也趋于无穷大。乘积、差和幂式须先化成适用的商式，每次再次使用法则时都须重新核对条件。导数之比没有极限，并不能反推原商没有极限，例如 $(x+\sin x)/x\to1$，而导数之比 $1+\cos x$ 不收敛。

若 $f$ 在 $a$ 的邻域内有连续的 $m$ 阶导数，则带 Peano 余项的\emph{Taylor 公式}为
\[
    f(a+h)=\sum_{k=0}^m\frac{f^{(k)}(a)}{k!}h^k+o(h^m).
\]
当低阶项相消时，应展开到首个未抵消的阶次。例如
\[
    n\log\left(1+\frac1n\right)
    =1-\frac1{2n}+O(n^{-2}),\qquad
    n\left[e-\left(1+\frac1n\right)^n\right]\longrightarrow\frac e2.
\]
后一结论来自 $(1+1/n)^n=e\exp[-1/(2n)+O(n^{-2})]$，随后使用 $e^t-1\sim t$，所以余项在乘以 $n$ 后仍趋于零。

若 $f$ 在 $[a,b]$ 上 Riemann 可积，则等距分割的任意取样积分和趋于积分；特别，$[0,1]$ 上连续函数满足
\[
    \frac1n\sum_{k=1}^nf(k/n)\longrightarrow\int_0^1f(t)\,dt.
\]
取 $f(t)=1/(1+t)$，再得式~\eqref{eq:pei-harmonic-block}。对连乘先取对数，例如
\[
    \left[\prod_{k=1}^n\left(1+\frac kn\right)\right]^{1/n}
    \longrightarrow\exp\left(\int_0^1\log(1+t)\,dt\right)=\frac4e.
\]
\emph[caution]{若取样函数在端点无界，不能直接套用上述有界区间上的 Riemann 和结论}，必须另外控制端点贡献或改用其他方法。

设 $f\ge0$ 在 $[0,1]$ 上连续，$M=\max f$，则
\begin{equation}
    \lim_{p\to+\infty}\left(\int_0^1f(t)^p\,dt\right)^{1/p}=M.
    \label{eq:pei-lp-max}
\end{equation}
若 $M=0$，结论直接成立；若 $M>0$，对 $0<\eta<M$，连续性保证在某个正长度 $\ell_\eta$ 的子区间上 $f>M-\eta$，于是 $(M-\eta)\ell_\eta^{1/p}\le(\int f^p)^{1/p}\le M$。先令 $p\to\infty$，再令 $\eta\downarrow0$ 即得结论，端点处取到最大值也可使用单侧子区间。

收敛级数可用于控制通项或尾项。若 $\sum u_k$ 收敛，则 $u_k\to0$，且任意有限尾段 $\sum_{k=n+1}^{m(n)}u_k\to0$，其中 $m(n)\ge n$；后者由级数的 Cauchy 准则保证。若 $\sum_{n=2}^\infty|x_n-x_{n-1}|<+\infty$，则对 $m>n$ 有 $|x_m-x_n|\le\sum_{k=n+1}^m|x_k-x_{k-1}|\to0$，所以 $x_n$ 收敛。对于正数列，若 $u_{n+1}/u_n\to q<1$，则从某项起有 $u_{n+1}\le r u_n$，其中 $q<r<1$，几何估计已足以证明 $u_n\to0$，不必先证明整个级数收敛。

渐近增长比较也常用于选择估计。对固定 $\alpha>0$、$a>1$，当 $n\to\infty$ 时，
\[
    \log\log n\ll\log n\ll n^\alpha\ll a^n\ll n!\ll n^n,
\]
其中 $u_n\ll v_n$ 表示 $u_n/v_n\to0$。对数与幂、幂与指数的比较可由 L'Hospital 法则得到；$a^n/n!\to0$ 可用相邻比值 $a/(n+1)\to0$，而 $n!/n^n\le2^{-\lfloor n/2\rfloor}\to0$。各参数固定这一条件不可省略。

\subsection{周期平均与辅助函数}

\begin{example}
\label{ex:pei-periodic-average}
若 $f$ 是周期为 $T>0$ 的连续函数，证明
\[
    \lim_{x\to+\infty}\frac1x\int_0^xf(t)\,dt
    =\frac1T\int_0^Tf(t)\,dt.
\]
\end{example}
\begin{solve}
写成 $x=nT+r$，其中 $0\le r<T$。连续周期函数有界，设 $|f|\le M$，则
\[
    \int_0^xf=n\int_0^Tf+\int_0^rf,
    \qquad\left|\int_0^rf\right|\le MT.
\]
除以 $x$，由 $n/x\to1/T$ 及有界余项除以 $x$ 趋零，得到结论。这里长时间平均由完整周期决定，末尾不足一周期的贡献有统一上界。
\end{solve}

\begin{example}
\label{ex:pei-integrating-factor}
设 $f\in C^1([0,+\infty))$，且 $f(x)+f'(x)\to0$，证明 $f(x)\to0$。
\end{example}
\begin{solve}
令 $h=f+f'$。由 $(e^xf(x))'=e^xh(x)$，对任意固定 $A\ge0$ 和 $x\ge A$，有
\[
    f(x)=e^{A-x}f(A)+\int_A^xe^{t-x}h(t)\,dt.
\]
给定 $\varepsilon>0$，先选 $A$ 使 $t\ge A$ 时 $|h(t)|<\varepsilon/2$，积分的绝对值便小于 $\varepsilon/2$；再固定此 $A$，令 $x$ 足够大，使首项绝对值小于 $\varepsilon/2$，即得结论。
\end{solve}

类似地，若 $f\in C^1(\mathbb{R})$ 有界且 $|f-f'|\le1$，令 $h=f-f'$，对 $R>x$ 积分 $ (e^{-t}f(t))'=-e^{-t}h(t)$，得到
\[
    f(x)=e^{x-R}f(R)+\int_x^Re^{x-t}h(t)\,dt.
\]
令 $R\to+\infty$，有界性消去首项，故 $|f(x)|\le\int_x^\infty e^{x-t}\,dt=1$。辅助函数的选择来自希望把 $f$ 与 $f'$ 合并成一个乘积的导数，而所需无穷积分由指数核的可积性保证。

% 第 1.4 节：PDF 第 46--55 页。
\section{Stolz 公式}
\label{sec:pei-stolz}

\subsection{两种数列形式及其证明}

\begin{theorem}[Stolz 公式]
\label{thm:pei-stolz}
设 $\{a_n\}$、$\{b_n\}$ 为实数列。如果 $b_n$ 严格递增且 $b_n\to+\infty$，并且
\begin{equation}
    \frac{a_{n+1}-a_n}{b_{n+1}-b_n}\longrightarrow L,
    \label{eq:pei-stolz-differences}
\end{equation}
其中 $L\in\mathbb{R}\cup\{\pm\infty\}$，则 $a_n/b_n\to L$。另一个版本是：若 $b_n>0$ 严格递减到零，$a_n\to0$，且式~\eqref{eq:pei-stolz-differences} 成立，则仍有 $a_n/b_n\to L$。
\end{theorem}
\begin{proof}
先设第一种情形下 $L$ 有限。给定 $\varepsilon>0$，存在 $N$ 使 $k\ge N$ 时
\[
    (L-\varepsilon)(b_{k+1}-b_k)
    \le a_{k+1}-a_k
    \le(L+\varepsilon)(b_{k+1}-b_k).
\]
从 $k=N$ 累加到 $n-1$，得 $|(a_n-a_N)-L(b_n-b_N)|\le\varepsilon(b_n-b_N)$。除以充分后期为正的 $b_n$，再令 $n\to\infty$，有限首段的贡献趋零，故 $\limsup|a_n/b_n-L|\le\varepsilon$；因 $\varepsilon$ 任意，结论成立。若 $L=+\infty$，对任意 $M>0$ 从某项起累加 $a_{k+1}-a_k\ge M(b_{k+1}-b_k)$，得到 $\liminf a_n/b_n\ge M$；$L=-\infty$ 时改看 $-a_n$。

第二种情形把差写成 $a_k-a_{k+1}$、$b_k-b_{k+1}>0$。有限 $L$ 时从 $k=n$ 累加到 $m-1$，再利用 $a_m,b_m\to0$，得到 $|a_n-Lb_n|\le\varepsilon b_n$，除以 $b_n$ 即可。无穷极限时作相同的单侧估计。
\end{proof}

第一种形式不要求分子趋于无穷大；第二种形式则必须同时要求分子、分母趋零。公式提供充分条件，不提供逆命题：例如 $a_n=(-1)^n$、$b_n=n$ 时 $a_n/b_n\to0$，但相邻差之比不收敛。使用时应先确认分母的单调性及去向，再计算相邻差，不能只凭式子看起来像不定式就套用。

\subsection{求和、连乘与比值的应用}

对固定实数 $p>-1$，令 $a_n=\sum_{k=1}^nk^p$、$b_n=n^{p+1}$，由中值定理有 $(n+1)^{p+1}-n^{p+1}\sim(p+1)n^p$，故
\begin{equation}
    \frac{1^p+2^p+\cdots+n^p}{n^{p+1}}\longrightarrow\frac1{p+1}.
    \label{eq:pei-power-sums}
\end{equation}
加权平均定理也可由 Stolz 公式验证：若 $p_n>0$，则 $P_n$ 严格递增，且相邻分子、分母差的比正好是 $a_{n+1}$；若部分 $p_n=0$，则直接使用定理~\ref{thm:pei-weighted-cesaro}，或去掉零权重后再应用公式。

\begin{example}
\label{ex:pei-factorial-root}
证明 $\displaystyle\frac{(n!)^{1/n}}n\to e^{-1}$。
\end{example}
\begin{solve}
对原式取对数，得到 $A_n/n$，其中 $A_n=\log(n!)-n\log n$。其相邻差为
\[
    A_{n+1}-A_n=-n\log\left(1+\frac1n\right)\longrightarrow-1.
\]
以 $b_n=n$ 应用 Stolz 公式，得到 $A_n/n\to-1$，再取指数即得结论。此法避免了把端点处无界的 $\log t$ 直接当作普通 Riemann 和函数使用，也不需要预先引用 Stirling 公式。
\end{solve}

若 $u_n>0$ 且 $u_{n+1}/u_n\to q>0$，对 $\log u_n/n$ 使用 Stolz 公式，得到 $u_n^{1/n}\to q$。若相邻比值趋于零，则对任意 $r>0$，充分后期有 $u_{n+1}\le r u_n$，故 $\limsup u_n^{1/n}\le r$，从而根式极限为零；相邻比值趋于 $+\infty$ 时类似得到根式趋于 $+\infty$。第~\ref{sec:pei-sequence-limsup} 节将把这些结论推广到比值不收敛的情形。

\subsection{函数极限的差分版本}

\begin{proposition}
\label{prop:pei-function-stolz}
设 $T>0$，$f$ 在 $[a,+\infty)$ 的每个有界闭子区间上有界，$g$ 严格递增且 $g(x)\to+\infty$。若
\[
    \frac{f(x+T)-f(x)}{g(x+T)-g(x)}\longrightarrow L
    \qquad(x\to+\infty),
\]
则 $f(x)/g(x)\to L$，其中 $L$ 可为有限实数或 $\pm\infty$。另若 $g(x)>0$、$g(x+T)<g(x)$，$f(x),g(x)\to0$，且同一差分比趋于 $L$，则也有 $f(x)/g(x)\to L$。
\end{proposition}
\begin{proof}
第一种情形下，选取充分大的 $A$，把每个足够大的 $x$ 唯一写成 $x=t+nT$，其中 $t\in[A,A+T)$。沿 $t,t+T,\ldots,t+nT$ 累加差分不等式，便得数列版本中的估计。关键是它对 $t$ 一致成立：$f(t)$ 在该初始区间上有界，单调的 $g(t)$ 也有界，而 $g(x)\to+\infty$，故所有首段误差除以 $g(x)$ 后一致趋零。第二种情形沿每条等距点列把估计累加至无穷，利用 $f,g\to0$ 消去尾端，即得与数列版本相同的界。
\end{proof}

特别，若 $f$ 局部有界，且 $f(x+1)-f(x)\to A$，则取 $g(x)=x$ 得 $f(x)/x\to A$。局部有界保证不同小数部分对应的首段值受到统一控制，不能仅对整数点列证明后就把结论扩展到全体实变量。

% 第 1.5 节：PDF 第 55--78 页。
\section{递推形式的极限}
\label{sec:pei-recursive-limits}

\subsection{不动点、单调性与收敛依据}

递推式 $x_{n+1}=f(x_n)$ 中满足 $f(c)=c$ 的点称为\emph{不动点}（fixed point）。若已知 $x_n\to A$ 且 $f$ 在 $A$ 处连续，才能由递推式推出 $A=f(A)$；解出这个方程只得到候选值。实际解题时可以先求候选点，以便建立界或估计误差，但随后必须独立验证收敛。还要注意，方程即使只有一个不动点，迭代也可能不收敛，例如 $f(x)=1-x$ 在 $[0,1]$ 上的不动点唯一，而从 $x_0=0$ 出发得到 $0,1,0,1,\ldots$。

\begin{example}
\label{ex:pei-square-root-recursion}
设 $x_0=1$、$x_{n+1}=\sqrt{2x_n}$，求其极限。
\end{example}
\begin{solve}
由归纳法，$1\le x_n<2$。在这个范围内 $2x_n\ge x_n^2$，故 $x_{n+1}\ge x_n$。由单调有界原理，$x_n\to A\in[1,2]$；再由连续性得到 $A=\sqrt{2A}$，候选根为 $0,2$，其中只有 $2$ 位于该范围，所以极限为 $2$。也可直接写出 $x_n=2^{1-2^{-n}}$，独立核查结果。
\end{solve}

对单调不减的连续函数，若已找到不动点 $c$，且 $x<c$ 时 $x\le f(x)\le c$、$x>c$ 时 $c\le f(x)\le x$，则从相应一侧出发的迭代单调有界，极限是该侧范围内的某个不动点；若该范围内不动点唯一，则极限必为 $c$。例如，已知存在 $c$ 且 $0<f'(x)<1$ 时，$f$ 递增、$f(x)-x$ 严格递减，因而满足上述指向 $c$ 的关系。这里不需要把逐点的 $f'(x)<1$ 错当成存在统一常数 $q<1$ 的导数上界。

单调性也允许可求和的扰动。若 $x_n\ge m$，且 $x_{n+1}\le x_n+\varepsilon_n$，其中 $\varepsilon_n\ge0$、$\sum_n\varepsilon_n<+\infty$，令 $r_n=\sum_{k=n}^\infty\varepsilon_k$，则 $x_n+r_n$ 单调不增且有下界，所以收敛；又 $r_n\to0$，故 $x_n$ 收敛。它把“每一步略有增加”的问题转成真正的单调数列，条件是全部向上扰动的总量有限。

\subsection{压缩估计与不变区间}

\begin{theorem}[相邻差的压缩判别]
\label{thm:pei-contractive-differences}
若存在固定 $q\in[0,1)$，使所有 $n\ge1$ 满足 $|x_{n+1}-x_n|\le q|x_n-x_{n-1}|$，则 $x_n$ 收敛。若其极限为 $A$，则
\begin{equation}
    |A-x_n|\le\frac{q^n}{1-q}|x_1-x_0|\qquad(n\ge1).
    \label{eq:pei-contraction-error}
\end{equation}
\end{theorem}
\begin{proof}
归纳得 $|x_{k+1}-x_k|\le q^k|x_1-x_0|$。对 $m>n$ 求和，
\[
    |x_m-x_n|\le\sum_{k=n}^{m-1}|x_{k+1}-x_k|
    \le\frac{q^n}{1-q}|x_1-x_0|.
\]
右端随 $n$ 趋零且与 $m$ 无关，故由 Cauchy 准则得收敛；再令 $m\to\infty$ 即得误差界。$q=0$ 时从 $x_1$ 起数列恒定，结论同样成立。
\end{proof}

\begin{theorem}[闭区间上的压缩迭代]
\label{thm:pei-contraction-map}
设 $I$ 为非空闭区间，$f:I\to I$，且存在 $q\in[0,1)$ 使
\[
    |f(x)-f(y)|\le q|x-y|\qquad(x,y\in I).
\]
则 $f$ 在 $I$ 上有唯一不动点 $c$，任取 $x_0\in I$ 产生的迭代 $x_{n+1}=f(x_n)$ 都趋于 $c$，且 $|x_n-c|\le q^n|x_0-c|$。
\end{theorem}
\begin{proof}
$f(I)\subseteq I$ 保证所有迭代有定义并留在 $I$ 中，相邻差满足定理~\ref{thm:pei-contractive-differences} 的条件，故收敛。闭性保证极限 $c\in I$，压缩条件保证 $f$ 连续，因而 $f(c)=c$。若 $d$ 也是不动点，则 $|c-d|\le q|c-d|$，只能有 $c=d$。把 $y=c$ 代入压缩条件并归纳，即得最后的误差界。
\end{proof}

若 $f$ 在区间内可微，且 $|f'|\le q<1$，中值定理便给出压缩条件，但仍须先验证 $f(I)\subseteq I$。例如取 $I=[a-r,a+r]$，若 $|f(a)-a|\le(1-q)r$ 且 $|f'|\le q$，则 $|f(x)-a|\le q|x-a|+(1-q)r\le r$，不变性成立。仅有 $|f(x)-f(y)|\le|x-y|$，即\emph{非扩张性}（nonexpansiveness），不能使用上述严格压缩结论。

\begin{example}
\label{ex:pei-averaged-nonexpansive}
设 $f:[a,b]\to[a,b]$ 满足 $|f(x)-f(y)|\le|x-y|$。证明平均迭代 $x_{n+1}=[x_n+f(x_n)]/2$ 收敛到 $f$ 的一个不动点。
\end{example}
\begin{solve}
令 $G(x)=[x+f(x)]/2$，则 $G$ 连续并把 $[a,b]$ 映入自身。若 $x<y$，有 $G(y)-G(x)=[y-x+f(y)-f(x)]/2\ge0$，所以 $G$ 单调不减。由 $x_1\ge x_0$ 或 $x_1\le x_0$ 分别归纳出整个迭代单调，且始终有界，故收敛到某个 $c$。由 $c=G(c)$ 得 $f(c)=c$。本题使用的是一维次序和平均后的单调性，不需要严格压缩，也不保证不动点唯一。
\end{solve}

\subsection{交替迭代、通项与变量替换}

当 $f$ 连续、单调递减且把有界闭区间映入自身时，$F=f\circ f$ 单调不减。于是偶数子列和奇数子列分别是 $F$ 的迭代，均单调有界，极限 $\alpha,\beta$ 满足 $\alpha=f(\beta)$、$\beta=f(\alpha)$。只有进一步证明 $\alpha=\beta$，才能断定原列收敛；否则可能形成二周期。几何上，交替把 $x_n$ 送到 $f(x_n)$，再把函数值视为下一次自变量，可以帮助猜测两子列的方向，但图形观察不能替代这些不等式与极限论证。

\begin{example}
\label{ex:pei-golden-recursion}
设 $x_0=1$，$x_{n+1}=1/(1+x_n)$，求极限。
\end{example}
\begin{solve}
所有项都在 $[1/2,1]$ 中，且 $f(x)=1/(1+x)$ 在此区间上满足 $|f'(x)|\le4/9<1$，故定理~\ref{thm:pei-contraction-map} 保证收敛。极限满足 $A^2+A-1=0$，由正性得 $A=(\sqrt5-1)/2$。若用奇、偶子列法，则先得到 $\alpha=1/(1+\beta)$、$\beta=1/(1+\alpha)$，两式相减给出 $\alpha=\beta$，随后求出相同的根。若只在较大的 $[0,1]$ 上估计导数，所得上界为一；应缩小到实际不变区间，或改用子列论证。
\end{solve}

对于线性递推，写出通项或找不变量通常更合适。例如 $x_0=a$、$x_1=b$、$x_{n+2}=(x_{n+1}+x_n)/2$ 满足 $x_{n+2}-x_{n+1}=-(x_{n+1}-x_n)/2$，所以
\[
    x_n=\frac{a+2b}{3}+\frac{2(a-b)}3\left(-\frac12\right)^n
    \longrightarrow\frac{a+2b}{3}.
\]
在递推式中直接代入极限只得到 $A=A$，无法识别这个依赖初值的数值；差分通项保留了所需信息。

\begin{proposition}[二阶线性递推的比值]
\label{prop:pei-second-order-ratio}
设 $p,q>0$，$u_{n+2}=pu_{n+1}+qu_n$，$u_0,u_1\in\mathbb{R}$。记
\[
    \alpha=\frac{p+\sqrt{p^2+4q}}2,\qquad
    \beta=\frac{p-\sqrt{p^2+4q}}2.
\]
若 $u_1\ne\beta u_0$，则从某项起 $u_n\ne0$，且 $u_{n+1}/u_n\to\alpha$。若 $u_1=\beta u_0$、$u_0\ne0$，则该比值恒为 $\beta$；若 $u_0=u_1=0$，则比值无定义。
\end{proposition}
\begin{proof}
两个特征根不同，直接代入并核对初值可得 $u_n=C\alpha^n+D\beta^n$，其中 $C=(u_1-\beta u_0)/(\alpha-\beta)$、$D=u_0-C$。因 $\alpha>0>\beta$ 且 $|\beta|<\alpha$，当 $C\ne0$ 时除以 $C\alpha^n$，余项 $(D/C)(\beta/\alpha)^n\to0$，所以分母最终非零且比值趋于 $\alpha$。$C=0$ 时则只剩 $D\beta^n$。
\end{proof}

变量替换可以把陌生递推化为已有类型。例如 $x_{n+1}=2+1/x_n$、$x_0\ge2$ 时，$[2,5/2]$ 从第一次迭代后是不变区间，导数绝对值不超过 $1/4$，故极限为 $1+\sqrt2$。一般的分式递推也可先寻找不动点，再考察 $x_n-c$ 或 $1/x_n$ 的递推；替换时须核对分母始终不为零。

\begin{example}
\label{ex:pei-conjugate-fractional-part}
以 $\{x\}=x-\lfloor x\rfloor$ 表示小数部分，求 $\displaystyle\lim_{n\to\infty}\{(2+\sqrt3)^n\}$。
\end{example}
\begin{solve}
由二项式展开，$s_n=(2+\sqrt3)^n+(2-\sqrt3)^n$ 是整数，因为奇数次的根式项相消，偶数次项为整数的两倍。又 $r_n=(2-\sqrt3)^n\in(0,1)$ 且趋零，所以 $(2+\sqrt3)^n=s_n-r_n$，其小数部分为 $1-r_n\to1$。小数部分始终小于一，不妨碍其极限等于一；这里利用共轭表达式识别了相邻整数。
\end{solve}

\subsection{求根迭代与耦合递推}

\begin{example}
\label{ex:pei-newton-pth-root}
设 $a>0$、整数 $p\ge2$、$x_0\ge a^{1/p}$，定义
\[
    x_{n+1}=\frac1p\left((p-1)x_n+\frac{a}{x_n^{p-1}}\right).
\]
证明 $x_n\to a^{1/p}$。
\end{example}
\begin{solve}
对 $p-1$ 个 $x_n$ 与一个 $a/x_n^{p-1}$ 应用算术平均与几何平均不等式，得 $x_{n+1}\ge a^{1/p}$。若 $x_n\ge a^{1/p}$，则
\[
    x_{n+1}-x_n=\frac{a-x_n^p}{p x_n^{p-1}}\le0.
\]
故归纳得数列单调不增且有正下界。令极限为 $A$，代入递推式可得 $A^p=a$，由 $A>0$ 得所求结论。正下界同时保证取极限时分母不趋零。
\end{solve}

更一般的\emph{Newton 迭代}（Newton iteration）为 $x_{n+1}=x_n-f(x_n)/f'(x_n)$。一个可核查的充分条件是：$f\in C^2([a,b])$，$f'>0$、$f''\ge0$，$f(a)\le0\le f(b)$，并从 $x_0=b$ 出发。此时连续性和严格单调性保证唯一零点 $c\in[a,b]$。对 $x\ge c$，由导数不减有 $0\le f(x)=\int_c^xf'(t)\,dt\le f'(x)(x-c)$，所以 $c\le x-f(x)/f'(x)\le x$。迭代因而单调下降且留在 $[c,b]$；若极限为 $A$，连续性及 $f'(A)>0$ 给出 $f(A)=0$，所以 $A=c$。离开这些条件时，切线图示本身不构成收敛保证。

\begin{example}
\label{ex:pei-agm-recursion}
设 $a_0\ge b_0>0$，定义 $a_{n+1}=(a_n+b_n)/2$、$b_{n+1}=\sqrt{a_nb_n}$。证明两列收敛到同一个数。
\end{example}
\begin{solve}
由算术平均与几何平均不等式，
\[
    b_n\le b_{n+1}\le a_{n+1}\le a_n.
\]
故两列单调有界，分别趋于 $A,B$，且 $A\ge B\ge b_0>0$。在算术平均递推中取极限，$A=(A+B)/2$，于是 $A=B$。这个公共值称为两初值的\emph{算术几何平均}（arithmetic-geometric mean），一般不能仅凭递推关系把它写成初值的一个简单代数式。
\end{solve}

二元递推也可用同样的次序思想。设连续正函数 $F(x,y)$ 分别对两个变量单调不减，并存在 $c>0$，使 $F(c,c)=c$，$0<t<c$ 时 $t<F(t,t)$，$t>c$ 时 $F(t,t)<t$。对 $x_0=x_1=d>0$、$x_{n+2}=F(x_{n+1},x_n)$，若 $d<c$，则由归纳法得到数列单调不减且不超过 $c$；若 $d>c$，则单调不增且不低于 $c$。极限 $A>0$ 满足 $A=F(A,A)$，上述严格符号条件又排除了 $A\ne c$，故两种情形都趋于 $c$。分别单调性保证对角线以外的递推也服从这些界，不能只检查 $F(t,t)$ 的图像。

\subsection{直接误差法与收敛速度}

如果已猜出候选极限 $c$，也可直接证明误差 $e_n=x_n-c$ 趋零，无须先证明 $x_n$ 单调。例如若 $|e_{n+1}|\le q|e_n|$、$0\le q<1$，则 $|e_n|\le q^n|e_0|$。对于二阶递推，若 $|e_{n+1}|\le q\max\{|e_n|,|e_{n-1}|\}$，令 $M_k=\max\{|e_{2k}|,|e_{2k+1}|\}$，连续使用两次估计便有 $M_{k+1}\le qM_k$，仍可得到收敛。关键是找到一个固定的严格小于一的因子，而不是随指标接近一的数。

\begin{proposition}[趋零递推的渐近速度]
\label{prop:pei-recursion-rate}
设 $x_n>0$、$x_n\to0$，存在常数 $c,\alpha>0$ 使
\[
    x_{n+1}=x_n-cx_n^{1+\alpha}+o(x_n^{1+\alpha}).
\]
则 $n x_n^\alpha\to1/(\alpha c)$。
\end{proposition}
\begin{proof}
令 $y_n=x_n^{-\alpha}$，则
\[
    y_{n+1}-y_n
    =x_n^{-\alpha}\left[
       (1-cx_n^\alpha+o(x_n^\alpha))^{-\alpha}-1\right]
    \longrightarrow\alpha c.
\]
这里使用零点处的一阶展开；对 $y_n/n$ 应用 Stolz 公式，得 $y_n/n\to\alpha c$，取倒数即得结论。
\end{proof}

\begin{example}
\label{ex:pei-sine-recursion-rate}
设 $0<x_0<\pi/2$、$x_{n+1}=\sin x_n$，证明 $x_n\to0$，并求其渐近速度。
\end{example}
\begin{solve}
由于 $0<\sin x<x$（$0<x<\pi/2$），数列递减且为正，所以有极限 $A\ge0$。由 $A=\sin A$，在 $[0,x_0]$ 上只能有 $A=0$。再使用 $\sin x=x-x^3/6+o(x^3)$，在命题~\ref{prop:pei-recursion-rate} 中取 $\alpha=2$、$c=1/6$，得到 $nx_n^2\to3$，即 $\sqrt{n/3}\,x_n\to1$。本题的不动点处导数为一，严格压缩估计不能直接给出这一幂次速度。
\end{solve}

同理，若 $0<a_0<1$、$a_{n+1}=a_n(1-a_n)$，先由单调有界及不动点方程证明 $a_n\to0$，再取 $\alpha=c=1$，得到 $na_n\to1$。收敛性和速度分成两个步骤：前者排除非零极限，后者利用首个非零误差项选择合适的倒数幂，再交给 Stolz 公式。

% 第 1.6 节：PDF 第 78--89 页。
\section{序列的上、下极限}
\label{sec:pei-sequence-limsup}

\subsection{尾部确界与量词刻画}

对不收敛的数列，仅说“极限不存在”会丢掉振荡范围的信息。为统一处理无界情形，以下允许上、下极限取扩展实数 $\overline{\mathbb{R}}=\mathbb{R}\cup\{-\infty,+\infty\}$，但不把 $+\infty-\infty$ 或 $0\cdot\infty$ 视为已有定义的运算。

\begin{definition}
\label{def:pei-sequence-limsup}
对实数列 $\{x_n\}$，令 $M_n=\sup_{k\ge n}x_k$、$m_n=\inf_{k\ge n}x_k$，允许无界集合的确界为 $\pm\infty$。定义其\emph{上极限}（limit superior）和\emph{下极限}（limit inferior）为
\begin{equation}
    \limsup_{n\to\infty}x_n=\lim_{n\to\infty}M_n=\inf_nM_n,
    \qquad
    \liminf_{n\to\infty}x_n=\lim_{n\to\infty}m_n=\sup_nm_n.
    \label{eq:pei-tail-extrema}
\end{equation}
因 $M_n$ 单调不增、$m_n$ 单调不减，上述扩展实数极限总存在。
\end{definition}

记 $\ell=\liminf x_n$、$u=\limsup x_n$，总有 $\ell\le u$。若 $u$ 有限，它等价地由两个条件刻画：任意 $\varepsilon>0$ 时最终有 $x_n<u+\varepsilon$；任意 $\varepsilon>0$、任意 $N$ 之后仍能找到 $n>N$ 使 $x_n>u-\varepsilon$。第一条限制最终上界，第二条保证数列不断逼近这个边界。下极限的对应条件是最终有 $x_n>\ell-\varepsilon$，且任意后段都存在 $x_n<\ell+\varepsilon$。这些条件中的严格不等号与非严格不等号可通过缩小 $\varepsilon$ 相互转换。

例如 $x_n=(-1)^n+1/n$ 的上极限为一，下极限为负一。若 $x_{2n}=0$、$x_{2n-1}=n$，则上极限为 $+\infty$、下极限为零；上极限为 $+\infty$ 只说明没有有限上界，不等于整个数列趋于 $+\infty$。后者要求任意 $M$ 之上包含全部充分后期的项，也即 $\liminf x_n=+\infty$。

\subsection{子列解释与普通极限}

\begin{theorem}
\label{thm:pei-limsup-subsequences}
实数列的上极限是其全部广义子列极限的最大者，下极限是最小者；“广义”允许极限为 $\pm\infty$。特别，存在子列分别趋于上极限与下极限，任意子列 $x_{n_k}$ 都满足
\[
    \liminf_nx_n\le\liminf_kx_{n_k}
       \le\limsup_kx_{n_k}\le\limsup_nx_n.
\]
\end{theorem}
\begin{proof}
先设上极限 $u$ 有限。由尾部逼近性质，可递归选取 $n_j>n_{j-1}$ 且 $x_{n_j}>u-1/j$；最终上界又保证 $x_{n_j}\to u$。任何收敛子列的极限均不超过 $u$，否则它的充分后期会超过原列的某个最终上界。若 $u=+\infty$，逐次选取 $n_j>n_{j-1}$ 使 $x_{n_j}>j$；若 $u=-\infty$，则全列已趋于 $-\infty$。下极限对 $-x_n$ 使用同样论证。子列尾部包含在原列的相应尾部中，故其确界范围不会超出原列的最终范围，得到所列不等式。
\end{proof}

\begin{corollary}
\label{cor:pei-limsup-convergence}
对 $A\in\overline{\mathbb{R}}$，有 $x_n\to A$ 当且仅当 $\liminf x_n=\limsup x_n=A$；普通收敛还要求 $A$ 有限。
\end{corollary}
\begin{proof}
有限 $A$ 时，最终上、下界把 $x_n$ 夹在 $(A-\varepsilon,A+\varepsilon)$ 中；$A=\pm\infty$ 时对应单侧最终界即为趋于无穷大的定义。反向直接由极限定义和尾部确界得到。
\end{proof}

\subsection{运算不等式及符号条件}

\begin{proposition}
\label{prop:pei-limsup-sum}
设 $x_n,y_n$ 有界，记 $\ell_x=\liminf x_n$、$u_x=\limsup x_n$，$\ell_y,u_y$ 类似定义，则
\begin{equation}
    \begin{aligned}
        \ell_x+\ell_y&\le\liminf(x_n+y_n)
          \le\min\{\ell_x+u_y,u_x+\ell_y\},\\
        \max\{\ell_x+u_y,u_x+\ell_y\}
          &\le\limsup(x_n+y_n)\le u_x+u_y.
    \end{aligned}
    \label{eq:pei-limsup-sum-bounds}
\end{equation}
\end{proposition}
\begin{proof}
任意 $\varepsilon>0$ 时，充分后期有 $\ell_x-\varepsilon<x_n<u_x+\varepsilon$，$\ell_y-\varepsilon<y_n<u_y+\varepsilon$；相加并令 $\varepsilon\downarrow0$，得到最外侧的界。取 $x_{n_k}\to\ell_x$ 的子列，利用 $y_{n_k}$ 的最终上界，可得 $\liminf(x_n+y_n)\le\ell_x+u_y$；交换两列得到另一个上界。对两列取负并用 $\limsup(-z_n)=-\liminf z_n$，得到第二行的内侧界。
\end{proof}

一般不能把这些界改成等式。例如 $x_n=(-1)^n$、$y_n=-(-1)^n$ 时，和恒为零，而两个上极限之和为二。若其中一列 $x_n\to a\in\mathbb{R}$，则上下夹界合并，得到
\[
    \limsup(x_n+y_n)=a+\limsup y_n,\qquad
    \liminf(x_n+y_n)=a+\liminf y_n.
\]
此处也允许 $y_n$ 无界，因为加上趋于零的误差不会改变扩展实数的最终界。

若 $x_n,y_n$ 均非负且有界，同样的确界估计给出
\[
    \ell_x\ell_y\le\liminf(x_ny_n)
      \le\min\{\ell_xu_y,u_x\ell_y\},\qquad
    \max\{\ell_xu_y,u_x\ell_y\}
      \le\limsup(x_ny_n)\le u_xu_y.
\]
证明内侧界时选取达到相应上、下极限的子列，非负性保证乘法保持不等号方向。若允许负值，则必须分别处理符号；特别，当 $x_n\to a$ 且 $y_n$ 有界时，$a>0$ 保持上、下极限，$a<0$ 交换上、下极限，$a=0$ 时乘积趋零。若 $\liminf y_n>0$ 且 $y_n$ 有界，则 $\limsup(1/y_n)=1/\liminf y_n$、$\liminf(1/y_n)=1/\limsup y_n$，可据此处理商。无界情形应另行检查是否出现未定义的扩展实数运算。

\subsection{相邻比值、根式与次可加性}

\begin{proposition}
\label{prop:pei-ratio-root-limsup}
对任意正数列 $a_n$，有
\begin{equation}
    \liminf\frac{a_{n+1}}{a_n}
    \le\liminf a_n^{1/n}
    \le\limsup a_n^{1/n}
    \le\limsup\frac{a_{n+1}}{a_n}.
    \label{eq:pei-ratio-root-limsup}
\end{equation}
各量允许取零或 $+\infty$。
\end{proposition}
\begin{proof}
记右端为 $U$。若 $U<+\infty$，任取 $r>U$，充分后期有 $a_{n+1}\le r a_n$，故 $a_n\le a_Nr^{n-N}$。开 $n$ 次方并取上极限，得 $\limsup a_n^{1/n}\le r$，再令 $r\downarrow U$；若 $U=+\infty$，该侧不等式直接成立。下界类似：若 $L=\liminf a_{n+1}/a_n>0$，对每个 $0<r<L$ 作反向的几何估计，得到 $\liminf a_n^{1/n}\ge r$，再令 $r\uparrow L$；$L=0$ 时由正性得到结论，$L=+\infty$ 时可取任意 $r>0$。
\end{proof}

\begin{theorem}[次可加数列的平均极限]
\label{thm:pei-subadditive}
设实数列 $a_n$ 满足\emph{次可加性}（subadditivity），即 $a_{m+n}\le a_m+a_n$（$m,n\ge1$）。则
\begin{equation}
    \lim_{n\to\infty}\frac{a_n}{n}
       =\inf_{m\ge1}\frac{a_m}{m},
    \label{eq:pei-subadditive}
\end{equation}
右端允许为 $-\infty$。
\end{theorem}
\begin{proof}
固定 $m\ge1$，把 $n=qm+r$，$0\le r<m$，并约定 $a_0=0$。反复使用次可加性，得 $a_n\le qa_m+a_r$，其中 $r=0$ 时右端就是 $qa_m$。因有限个余数项有统一上界，除以 $n$ 后取上极限，得到 $\limsup a_n/n\le a_m/m$。若 $I=\inf_ma_m/m$ 有限，再对 $m$ 取下确界，并结合每项 $a_n/n\ge I$，即得上下极限相等；若 $I=-\infty$，上述上极限估计可小于任意实数，说明 $a_n/n\to-\infty$。
\end{proof}

这个证明的核心是固定块长后控制有限种余数，再让总长度趋于无穷。它也展示了上、下极限的一般用法：不预设收敛，先分别估计，再证明两者必须重合。

% 第 1.7 节：PDF 第 89--95 页。
\section{函数的上、下极限}
\label{sec:pei-function-limsup}

\subsection{子极限与去心邻域确界}

设 $f:E\to\mathbb{R}$，$a$ 是 $E$ 的聚点。若存在 $t_n\in E\setminus\{a\}$，$t_n\to a$ 且 $f(t_n)\to L\in\overline{\mathbb{R}}$，则称 $L$ 为 $f$ 在 $a$ 处的\emph{子极限}（cluster value）。它记录沿某条趋近途径得到的极限；上、下极限则分别记录所有这些途径中的最大、最小边界。

\begin{definition}
\label{def:pei-function-limsup}
令 $E_\delta=\{x\in E:0<|x-a|<\delta\}$，它对每个 $\delta>0$ 都非空。定义
\begin{equation}
    \begin{aligned}
        \limsup_{\substack{x\to a\\x\in E}}f(x)
          &=\lim_{\delta\downarrow0}\sup_{x\in E_\delta}f(x)
           =\inf_{\delta>0}\sup_{x\in E_\delta}f(x),\\
        \liminf_{\substack{x\to a\\x\in E}}f(x)
          &=\lim_{\delta\downarrow0}\inf_{x\in E_\delta}f(x)
           =\sup_{\delta>0}\inf_{x\in E_\delta}f(x).
    \end{aligned}
    \label{eq:pei-function-tail-extrema}
\end{equation}
确界和极限允许取扩展实数值。邻域缩小时上确界不增、下确界不减，所以两者总存在。
\end{definition}

有限上极限 $U$ 的量词刻画是：对任意 $\varepsilon>0$，存在 $\delta>0$，使全部 $x\in E_\delta$ 满足 $f(x)<U+\varepsilon$；同时对任意 $\varepsilon>0$ 和任意 $\delta>0$，总有某个 $x\in E_\delta$ 满足 $f(x)>U-\varepsilon$。下极限把两种不等号反向。由第二条递归选取 $|t_n-a|<1/n$ 的点，并结合第一条，可使 $f(t_n)\to U$；其余子极限不超过 $U$。对无穷边界以 $n$、$-n$ 替代 $U\pm1/n$ 作同样选取，因而上、下极限确实分别是最大、最小广义子极限。

设 $t_n\to a$、$t_n\in E\setminus\{a\}$，则
\[
    \liminf_{x\to a}f(x)\le\liminf_nf(t_n)
       \le\limsup_nf(t_n)\le\limsup_{x\to a}f(x).
\]
这给出函数与数列两种上、下极限之间的联系，且两端均可由适当选取趋近数列达到。普通函数极限存在的充要条件是这两个边界重合且有限；若重合于 $\pm\infty$，则对应趋于该无穷大。

\subsection{单侧形式与振荡例子}

若 $a$ 同时是 $E$ 左侧和右侧的聚点，在 $E_\delta$ 中分别增加 $x<a$ 或 $x>a$ 的限制，便得到左、右上极限和下极限。去心邻域是两侧邻域的并，因此
\[
    \begin{aligned}
        \limsup_{x\to a}f(x)
          &=\max\left\{\limsup_{x\to a-}f(x),
                       \limsup_{x\to a+}f(x)\right\},\\
        \liminf_{x\to a}f(x)
          &=\min\left\{\liminf_{x\to a-}f(x),
                       \liminf_{x\to a+}f(x)\right\}.
    \end{aligned}
\]
若定义域只从一侧积聚，应只讨论该侧，不能给另一侧虚设一个极限。$x\to+\infty$ 时，把去心邻域换成尾部 $E\cap(R,+\infty)$ 即得对应定义；$x\to-\infty$ 类似。

\begin{example}
\label{ex:pei-oscillation-cluster-values}
求 $f(x)=\sin(1/x)$ 在 $x\to0$ 时，以及 $D(x)=\mathbf{1}_{\mathbb{Q}}(x)$ 在任意点处的上、下极限和子极限。
\end{example}
\begin{solve}
前者始终在 $[-1,1]$ 中。任取 $c\in[-1,1]$，选 $\theta$ 使 $\sin\theta=c$，再令 $t_n=1/(2\pi n+\theta)$（去掉有限个分母为零的指标），则 $t_n\to0$ 且 $f(t_n)=c$。所以全部子极限组成 $[-1,1]$，上、下极限分别为 $1,-1$。有理数与无理数在每个开区间内都存在，故 $D$ 在任意点的上极限为一，下极限为零；其取值只有零与一，因此子极限只有这两个数，并不包含它们之间的其他数。
\end{solve}

这两个例子说明，子极限集的最大、最小值不能决定中间是否填满；后者还需要函数的其他性质。对连续振荡函数，可以通过适当的趋近点得到整段子极限；对不连续函数，子极限集可以只包含有限个点。

\subsection{估计规则与典型应用}

若 $f\le g$ 在某个充分小的去心邻域内成立，则 $\limsup f\le\limsup g$、$\liminf f\le\liminf g$。局部有界函数的和、非负函数的积满足与式~\eqref{eq:pei-limsup-sum-bounds} 及其乘法版本完全相同的不等式；证明时可直接对同一去心邻域中的确界作估计，或选取逼近边界的数列。若 $h(x)\to c\in\mathbb{R}$，则
\[
    \limsup(f+h)=\limsup f+c,\qquad
    \liminf(f+h)=\liminf f+c.
\]
特别，趋零加性误差不会改变上、下极限。乘法仍需核对极限因子的符号，并在负号出现时交换上、下极限。

\begin{example}
\label{ex:pei-rational-oscillation}
对 $x>1$，令 $f(x)=(x\sin^2x-1)/(x-\sin x)$，求 $x\to+\infty$ 时的上、下极限。
\end{example}
\begin{solve}
先隔离振荡主项：
\[
    f(x)-\sin^2x=\frac{\sin^3x-1}{x-\sin x},\qquad
    |f(x)-\sin^2x|\le\frac2{x-1}\longrightarrow0.
\]
又 $0\le\sin^2x\le1$，在 $x_n=\pi/2+n\pi$ 处取一，在 $y_n=n\pi$ 处取零，所以 $\limsup f=1$、$\liminf f=0$。分母最终离开零点的界保证了误差估计适用于所有足够大的实数，而非只在两条选定点列上成立。
\end{solve}

\begin{example}
\label{ex:pei-dyadic-difference}
设 $f$ 在零点的某个去心邻域中有定义，$F(x)=f(x)/x$ 在那里有界，并且
\[
    \frac{f(x)-f(x/2)}x\longrightarrow A\in\mathbb{R}
    \qquad(x\to0).
\]
证明 $f(x)/x\to2A$。
\end{example}
\begin{solve}
令 $H(x)=[f(x)-f(x/2)]/x$，则恒等式为 $F(x)=H(x)+F(x/2)/2$。记 $U=\limsup_{x\to0}F(x)$、$L=\liminf_{x\to0}F(x)$，有界性保证它们有限；变量替换 $t=x/2$ 不改变零点处的上、下极限。因 $H\to A$，分别取上、下极限得 $U=A+U/2$、$L=A+L/2$，所以 $U=L=2A$，结论成立。这里假设的是商 $f(x)/x$ 有界，不能把它替换成较弱的 $f(x)$ 有界。
\end{solve}

% 第 1.8 节：PDF 第 95--100、102--106 页；第 101 页为夹入的分册目录。
\section{实数及其基本定理}
\label{sec:pei-real-completeness}

\subsection{实数构造与完备性的作用}

有理数具有加、减、乘、除和大小比较，但并不能保证每个有界单调列或每个 Cauchy 列都在有理数中有极限。例如逐位截取 $\sqrt2$ 所得到的有理数近似列满足 Cauchy 条件，却不可能收敛到有理数。引入实数的目标是补足这些由有理近似逼近的位置，同时保留原有运算与次序。

教材介绍了小数、分割和基本列三种表述。用无限小数描述实数时，需要识别 $0.999\ldots=1.000\ldots$ 等不同表示。用有理数的\emph{Dedekind 分割}（Dedekind cut）构造时，以一个非空、真包含于 $\mathbb{Q}$、向下封闭且没有最大元的有理数子集表示一个数；有理数 $r$ 对应 $\{q\in\mathbb{Q}:q<r\}$，其余合适的分割表示无理数。用基本列构造时，则把两个有理 Cauchy 列 $\{q_n\}$、$\{r_n\}$ 在 $|q_n-r_n|$ 最终小于任意给定正有理数时视为等价，并以等价类作为实数。这种定义无需事先假定它们已有实数极限；运算由代表列逐项进行，再证明不依赖代表列的选择。这里给出的是构造思路，完整的域与序公理验证需另行展开。

本章采用公理观点：$\mathbb{R}$ 是满足上确界性质的有序域。其\emph{Archimedes 性质}为：对任意 $x,y>0$，存在 $n\in\mathbb{N}$ 使 $nx>y$。这也能由上确界性质推出：若 $\mathbb{N}$ 有上确界 $s$，则存在整数 $n>s-1$，于是整数 $n+1>s$，矛盾；再对 $y/x$ 选取足够大的整数即可。与有理数相比，新增的实质是\emph{完备性}（completeness），而不是增加另一套形式运算法则。

\subsection{八个基本定理的准确陈述}

下面把本章使用的存在性依据集中列出。为避免把“序列中重复出现的同一个值”与“数集中不同点的积聚”混同，先给出两种对象的定义。

\begin{definition}
\label{def:pei-cluster-point}
若对每个 $\varepsilon>0$ 和每个 $N$，都存在 $n>N$ 使 $|x_n-c|<\varepsilon$，则称 $c$ 为数列 $\{x_n\}$ 的一个\emph{子列极限点}，简称本章的数列极限点；它等价于存在子列 $x_{n_k}\to c$。若对每个 $\varepsilon>0$，数集 $E$ 都在 $0<|x-c|<\varepsilon$ 内有点，则称 $c$ 为 $E$ 的\emph{聚点}（accumulation point）。聚点不要求属于 $E$。若 $c\in E$ 且某个去心邻域内没有 $E$ 的点，则称 $c$ 为\emph{孤立点}（isolated point）。
\end{definition}

数集聚点的定义等价于每个邻域内有无穷多个不同的 $E$ 中点：若只有有限个，可取小于它们到 $c$ 的全部正距离的邻域，排除所有点而产生矛盾。数列的极限点则按指标出现次数计数；例如常数列 $x_n=1$ 的极限点为一，但其值域集合 $\{1\}$ 没有聚点。

\begin{theorem}[确界定理]
\label{thm:pei-supremum}
非空实数集有上界则有实数上确界；有下界则有实数下确界。
\end{theorem}

\begin{theorem}[单调有界定理]
\label{thm:pei-completeness-monotone}
实数列单调不减且有上界，或单调不增且有下界时，必有有限极限。
\end{theorem}

\begin{theorem}[基本列收敛定理]
\label{thm:pei-completeness-cauchy}
每个实数 Cauchy 列都收敛到实数。结合 Cauchy 条件的必要性，这正是定理~\ref{thm:pei-cauchy}。
\end{theorem}

\begin{theorem}[Bolzano--Weierstrass 定理]
\label{thm:pei-bolzano-weierstrass}
每个有界实数列都有收敛子列，其极限仍为实数。这是教材所称的致密性定理。
\end{theorem}

\begin{theorem}[聚点原理]
\label{thm:pei-accumulation}
每个有界无限实数集至少有一个实数聚点。
\end{theorem}

\begin{theorem}[闭区间套定理]
\label{thm:pei-nested-intervals}
设非空有界闭区间 $I_n=[a_n,b_n]$ 满足 $I_{n+1}\subseteq I_n$。则 $\bigcap_nI_n\ne\varnothing$；若还满足 $b_n-a_n\to0$，则公共点唯一。
\end{theorem}

\begin{theorem}[有限覆盖定理]
\label{thm:pei-finite-cover}
设开区间族 $\mathcal{U}$ 覆盖有限闭区间 $[a,b]$，即 $[a,b]\subseteq\bigcup_{U\in\mathcal{U}}U$。则存在有限个 $U_1,\ldots,U_m\in\mathcal{U}$，仍有 $[a,b]\subseteq\bigcup_{j=1}^mU_j$。这些区间称为一个有限子覆盖。
\end{theorem}

\begin{theorem}[Dedekind 分割定理]
\label{thm:pei-dedekind}
设非空集合 $A,B$ 把 $\mathbb{R}$ 分为两部分，$A\cap B=\varnothing$、$A\cup B=\mathbb{R}$，且每个 $a\in A$、$b\in B$ 都满足 $a<b$。则存在唯一实数 $c$ 分隔两组：或者 $c=\max A$ 且 $B$ 无最小值，或者 $c=\min B$ 且 $A$ 无最大值。
\end{theorem}

\emph[caution]{区间套的闭性和长度条件承担不同作用。}没有闭性时，即使长度趋零也可能没有公共点，例如 $(0,1/n)$；有闭性而没有长度趋零时，公共部分可含多个点，例如恒定区间 $[0,1]$。有限覆盖定理同样依赖闭且有界的对象：$(0,1)$ 被 $\{(1/n,1):n\ge2\}$ 覆盖，却不存在有限子覆盖。

\subsection{等价性的主证明链}

这些定理在有序域及 Archimedes 性质的基础上从不同角度表述同一完备性。下面给出一条闭合的证明链，并把聚点、覆盖和分割接入其中；引用前面定理时，只使用当前箭头已经建立的结论。

\begin{proof}[确界、单调列、区间套、子列与基本列之间的推导]
确界定理推出单调有界定理，证明已在定理~\ref{thm:pei-monotone} 中给出。单调有界定理推出区间套定理：$a_n$ 递增有上界、$b_n$ 递减有下界，故分别趋于 $a,b$，且对任意 $n,m$ 有 $a_n\le b_m$，所以 $a\le b$。逐项取极限可知公共部分恰为 $[a,b]$；若长度趋零，则 $a=b$。

区间套定理推出有界列的收敛子列：取一个含所有项的闭区间，将其二等分，至少一半含有无穷多个指标对应的项；保留该半区间并反复进行，得到长度趋零的闭区间套 $I_k$。设唯一公共点为 $c$，每次选择比前次更大的指标 $n_k$，使 $x_{n_k}\in I_k$，便有 $|x_{n_k}-c|\le|I_k|\to0$。这里“无穷多项”按指标计数，允许项值重复。

有界列的收敛子列性质推出 Cauchy 列收敛：基本列先必有界，因为在误差取一后，其尾部落在某个固定项的一邻域内，再并入有限首段即可。取收敛子列 $x_{n_k}\to c$；给定 $\varepsilon>0$，用 Cauchy 条件控制两项之差小于 $\varepsilon/2$，再选一个充分后期且距离 $c$ 小于 $\varepsilon/2$ 的子列项，三角不等式便使全部尾项距离 $c$ 小于 $\varepsilon$。

最后仅用基本列收敛性质反推确界定理。设 $E\ne\varnothing$ 且有上界 $b_0$；若某个 $a_0\in E$ 已是上界，它就是最大值及上确界。否则 $a_0$ 不是上界，$a_0<b_0$。逐次二分 $[a_n,b_n]$：中点若为上界，就替换右端；若不是上界，就替换左端。始终保持 $a_n$ 不是上界、$b_n$ 是上界，且长度为 $2^{-n}(b_0-a_0)$。两端点列都是 Cauchy 列，故由所假定的基本列收敛性质有共同极限 $c$。若某个 $x\in E$ 大于 $c$，则充分后期有 $b_n<x$，与上界性质矛盾；若 $y<c$，则充分后期有 $a_n>y$，而 $a_n$ 不是上界，故存在 $x\in E$ 使 $x>a_n>y$。因此 $c$ 是最小上界，闭合证明链完成。
\end{proof}

聚点原理与收敛子列性质的等价需要注意重复值。由有界无限集选出两两不同的点列，再选其收敛子列，极限即为原集聚点。反过来，若有界数列的某个值出现无穷多次，直接取常值子列；否则值域必为无限有界集。取其聚点 $c$，每个邻域有无穷多个不同值，因而有任意大的原指标，可递归选出趋于 $c$ 的子列。这两种情况涵盖全部有界数列。

\begin{proof}[区间套推出有限覆盖，有限覆盖推出收敛子列]
若某个覆盖 $[a,b]$ 的开区间族没有有限子覆盖，将 $[a,b]$ 二等分，至少有一半仍没有有限子覆盖；不断保留这样的半区间，得到长度趋零的闭区间套。其公共点 $c$ 属于某个原覆盖区间 $U$，而 $U$ 是开集，故存在 $r>0$ 使 $(c-r,c+r)\subseteq U$。充分后期的套区间长度小于 $r$，又包含 $c$，所以整个套区间都在 $U$ 内，竟由单个区间覆盖，矛盾。

反向，设有界列 $x_n\in[a,b]$ 没有收敛子列。对每个 $c\in[a,b]$，必存在一个开邻域 $U_c$，其中只含有限多个指标对应的项；否则按半径 $1/k$ 逐次选取任意大指标，就能得到趋于 $c$ 的子列。这些 $U_c$ 覆盖 $[a,b]$，有限覆盖定理给出有限子覆盖，因而全部指标只能落在有限个有限指标集中，矛盾。于是有限覆盖性质也接入了上述等价链。
\end{proof}

\begin{proof}[确界与分割之间的推导]
给定定理~\ref{thm:pei-dedekind} 的分割 $A,B$，集合 $A$ 非空并被 $B$ 的任意元素界住，令 $c=\sup A$。每个 $a\in A$ 都不超过 $c$，每个 $b\in B$ 都不小于 $c$。由于 $A,B$ 分割全部实数，$c$ 必属于其中一组；属于 $A$ 时是其最大值，属于 $B$ 时是其最小值。两组不能同时有最大值和最小值，因为若两者不同，它们的中点将无处归属；若相同，则违反两组不交。分割点的唯一性也由同一中点论证得到。

反向，设非空集 $E$ 有上界，令 $B$ 为全部上界组成的集合，$A$ 为全部非上界组成的集合。两组均非空，且任意 $a\in A$ 之后存在 $x\in E$，从而 $a<x\le b$ 对每个 $b\in B$ 成立，所以它们构成分割。$A$ 不可能有最大值：对 $a\in A$，取 $x\in E$ 且 $x>a$，则 $(a+x)/2$ 仍不是上界且大于 $a$。分割定理因而保证 $B$ 有最小值，它正是 $E$ 的上确界。
\end{proof}

\subsection{局部性质、整体结论与极限点集}

\begin{example}
\label{ex:pei-local-boundedness-compact}
若函数 $f:[a,b]\to\mathbb{R}$ 在每点附近相对于 $[a,b]$ 局部有界，证明它在整个区间上有界。
\end{example}
\begin{solve}
对每个 $x\in[a,b]$，取开邻域 $U_x$ 与有限常数 $M_x$，使 $t\in U_x\cap[a,b]$ 时 $|f(t)|\le M_x$。这些邻域构成开覆盖，取有限子覆盖 $U_{x_1},\ldots,U_{x_m}$，便有 $|f(t)|\le\max_jM_{x_j}$ 对全区间成立。因此每点存在有限双侧极限、端点存在相应有限单侧极限的函数，即使不连续，也在 $[a,b]$ 上有界：极限控制去心邻域，再把该点有限的函数值并入上界即可。若区间改为 $(0,1)$，$f(t)=1/t$ 在每点局部有界却整体无界，说明闭区间条件的作用。
\end{solve}

\begin{example}
\label{ex:pei-bisection-fixed-point}
设 $f:[0,1]\to\mathbb{R}$ 连续，且 $f(0)>0$、$f(1)<1$。用区间套构造 $c\in(0,1)$ 使 $f(c)=c$。
\end{example}
\begin{solve}
令 $g(x)=f(x)-x$，则 $g(0)>0$、$g(1)<0$。二分 $[0,1]$，中点若使 $g=0$，结论已得；否则保留左端为正、右端为负的半区间。反复进行得到闭区间套 $[a_n,b_n]$，长度趋零且 $g(a_n)>0$、$g(b_n)<0$。其唯一公共点为 $c$，连续性给出 $g(c)\ge0$ 及 $g(c)\le0$，所以 $g(c)=0$；端点处符号严格，故 $c\ne0,1$。本例用连续性把区间内的逼近转化为方程的精确解，与例~\ref{ex:pei-monotone-fixed-point} 的单调映射证明使用了不同假设。
\end{solve}

\begin{proposition}
\label{prop:pei-closed-cluster-set}
设 $C$ 是一个实数列全部有限子列极限点组成的集合，则 $C$ 是闭集；若原数列有界，$C$ 还是非空有界闭集，且 $C$ 只有一个元素当且仅当原数列收敛。
\end{proposition}
\begin{proof}
设 $c_j\in C$ 且 $c_j\to c\in\mathbb{R}$。因每个 $c_j$ 都是原数列的子列极限，可递归选取 $n_j>n_{j-1}$，使 $|x_{n_j}-c_j|<1/j$。于是 $|x_{n_j}-c|\le1/j+|c_j-c|\to0$，故 $c\in C$，证明了闭性。有界性由原列传给各子列极限，非空性由 Bolzano--Weierstrass 定理保证。

若有界原列仅有一个子列极限点 $c$ 而不趋于 $c$，则存在 $\varepsilon_0>0$ 及子列使 $|x_{n_j}-c|\ge\varepsilon_0$。该子列仍有界，故存在收敛子列，其极限 $d$ 满足 $|d-c|\ge\varepsilon_0$，与唯一性矛盾。反向由收敛列的全部子列都具有相同极限得到。
\end{proof}

例如，不存在有限子列极限点恰好为 $\{1,1/2,1/3,\ldots\}$ 的实数列，因为该集合不闭，漏掉了零。另一方面，“唯一有限子列极限点”若没有有界性并不足以保证收敛：$x_{2n}=0$、$x_{2n-1}=n$ 的唯一有限子列极限点为零，原列仍不收敛。与第~\ref{sec:pei-sequence-limsup} 节允许无穷子列极限的表述相对照，可以看出有界性正是在排除逃向无穷大的可能。


% 教材来源：数学分析/textbooks/数分习题（分章原文）/数分习题-2.pdf。
% PDF 第 1--19 页对应第 2.1 节，第 19--32 页对应第 2.2 节，
% 第 32--33 页对应第 2.3 节，第 33--42 页对应第 2.4 节。
\chapter{一元函数的连续性}
\label{chap:pei-continuity}

本章依据所给教材第二章“一元函数的连续性”整理，围绕连续性的证明与应用、一致连续性、上与下半连续、函数方程四个主题展开。连续性把自变量的逼近传递为函数值的逼近；闭区间的紧性则把逐点的局部信息转化为有界性、最值和统一的误差控制。学习时应区分三种不同任务：由已知条件推出连续，由连续性推出存在性或结构性质，以及用连续性排除函数方程中的非标准解。以下保留教材的主题次序，合并同类方法，并选取代表例题及单元练习补足论证。

\section{连续性的证明与应用}
\label{sec:pei-continuity-proofs}

\subsection{从局部定义到集合刻画}

\begin{definition}
设 $E\subseteq\mathbb R$，$f:E\to\mathbb R$，$x_0\in E$。若对每个 $\varepsilon>0$，存在 $\delta>0$，使得对所有 $x\in E$，
\[
|x-x_0|<\delta\quad\Longrightarrow\quad |f(x)-f(x_0)|<\varepsilon,
\]
则称 $f$ 在 $x_0$ 处\emph{连续（continuous）}。这里的邻域应与定义域 $E$ 相交；因此区间端点采用相应的单侧条件，孤立点处则自动连续。
\end{definition}

证明连续性时，首先检查已有的和、积、商、复合运算能否直接使用；商的分母在所讨论的点不能为零，复合函数的内层取值必须落在外层定义域中。若条件以数列形式给出，通常使用序列判别；若条件涉及单调性、介值性或集合的开闭性，则应把函数值的限制转成邻域或水平集的限制。无论采用哪种形式，待证结论始终是在一个任意但固定的点控制所有充分接近的自变量。

\begin{proposition}
\label{prop:pei-continuity-characterizations}
对 $f:E\to\mathbb R$，以下三个条件等价：$f$ 在 $E$ 上连续；每当 $x_n\in E$ 且 $x_n\to x_0\in E$ 时，都有 $f(x_n)\to f(x_0)$；每个开集 $U\subseteq\mathbb R$ 的逆像 $f^{-1}(U)$ 都是 $E$ 中的相对开集。最后一个条件又等价于：对每个 $c\in\mathbb R$，集合 $\{x\in E:f(x)<c\}$ 与 $\{x\in E:f(x)>c\}$ 都相对开；也等价于每个闭集的逆像都相对闭。
\end{proposition}

\begin{proof}
连续性直接保证收敛数列的像收敛。反之，若 $f$ 在 $x_0$ 不连续，则有 $\varepsilon_0>0$，使每个 $n$ 都能选到 $x_n\in E$，满足 $|x_n-x_0|<1/n$ 而 $|f(x_n)-f(x_0)|\ge\varepsilon_0$，与序列条件矛盾。

若 $f$ 连续，$x_0\in f^{-1}(U)$，则可先取 $\varepsilon>0$ 使 $(f(x_0)-\varepsilon,f(x_0)+\varepsilon)\subseteq U$，再用连续性找到 $x_0$ 的相对邻域包含在 $f^{-1}(U)$ 中。反之，将逆像条件用于这个以 $f(x_0)$ 为中心的开区间，便得到连续性的定义。两个严格水平集的交给出开区间的逆像，任意开集是开区间的并，故水平集条件也充分。闭集形式由取补集得到。
\end{proof}

\begin{example}
\label{ex:pei-thomae}
教材的 Riemann 函数也常称为 Thomae 函数，定义为
\[
R(x)=\begin{cases}
1/q,&x=p/q\in\mathbb Q,\ \gcd(p,q)=1,\ q>0,\\
0,&x\notin\mathbb Q.
\end{cases}
\]
证明它在无理点连续、在有理点不连续。
\end{example}

\begin{solve}
若 $x_0=p/q$ 为既约有理数，取无理数列 $x_n\to x_0$，则 $R(x_n)=0$ 而 $R(x_0)=1/q>0$，故不连续。若 $x_0$ 无理，给定 $\varepsilon>0$，在有界区间 $[x_0-1,x_0+1]$ 内使 $R(x)\ge\varepsilon$ 的点只能是分母 $q\le1/\varepsilon$ 的既约分数。分母只有有限种选择，而每个分母对应的分子在该区间内也只有有限个，故这类点构成有限集 $F$。由于 $x_0\notin F$，可以取 $0<\delta<1$，使 $(x_0-\delta,x_0+\delta)$ 与 $F$ 不交；于是 $|x-x_0|<\delta$ 时有 $|R(x)-R(x_0)|=R(x)<\varepsilon$。注意 $R(0)=1$，因为 $0$ 的既约表示为 $0/1$。
\end{solve}

\begin{example}
设 $u,v:E\to\mathbb R$ 连续，证明 $\max\{u,v\}$、$\min\{u,v\}$ 连续，并说明三个连续函数的中间值也连续。
\end{example}

\begin{solve}
利用恒等式
\[
\max\{u,v\}=\frac{u+v+|u-v|}{2},\qquad
\min\{u,v\}=\frac{u+v-|u-v|}{2},
\]
即可归结为连续函数的运算。三个实数的中间值等于三者之和减去最大值与最小值，故三个连续函数逐点取中间值也连续。类似地，对 $n\in\mathbb N$ 定义截断函数 $T_n(t)=\max\{-n,\min\{t,n\}\}$，则 $f$ 连续当且仅当每个 $T_n\circ f$ 连续。充分性的关键是局部恢复原函数：固定 $x_0$，取整数 $n>|f(x_0)|$，由 $T_n\circ f$ 的连续性，使其在 $x_0$ 附近的值严格位于 $(-n,n)$ 内；此时截断没有发生，因而 $T_n(f(x))=f(x)$。
\end{solve}

\subsection{介值性、单调性与紧性}

\begin{theorem}
\label{thm:pei-compact-continuity}
若 $f$ 在非退化闭区间 $[a,b]$ 上连续，则 $f$ 有界且达到最大值与最小值；其像为闭区间
\[
f([a,b])=[m,M],\qquad m=\min_{[a,b]}f,\quad M=\max_{[a,b]}f.
\]
特别地，若 $f(a)f(b)<0$，则存在 $c\in(a,b)$ 使 $f(c)=0$；若 $f(x)>0$ 对每个 $x\in[a,b]$ 成立，则存在统一的 $m>0$ 使 $f(x)\ge m$。
\end{theorem}

\begin{proof}
连续性保证每点 $x$ 附近都有 $|f(t)|<|f(x)|+1$。这些邻域覆盖 $[a,b]$，取有限子覆盖后，将有限个上界取最大值便得到整体有界性。设 $M=\sup f([a,b])$，按确界定义选 $x_n$ 使 $f(x_n)\to M$。闭区间的序列紧性给出子列 $x_{n_k}\to c\in[a,b]$，故 $f(c)=M$；最小值同理。

零点结论也可直接由确界原理证明。不妨 $f(a)<0<f(b)$，令 $A=\{x\in[a,b]:f(x)\le0\}$，$c=\sup A$。端点处的局部保号性保证 $a<c<b$。由 $A$ 中趋于 $c$ 的数列和连续性有 $f(c)\le0$；若 $f(c)<0$，则 $c$ 右侧仍有 $A$ 中的点，违反上确界定义，故 $f(c)=0$。对 $f-y$ 使用这一结论可得介值性，进而得到像区间的表达式。最后，正值函数所达到的最小值必严格为正。
\end{proof}

上述结论中的闭区间条件不能遗漏。例如 $f(x)=1/x$ 在 $(0,1)$ 上连续却无界，$f(x)=x$ 在 $(0,1)$ 上有界却不达到两端确界，而同一个 $f(x)=x$ 处处为正却没有正的统一下界。若 $f$ 在 $[a,+\infty)$ 上连续且在 $+\infty$ 有有限极限，则先在尾部利用极限控制函数值，再在剩余的闭区间应用有界性定理，即知 $f$ 整体有界；若 $f$ 在 $(a,b)$ 上连续，且两端极限均为 $-\infty$，则可先排除端点附近比某个内点函数值更小的部分，再在中间闭区间找到全局最大值。

\begin{proposition}
\label{prop:pei-injective-ivt}
设 $I$ 是区间，$f:I\to\mathbb R$ 具有\emph{介值性（intermediate value property）}且为单射，则 $f$ 严格单调，并且 $f$ 与逆函数 $f^{-1}:f(I)\to I$ 都连续。特别地，连续单射在区间上必严格单调。
\end{proposition}

\begin{proof}
若单射 $f$ 不严格单调，就能在 $I$ 中找到 $x_1<x_2<x_3$，使 $f(x_2)$ 严格大于两端函数值或严格小于两端函数值。在前一种情形，在 $\max\{f(x_1),f(x_3)\}$ 与 $f(x_2)$ 之间取一数，由介值性，它在 $(x_1,x_2)$ 与 $(x_2,x_3)$ 内各有一个原像，违反单射性；另一情形同理。因此 $f$ 严格单调。单调函数若在某点不连续，函数值与单侧极限之间就有未被取到的值，违反介值性，故 $f$ 连续。严格单调函数的逆函数仍严格单调，而任取 $y_1<y<y_2$ 属于像区间，其原像也夹在两端原像之间，故逆函数具有介值性，从而也连续。端点按单侧条件处理。
\end{proof}

\begin{notealert}[type=caution,title={严格单调与非严格单调}]
若只知 $f$ 单调不减且像集闭，并不能推出连续：阶跃函数的像集 $\{0,1\}$ 就是反例。若改为 $f$ 在 $[0,1]$ 上严格递增且像集闭，则可以推出连续。事实上，左极限 $L=\lim_{x\uparrow x_0}f(x)$ 属于闭像集；若 $L<f(x_0)$，则 $L=f(t)$ 的原像只能满足 $t<x_0$，但取 $t<s<x_0$ 又有 $f(t)<f(s)\le L$，矛盾。右极限同理。使用教材的“单调递增”条件时，应先明确其是否表示严格递增。
\end{notealert}

\begin{example}
\label{ex:pei-running-extrema}
设 $f\in C([a,b])$，定义
\[
M(x)=\max_{a\le t\le x}f(t),\qquad m(x)=\min_{a\le t\le x}f(t).
\]
证明 $M,m$ 在 $[a,b]$ 连续。
\end{example}

\begin{solve}
$M$ 随区间扩大而单调不减。固定 $x_0$，由 $f$ 在 $x_0$ 连续，对充分接近且大于 $x_0$ 的 $x$，新增区间上的函数值均不超过 $f(x_0)+\varepsilon\le M(x_0)+\varepsilon$，所以 $M(x_0)\le M(x)\le M(x_0)+\varepsilon$。再看左侧，在 $[a,x_0]$ 选最大值点 $t_0$。若 $t_0<x_0$，则 $t_0<x<x_0$ 时已有 $M(x)=M(x_0)$；若最大值只能在 $x_0$ 取得，则 $f(x)\le M(x)\le M(x_0)=f(x_0)$，由连续性仍有 $M(x)\to M(x_0)$。端点只需相应一侧。对 $-f$ 使用同一结论，并注意 $m=-\max_{a\le t\le x}(-f(t))$，便得 $m$ 连续。
\end{solve}

\begin{example}
设 $f:[0,1]\to[0,1]$ 的图像在 $[0,1]^2$ 中为闭集，证明 $f$ 连续。
\end{example}

\begin{solve}
若不连续，则存在 $x_n\to x_0$ 与 $\varepsilon_0>0$，使 $|f(x_n)-f(x_0)|\ge\varepsilon_0$。像序列位于紧区间 $[0,1]$，故有子列 $f(x_{n_k})\to y$。图像闭保证 $(x_0,y)$ 仍在图像上，即 $y=f(x_0)$，与距离下界矛盾。这里有界的像空间不可随意去掉：在 $[0,1]$ 上定义 $f(0)=0$、$f(x)=1/x$（$x>0$），图像在 $[0,1]\times\mathbb R$ 中闭，但函数在 $0$ 处不连续。
\end{solve}

\subsection{构造辅助函数与利用连续的极限传递}

存在性问题通常不直接求未知点，而是构造连续函数 $F$，使目标等式等价于 $F(x)=0$。关键是找到能保证 $F$ 异号的两点；它们可能是区间端点、原函数的最值点，也可能来自一串差值的望远镜求和。若条件含平移、周期或递推关系，则应尝试构造趋于已知连续点的自变量序列。

\begin{example}
若 $f\in C([0,1])$，$f\ge0$，且 $f(0)=f(1)=0$，证明对每个 $h\in(0,1)$，存在 $x\in[0,1-h]$ 使 $f(x+h)=f(x)$。再证明：仅假定 $f(0)=f(1)$ 时，对每个整数 $n\ge2$，步长 $h=1/n$ 仍有这一结论。
\end{example}

\begin{solve}
令 $F(x)=f(x+h)-f(x)$。在第一个问题中，$F(0)=f(h)\ge0$、$F(1-h)=-f(1-h)\le0$，故有零点。在第二个问题中，
\[
\sum_{k=0}^{n-1}F(k/n)=f(1)-f(0)=0.
\]
若所有差值为零，则已经得解；否则这些差值中必同时出现正数与负数，再由 $F$ 的连续性得到零点。这说明“任意步长”与“等分步长”依赖的条件不同。
\end{solve}

\begin{example}
连续函数在圆周上取实值，证明存在一对对径点函数值相同。另设 $f,g\in C([a,b])$ 且 $\max f=\max g$，证明存在 $x_0$ 使 $e^{f(x_0)}+3f(x_0)=e^{g(x_0)}+3g(x_0)$。
\end{example}

\begin{solve}
将圆周参数化，记对应的 $2\pi$ 周期连续函数为 $u(\theta)$，取 $F(\theta)=u(\theta)-u(\theta+\pi)$，则 $F(\pi)=-F(0)$，由介值性得零点。第二问令 $H(t)=e^t+3t$，它严格递增，故只需证明某点 $f=g$。在 $f$ 的最大值点 $x_1$ 有 $f(x_1)\ge g(x_1)$，在 $g$ 的最大值点 $x_2$ 有 $f(x_2)\le g(x_2)$，对 $f-g$ 用介值性即可。
\end{solve}

\begin{example}
\label{ex:pei-minimal-period}
证明：非常数的连续周期函数 $f:\mathbb R\to\mathbb R$ 存在最小正周期 $T_0$，所有正周期恰为 $T_0$ 的正整数倍。
\end{example}

\begin{solve}
设 $P$ 为所有正周期的集合，$T_0=\inf P\ge0$。取 $T_n\in P$ 使 $T_n\to T_0$，则对每个 $x$，由连续性有 $f(x+T_0)=\lim_n f(x+T_n)=f(x)$。还须证明 $T_0>0$。若 $T_n\to0$，对任意 $x\in\mathbb R$，取整数 $k_n=\lfloor x/T_n\rfloor$，则 $k_nT_n\to x$，而 $f(k_nT_n)=f(0)$，所以 $f(x)=f(0)$，违反非常数条件。

因此 $T_0$ 是最小正周期。任意 $T\in P$ 可写成 $T=kT_0+r$，其中 $k\in\mathbb N$、$0\le r<T_0$。周期的整数倍及两个周期之差仍满足平移不变关系，故若 $r>0$，便得到更小的正周期，矛盾。因此 $r=0$。特别地，若一个非常数连续函数有两个正周期，则它们的比必为有理数。
\end{solve}

\begin{example}
证明不存在连续函数 $f:\mathbb R\to\mathbb R$ 满足 $f(f(x))=e^{-x}$；也不存在在非退化区间上连续、在每个有理点取无理值而在每个无理点取有理值的函数。
\end{example}

\begin{solve}
第一问中，$e^{-x}$ 是单射，故 $f$ 必为单射；由命题~\ref{prop:pei-injective-ivt}，$f$ 严格单调。但无论 $f$ 递增还是递减，$f\circ f$ 都严格递增，与 $e^{-x}$ 严格递减矛盾。

第二问中，有理点只有可数多个，其像至多可数；无理点的像包含在 $\mathbb Q$ 中，也至多可数。因此整个像集可数。若 $f$ 不是常数，连续性和介值性使其像包含一个非退化区间，这个区间不可数，矛盾；若 $f$ 为常数，则由于定义域同时含有有理点与无理点，同一个常数又不可能同时是无理数和有理数，仍矛盾。
\end{solve}

\subsection{连续映射的迭代}

\begin{proposition}
若连续函数 $f:\mathbb R\to\mathbb R$ 没有不动点，则从任意 $x_0$ 出发的迭代 $x_{n+1}=f(x_n)$ 严格单调，且趋于 $+\infty$ 或 $-\infty$。
\end{proposition}

\begin{proof}
连续函数 $f(x)-x$ 无零点，由介值性，它在实轴上恒正或恒负。因此每次迭代都向同一方向严格移动。若单调轨道有有限极限 $\ell$，则连续性给出 $\ell=f(\ell)$，矛盾；所以只能向相应方向无界发散。这里单调的是迭代数列，结论并未要求 $f$ 本身单调。
\end{proof}

\begin{theorem}
\label{thm:pei-iteration-small-step}
设 $f:[a,b]\to[a,b]$ 连续，$x_{n+1}=f(x_n)$。则 $\{x_n\}$ 收敛当且仅当 $x_{n+1}-x_n\to0$。
\end{theorem}

\begin{proof}
必要性直接成立。反之，设步差趋于零而数列不收敛，令 $\alpha=\liminf x_n$、$\beta=\limsup x_n$，则 $a\le\alpha<\beta\le b$。先证明每个 $c\in(\alpha,\beta)$ 都是聚点：取充分小的 $\eta>0$，由于数列无限多次出现在 $c-\eta$ 左侧和 $c+\eta$ 右侧，且充分靠后的每一步长度小于 $\eta$，在这些穿越中必有任意靠后的项落入 $(c-\eta,c+\eta)$。令 $\eta\downarrow0$ 并依次选指标，便得到趋于 $c$ 的子列。

对任意聚点 $c$，若 $x_{n_k}\to c$，则步差趋零给出 $x_{n_k+1}\to c$，另一方面连续性给出 $x_{n_k+1}=f(x_{n_k})\to f(c)$，所以 $f(c)=c$。因而 $(\alpha,\beta)$ 内每一点都是不动点。数列既然以其中一点为聚点，就必有某项进入该开区间，之后立即保持常数，与 $\alpha<\beta$ 矛盾。故数列收敛。
\end{proof}

\emph[caution]{步差趋零本身不是一般数列收敛的充分条件}。例如 $\sin(\log(n+1))$ 有界、相邻项之差趋零却不收敛；定理依赖紧区间上的连续自治迭代结构，不能仅凭步差小就对任意递推式断言收敛。

\section{一致连续性}
\label{sec:pei-uniform-continuity}

\subsection{量词、序列判别与基本估计}

\begin{definition}
设 $E\subseteq\mathbb R$，$f:E\to\mathbb R$。若
\begin{equation}
\label{eq:pei-uniform-definition}
\forall\varepsilon>0\ \exists\delta>0\ \forall x,y\in E:
\quad |x-y|<\delta\Longrightarrow |f(x)-f(y)|<\varepsilon,
\end{equation}
则称 $f$ 在 $E$ 上\emph{一致连续（uniformly continuous）}。这里 $\delta$ 只依赖误差 $\varepsilon$，不依赖点的位置。否定形式为
\[
\exists\varepsilon_0>0\ \forall\delta>0\ \exists x,y\in E:
\quad |x-y|<\delta,\qquad |f(x)-f(y)|\ge\varepsilon_0.
\]
\end{definition}

\begin{proposition}
\label{prop:pei-uniform-sequences}
$f$ 在 $E$ 上一致连续，当且仅当对任意 $x_n,y_n\in E$，$|x_n-y_n|\to0$ 都能推出 $|f(x_n)-f(y_n)|\to0$。特别地，要证明不一致连续，只需找到两列距离趋零而函数值之差保持正下界的点。
\end{proposition}

\begin{proof}
必要性将式~\eqref{eq:pei-uniform-definition} 用于每对 $x_n,y_n$ 即得。若不一致连续，在否定形式中依次取 $\delta=1/n$，便得到 $|x_n-y_n|<1/n$ 而 $|f(x_n)-f(y_n)|\ge\varepsilon_0$，违反所给序列条件。
\end{proof}

\begin{definition}
若存在常数 $L\ge0$ 使所有 $x,y\in E$ 都满足 $|f(x)-f(y)|\le L|x-y|$，则称 $f$ 满足 \emph{Lipschitz 条件}。更一般地，若对某个 $0<\gamma\le1$ 及常数 $C\ge0$ 有 $|f(x)-f(y)|\le C|x-y|^\gamma$，则称其满足 $\gamma$ 阶 \emph{H\"older 条件}。两种条件均推出一致连续。
\end{definition}

若 $f$ 在区间上可导且 $|f'|\le L$，由中值定理可推出 Lipschitz 条件；在尚未学习微分学时，可直接对差值作代数估计。反向蕴含一般不成立：$\sqrt{x}$ 在 $[0,+\infty)$ 上由
\[
|\sqrt{x}-\sqrt{y}|\le\sqrt{|x-y|}
\]
一致连续，却不满足全局 Lipschitz 条件，因为 $\sqrt{x}/x\to+\infty$（$x\downarrow0$）。有界导数是方便的充分条件，不能作为一致连续性的定义或必要条件。

\begin{example}
判断 $1/x$、$x^2$ 与 $\sin(x^2)$ 在相应区间上的一致连续性。
\end{example}

\begin{solve}
在 $(0,1)$ 上取 $x_n=1/(n+1)$、$y_n=1/(n+2)$，两点距离趋零，但 $|1/x_n-1/y_n|=1$，故 $1/x$ 不一致连续；在 $[a,+\infty)$（$a>0$）上则有
\[
\left|\frac1x-\frac1y\right|=\frac{|x-y|}{xy}\le\frac{|x-y|}{a^2},
\]
因而一致连续。对 $x^2$ 在 $[0,+\infty)$ 或 $\mathbb R$ 上取 $x_n=n$、$y_n=n+1/n$，平方之差为 $2+1/n^2$，故不一致连续。对 $\sin(x^2)$ 取
\[
x_n=\sqrt{2\pi n+\pi/2},\qquad y_n=\sqrt{2\pi n+3\pi/2}.
\]
两点距离等于 $\pi/(x_n+y_n)\to0$，函数值之差却恒为 $2$。这个例子同时说明，有界、连续乃至处处可导，都不能单独保证无穷区间上的一致连续性。
\end{solve}

\begin{proposition}
一致连续函数的和、差及一致连续映射的复合仍一致连续。有界且一致连续的两个函数之积一致连续；若其中作为分母的函数绝对值还有正下界，则商也一致连续。
\end{proposition}

\begin{proof}
和、差由三角不等式分别分配误差；复合时先对外层选择自变量误差，再由内层将该误差转成原自变量的距离限制。若 $|f|\le M$、$|g|\le N$，则
\[
|f(x)g(x)-f(y)g(y)|\le M|g(x)-g(y)|+N|f(x)-f(y)|.
\]
若另有 $|g|\ge c>0$，则 $|1/g(x)-1/g(y)|\le c^{-2}|g(x)-g(y)|$，再用积的结论。无穷区间上不能删去积的有界条件：$f(x)=g(x)=x$ 都一致连续，但乘积 $x^2$ 不一致连续。分母处处非零也不能代替正下界，例如 $(0,1)$ 上的 $1$ 与 $x$ 都一致连续，商 $1/x$ 却不是。
\end{proof}

\subsection{紧区间、端点延拓与 Cauchy 序列}

\begin{theorem}[Heine--Cantor 定理]
\label{thm:pei-heine-cantor}
闭区间上的连续函数一致连续。
\end{theorem}

\begin{proof}
若 $f\in C([a,b])$ 不一致连续，则由命题~\ref{prop:pei-uniform-sequences} 的否定形式，可取 $x_n,y_n\in[a,b]$，使 $|x_n-y_n|\to0$ 而函数值之差不小于某个 $\varepsilon_0>0$。取子列 $x_{n_k}\to c$，则 $y_{n_k}\to c$，由连续性两列函数值都趋于 $f(c)$，矛盾。
\end{proof}

\begin{theorem}
\label{thm:pei-endpoint-extension}
设 $a<b$ 为有限实数，$f\in C((a,b))$。则 $f$ 在 $(a,b)$ 上一致连续，当且仅当有限极限 $f(a+)$ 与 $f(b-)$ 都存在；等价地，$f$ 能连续延拓至 $[a,b]$。因此有限区间上的一致连续函数必有界。
\end{theorem}

\begin{proof}
若 $f$ 一致连续，给定 $\varepsilon>0$，取定义中的 $\delta>0$。当 $a<x,y<a+\delta$ 时有 $|x-y|<\delta$，故 $|f(x)-f(y)|<\varepsilon$。由函数极限的 Cauchy 准则，$f(a+)$ 存在且有限；右端同理。以这两个极限补定义端点值，便得到闭区间上的连续函数，再由定理~\ref{thm:pei-heine-cantor} 得到充分性及有界性。其他有限区间按其缺失的端点作同样处理。
\end{proof}

例如 $\sin x/x$ 在 $(0,1)$ 上可于 $0$ 补值为 $1$，故一致连续；$\sin(1/x)$ 在 $(0,1)$ 上虽有界连续，却无 $0$ 处的右极限，故不一致连续。有限开区间上，两个一致连续函数都自动有界，因而乘积一致连续；这一推论不能直接推广到无穷区间。

\begin{proposition}
\label{prop:pei-cauchy-preserving}
一致连续映射将 Cauchy 序列映为 Cauchy 序列。若 $E\subseteq\mathbb R$ 有界，则反之也成立：$f:E\to\mathbb R$ 将每个 $E$ 中的 Cauchy 序列映为 Cauchy 序列，当且仅当 $f$ 一致连续。
\end{proposition}

\begin{proof}
必要性直接将输入序列的 Cauchy 条件与一致连续的 $\delta$ 相接。反之，若 $f$ 不一致连续，则有距离趋零但像距离保持正下界的两列 $x_n,y_n\in E$。有界性使某子列 $x_{n_k}\to c$，同时 $y_{n_k}\to c$。交错排列 $x_{n_1},y_{n_1},x_{n_2},y_{n_2},\ldots$ 得到 $E$ 中的 Cauchy 序列，但其像中每对相邻指定项之差均不小于 $\varepsilon_0$，不是 Cauchy 序列，矛盾。这里不要求极限 $c$ 属于 $E$。
\end{proof}

在无界集合上，保持 Cauchy 序列未必推出一致连续。实轴上的 $x^2$ 就是反例：每个实数 Cauchy 序列都收敛，平方后的序列仍收敛，但平方函数在实轴上不一致连续。这一限制来自充分性证明所用的有界子列抽取，不能省略。

\begin{theorem}
\label{thm:pei-dense-extension}
设 $D$ 在 $E\subseteq\mathbb R$ 中稠密。若 $u:D\to\mathbb R$ 一致连续，则存在唯一的连续延拓 $f:E\to\mathbb R$，且 $f$ 仍一致连续。更一般地，若 $u$ 在 $D$ 的每个有界子集上一致连续，则它也有唯一的连续延拓，并且延拓在 $E$ 的每个有界子集上一致连续。
\end{theorem}

\begin{proof}
对 $x\in E$，由稠密性取 $x_n\in D$，$x_n\to x$。这些点构成有界 Cauchy 序列，由相应的一致连续性，$u(x_n)$ 是实数 Cauchy 序列，故可定义 $f(x)=\lim_nu(x_n)$。若另取 $y_n\to x$，将两列交错排列便知两列像有相同极限，故定义与选择无关。对 $x\in D$ 取常值列即知 $f(x)=u(x)$。

先考虑全局一致连续的情形。对 $\varepsilon/2$ 取 $u$ 的统一尺度 $\delta$。若 $x,y\in E$ 且 $|x-y|<\delta/2$，分别选 $D$ 中的逼近列，则充分靠后的两列点距离小于 $\delta$，从而 $|u(x_n)-u(y_n)|<\varepsilon/2$；取极限有 $|f(x)-f(y)|\le\varepsilon/2<\varepsilon$。若仅有有界子集上的一致连续性，对任意有界 $B\subseteq E$，可在一个固定的更大有界区间内选择所有逼近列，重复论证便知 $f|_B$ 一致连续。因此 $f$ 在各点连续。最后，任意两个连续延拓在稠密集 $D$ 上相同，沿逼近列取极限便在 $E$ 上处处相同。
\end{proof}

这一定理特别给出：$D$ 上的一致连续实函数能唯一延拓至闭包 $\overline D$。值域取 $\mathbb R$ 的完备性保证像 Cauchy 序列有极限；若值域不是完备空间，就不能直接照搬这一存在性结论。

\subsection{无穷区间上的尾部控制}

\begin{proposition}
\label{prop:pei-uniform-gluing}
若 $f$ 在 $[a,c]$ 与 $[c,b]$ 上分别一致连续，且在公共点采用同一函数值，则它在 $[a,b]$ 上一致连续；任一端点允许为无穷远，此时按对应的半无限区间理解。
\end{proposition}

\begin{proof}
对 $\varepsilon/2$ 分别取两个尺度，令 $\delta$ 为二者的最小值。若两点属于同一段，直接使用相应估计；若 $x<c<y$ 且 $|x-y|<\delta$，则 $|x-c|,|y-c|<\delta$，通过公共点 $c$ 使用三角不等式，有 $|f(x)-f(y)|<\varepsilon$。
\end{proof}

\begin{theorem}
\label{thm:pei-uniform-tail}
若 $f\in C([a,+\infty))$ 且 $\lim_{x\to+\infty}f(x)=L\in\mathbb R$，则 $f$ 一致连续。更一般地，若 $g$ 在 $[a,+\infty)$ 上一致连续，$f$ 连续且 $f(x)-g(x)\to0$，则 $f$ 也一致连续。
\end{theorem}

\begin{proof}
对第一种情形，给定 $\varepsilon>0$，取 $A>a$，使 $x\ge A$ 时有 $|f(x)-L|<\varepsilon/2$；于是两点均在尾部时，函数值之差小于 $\varepsilon$。又由 Heine--Cantor 定理，$f$ 在固定于本次误差的闭区间 $[a,A+1]$ 上一致连续，取其尺度 $\delta_0$，再令 $\delta=\min\{\delta_0,1\}$。距离小于 $\delta$ 的两点，若较小者不超过 $A$，两者都在 $[a,A+1]$；否则两者都在尾部。这就覆盖了全部情形。

对一般情形，将误差拆成三份。先取 $A$ 使 $|f-g|<\varepsilon/3$ 在尾部成立，再取 $g$ 对 $\varepsilon/3$ 的全局尺度，以及 $f$ 在 $[a,A+1]$ 上对 $\varepsilon$ 的尺度，并使最终尺度不超过 $1$。尾部使用
\[
|f(x)-f(y)|\le|f(x)-g(x)|+|g(x)-g(y)|+|g(y)-f(y)|,
\]
其余情形使用闭区间估计，结论即得。
\end{proof}

此证明需要重叠区间 $[a,A+1]$ 来处理跨过分界点的两点。还要注意，由“对每个误差可以选出一个尾部，使其中任意两点的像很接近”不能断言存在一个固定尾区间在所有误差下都满足同一性质，因为 $A$ 本身依赖误差。两端均有有限极限的实轴连续函数也一致连续；连续周期函数则可把任意两点同时平移到一个固定的、长度略大于周期的闭区间，在该区间应用 Heine--Cantor 定理。具体地，若周期为 $T>0$，先要求 $|x-y|<T$，再同时平移使 $x\in[0,T]$，则 $y\in[-T,2T]$，从而有统一尺度。

\begin{proposition}
\label{prop:pei-uniform-linear-growth}
若 $f:\mathbb R\to\mathbb R$ 一致连续，则存在 $A,B\ge0$，使 $|f(x)|\le A|x|+B$ 对所有实数 $x$ 成立。
\end{proposition}

\begin{proof}
对误差 $1$ 取统一尺度 $\delta>0$，令 $h=\delta/2$。对 $x\ne0$，取 $N=\lceil |x|/h\rceil$，将 $0$ 到 $x$ 等分成 $N$ 段，每段长度不超过 $h<\delta$，故
\[
|f(x)-f(0)|\le\sum_{j=1}^N|f(jx/N)-f((j-1)x/N)|<N\le |x|/h+1.
\]
取 $A=1/h$、$B=|f(0)|+1$ 即可。$x=0$ 时也成立。
\end{proof}

增长至多线性只是必要条件，不能反推一致连续，前面的有界反例 $\sin(x^2)$ 已经说明这一点。对可导函数还有一种尾部判别：若 $f\in C([0,+\infty))$，在 $(0,+\infty)$ 可导，且 $\lim_{x\to\infty}|f'(x)|=A$ 存在，则 $A<\infty$ 时可用尾部有界导数与头部紧区间拼接，得到一致连续；若 $A=+\infty$，可选 $x_n\ge n$，使 $x\ge x_n$ 时 $|f'(x)|\ge n$，再令 $y_n=x_n+1/n$，由中值定理有 $|f(y_n)-f(x_n)|\ge1$，故不一致连续。该判别以导数绝对值的极限存在为前提。

\begin{example}
设 $f$ 在 $[0,+\infty)$ 上一致连续，且每个固定的 $t\ge0$ 都满足 $f(n+t)\to0$（$n\to\infty$，$n\in\mathbb N$）。证明 $f(x)\to0$（$x\to+\infty$）。
\end{example}

\begin{solve}
给定 $\varepsilon>0$，取一致连续性对 $\varepsilon/2$ 的尺度 $\delta$，再在 $[0,1]$ 内选择有限个网格点 $t_1,\ldots,t_m$，使每个 $r\in[0,1]$ 与某个 $t_i$ 距离小于 $\delta$。对每个 $i$，存在 $N_i$，使 $n\ge N_i$ 时 $|f(n+t_i)|<\varepsilon/2$。令 $N=\max_iN_i$。任意充分大的 $x$ 可写成 $x=n+r$，其中 $n=\lfloor x\rfloor\ge N$、$r\in[0,1)$，选相应网格点便有
\[
|f(x)|\le|f(n+r)-f(n+t_i)|+|f(n+t_i)|<\varepsilon.
\]
有限网格使有限个收敛起点可以取最大值，一致连续性则把网格控制传递到全部实数。
\end{solve}

\subsection{连续模数}

\begin{definition}
设 $I$ 是区间，$f:I\to\mathbb R$。对 $t\ge0$，定义\emph{连续模数（modulus of continuity）}
\begin{equation}
\label{eq:pei-modulus}
\omega_f(t)=\sup\{|f(x)-f(y)|:x,y\in I,\ |x-y|\le t\}.
\end{equation}
一般允许其值为 $+\infty$。它非负、单调不减，且 $\omega_f(0)=0$；$f$ 一致连续当且仅当 $\omega_f(t)\to0$（$t\downarrow0$）。
\end{definition}

\begin{proof}
若 $f$ 一致连续，先对 $\varepsilon/2$ 取 $\delta$，当 $0<t<\delta$ 时，定义中的每个差值均小于 $\varepsilon/2$，故上确界不超过 $\varepsilon/2<\varepsilon$。反之，选 $\delta>0$ 使 $\omega_f(\delta)<\varepsilon$，则 $|x-y|<\delta$ 时的差值都小于 $\varepsilon$。
\end{proof}

在区间上还可沿两点之间分段，得到 $\omega_f(s+t)\le\omega_f(s)+\omega_f(t)$。若 $f$ 一致连续，先找一个小尺度使其模数有限，再把任意有限距离分成有限段，即知所有有限 $t$ 的模数都有限。例如在 $[0,+\infty)$ 上，$\omega_{\sqrt{\cdot}}(t)=\sqrt t$；在 $[a,+\infty)$（$a>0$）上，
\[
\omega_{1/\cdot}(t)=\frac{t}{a(a+t)};
\]
在实轴上，$\omega_{\sin}(t)\le\min\{t,2\}$。连续模数记录的是给定水平尺度下的最大振幅，比只看某点的斜率更适合判断一致性。

\section{上、下半连续}
\label{sec:pei-semicontinuity}

\subsection{定义与等价描述}

\begin{definition}
设 $f:E\to\mathbb R$、$x_0\in E$。若对每个 $\varepsilon>0$，存在 $\delta>0$，使 $x\in E$、$|x-x_0|<\delta$ 时有 $f(x)<f(x_0)+\varepsilon$，则称 $f$ 在 $x_0$ \emph{上半连续（upper semicontinuous）}。若相应地有 $f(x)>f(x_0)-\varepsilon$，则称其\emph{下半连续（lower semicontinuous）}。两者同时成立当且仅当 $f$ 在该点连续。
\end{definition}

\begin{proposition}
\label{prop:pei-semicontinuity-characterization}
$f$ 在 $E$ 上上半连续，当且仅当对每个收敛于 $x_0\in E$ 的序列 $x_n\in E$，有 $\limsup_nf(x_n)\le f(x_0)$；也等价于每个严格下水平集 $\{f<c\}$ 在 $E$ 中相对开，或每个上水平集 $\{f\ge c\}$ 在 $E$ 中相对闭。下半连续的对应条件为 $\liminf_nf(x_n)\ge f(x_0)$、$\{f>c\}$ 相对开、$\{f\le c\}$ 相对闭。若 $E$ 本身闭，这些相对闭集也都是实轴中的闭集。
\end{proposition}

\begin{proof}
从上半连续定义可直接取序列的上极限。若定义不成立，则存在 $\varepsilon_0>0$，对每个 $n$ 能取 $|x_n-x_0|<1/n$ 而 $f(x_n)\ge f(x_0)+\varepsilon_0$，违反序列条件。若上半连续且 $f(x_0)<c$，对小于 $c-f(x_0)$ 的误差使用定义，可得 $x_0$ 的相对邻域仍在 $\{f<c\}$ 中。反之，在 $c=f(x_0)+\varepsilon$ 时使用水平集的开性即得定义。闭集形式由补集得到；下半连续的结论通过将 $f$ 换为 $-f$ 得到。
\end{proof}

上半连续限制函数从附近逼近时不能超过点值，却允许点值向上跳；下半连续具有对偶方向。因此上半连续在负值点保负，下半连续在正值点保正；反方向一般不成立，介值性也不随半连续性保留。Dirichlet 函数 $D=\boldsymbol{1}_{\mathbb Q}$ 在有理点上半连续而不下半连续，在无理点下半连续而不上半连续。例~\ref{ex:pei-thomae} 的 Riemann 函数处处上半连续：无理点已经连续；在有理点，对给定的高度阈值排除附近有限个分母较小的其他有理点，即得上半连续，而无理数列趋近该点说明它不下半连续。

\subsection{运算、最值与单调极限}

\begin{proposition}
有限个上半连续实函数的和、最大值与最小值仍上半连续，乘以非负常数也保持上半连续；取负号则变成下半连续。非负上半连续函数的乘积上半连续。任意一族上半连续函数的逐点下确界若为有限实值，仍上半连续；任意一族下半连续函数的逐点上确界若为有限实值，仍下半连续。
\end{proposition}

\begin{proof}
和的结论将误差分配给各项即可。最大值的严格下水平集是有限个相应开集的交，最小值的严格下水平集是相应开集的并，故都上半连续。对非负乘积，设点值为 $A,B\ge0$，在邻域中分别用 $A+\eta$、$B+\eta$ 从上控制两因子，再令 $\eta$ 足够小，使 $(A+\eta)(B+\eta)<AB+\varepsilon$。最后，若 $h=\inf_{\lambda}f_\lambda$，则
\[
\{h<c\}=\bigcup_\lambda\{f_\lambda<c\},
\]
为相对开集。下半连续的对偶结论由取负号得到。
\end{proof}

乘积的符号条件有实质作用。例如 $\boldsymbol{1}_{\{0\}}$ 与常函数 $-1$ 都上半连续，但二者之积在 $0$ 处不上半连续。由下确界的结论可知，递减上半连续函数列的有限逐点极限仍上半连续；递增下半连续函数列的有限逐点极限仍下半连续。这里单调方向不可随意交换。

\begin{theorem}
\label{thm:pei-semicontinuity-extrema}
闭区间上的有限实值上半连续函数有上界并达到最大值；有限实值下半连续函数有下界并达到最小值。
\end{theorem}

\begin{proof}
对每个点 $x$，由上半连续性取相对邻域，使其中 $f(t)<f(x)+1$；有限子覆盖便给出统一上界。记 $M=\sup f$，取 $f(x_n)\to M$，再抽取 $x_{n_k}\to c\in[a,b]$。由命题~\ref{prop:pei-semicontinuity-characterization}，有 $M\le f(c)$，而上确界定义又给出 $f(c)\le M$，所以 $f(c)=M$。下半连续情形对 $-f$ 使用结论。
\end{proof}

上半连续并不保证下界。例如在 $[0,1]$ 上令 $f(0)=0$，$f(x)=-1/x$（$x>0$），它上半连续、最大值为 $0$，却无下界。若只在开区间上上半连续，连上界也未必存在。不过仍有一个较弱的局部结论：任意有限实值半连续函数都能在某个非退化内闭区间上同时有上、下界。

\begin{proof}
只证上半连续情形。若每个非退化内闭区间上都无下界，从任一 $I_0\subset(a,b)$ 出发，可在 $I_0$ 的内部取 $x_1$ 使 $f(x_1)<-1$，再由上半连续性取非退化闭区间 $I_1\subset\operatorname{int}I_0$，使 $f<-1$ 在 $I_1$ 上成立。依次在 $\operatorname{int}I_{n-1}$ 内取 $f<-n$ 的点，并缩成长度不超过 $2^{-n}$、其上恒有 $f<-n$ 的闭区间 $I_n$。闭区间套定理给出 $c\in\bigcap_nI_n$，于是 $f(c)<-n$ 对所有 $n$ 成立，与 $f(c)$ 为有限实数矛盾。因此某个内闭区间上有下界；再由定理~\ref{thm:pei-semicontinuity-extrema} 得到上界。下半连续的情形对 $-f$ 处理。
\end{proof}

\subsection{用连续函数单调逼近}

\begin{theorem}
\label{thm:pei-semicontinuity-approximation}
设 $f:[a,b]\to\mathbb R$ 上半连续。则存在递减的连续函数列 $f_n$，使 $f_n(x)\downarrow f(x)$ 对每个 $x\in[a,b]$ 成立。下半连续函数则可由递增连续函数列逐点逼近。
\end{theorem}

\begin{proof}
由定理~\ref{thm:pei-semicontinuity-extrema}，$f$ 有有限上界 $M$。对 $n\in\mathbb N$ 定义
\begin{equation}
\label{eq:pei-semicontinuity-envelope}
f_n(x)=\sup_{y\in[a,b]}\bigl(f(y)-n|x-y|\bigr).
\end{equation}
取 $y=x$ 得 $f_n(x)\ge f(x)$，而 $f_n(x)\le M$，故函数有限；随 $n$ 增大它单调不增。三角不等式给出 $f(y)-n|x-y|\le f(y)-n|z-y|+n|x-z|$，取上确界并交换 $x,z$ 后得 $|f_n(x)-f_n(z)|\le n|x-z|$，因此 $f_n$ 连续。

固定 $x$ 与 $\varepsilon>0$，由上半连续性选 $\delta>0$，使 $|y-x|<\delta$ 时有 $f(y)<f(x)+\varepsilon$。对于其余 $y$，有 $f(y)-n|x-y|\le M-n\delta$，在充分大的 $n$ 时不超过 $f(x)+\varepsilon$。分开这两类点取上确界，就有 $f(x)\le f_n(x)\le f(x)+\varepsilon$，故 $f_n(x)\downarrow f(x)$。对下半连续函数 $g$，先逼近 $-g$ 再取负号即可。
\end{proof}

式~\eqref{eq:pei-semicontinuity-envelope} 给出的是逐点收敛；若 $f$ 不连续，就不可能是一致收敛，因为连续函数列的一致极限连续。这个构造也说明，上半连续函数可以看成一列连续上界逐步下降所形成的极限，而下半连续函数具有对偶的连续下界逼近。

\section{函数方程}
\label{sec:pei-functional-equations}

\subsection{加法方程与从有理数到实数的推归}

函数方程同时约束不同自变量处的函数值。有效的起点通常是令某个变量为 $0$、$1$、另一个变量或其相反数，先确定特殊值、奇偶性和整数倍关系，再推到有理倍，最后利用连续性或单调性补足实数。连续性等正则条件是方程假设的一部分，不能在推导完成前默认所有代数解都连续。

\begin{theorem}[Cauchy 加法方程]
\label{thm:pei-cauchy-additive}
设 $f:\mathbb R\to\mathbb R$ 满足
\begin{equation}
\label{eq:pei-cauchy-additive}
f(x+y)=f(x)+f(y)\qquad(x,y\in\mathbb R).
\end{equation}
若 $f$ 在 $0$ 处连续，则存在 $c\in\mathbb R$ 使 $f(x)=cx$ 对所有实数成立。在任意一点连续、在某个非退化区间上有界、在实轴上单调，都分别足以替代在 $0$ 处连续的条件。
\end{theorem}

\begin{proof}
令 $x=y=0$ 得 $f(0)=0$，再令 $y=-x$ 得 $f(-x)=-f(x)$。归纳得 $f(nx)=nf(x)$（$n\in\mathbb Z$），对 $nx/n=x$ 有 $f(x/n)=f(x)/n$，所以 $f(qx)=qf(x)$ 对有理数 $q$ 成立。记 $c=f(1)$，则 $f(q)=cq$。任给实数 $x$，取有理数列 $q_n\to x$，由在 $0$ 处连续及
\[
f(x)-f(q_n)=f(x-q_n)\to0,
\]
得到 $f(x)=\lim_ncq_n=cx$，并直接代回检验。

若 $f$ 在 $x_0$ 连续，则 $f(h)=f(x_0+h)-f(x_0)\to0$，故在 $0$ 连续。若 $|f|\le M$ 在某个非退化区间上成立，在其中选内点 $x_0$ 和 $r>0$，使 $x_0+h$ 在该区间内（$|h|<r$），则 $|f(h)|\le2M$。给定 $\varepsilon>0$，取整数 $N$ 使 $2M/N<\varepsilon$，当 $|h|<r/N$ 时有
\[
|f(h)|=\frac{|f(Nh)|}{N}\le\frac{2M}{N}<\varepsilon.
\]
因此仍在 $0$ 连续。单调函数在任意固定的有限闭区间上有界，故也适用。若 $f$ 在某个非退化闭区间上 Riemann 可积，也由其有界性归入同一情形。
\end{proof}

\begin{example}
求连续的 Jensen 中点方程
\[
f\left(\frac{x+y}{2}\right)=\frac{f(x)+f(y)}2\qquad(x,y\in\mathbb R)
\]
的全部解，并讨论二阶差分为零的条件。
\end{example}

\begin{solve}
令 $b=f(0)$、$g=f-b$，则 $g(0)=0$。在中点方程中取 $(x,y)=(2t,0)$，得到 $g(2t)=2g(t)$，因此 $g(x+y)=g(x)+g(y)$，由定理~\ref{thm:pei-cauchy-additive} 知 $f(x)=cx+b$。反之，所有仿射函数都满足方程。

若以 $\Delta_hf(x)=f(x+h)-f(x)$ 记差分，则要求 $\Delta_h^2f(x)=0$ 对\emph{所有 $x,h\in\mathbb R$} 成立，等价于 $f(x+2h)+f(x)=2f(x+h)$，也就是任意两点的中点方程，所以连续解仍为仿射函数。若只固定一个非零步长 $h$，结论就不成立，例如 $f(x)=x+\sin(2\pi x/h)$ 的该步长二阶差分恒为零，却不是仿射函数。
\end{solve}

\subsection{通过换元化为加法方程}

\begin{proposition}
\label{prop:pei-exponential-equations}
在各自定义域上连续的下列方程，其全部实值解分别为
\[
\begin{aligned}
f(x+y)=f(x)f(y)\quad &(x,y\in\mathbb R):
&&f\equiv0\ \text{或}\ f(x)=e^{cx};\\
f(xy)=f(x)+f(y)\quad &(x,y>0):
&&f(x)=c\log x;\\
f(xy)=f(x)f(y)\quad &(x,y>0):
&&f\equiv0\ \text{或}\ f(x)=x^c,
\end{aligned}
\]
其中 $c\in\mathbb R$，$\log$ 表示自然对数。
\end{proposition}

\begin{proof}
第一式若非零解，则取 $y=0$ 得 $f(0)=1$，再由 $f(x)f(-x)=1$ 知处处非零；同时 $f(x)=f(x/2)^2>0$。于是 $F(x)=\log f(x)$ 连续且可加，故 $F(x)=cx$。第二式令 $F(t)=f(e^t)$，则 $F(s+t)=F(s)+F(t)$，故 $F(t)=ct$。第三式若非零解，先由 $f(1)=1$、$f(x)f(1/x)=1$ 排除零点，再由 $f(x)=f(\sqrt{x})^2>0$ 得到正值；令 $F(t)=\log f(e^t)$，便又化为连续加法方程。这些候选函数代回原式均成立。
\end{proof}

对数换元必须先证明取对数的对象为正，不能从“连续”直接推出正值。乘法方程的定义域也影响解的形式：若要求 $f:\mathbb R\to\mathbb R$ 连续且满足 $f(xy)=f(x)f(y)$，则全部解为 $f\equiv0$、$f\equiv1$，以及对 $c>0$ 的
\[
f(0)=0,\qquad f(x)=|x|^c\ (x\ne0),
\]
或 $f(0)=0$、$f(x)=\operatorname{sgn}(x)|x|^c$（$x\ne0$）。事实上，正半轴上的非零解是 $x^c$，$f(-1)^2=1$ 决定负半轴的两种符号；若不是常函数 $1$，方程迫使 $f(0)=0$，原点连续又迫使 $c>0$。若去掉原点连续，$\operatorname{sgn}x$ 也满足乘法方程，这说明检查特殊点不能省略。

\subsection{余弦方程与联立加法公式}

\begin{theorem}
\label{thm:pei-dalembert}
连续函数 $f:\mathbb R\to\mathbb R$ 满足 d'Alembert 方程
\begin{equation}
\label{eq:pei-dalembert}
f(x+y)+f(x-y)=2f(x)f(y)
\end{equation}
当且仅当它属于 $f\equiv0$、$f(x)=\cos(cx)$、$f(x)=\cosh(cx)$ 这三类，其中 $c\ge0$，$\cosh t=(e^t+e^{-t})/2$。
\end{theorem}

\begin{proof}
零函数可直接检验。对非零解，取 $y=0$ 得 $f(0)=1$，取 $x=0$ 得 $f(-y)=f(y)$。再令 $x=y=t$ 得
\begin{equation}
\label{eq:pei-dalembert-halving}
f(2t)=2f(t)^2-1.
\end{equation}
由连续性可取 $r>0$，使 $f>0$ 在 $[0,r]$ 上成立。若 $0<f(r)\le1$，写成 $f(r)=\cos\theta$，$0\le\theta<\pi/2$。由式~\eqref{eq:pei-dalembert-halving} 及正值条件，每次开方只能取正根，因此
\[
f(r/2^n)=\cos(\theta/2^n).
\]
固定 $n$、令 $h=r/2^n$，原方程给出递推
\[
f((m+1)h)=2f(mh)f(h)-f((m-1)h).
\]
由 $f(0)=1$、$f(h)=\cos(\theta/2^n)$ 归纳可得 $f(mr/2^n)=\cos(m\theta/2^n)$（$m\ge0$）。这些非负二进有理网格点在 $[0,+\infty)$ 稠密，由连续性得到 $f(x)=\cos(\theta x/r)$，再由偶性扩至实轴。

若 $f(r)>1$，写成 $f(r)=\cosh\theta$（$\theta>0$）。同一正根选择与递推，将上式中的 $\cos$ 换为 $\cosh$，得到 $f(x)=\cosh(\theta x/r)$。反过来，用余弦或双曲余弦加法公式检验即可。这一证明只用连续性、二分与递推，不预设可导性。
\end{proof}

\begin{example}
设连续有界实函数 $f,g$ 满足
\[
\begin{aligned}
f(x+y)&=f(x)f(y)-g(x)g(y),\\
g(x+y)&=f(x)g(y)+g(x)f(y),
\end{aligned}
\qquad f(0)=1,\quad g(0)=0.
\]
证明它们恰为 $f(x)=\cos(cx)$、$g(x)=\sin(cx)$（$c\in\mathbb R$）。
\end{example}

\begin{solve}
令 $H=f^2+g^2$，两式平方相加得到 $H(x+y)=H(x)H(y)$，且 $H(0)=1$，由命题~\ref{prop:pei-exponential-equations} 知 $H(x)=e^{kx}$。因为 $H$ 在整个实轴有界，只能 $k=0$，故 $f^2+g^2=1$。将 $y=-x$ 代入原方程，解二元线性方程可得 $f(-x)=f(x)$、$g(-x)=-g(x)$。把 $y$ 与 $-y$ 的第一式相加，即得 $f$ 满足式~\eqref{eq:pei-dalembert}。又 $f(0)=1$、$|f|\le1$，所以 $f(x)=\cos(ax)$，其中 $a\ge0$。

若 $a=0$，则 $g\equiv0$。若 $a>0$，取 $t$ 使 $\sin(at)\ne0$，则 $g(t)^2=\sin^2(at)$，且第一式给出
\[
2g(x)g(t)=f(x-t)-f(x+t)=2\sin(ax)\sin(at).
\]
因此 $g$ 的符号选择在整个实轴一致，$g(x)=\sin(ax)$ 或 $g(x)=-\sin(ax)$；将后一种情形吸收到参数 $c=-a$ 中，即得结论。
\end{solve}

\subsection{缩放迭代与连续点的作用}

\begin{example}
若 $f:(0,+\infty)\to\mathbb R$ 满足 $f(2x)=f(x)$ 且 $\lim_{x\downarrow0}f(x)=A$，证明 $f\equiv A$。若改为 $f:\mathbb R\to\mathbb R$ 在 $0$ 连续，$f(0)=1$，并满足 $f(2x)=f(x)\cos x$，求全部解。
\end{example}

\begin{solve}
第一问反复将自变量减半，有 $f(x)=f(x/2^n)\to A$。第二问同样迭代得
\[
f(x)=f(x/2^n)\prod_{k=1}^n\cos(x/2^k).
\]
由反复使用 $\sin(2t)=2\sin t\cos t$，对 $x\ne0$、充分大的 $n$ 有
\[
\prod_{k=1}^n\cos(x/2^k)=\frac{\sin x}{2^n\sin(x/2^n)}\longrightarrow\frac{\sin x}{x}.
\]
这里分母的极限为 $x$，而 $f(x/2^n)\to1$，所以唯一可能的解为 $f(x)=\sin x/x$（$x\ne0$）、$f(0)=1$。补值后的函数连续且代回满足方程，故确为唯一解。这类方法的核心是把任意点沿迭代送到一个已知极限或连续点。
\end{solve}

\subsection{转化为微分方程时的定义域检查}

以下方法需要导数及一阶微分方程知识，可在学习微分学后再读。若函数方程能把任意 $x$ 处的差商写成 $0$ 处的差商乘以一个连续因子，那么在 $0$ 处可导常能推出处处可导；随后解微分方程只是必要步骤，仍须把所得函数放回原来的定义域检验。

\begin{example}
\label{ex:pei-tangent-domain}
设 $\lambda>0$，实函数 $f$ 连续、在 $0$ 处可导，且满足
\begin{equation}
\label{eq:pei-tangent-equation}
f(x+y)=\frac{f(x)+f(y)}{1-\lambda f(x)f(y)}.
\end{equation}
在原点附近求其局部形式，并判断若该等式要求对所有 $x,y\in\mathbb R$ 成立且分母始终非零，有哪些全实轴解。
\end{example}

\begin{solve}
令 $x=y=0$，由实值及 $\lambda>0$ 得 $f(0)=0$。记 $c=f'(0)$，将 $y=h$ 代入并移项，有
\[
\frac{f(x+h)-f(x)}h=
\frac{f(h)}h\,\frac{1+\lambda f(x)^2}{1-\lambda f(x)f(h)}.
\]
因 $f(h)\to0$，每个所讨论的内点均可取极限，从而 $f'(x)=c(1+\lambda f(x)^2)$。利用链式法则积分并代入初值，得到
\[
\arctan(\sqrt\lambda f(x))=c\sqrt\lambda\,x,
\qquad
f(x)=\frac1{\sqrt\lambda}\tan(c\sqrt\lambda\,x).
\]
这是包含 $0$ 且避开极点的区间内的局部解，原方程只在 $x,y,x+y$ 均属于相应区间、分母非零时使用。若要求 $f$ 定义于整个实轴，则左边的反正切有界，右边只有在 $c=0$ 时才能对全部实数成立，所以全实轴解只有 $f\equiv0$。不能把有极点的正切函数直接称为全实轴连续解。
\end{solve}

\begin{example}
设 $f:\mathbb R\to\mathbb R$ 连续，在 $0$ 处可导，$f(0)=0$，并对所有实数 $x,y$ 满足
\[
f(x+y)=\frac{f(x)+f(y)}{1+f(x)f(y)},
\]
且分母非零。求 $f$。
\end{example}

\begin{solve}
令 $c=f'(0)$，与例~\ref{ex:pei-tangent-domain} 相同的差商运算给出 $f'(x)=c(1-f(x)^2)$、$f(0)=0$。在原点附近 $|f|<1$，分离变量得到
\[
\frac12\log\frac{1+f(x)}{1-f(x)}=cx,
\qquad f(x)=\tanh(cx)=\frac{e^{2cx}-1}{e^{2cx}+1}.
\]
这一表达式在实轴上始终位于 $(-1,1)$；由微分方程右端对未知函数值局部 Lipschitz 所保证的初值解唯一性，原函数在全实轴上都与它相同。直接代入双曲正切加法公式可验算原方程，且 $1+f(x)f(y)>0$，所以所有 $c\in\mathbb R$ 均给出解。这里假设 $f(0)=0$ 已排除常函数 $1$ 与 $-1$；若不加该初值，它们也满足原方程。
\end{solve}

函数方程的完整解答因此需要同时检查四件事：代数推导是否涵盖零解或常值解，连续等正则条件在何处使用，换元与取对数是否保持定义域，以及求得的公式是否在原先要求的全部自变量上有意义。这些检查既决定解的唯一性，也决定某个公式究竟只是局部解还是全局解。

% 教材来源：数学分析/textbooks/数分习题（分章原文）/数分习题-3.pdf。
% 对应原书第三章“一元微分学”，本 PDF 第 1--102 页；按主题整合要点与代表例题，不逐题誊录。
% 扫描原文通过嵌入扫描数据 OCR 核对；无法可靠辨识的个别题目不按原题转写。
\chapter{一元微分学}
\label{chap:pei-diff}

一元微分学以函数的局部线性近似为出发点，通过中值定理把局部变化率与区间上的函数增量联系起来，再用 Taylor 公式描述更高阶的局部变化。本章依据教材第三章，依次讨论导数、微分中值定理、Taylor 公式、不等式与凸函数以及导数的综合应用。学习的重点是从题目条件中辨认可以使用的定理，并把目标式转化为适当的函数值、导数或余项；计算中的技巧必须建立在可导性、区间条件和极限存在性已经核实的基础上。

\section{导数}
\label{sec:pei-diff-derivative}
% 教材对应：PDF 第 1--17 页。

\subsection{导数定义与可微性的判定}

\begin{definition}
设实函数 $f$ 在 $x_0$ 的某邻域内有定义。若有限极限
\begin{equation}
  f'(x_0)=\lim_{h\to0}\frac{f(x_0+h)-f(x_0)}{h}
  \label{eq:pei-diff-derivative-definition}
\end{equation}
存在，则称 $f$ 在 $x_0$ 处\emph{可导（differentiable）}。若存在常数 $A$，使
\begin{equation}
  f(x_0+h)=f(x_0)+Ah+o(h),\qquad h\to0,
  \label{eq:pei-diff-differentiability}
\end{equation}
则称 $f$ 在 $x_0$ 处\emph{可微}，并记 $\mathrm{d}f=A\,\mathrm{d}x$。一元实函数在一点可导与可微等价，且 $A=f'(x_0)$；这里 $r(h)=o(h)$ 表示 $r(h)/h\to0$。
\end{definition}

上述等价性直接来自差商极限：若差商趋于 $A$，则余项为 $h$ 乘以趋于零的量；反之，将式~\eqref{eq:pei-diff-differentiability} 除以 $h\ne0$ 即得导数。它也说明可导必连续，因为 $Ah+o(h)\to0$。在分段点，应先检验函数值与两侧极限是否一致，再分别计算左右差商；左右导数都存在且相等，才有双侧导数。只计算两段表达式的导数并令其极限相等，不能替代对分段点本身的检验。

\begin{proposition}
\label{prop:pei-diff-caratheodory}
函数 $f$ 在 $x_0$ 处可导，当且仅当存在一个在 $x_0$ 连续的函数 $g$，使其在 $x_0$ 的某邻域内满足
\[
  f(x)-f(x_0)=(x-x_0)g(x).
\]
此时 $g(x_0)=f'(x_0)$。
\end{proposition}
\begin{proof}
必要性只需对 $x\ne x_0$ 取 $g(x)=[f(x)-f(x_0)]/(x-x_0)$，并定义 $g(x_0)=f'(x_0)$。充分性则由差商等于 $g(x)$ 及其在 $x_0$ 的连续性得到。这一形式把可导性转化为连续因子的存在，尤其适合处理含 $(x-x_0)$ 的乘积。
\end{proof}

\begin{example}
\label{ex:pei-diff-oscillation-derivative}
对正整数 $m$，令
\[
 f_m(x)=\begin{cases}x^m\sin(1/x),&x\ne0,\\0,&x=0.\end{cases}
\]
讨论它在原点的连续性、可导性及导函数的连续性。
\end{example}
\begin{solve}
由 $|f_m(x)|\le |x|^m$，所有这些函数在原点都连续。原点处的差商为 $x^{m-1}\sin(1/x)$，故 $m\ge2$ 时可导且 $f_m'(0)=0$；$m=1$ 时差商振荡，导数不存在。对 $x\ne0$，
\[
 f_m'(x)=mx^{m-1}\sin(1/x)-x^{m-2}\cos(1/x).
\]
所以 $m\ge3$ 时 $f_m'(x)\to0=f_m'(0)$，而 $m=2$ 时导函数不连续。由此可见，\emph[caution]{可导不蕴含导函数连续}；在原点附近可用求导公式，并不意味着该公式可以不经检验地代入原点。
\end{solve}


\begin{example}
\label{ex:pei-diff-riemann-nondifferentiable}
令 $R$ 为 $[0,1]$ 上的 Riemann 函数，也称 Thomae 函数：若 $x=p/q\in(0,1)$ 为既约分数，取 $R(x)=1/q$；若 $x$ 为无理数或 $x=0,1$，取 $R(x)=0$。证明它在每个内点都不可导，且两个端点的相应单侧导数也不存在。
\end{example}
\begin{solve}
在内部有理点 $p/q$，沿无理点趋近时函数值恒为零，与 $R(p/q)>0$ 不符，故函数不连续，从而不可导。设 $x_0\in(0,1)$ 为无理数，并取十进制截断 $r_n=\lfloor10^n x_0\rfloor/10^n$。当 $n$ 充分大时 $r_n\in(0,1)$，其既约分母至多为 $10^n$，且 $0<x_0-r_n<10^{-n}$，因而
\[
 \left|\frac{R(r_n)-R(x_0)}{r_n-x_0}\right|\ge1.
\]
另一方面，沿无理点趋近 $x_0$ 的差商恒为零，因此双侧导数不存在。在 $0$ 处沿 $1/q\to0$ 的右差商为 $1$，沿无理点的右差商为零；在 $1$ 处沿 $(q-1)/q\to1$ 的左差商为 $-1$，沿无理点则为零。两端的单侧导数也不存在。
\end{solve}

\begin{notealert}[type=caution,title={对称差商不能替代导数}]
可导必推出 $[f(x_0+h)-f(x_0-h)]/(2h)\to f'(x_0)$，逆命题不成立。例如 $f(x)=|x|$ 在原点的对称差商恒为零，但左右导数分别为 $-1$ 与 $1$。证明不可导时，只需找到两列趋近方式给出不同的差商极限，或证明连续性等必要条件不成立。
\end{notealert}

\subsection{由差分条件恢复导数}

\begin{proposition}
\label{prop:pei-diff-dyadic-derivative}
设 $f$ 在原点连续，且
\[
 \lim_{x\to0}\frac{f(2x)-f(x)}{x}=A\in\mathbb R.
\]
则 $f'(0)$ 存在且等于 $A$。
\end{proposition}
\begin{proof}
任给 $\varepsilon>0$，取 $\delta>0$，使 $0<|t|<\delta$ 时有
$|f(2t)-f(t)-At|\le\varepsilon|t|$。固定 $0<|x|<\delta$，依次取 $t=x/2^k$，将 $k=1,\ldots,n$ 对应的不等式相加，得
\[
 \left|f(x)-f(x/2^n)-Ax(1-2^{-n})\right|
 \le\varepsilon|x|(1-2^{-n}).
\]
先让 $n\to\infty$，利用原点连续性得到
$|f(x)-f(0)-Ax|\le\varepsilon|x|$，再按导数定义得结论。这里对固定的 $x$ 先作有限次伸缩与求和，避免了预先假定待证差商有极限。
\end{proof}

\begin{remark}
原点连续性不能删除。例如令 $f(0)=0$，而 $x\ne0$ 时 $f(x)=1$，则 $f(2x)-f(x)=0$，但原点不可导。即使已经知道 $f(2x)-f(x)=Ax+o(x)$，也不能直接设 $[f(x)-f(0)]/x$ 有极限后解方程；只有独立证明极限存在，这种运算才成立。
\end{remark}

\begin{proposition}
\label{prop:pei-diff-uniform-differentiability}
设 $f$ 在 $[a,b]$ 上可导，端点按向内的单侧导数理解。则 $f'\in C([a,b])$ 当且仅当 $f$ 在该区间上\emph{一致可微（uniformly differentiable）}，即对每个 $\varepsilon>0$，存在 $\delta>0$，使对所有满足 $x,x+h\in[a,b]$、$0<|h|<\delta$ 的 $x,h$，都有
\[
 \left|\frac{f(x+h)-f(x)}h-f'(x)\right|<\varepsilon.
\]
\end{proposition}
\begin{proof}
若 $f'$ 连续，则它在紧区间上一致连续；对 $[x,x+h]$ 应用 Lagrange 中值定理，差商等于两端之间某点的导数，故上述估计对所有 $x$ 一致成立。反之，固定 $x_0$ 和 $\varepsilon$，选一个足够小且使 $x_0+h\in(a,b)$ 的非零 $h$；在 $x_0$ 的相对邻域内，$x+h$ 仍在区间中。一致可微性使 $f'(x)$ 与连续函数 $[f(x+h)-f(x)]/h$ 的差一致小于 $\varepsilon$。再令 $x\to x_0$，由两次这样的近似得 $\limsup_{x\to x_0}|f'(x)-f'(x_0)|\le2\varepsilon$，遂得连续性。
\end{proof}

\subsection{高阶导数与求导方法的选择}

\emph{高阶导数（higher derivatives）}递归定义为 $f^{(n)}=(f^{(n-1)})'$，其中 $f^{(0)}=f$；声称某点有 $n$ 阶导数时，前一阶导函数须在该点附近有定义。求高阶导数通常先整理结构：三角乘积用积化和差降为基本频率的和，真分式尽量分解为简单分式，复合一次函数则保留链式法则产生的幂因子。例如
\[
 \frac{\mathrm d^n}{\mathrm dx^n}\sin(ax+b)
 =a^n\sin\left(ax+b+\frac{n\pi}{2}\right),\qquad
 \frac{\mathrm d^n}{\mathrm dx^n}\frac1{x-c}
 =\frac{(-1)^n n!}{(x-c)^{n+1}}.
\]
这些公式分别在全部实数与 $x\ne c$ 的范围内使用；先拆项可以避免把简单结构掩盖在多次商法则之中。

\begin{theorem}[Leibniz 公式]
\label{thm:pei-diff-leibniz}
若 $u,v$ 在某区间内具有直到 $n$ 阶导数，则
\begin{equation}
 (uv)^{(n)}=\sum_{k=0}^n\binom nk u^{(k)}v^{(n-k)}.
 \label{eq:pei-diff-leibniz}
\end{equation}
更一般地，$m$ 个具有这些导数的函数满足
\[
 (u_1\cdots u_m)^{(n)}
 =\sum_{k_1+\cdots+k_m=n}
 \frac{n!}{k_1!\cdots k_m!}\prod_{j=1}^m u_j^{(k_j)},
 \qquad k_j\ge0.
\]
\end{theorem}
\begin{proof}
$n=1$ 时为乘积法则。对 $n$ 阶公式逐项求导，并把相同的 $u^{(k)}v^{(n+1-k)}$ 合并，系数由
$\binom nk+\binom n{k-1}=\binom{n+1}k$ 给出，完成归纳。多因子形式再按因子个数归纳，系数化简为多项式系数即可。
\end{proof}

直接反复求导时，应把前几阶的结构猜测写成命题，再用归纳法核实；若函数满足微分恒等式，则可对恒等式求导，将高阶导数变成递推问题。后一种方法在只求某点的导数时尤其有效，不必求出一般的 $f^{(n)}(x)$。与之不同，通过 Taylor 系数读取导数时，应先证明相应展开具有足够阶数的余项，不能只凭一个形式级数作结论。

\begin{example}
\label{ex:pei-diff-arcsin-recurrence}
求 $y(x)=(\arcsin x)^2$ 在原点的各阶导数。
\end{example}
\begin{solve}
在 $(-1,1)$ 内，直接求导得到
\[
 (1-x^2)y''-xy'=2.
\]
令 $a_n=y^{(n)}(0)$。由 $a_0=a_1=0$ 得 $a_2=2$；对恒等式求 $n$ 阶导数并代入原点，由 Leibniz 公式可得
\[
 a_{n+2}=n^2a_n,\qquad n\ge1.
\]
于是所有奇数阶导数为零，而
\[
 y^{(2m)}(0)=2\bigl((2m-2)!!\bigr)^2,\qquad m\ge1,
\]
其中 $0!!=1$。递推的起始值与奇偶两条链缺一不可。
\end{solve}

\begin{example}
\label{ex:pei-diff-flat-function}
令 $\phi(x)=e^{-1/x^2}$（$x\ne0$），$\phi(0)=0$。证明它无穷次可导，且每阶导数在原点都等于零。
\end{example}
\begin{solve}
对任意 $N\ge0$，令 $t=1/x^2$，由 $t^{N/2}e^{-t}\to0$ 得
$|x|^{-N}e^{-1/x^2}\to0$。在 $x\ne0$ 处反复求导，归纳得到 $\phi^{(n)}(x)=P_n(1/x)e^{-1/x^2}$，其中 $P_n$ 是多项式。若已知 $\phi^{(n)}(0)=0$，则 $\phi^{(n)}(h)/h\to0$，由定义得 $\phi^{(n+1)}(0)=0$；同一估计也说明该阶导函数在原点连续。逐阶归纳即得结论。这个例子说明，无穷次可导并不足以保证函数等于自己的 Taylor 级数。
\end{solve}

对不同方式给出的函数，应先确定求导变量及局部可解性。参数曲线 $x=x(t),y=y(t)$ 在 $x'(t)\ne0$ 处满足 $\mathrm dy/\mathrm dx=y'(t)/x'(t)$，再次求导须使用 $\mathrm d/\mathrm dx=[1/x'(t)]\mathrm d/\mathrm dt$。若 $F$ 在相关点的邻域内连续可微，且已知可微函数 $y=y(x)$ 满足 $F(x,y(x))=0$、$F_y\ne0$，则链式法则给出 $y'=-F_x/F_y$；函数 $y$ 的局部存在性应由题设或隐函数定理另行保证。反函数的求导同样要求局部可逆，并在相应点满足原函数导数非零。

\section{微分中值定理}
\label{sec:pei-diff-mean-value}
% 教材对应：PDF 第 17--50 页。

\subsection{Rolle 定理、零点与辅助函数}

\begin{theorem}[Fermat 定理与 Rolle 定理]
若 $f$ 在内点 $c$ 取得局部极值且在 $c$ 可导，则 $f'(c)=0$。若 $f\in C([a,b])$，在 $(a,b)$ 内可导，且 $f(a)=f(b)$，则存在 $c\in(a,b)$ 使 $f'(c)=0$。
\end{theorem}
\begin{proof}
以内点极大值为例，充分小的正、负 $h$ 对应的差商分别非正、非负；若它们有共同极限，该极限只能为零。对 Rolle 定理，常函数的结论直接成立；非常函数在紧区间上取得最大值与最小值，由两端值相等，至少有一个非平凡极值在内点取得，遂应用前一结论。
\end{proof}

Rolle 定理把两个函数零点转化为其间的一个导数零点。反复使用可知：若 $f$ 在区间内 $n$ 次可导且有 $n+1$ 个不同零点，则 $f^{(n)}$ 至少有一个零点。需要证明零点唯一时，则可证明函数严格单调，或反设有两个零点再推出导数出现不可能的零点。\emph[tip]{目标中若出现 $f'-\lambda f$，可先乘以积分因子}，因为
\[
 (e^{-\lambda x}f(x))'=e^{-\lambda x}(f'(x)-\lambda f(x)).
\]
因此，$f$ 的两个零点之间至少有一个 $f'-\lambda f$ 的零点；对若干实数 $\lambda_j$ 依次使用，就能处理 $(D-\lambda_1)\cdots(D-\lambda_n)f$，其中 $D$ 表示求导算子。

\begin{example}
\label{ex:pei-diff-real-roots}
设实系数 $n$ 次多项式 $P$ 的全部根均为实数。证明 $P',\ldots,P^{(n-1)}$ 的全部根也都是实数。
\end{example}
\begin{solve}
设 $P$ 的不同实根为 $r_1<\cdots<r_m$，重数分别为 $k_1,\ldots,k_m$，且 $\sum k_j=n$。在 $r_j$ 处，$P'$ 至少保留 $k_j-1$ 重根；相邻不同根之间，Rolle 定理另给出一个实根。合计实根重数至少为
\[
 \sum_{j=1}^m(k_j-1)+(m-1)=n-1.
\]
而 $P'$ 的次数正是 $n-1$，故其根全部为实数。对各阶导数重复论证即可。计数时必须同时考虑原有重根及根间新增的根。
\end{solve}

\begin{example}
\label{ex:pei-diff-legendre-roots}
对 $n\ge1$，证明 Legendre 多项式
\[
 P_n(x)=\frac1{2^n n!}\frac{\mathrm d^n}{\mathrm dx^n}(x^2-1)^n
\]
在 $(-1,1)$ 内有 $n$ 个不同的实根。
\end{example}
\begin{solve}
令 $Q(x)=(x^2-1)^n$。它在 $-1,1$ 各有 $n$ 重零点。对 $k=0,\ldots,n$ 归纳：$Q^{(k)}$ 在 $(-1,1)$ 至少有 $k$ 个不同零点，而 $k<n$ 时两个端点仍为零点。将这些零点依次排列，在相邻零点间使用 Rolle 定理，得到下一阶导数的至少 $k+1$ 个内部零点。故 $Q^{(n)}$ 有至少 $n$ 个内部零点；它本身只有 $n$ 次，所以恰有这些根且均为单根。
\end{solve}

需要证明一个复杂的中值等式时，可先把目标写成 $F'(c)=0$，再检查能否让 $F(a)=F(b)$。例如 $f,g,h$ 在 $[a,b]$ 连续、在内部可导时，取
\[
 F(x)=\begin{vmatrix}
 f(a)&g(a)&h(a)\\f(b)&g(b)&h(b)\\f(x)&g(x)&h(x)
 \end{vmatrix}.
\]
两端均有重复行，故 $F(a)=F(b)=0$，于是存在 $c\in(a,b)$，使把第三行换成 $(f'(c),g'(c),h'(c))$ 后的行列式为零。这是把端点条件直接编码到辅助函数中的一种方法。

\subsection{Lagrange 定理与增量估计}

\begin{theorem}[Lagrange 中值定理]
\label{thm:pei-diff-lagrange}
若 $f\in C([a,b])$ 且在 $(a,b)$ 内可导，则存在 $c\in(a,b)$ 使
\begin{equation}
 f(b)-f(a)=f'(c)(b-a).
 \label{eq:pei-diff-lagrange}
\end{equation}
\end{theorem}
\begin{proof}
对 $F(x)=f(x)-f(a)-[f(b)-f(a)](x-a)/(b-a)$ 应用 Rolle 定理即可。几何上，割线的斜率等于区间内某条切线的斜率，但定理一般不确定 $c$ 的具体位置。
\end{proof}

若 $|f'|\le M$，式~\eqref{eq:pei-diff-lagrange} 给出 $|f(y)-f(x)|\le M|y-x|$，即 \emph{Lipschitz 条件（Lipschitz condition）}，从而一致连续。若 $f'\equiv0$，则 $f$ 为常数；若 $f'\ge0$，则 $f$ 单调不减。严格递增不要求每一点导数都为正，例如 $x^3$ 在原点的导数为零；而 $f'>0$ 是严格递增的充分条件。比较两条割线时，可以在相邻区间分别使用中值定理，再用两个中值点的位置关系比较导数，这比试图求出中值点更有效。

\begin{example}
\label{ex:pei-diff-increasing-derivative-chord}
若 $f$ 在 $[a,b]$ 连续、在内部可导，$f'$ 严格递增且 $f(a)=f(b)$，证明 $a<x<b$ 时 $f(x)<f(a)$。
\end{example}
\begin{solve}
分别对 $[a,x]$ 与 $[x,b]$ 使用中值定理，得到 $u\in(a,x)$、$v\in(x,b)$，使
\[
 \frac{f(x)-f(a)}{x-a}=f'(u)<f'(v)
 =\frac{f(b)-f(x)}{b-x}.
\]
令 $f(a)=f(b)=A$，移项得 $(b-a)(f(x)-A)<0$，即得结论。这也预示了导数单调性与凸性的联系。
\end{solve}

\begin{example}
\label{ex:pei-diff-sublinear}
设 $f$ 在某半直线上可导。证明：若 $f'(x)\to0$（$x\to+\infty$），则 $f(x)/x\to0$；反过来，若 $f(x)/x\to0$，则至少有一列 $c_n\to+\infty$，使 $f'(c_n)\to0$。
\end{example}
\begin{solve}
任给 $\varepsilon>0$，取 $A$ 使 $t>A$ 时 $|f'(t)|<\varepsilon$。对 $x>A$，中值定理给出
$|f(x)|\le|f(A)|+\varepsilon(x-A)$；除以 $x$ 再令 $x\to\infty$，得所需极限。反向命题中，对充分大的正整数 $n$，在 $[n,2n]$ 上取中值点 $c_n$，则
\[
 f'(c_n)=\frac{f(2n)-f(n)}n
 =2\frac{f(2n)}{2n}-\frac{f(n)}n\longrightarrow0.
\]
反向只能保证这样一列点，不能推出整个导函数趋于零，例如有界函数 $f(x)=\sin x$ 的导函数 $\cos x$ 没有极限。
\end{solve}

\begin{example}
\label{ex:pei-diff-differential-uniqueness}
设 $f$ 在 $[0,+\infty)$ 上连续、在 $(0,+\infty)$ 内可导，$f(0)=0$，并存在常数 $A\ge0$ 使 $|f'(x)|\le A|f(x)|$ 对所有 $x>0$ 成立。证明 $f\equiv0$。
\end{example}
\begin{solve}
$A=0$ 时由中值定理即可。设 $A>0$，取 $\ell>0$ 使 $A\ell<1$。在 $[0,\ell]$ 上令 $M=\max|f|$，对每个 $x$ 使用中值定理，有 $|f(x)|\le AxM\le A\ell M$，于是 $M\le A\ell M$，只能有 $M=0$。已知 $f(\ell)=0$ 后，在 $[\ell,2\ell]$ 重复同一论证，逐段推进即得结论。此处的有限区间最大值来自连续性，推进步长由 $A$ 统一控制。
\end{solve}

\subsection{导数的介值性与间断点}

\begin{theorem}[Darboux 定理]
\label{thm:pei-diff-darboux}
设 $f$ 在一个区间内处处可导。若 $u<v$ 且 $\lambda$ 严格介于 $f'(u)$ 与 $f'(v)$ 之间，则存在 $c\in(u,v)$ 使 $f'(c)=\lambda$。这一性质称为导数的\emph{介值性（intermediate value property）}。
\end{theorem}
\begin{proof}
先设 $f'(u)<\lambda<f'(v)$，令 $F(x)=f(x)-\lambda x$。由 $F'(u)<0$ 可知，$u$ 右侧有点使函数值小于 $F(u)$；由 $F'(v)>0$ 可知，$v$ 左侧也有点使函数值小于 $F(v)$。连续函数 $F$ 在 $[u,v]$ 达到最小值，而两端均不能为最小值点，故在内点 $c$ 有 $F'(c)=0$。另一种次序对 $-f$ 使用同一论证即可。证明没有假定 $f'$ 连续。
\end{proof}

\begin{proposition}
\label{prop:pei-diff-derivative-no-jump}
若 $f$ 在 $x_0$ 的某邻域内处处可导，且有限单侧极限 $\lim_{x\to x_0+}f'(x)=L$ 存在，则 $L=f'(x_0)$；左侧同理。因此导函数没有可去间断点或跳跃间断点，它的间断点只能是第二类间断点。
\end{proposition}
\begin{proof}
对 $h>0$，中值定理给出
$[f(x_0+h)-f(x_0)]/h=f'(c_h)$，其中 $x_0<c_h<x_0+h$。当 $h\to0+$ 时，左端趋于 $f'(x_0)$，右端趋于 $L$。若左右有限极限都存在，它们都必须等于函数值，故导函数在该点连续。
\end{proof}

\begin{example}
\label{ex:pei-diff-bounded-second-zero}
设 $f$ 在 $\mathbb R$ 上二次可导且有界，证明存在 $c\in\mathbb R$ 使 $f''(c)=0$。
\end{example}
\begin{solve}
反设 $f''$ 处处不为零。因它也是导函数，由 Darboux 定理，$f''$ 只能恒正或恒负。设恒正，则 $f'$ 严格递增。若某点 $x_0$ 有 $f'(x_0)>0$，则在其右侧 $f'(x)\ge f'(x_0)>0$，由中值定理得 $f$ 向右无界；若 $f'(x_0)<0$，向左作相同估计也使 $f$ 无界。若每点的 $f'$ 都为零，又与严格递增矛盾。恒负的情形对 $-f$ 处理。故反设不成立。
\end{solve}

\subsection{Cauchy 中值定理与 L'Hospital 法则}

\begin{theorem}[Cauchy 中值定理]
\label{thm:pei-diff-cauchy}
设 $F,G\in C([a,b])$，在 $(a,b)$ 内可导，则存在 $c\in(a,b)$，使
\[
 [F(b)-F(a)]G'(c)=[G(b)-G(a)]F'(c).
\]
若进一步有 $G'(x)\ne0$ 对所有 $x\in(a,b)$ 成立，则 $G(b)\ne G(a)$，且
\begin{equation}
 \frac{F(b)-F(a)}{G(b)-G(a)}=\frac{F'(c)}{G'(c)}.
 \label{eq:pei-diff-cauchy}
\end{equation}
\end{theorem}
\begin{proof}
对 $[F(b)-F(a)]G(x)-[G(b)-G(a)]F(x)$ 使用 Rolle 定理可得第一式。如果 $G(b)=G(a)$，Rolle 定理会迫使某处 $G'=0$，故在附加条件下可除以相应的两个量，得到比值形式。分别对 $F,G$ 使用 Lagrange 定理所产生的中值点一般不同，不能因此直接得到同一点 $c$ 的比值公式。
\end{proof}

\begin{theorem}[L'Hospital 法则的常用形式]
\label{thm:pei-diff-lhospital}
设 $f,g$ 在趋近点的某个单侧去心邻域内可导，$g'$ 处处非零，且 $f'/g'\to L$，其中 $L$ 可以是有限数或正、负无穷。若满足 $f\to0$、$g\to0$，或满足 $|g|\to\infty$，则 $f/g\to L$。趋近点也可为正、负无穷；在第二种情形中不必另假定 $f$ 本身趋于无穷。
\end{theorem}
\begin{proof}
先说明有限趋近点、$0/0$ 型的情形。把端点函数值补为零，对端点与邻近点之间应用式~\eqref{eq:pei-diff-cauchy}，所得中值点趋于该端点，故比值趋于 $L$。对 $|g|\to\infty$ 型，固定一个足够靠近趋近端的点 $t$，有
\[
 \frac{f(x)}{g(x)}
 =\frac{f(x)-f(t)}{g(x)-g(t)}
   \left(1-\frac{g(t)}{g(x)}\right)+\frac{f(t)}{g(x)}.
\]
若 $L$ 有限，先选 $t$ 使其与趋近端之间的全部导数比都与 $L$ 相差小于给定的 $\varepsilon$；Cauchy 定理使上式中的差分比也满足此估计。再令 $x$ 趋近该端，另两项中的固定点比值趋于零，便得结论。若 $L$ 为无穷，则以任意给定实数为界，对导数比作相应的单侧估计，沿用同一式子即可；无穷远端点的 $0/0$ 型可先对 $x,t$ 使用 Cauchy 定理，再令 $t$ 趋向无穷远端点。
\end{proof}

\begin{notealert}[type=caution,title={使用法则前先辨认条件}]
L'Hospital 法则的逻辑是从导数比的极限推出原比值的极限，导数比没有极限时不能反推原比值没有极限。例如 $(x+\sin x)/x\to1$，而导数比 $1+\cos x$ 无极限。遇到 $0\cdot\infty$、$\infty-\infty$ 或幂型未定式，应先作代数变形或取对数，转成满足条件的比值；每次重复使用，都要重新核实可导性、分母导数非零及未定式类型。
\end{notealert}

\begin{example}
\label{ex:pei-diff-f-plus-derivative}
设 $f$ 在某半直线上可导，且 $f(x)+f'(x)\to A\in\mathbb R$。证明 $f(x)\to A$ 且 $f'(x)\to0$。
\end{example}
\begin{solve}
将 $f(x)$ 写成 $[e^x f(x)]/e^x$。分母趋于无穷，其导数非零，而导数比为
\[
 \frac{(e^xf(x))'}{(e^x)'}=f(x)+f'(x)\longrightarrow A.
\]
由定理~\ref{thm:pei-diff-lhospital} 的无穷分母形式得 $f(x)\to A$，再由 $f'=(f+f')-f$ 得 $f'(x)\to0$。使用此法不需要预先知道分子 $e^xf(x)$ 的极限。
\end{solve}

\subsection{对称导数与连续性条件}

教材附录讨论了导数的推广。对内点 $x$，若极限存在，可分别定义\emph{对称导数（symmetric derivative）}与\emph{二阶对称差商极限}
\[
 D_s f(x)=\lim_{h\to0}\frac{f(x+h)-f(x-h)}{2h},\qquad
 D_s^2 f(x)=\lim_{h\to0}\frac{f(x+h)-2f(x)+f(x-h)}{h^2}.
\]
这里的 $D_s^2$ 表示二阶对称差商的极限，不是对 $D_s f$ 再求一次对称导数。普通一、二阶导数存在时，相应的对称差商极限分别等于 $f'$、$f''$，二阶情形可由定理~\ref{thm:pei-diff-taylor-peano} 的展开式看出；逆命题一般不成立。特别地，一阶对称差商不含中心点的函数值，所以它存在甚至不保证函数在该点连续。连续性等正则条件在推广结论中须单独核实。

\begin{proposition}[Schwarz 型结论]
\label{prop:pei-diff-symmetric-affine}
设 $f\in C([a,b])$，且对所有 $x\in(a,b)$，$D_s^2 f(x)$ 都存在并等于零。则 $f$ 在 $[a,b]$ 上为仿射函数，即 $f(x)=\alpha x+\beta$。
\end{proposition}
\begin{proof}
设 $L$ 为连接两个端点的仿射函数，令 $K=f-L$，于是 $K(a)=K(b)=0$ 且 $D_s^2K=0$。任给 $\varepsilon>0$，函数
$K_\varepsilon(x)=K(x)+\varepsilon(x-a)(x-b)$ 的二阶对称差商极限为 $2\varepsilon>0$。若它有正值，由连续性及端点值为零，必在内点取得正的最大值；但最大值点的二阶对称差商均非正，矛盾。因此 $K_\varepsilon\le0$。对 $K(x)-\varepsilon(x-a)(x-b)$ 同理可知它非负。令 $\varepsilon\to0$，得到 $K\equiv0$。
\end{proof}

\section{Taylor 公式}
\label{sec:pei-diff-taylor}
% 教材对应：PDF 第 50--68 页。

\subsection{两种余项及其使用范围}

\begin{theorem}[带 Lagrange 余项的 Taylor 公式]
\label{thm:pei-diff-taylor-lagrange}
设 $n\ge0$，$f$ 在连接 $x_0,x$ 的闭区间上有连续的直到 $n$ 阶导数，并在内部有 $n+1$ 阶导数。则存在严格位于 $x_0,x$ 之间的 $\xi$，使
\begin{equation}
 f(x)=\sum_{k=0}^n\frac{f^{(k)}(x_0)}{k!}(x-x_0)^k
       +\frac{f^{(n+1)}(\xi)}{(n+1)!}(x-x_0)^{n+1}.
 \label{eq:pei-diff-taylor-lagrange}
\end{equation}
$n=0$ 时连续性条件只要求 $f$ 本身连续。
\end{theorem}
\begin{proof}
设 $x\ne x_0$，记前面的 $n$ 次多项式为 $T_n(t)$，并取常数 $K$ 使
$F(t)=f(t)-T_n(t)-K(t-x_0)^{n+1}$ 在 $t=x$ 时为零。它在 $x_0$ 的直到 $n$ 阶导数都为零；反复使用 Rolle 定理，可得内点 $\xi$ 使 $F^{(n+1)}(\xi)=0$。所以 $K=f^{(n+1)}(\xi)/(n+1)!$，代回即得公式。$x=x_0$ 时等式直接成立，无须选取中值点。
\end{proof}

\begin{theorem}[带 Peano 余项的 Taylor 公式]
\label{thm:pei-diff-taylor-peano}
设 $n\ge1$，$f$ 在 $x_0$ 的某邻域内有直到 $n-1$ 阶导数，且 $f^{(n-1)}$ 在 $x_0$ 可导。则
\begin{equation}
 f(x)=\sum_{k=0}^n\frac{f^{(k)}(x_0)}{k!}(x-x_0)^k
       +o\bigl((x-x_0)^n\bigr),\qquad x\to x_0.
 \label{eq:pei-diff-taylor-peano}
\end{equation}
\end{theorem}
\begin{proof}
记差值为 $r(x)=f(x)-T_n(x)$。当 $n=1$ 时，结论就是可导的定义。当 $n\ge2$ 时，连续使用 $0/0$ 型 L'Hospital 法则到 $n-1$ 次，把 $r(x)/(x-x_0)^n$ 的极限归结为
\[
 \lim_{x\to x_0}
 \frac{f^{(n-1)}(x)-f^{(n-1)}(x_0)-f^{(n)}(x_0)(x-x_0)}{n!(x-x_0)}=0.
\]
最后一步用 $f^{(n-1)}$ 在 $x_0$ 可导，不需要假定 $f^{(n)}$ 在去心邻域存在或连续。
\end{proof}

Lagrange 余项在一段区间上用某个中值点的高阶导数控制误差，适合证明不等式、构造中值公式与估计函数或导数的界；Peano 余项只描述展开点附近的阶数，适合局部极限、极值判别及系数比较。\emph[caution]{不能对一般的 $o(x^n)$ 余项任意求导}：例如 $r(x)=x^3\sin(1/x^2)$ 在原点补为零后满足 $r(x)=o(x^2)$，但对 $x\ne0$，$r'(x)=3x^2\sin(1/x^2)-2\cos(1/x^2)$，并不趋于零。求导后的余项必须另有依据。

\subsection{中值公式、不等式与导数估计}

Taylor 公式的展开点应由条件决定：已知某点的若干阶导数为零，就围绕该点展开；含两个端点的函数值，可以向同一内点展开并消去一阶项；含对称的自变量，则分别在 $x_0\pm h$ 展开以消去奇次项或偶次项。若已知 $f^{(n+1)}$ 的符号，式~\eqref{eq:pei-diff-taylor-lagrange} 还给出了函数与 Taylor 多项式的大小关系，其中必须连同 $(x-x_0)^{n+1}$ 的符号一起判断。

\begin{example}
\label{ex:pei-diff-second-difference}
设 $f\in C^2([a,b])$，$m=(a+b)/2$。证明存在 $c\in(a,b)$，使
\[
 f(a)-2f(m)+f(b)=\frac{(b-a)^2}{4}f''(c).
\]
\end{example}
\begin{solve}
将两个端点在 $m$ 处展开到一次项，得
\[
 f(a)-2f(m)+f(b)
 =\frac{(b-a)^2}{8}\bigl(f''(\xi)+f''(\eta)\bigr),
 \quad \xi\in(a,m),\quad\eta\in(m,b).
\]
由 $f''$ 的连续性，它在 $\xi,\eta$ 之间取得两值的平均数，即得结论。两个展开式的中值点一般不同；将它们合并成同一个点，需要介值性这样的额外依据。
\end{solve}

\begin{example}
\label{ex:pei-diff-curvature-bound}
设 $f\in C^2([0,1])$，$f(0)=f(1)=0$，且 $\min_{[0,1]}f=-1$。证明存在 $c\in(0,1)$ 使 $f''(c)\ge8$。
\end{example}
\begin{solve}
取内点 $x_0$ 使 $f(x_0)=-1$，则 $f'(x_0)=0$。向两个端点展开，分别得到
\[
 1=\tfrac12f''(\xi)x_0^2,\qquad
 1=\tfrac12f''(\eta)(1-x_0)^2.
\]
因为 $\min\{x_0,1-x_0\}\le1/2$，相应一个中值点的二阶导数至少为 $8$。常数 $8$ 可由 $f(x)=4x(x-1)$ 达到，因此不能改大。
\end{solve}

\begin{example}
\label{ex:pei-diff-finite-interval-bound}
设 $f\in C^2([0,1])$，$|f(x)|\le1$，$|f''(x)|\le2$。证明 $|f'(x)|\le3$。
\end{example}
\begin{solve}
固定 $x\in(0,1)$，把 $f(0),f(1)$ 都在 $x$ 处展开。相减后有
\[
 f(1)-f(0)=f'(x)+\tfrac12f''(\xi)(1-x)^2
                         -\tfrac12f''(\eta)x^2.
\]
故 $|f'(x)|\le2+(1-x)^2+x^2\le3$。端点处由 $f'$ 连续取极限得到同一估计。这里把两端向同一点展开，正是为了得到系数为 $1$ 的 $f'(x)$。
\end{solve}

\begin{proposition}[全实轴上的导数界]
\label{prop:pei-diff-landau-bound}
设 $f$ 在 $\mathbb R$ 上二次可导，且
$M_0=\sup_{\mathbb R}|f|<\infty$、$M_2=\sup_{\mathbb R}|f''|<\infty$。则 $M_1=\sup_{\mathbb R}|f'|<\infty$，并且
\begin{equation}
 M_1^2\le2M_0M_2.
 \label{eq:pei-diff-landau-bound}
\end{equation}
\end{proposition}
\begin{proof}
对任意 $x\in\mathbb R$、$h>0$，向 $x+h$ 与 $x-h$ 展开再相减，得
\[
 2h|f'(x)|\le2M_0+M_2h^2,
 \qquad |f'(x)|\le\frac{M_0}{h}+\frac{M_2h}{2}.
\]
若 $M_0M_2>0$，取 $h=\sqrt{2M_0/M_2}$，右端为 $\sqrt{2M_0M_2}$，对 $x$ 取上确界即得结论。若 $M_0=0$，则 $f\equiv0$；若 $M_2=0$，则 $f$ 为仿射函数，再由全实轴上有界知它为常数。任意选择正步长的自由依赖于定义域为整个实轴，不能直接在有限区间套用同一常数。
\end{proof}

\subsection{无穷远处的极限与中间阶导数}

\begin{proposition}
\label{prop:pei-diff-limit-bounded-second}
设 $f$ 在某半直线上二次可导，$\lim_{x\to+\infty}f(x)$ 存在且有限，且 $|f''|\le M$。则 $f'(x)\to0$。
\end{proposition}
\begin{proof}
固定 $h>0$，由 Taylor 公式，
\[
 |f'(x)|\le\frac{|f(x+h)-f(x)|}{h}+\frac{Mh}{2}.
\]
先保持 $h$ 不变令 $x\to+\infty$，利用函数极限有限，得 $\limsup|f'(x)|\le Mh/2$。这一估计对任意 $h>0$ 都成立，再令 $h\to0+$，即得结论。这两个极限有先后顺序，并没有交换极限运算。
\end{proof}

\begin{proposition}
\label{prop:pei-diff-intermediate-derivative-limit}
设 $n\ge2$，$f$ 在某半直线上有 $n$ 阶导数，且 $f(x)\to A\in\mathbb R$、$f^{(n)}(x)\to0$。则 $f^{(j)}(x)\to0$ 对 $j=1,\ldots,n-1$ 都成立。
\end{proposition}
\begin{proof}
对 $m=1,\ldots,n-1$，在 $x$ 处将 $f(x+m)$ 展开到 $n-1$ 阶。因每个余项中的中值点位于 $[x,x+n-1]$，且 $f^{(n)}$ 趋于零，得到
\[
 f(x+m)-f(x)=\sum_{j=1}^{n-1}\frac{m^j}{j!}f^{(j)}(x)+o(1).
\]
左端都趋于零。这个关于 $f'(x),\ldots,f^{(n-1)}(x)$ 的线性方程组有常数系数矩阵 $(m^j/j!)$；从各行提出 $m$、各列提出 $1/j!$ 后，剩下的是节点 $1,\ldots,n-1$ 的 Vandermonde 矩阵，行列式非零。因此各个未知导数都是有限个趋零量的固定线性组合，结论成立。
\end{proof}

仅有 $f(x)$ 的极限有限，并不能断言 $f'(x)\to0$。例如 $f(x)=\sin(x^2)/x$ 在 $x\ge1$ 上趋于零，但 $f'(x)=2\cos(x^2)-\sin(x^2)/x^2$ 无极限；上面的有界高阶导数或高阶导数趋零条件，正是用来排除这种越来越快的振荡。

\subsection{中值点的极限}

\begin{proposition}
\label{prop:pei-diff-mean-point-limit}
设 $n\ge2$，$f\in C^n$ 于 $x_0$ 的某邻域，
$f^{(2)}(x_0)=\cdots=f^{(n-1)}(x_0)=0$、$f^{(n)}(x_0)\ne0$；$n=2$ 时中间的零值条件为空。对每个充分小的非零 $h$，任取满足
\[
 \frac{f(x_0+h)-f(x_0)}h=f'(x_0+\theta(h)h),
 \qquad 0<\theta(h)<1
\]
的中值参数。则
\[
 \lim_{h\to0}\theta(h)=n^{-1/(n-1)}.
\]
\end{proposition}
\begin{proof}
把左边的 $f$ 与右边的 $f'$ 都在 $x_0$ 展开，并消去共同的一阶导数项，得
\[
 \frac{f^{(n)}(x_0)}{n!}h^{n-1}+o(h^{n-1})
 =\frac{f^{(n)}(x_0)}{(n-1)!}\theta(h)^{n-1}h^{n-1}
  +o(h^{n-1}).
\]
右侧余项可统一写成 $o(h^{n-1})$，因为其原形式为 $o((\theta(h)h)^{n-1})$，且 $0<\theta(h)<1$。除以非零的 $f^{(n)}(x_0)h^{n-1}$，得到 $\theta(h)^{n-1}\to1/n$。由于 $\theta(h)>0$，便得所需极限。证明也说明，不必先假定 $\theta(h)$ 有极限或中值点唯一。
\end{proof}

\begin{example}
\label{ex:pei-diff-arctan-mean-point}
对 $f(x)=\arctan x$，若 $[f(h)-f(0)]/h=f'(\theta(h)h)$，求 $h\to0$ 时 $\theta(h)$ 的极限。
\end{example}
\begin{solve}
因为 $f''(0)=0$、$f'''(0)=-2\ne0$，命题~\ref{prop:pei-diff-mean-point-limit} 取 $n=3$ 给出 $\theta(h)\to1/\sqrt3$。也可由 $\arctan h/h=1-h^2/3+o(h^2)$ 与 $1/[1+\theta(h)^2h^2]$ 比较得到同一结果。只有首个非零的高阶项为二阶时，中值参数才由这一论证趋于 $1/2$。
\end{solve}

\subsection{函数方程与展开系数的唯一性}

\begin{example}
\label{ex:pei-diff-fixed-mean-parameter}
设 $f\in C^3(\mathbb R)$。假设存在固定的 $\theta\in(0,1)$，使对任意 $x$ 及充分小的 $h$ 恒有
\[
 f(x+h)-f(x)=h f'(x+\theta h).
\]
证明：$\theta\ne1/2$ 时 $f$ 为至多一次多项式；$\theta=1/2$ 时 $f$ 为至多二次多项式。
\end{example}
\begin{solve}
在固定的 $x$ 处展开并除以 $h$，得
\[
 f'(x)+\tfrac12h f''(x)+\tfrac16h^2f'''(x)+o(h^2)
 =f'(x)+\theta h f''(x)+\tfrac12\theta^2h^2f'''(x)+o(h^2).
\]
若 $\theta\ne1/2$，比较一次项得 $f''(x)=0$ 对每个 $x$ 成立。若 $\theta=1/2$，一次项抵消，比较二次项得 $(1/6-1/8)f'''(x)=0$，故 $f'''\equiv0$。逐次由中值定理得到所述多项式形式。这里的固定参数恒等式远强于中值定理中允许中值点随 $x,h$ 改变的存在性结论。
\end{solve}

\begin{proposition}[有限阶渐近展开的唯一性]
\label{prop:pei-diff-taylor-unique}
若在 $x\to x_0$ 时
\[
 f(x)=\sum_{k=0}^n a_k(x-x_0)^k+o((x-x_0)^n)
     =\sum_{k=0}^n b_k(x-x_0)^k+o((x-x_0)^n),
\]
则 $a_k=b_k$ 对所有 $k=0,\ldots,n$ 成立。若 $f$ 满足定理~\ref{thm:pei-diff-taylor-peano} 的条件，这些系数必为 $f^{(k)}(x_0)/k!$。
\end{proposition}
\begin{proof}
相减后得到一个多项式等于 $o((x-x_0)^n)$。若有系数不同，取最小的下标 $j$ 使 $a_j\ne b_j$，再除以 $(x-x_0)^j$ 并令 $x\to x_0$，左端趋于非零常数 $a_j-b_j$，右端趋于零，矛盾。后一结论由 Taylor 公式与唯一性比较得到，不能仅凭展开形式反推所有阶导数存在。
\end{proof}

\subsection{大小阶记号与求和中的一致控制}

固定一个趋近过程，并设比较函数 $g$ 在相应去心邻域内非零。记 $r=O(g)$，表示存在常数 $C$ 使最终有 $|r|\le C|g|$；记 $r=o(g)$，表示 $r/g\to0$。特别地，$O(1)$ 只表示有界，$o(1)$ 才表示趋于零。有限个同阶小量可以相加，但当项数也趋于无穷时，必须控制每一项的误差是否对下标一致；仅仅把每项写成 $o(1)$，不足以推出总和也是 $o(1)$。

\begin{example}
\label{ex:pei-diff-uniform-small-remainders}
设 $f(0)=0$ 且 $f'(0)=A$。求
$\displaystyle\lim_{n\to\infty}\sum_{k=1}^n f(k/n^2)$。
\end{example}
\begin{solve}
写 $f(t)=At+t\varepsilon(t)$，其中 $\varepsilon(t)\to0$。对充分大的 $n$，所有采样点都在 $(0,1/n]$ 内，令 $\eta_n=\sup_{0<t\le1/n}|\varepsilon(t)|$，则 $\eta_n\to0$。因此
\[
 \sum_{k=1}^n f(k/n^2)
 =A\frac{n+1}{2n}+E_n,
 \qquad |E_n|\le\eta_n\frac{n+1}{2n}\longrightarrow0.
\]
极限为 $A/2$。关键不在于每个固定下标的余项趋零，而在于全部采样点共同趋于展开点，且权重之和保持有界。
\end{solve}

\section{不等式与凸函数}
\label{sec:pei-diff-convexity}
% 教材对应：PDF 第 68--89 页。

\subsection{用单调性、中值定理和余项证明不等式}

将待证不等式写成 $F(x)\ge0$ 后，先寻找有利的基准点 $x_0$：若 $F(x_0)=0$，就可通过 $F'$ 的符号证明 $x_0$ 为最小值点；若 $F'$ 仍难判断，可继续求导，直到出现能确定符号的表达式，再逐层返回。中值定理适合把函数差转成导数的上下界，Taylor 公式适合让低阶项抵消后由高阶余项决定符号。对于“求所有参数，使不等式对所有自变量成立”，应将其转化为参数与某函数上确界或下确界的比较，再核实界是否达到以及原不等式是否严格。

\begin{example}
\label{ex:pei-diff-log-inequalities}
证明对所有 $x>-1$，
\[
 \frac{x}{1+x}\le\ln(1+x)\le x,
\]
且两个等号都仅在 $x=0$ 成立。
\end{example}
\begin{solve}
令 $F(x)=x-\ln(1+x)$、$G(x)=\ln(1+x)-x/(1+x)$，则
\[
 F'(x)=\frac{x}{1+x},\qquad G'(x)=\frac{x}{(1+x)^2}.
\]
二者在 $(-1,0)$ 上严格递减、在 $(0,+\infty)$ 上严格递增，且 $F(0)=G(0)=0$，所以都以零为唯一最小值。这个证明同时处理正、负 $x$；只在正半轴证明导数非负，并不足以得到整个定义域的结论。
\end{solve}

\begin{example}
\label{ex:pei-diff-power-sum}
设 $s\ge0$ 为实数、$n\ge1$ 为整数，证明
\[
 \frac{n^{s+1}}{s+1}\le\sum_{k=1}^n k^s
 \le\frac{(n+1)^{s+1}-1}{s+1}.
\]
\end{example}
\begin{solve}
对 $H(t)=t^{s+1}/(s+1)$，分别在 $[k-1,k]$ 与 $[k,k+1]$ 上用中值定理。因为 $H'(t)=t^s$ 在正半轴单调不减，得到
\[
 H(k)-H(k-1)\le k^s\le H(k+1)-H(k).
\]
对 $k=1,\ldots,n$ 求和，函数增量依次抵消，即得所求界。$s=0$ 时函数增量直接等于 $1$，同样成立；$s>0$ 时两端均为严格不等式。此法把幂和比较转化为原函数的单位区间增量比较。
\end{solve}

\begin{example}
\label{ex:pei-diff-best-constant}
求使
\[
 \left(1+\frac1x\right)^{x+a}>e\qquad\text{对所有 }x>0
\]
成立的最小实数 $a$。
\end{example}
\begin{solve}
取对数后，条件等价于
$a>F(x)=1/\ln(1+1/x)-x$ 对所有 $x>0$ 成立。先证明 $F$ 严格递增。令
$Q(t)=t/\sqrt{1+t}-\ln(1+t)$，则 $Q(0)=0$，且对 $t>0$，
\[
 Q'(t)=\frac{1+t/2-\sqrt{1+t}}{(1+t)^{3/2}}>0,
\]
因为 $(1+t/2)^2>1+t$。因此 $\ln(1+1/x)<1/\sqrt{x(x+1)}$，从而
\[
 F'(x)=\frac1{x(x+1)\ln^2(1+1/x)}-1>0.
\]
又由 $\ln(1+t)=t-t^2/2+o(t^2)$ 可得
$F(x)=1/2+o(1)$（$x\to+\infty$）。严格递增且有极限说明 $F(x)<1/2$ 对每个有限 $x$ 成立，所以 $a=1/2$ 已满足原严格不等式；任意 $a<1/2$ 都会在足够大的 $x$ 处失败。故所求最小值为 $1/2$。上确界未达到，使最优参数仍可满足严格不等式。
\end{solve}

\subsection{凸性的定义、中点条件与 Jensen 不等式}

\begin{definition}
\label{def:pei-diff-convex}
设 $I$ 为区间，$f:I\to\mathbb R$。若任意 $x,y\in I$ 及 $\lambda\in[0,1]$ 都满足
\begin{equation}
 f(\lambda x+(1-\lambda)y)
 \le\lambda f(x)+(1-\lambda)f(y),
 \label{eq:pei-diff-convex-definition}
\end{equation}
则称 $f$ 为\emph{凸函数（convex function）}。若 $x\ne y$、$0<\lambda<1$ 时不等式总是严格的，则称为\emph{严格凸函数（strictly convex function）}。把不等号反向得到凹函数与严格凹函数；等价地，$f$ 为凹函数当且仅当 $-f$ 为凸函数。
\end{definition}

这个约定表示函数图像位于任意弦的下方。若只要求
$f((x+y)/2)\le[f(x)+f(y)]/2$，称为\emph{中点凸（midpoint convex）}。完整凸性必推出中点凸性；反向则需要正则条件。例如，当 $f$ 连续时，反复二等分可先证明式~\eqref{eq:pei-diff-convex-definition} 对所有二进有理权重成立，再利用这些权重在 $[0,1]$ 中稠密及函数连续性，推广到任意权重。\emph[caution]{不能删去正则条件，直接把中点不等式当作完整凸性的定义}。

\begin{theorem}[Jensen 不等式]
\label{thm:pei-diff-jensen}
若 $f$ 在区间 $I$ 上为凸函数，$x_1,\ldots,x_n\in I$，$\lambda_i\ge0$ 且 $\sum_{i=1}^n\lambda_i=1$，则
\begin{equation}
 f\left(\sum_{i=1}^n\lambda_i x_i\right)
 \le\sum_{i=1}^n\lambda_i f(x_i).
 \label{eq:pei-diff-jensen}
\end{equation}
若 $f$ 严格凸，则等号当且仅当所有正权重所对应的 $x_i$ 相同。
\end{theorem}
\begin{proof}
零权重可以删除。对正权重的个数归纳：两个点时就是定义；若结论对 $n-1$ 个点成立，将前 $n-1$ 个点合并成加权均值
$z=(\sum_{i<n}\lambda_i x_i)/(1-\lambda_n)$，先对 $z,x_n$ 使用凸性，再对 $z$ 使用归纳假设。严格凸时，任意一次合并了不同的正权重点都会产生严格不等式；只有所有这些点都相同时才能始终取等号。
\end{proof}

\begin{example}
\label{ex:pei-diff-weighted-amgm}
设 $a_i>0$、$\lambda_i>0$、$\sum\lambda_i=1$。由 $\ln t$ 在 $(0,+\infty)$ 上严格凹，
\[
 \sum_i\lambda_i\ln a_i\le\ln\left(\sum_i\lambda_i a_i\right),
 \qquad
 \prod_i a_i^{\lambda_i}\le\sum_i\lambda_i a_i.
\]
两式等号都当且仅当 $a_1=\cdots=a_n$。这是加权算术平均与几何平均不等式；使用对数前的正值条件和严格凹性给出的等号条件都应保留。
\end{example}

\subsection{割线、单侧导数与支撑直线}

\begin{proposition}
\label{prop:pei-diff-convex-slopes}
函数 $f$ 在区间 $I$ 上凸，当且仅当任意 $u<v<w$ 都有
\begin{equation}
 \frac{f(v)-f(u)}{v-u}
 \le\frac{f(w)-f(u)}{w-u}
 \le\frac{f(w)-f(v)}{w-v}.
 \label{eq:pei-diff-convex-slopes}
\end{equation}
\end{proposition}
\begin{proof}
写 $v=\lambda u+(1-\lambda)w$，其中 $\lambda=(w-v)/(w-u)$。把定义中的凸性不等式乘以 $w-u$ 并移项，即得两条割线斜率不等式。反过来，任取两点之间的凸组合，可由相同变形还原定义。几何上，这表示从左到右的割线斜率依次增大。
\end{proof}

\begin{proposition}
\label{prop:pei-diff-convex-local-properties}
有限值凸函数在定义区间的每个内部紧子区间上满足 Lipschitz 条件，因而在每个内点连续；其左右导数在每个内点都存在且有限，并满足
\[
 f'_-(x)\le f'_+(x),\qquad
 f'_+(x)\le f'_-(y)\quad(x<y).
\]
对任一内点 $x_0$ 及任一
$s\in[f'_-(x_0),f'_+(x_0)]$，都有
\begin{equation}
 f(x)\ge f(x_0)+s(x-x_0),\qquad x\in I.
 \label{eq:pei-diff-supporting-line}
\end{equation}
相应直线称为 $f$ 在 $x_0$ 处的\emph{支撑直线（supporting line）}。
\end{proposition}
\begin{proof}
先取定义区间内的闭区间 $[u,v]$。凸性给出上界 $M=\max\{f(u),f(v)\}$；以 $c=(u+v)/2$ 为中心，对每个 $x\in[u,v]$ 取对称点 $x^*=u+v-x$，有 $2f(c)\le f(x)+f(x^*)$，故还有下界 $m=2f(c)-M$。对内部紧子区间 $[\alpha,\beta]\subset(u,v)$，用式~\eqref{eq:pei-diff-convex-slopes} 把任意两点间的割线斜率夹在通向外侧端点的割线斜率之间；分母分别至少为 $\alpha-u$、$v-\beta$，分子绝对值至多为 $M-m$，所以斜率绝对值有统一上界，得到 Lipschitz 条件。

固定内点后，左右差商分别随端点位置单调变化，并被通向更外侧固定点的割线斜率夹住，所以单侧极限有限且存在。对割线斜率取相应的单侧极限，即得上述导数次序。最后，$x>x_0$ 时割线斜率不小于 $f'_+(x_0)\ge s$；$x<x_0$ 时割线斜率不大于 $f'_-(x_0)\le s$，再注意乘以负的 $x-x_0$ 要反向，便得支撑不等式。
\end{proof}

\begin{notealert}[type=caution,title={凸性不保证端点连续或处处可导}]
函数 $f(x)=|x|$ 为凸函数，但原点只有左右导数，没有普通导数。函数 $f(0)=1$、$f(x)=0$（$0<x\le1$）在 $[0,1]$ 上仍凸，却在左端点不连续。因此上述连续性、局部 Lipschitz 性及有限单侧导数的结论，均须保留“内点”这一范围。
\end{notealert}

\subsection{导数判别、严格凸性与应用}

\begin{theorem}
\label{thm:pei-diff-convex-derivatives}
设 $f$ 在开区间 $I$ 内可导。下列条件等价：$f$ 凸；$f'$ 单调不减；任意 $x,x_0\in I$ 都有
\[
 f(x)\ge f(x_0)+f'(x_0)(x-x_0).
\]
在此条件下，$f$ 严格凸当且仅当 $f'$ 严格递增。若 $f$ 二次可导，则凸性等价于 $f''\ge0$；严格凸性等价于 $f''\ge0$ 且 $f''$ 在任一非退化子区间上都不恒为零。
\end{theorem}
\begin{proof}
凸性推出导数单调与支撑不等式，见命题~\ref{prop:pei-diff-convex-local-properties}。若 $f'$ 单调，则对 $u<v<w$ 的两个相邻区间使用中值定理，可得左割线斜率不大于右割线斜率，从而由命题~\ref{prop:pei-diff-convex-slopes} 得凸性。若已知切线支撑不等式，在两端 $u,w$ 分别以它们之间的点 $v$ 为支撑点，乘以使 $v=\lambda u+(1-\lambda)w$ 的权重再相加，导数项抵消，即得定义中的凸性。

当导数严格递增时，割线比较也严格，故严格凸。反过来，凸函数的导数已经单调；若不同两点有相同导数，它在两点间必恒定，由中值定理 $f$ 在该段为仿射函数，与严格凸矛盾。二阶判别来自 $f'$ 的单调性判别；在 $f''\ge0$ 时，导数在某段不严格增加恰意味着它在该段恒定，也恰意味着 $f''$ 在该段恒为零。
\end{proof}

\emph[caution]{$f''>0$ 是严格凸的充分条件，并非必要条件}。例如 $f(x)=x^4$ 的导数 $4x^3$ 严格递增，所以它严格凸，虽然 $f''(0)=0$。严格凸函数的极小值若存在则唯一：若在两个不同点取得相同的最小值，内部凸组合的函数值会严格更小。凸函数的任何内点局部最小值也是全局最小值，因为只要另一点函数值更低，靠近该局部最小点的凸组合就已违反局部最小性。凹函数的对应结论由取负得到。

\begin{example}
\label{ex:pei-diff-convex-ratio}
设 $f$ 在 $[0,+\infty)$ 上非负且严格凸，$f(0)=0$。证明 $f(x)/x$ 在 $(0,+\infty)$ 上严格递增，且 $f$ 本身在 $[0,+\infty)$ 上严格递增。
\end{example}
\begin{solve}
对 $0<x<y$，将 $x$ 写成 $(x/y)y+(1-x/y)0$，严格凸性给出 $f(x)<(x/y)f(y)$，即 $f(x)/x<f(y)/y$。又 $f$ 在正点不能为零，否则该点与原点之间的严格凸性会迫使函数出现负值。因此 $f(y)>0$，再由 $x/y<1$ 得 $f(x)<f(y)$；与原点比较也严格成立。
\end{solve}

\begin{example}
\label{ex:pei-diff-positive-composition}
是否存在 $C^1(\mathbb R)$ 函数，满足对所有 $x\in\mathbb R$ 都有 $f(x)>0$ 且 $f'(x)=f(f(x))$？
\end{example}
\begin{solve}
不存在。由等式及正值性得 $f'>0$，故 $f$ 严格递增。记 $c=f(0)>0$，因 $f(x)>0$，有 $f'(x)=f(f(x))>f(0)=c$ 对所有实数成立。取 $x<-1$，在 $[x,0]$ 上用中值定理，得到
$f(0)-f(x)>-cx$，即 $f(x)<c(1+x)<0$，矛盾。此题的关键是定义域为整个实轴及函数始终正值；仅凭局部微分关系不能推出同样的不存在性结论。
\end{solve}

\section{导数的综合应用}
\label{sec:pei-diff-derivative-applications}
% 教材对应：PDF 第 89--102 页。

\subsection{极值、最值与高阶判别}

\begin{definition}
设 $x_0$ 为定义域 $D$ 中的一点。若存在 $\delta>0$，使对所有 $x\in D\cap(x_0-\delta,x_0+\delta)$ 都有 $f(x)\le f(x_0)$，则称 $x_0$ 为\emph{局部极大值点（local maximum point）}；若对 $x\ne x_0$ 不等式严格，则称为严格局部极大值点。局部极小值点反向定义。若相应比较对所有 $x\in D$ 成立，则为全局最大值或全局最小值。可导且导数为零的点称为\emph{驻点（stationary point）}。
\end{definition}

求最值时，候选位置包括内部驻点、内部不可导点及包含在定义域中的端点；在开区间或无界区间上，还须考察边界极限及无穷远处的行为。Fermat 定理只说明可导的内部极值点一定是驻点，驻点未必有极值，例如 $x^3$ 在原点。若函数连续，且候选点左、右邻域内导数由正变负，则为严格局部极大值；由负变正，则为严格局部极小值。仅仅把候选点的函数值比较一遍，须以已经证明全局最值存在且候选位置完整为前提。

\begin{proposition}[首个非零高阶导数判别]
\label{prop:pei-diff-higher-extremum}
设 $f$ 满足在 $x_0$ 处使用 $n$ 阶 Peano 公式的条件，且
$f'(x_0)=\cdots=f^{(n-1)}(x_0)=0$、$f^{(n)}(x_0)\ne0$。若 $n$ 为偶数，则 $f^{(n)}(x_0)>0$ 时为严格局部极小值，$f^{(n)}(x_0)<0$ 时为严格局部极大值；若 $n$ 为奇数，则不是局部极值点。
\end{proposition}
\begin{proof}
由定理~\ref{thm:pei-diff-taylor-peano}，
\[
 f(x)-f(x_0)=(x-x_0)^n
 \left(\frac{f^{(n)}(x_0)}{n!}+o(1)\right).
\]
括号内在充分小的去心邻域中与非零的首项同号，结论便由 $(x-x_0)^n$ 在两侧的符号得到。这里无需高阶导函数在整个邻域中连续。
\end{proof}

偶函数在原点可导时导数为零，但偶性本身不能判断极值类型。若二阶导数非零，可用上述 $n=2$ 的判别；若各阶导数都为零，仍可能有严格极值，例如例~\ref{ex:pei-diff-flat-function} 的函数在原点为严格全局最小值。反过来，极值点也不必彼此孤立：连续函数 $f(x)=x\sin(1/x)+2|x|$（$x\ne0$）、$f(0)=0$ 在原点有严格最小值，同时有其他极值点不断趋近原点：在正半轴上，$f'(1/(2m\pi))=2-2m\pi<0$，而 $f'(1/((2m+1)\pi))=2+(2m+1)\pi>0$（$m\ge1$），两列点之间交替出现的导数符号保证了不断出现的极值。

\begin{proposition}
\label{prop:pei-diff-countable-strict-extrema}
任意实函数的严格局部极大值点至多可数，严格局部极小值点也至多可数；非严格极值点不满足这一结论。
\end{proposition}
\begin{proof}
对每个严格局部极大值点 $x$，选取有理数 $p<x<q$，使 $(p,q)$ 包含在它的严格极大值邻域中，则 $x$ 是 $(p,q)$ 上唯一的最大值点。两个不同的这类点不可能选到同一对 $(p,q)$，而有理数对只有可数多个，故这些点至多可数。极小值对 $-f$ 使用同一论证。常函数的每个点都是非严格极大值点和极小值点，说明不能删除“严格”。
\end{proof}

\begin{remark}
严格极值点在有界闭区间上仍可有无穷多个，不能用紧致性直接断言它们只有有限个。两个严格局部极大值点之间，连续函数在连接两点的闭区间上取得最小值；两个端点都不能是最小值点，故至少有一个内部局部极小值点。但在没有进一步条件时，不能保证得到的极小值点也是严格的。
\end{remark}

\subsection{利用导数判断方程实根}

零点问题应分别证明存在性、唯一性与根的个数。连续性及端点异号可以证明存在；严格单调可以证明唯一；若导数只有有限个零点，可将定义域分为若干单调区间，逐段比较端点值。若要证明没有实根，寻找正的全局最小值往往比直接解方程更有效。遇到多项式，Rolle 定理还可以将原函数的根数与导数的根数联系起来，但重根必须计入，不能只数相邻根之间的零点。

\begin{example}
\label{ex:pei-diff-alternating-polynomial-roots}
设 $n\ge1$，
\[
 P_n(x)=1+\sum_{k=1}^n\frac{(-1)^k}{k}x^k.
\]
证明 $n$ 为奇数时方程 $P_n(x)=0$ 恰有一个实根，$n$ 为偶数时没有实根。
\end{example}
\begin{solve}
对 $x\ne-1$，等比求和给出
\[
 P_n'(x)=\sum_{k=1}^n(-1)^kx^{k-1}
        =\frac{-1+(-1)^n x^n}{1+x};
\]
在 $x=-1$ 处由原多项式求导得 $P_n'(-1)=-n$。若 $n$ 为奇数，则这个导数在全部实轴上严格为负，且由最高次项知 $P_n(-\infty)=+\infty$、$P_n(+\infty)=-\infty$，所以恰有一根。若 $n=2m$，则 $P_n'<0$ 于 $x<1$，$P_n'>0$ 于 $x>1$，故全局最小值在 $x=1$；而
\[
 P_{2m}(1)=\sum_{j=1}^{m-1}\left(\frac1{2j}-\frac1{2j+1}\right)
                  +\frac1{2m}>0.
\]
因此偶数情形没有实根。分式导数表达式在 $x=-1$ 的表面奇点应回到原多项式处理。
\end{solve}

\subsection{切线、法线、曲率圆与函数形状}

对可微曲线 $y=f(x)$，在 $x=x_0$ 处记 $y_0=f(x_0)$、$p=f'(x_0)$，切线方程为 $Y-y_0=p(X-x_0)$，法线可写成 $X-x_0+p(Y-y_0)=0$，该形式也适用于 $p=0$。对正则参数曲线 $(x(t),y(t))$，即切向量 $(x'(t_0),y'(t_0))\ne(0,0)$，更统一的写法为
\[
 y'(t_0)(X-x(t_0))-x'(t_0)(Y-y(t_0))=0
\]
和
\[
 x'(t_0)(X-x(t_0))+y'(t_0)(Y-y(t_0))=0.
\]
前者为切线，后者为法线；不必为了使用斜率形式而除以可能为零的分量。

\begin{example}
\label{ex:pei-diff-tangent-triangle}
在曲线 $y=x^2-1$ 的 $x>0$ 部分取点 $P$，其切线与两坐标轴围成三角形。求面积的最小值及相应切点。
\end{example}
\begin{solve}
设 $P=(x,x^2-1)$，切线为 $Y=2xX-x^2-1$。两个截距分别为 $(x^2+1)/(2x)$ 与 $-(x^2+1)$，故面积为
\[
 S(x)=\frac{(x^2+1)^2}{4x},\qquad
 S'(x)=\frac{(x^2+1)(3x^2-1)}{4x^2}.
\]
它在 $(0,1/\sqrt3)$ 下降，在 $(1/\sqrt3,+\infty)$ 上升，所以唯一最小值为 $4\sqrt3/9$，切点为 $(1/\sqrt3,-2/3)$。面积应使用截距乘积的绝对值，避免把第二个负截距误当作负面积。
\end{solve}

\begin{proposition}[曲率圆]
\label{prop:pei-diff-osculating-circle}
设 $f\in C^2$ 于 $x_0$ 附近，记 $p=f'(x_0)$、$q=f''(x_0)\ne0$。与曲线在 $(x_0,f(x_0))$ 处相切且具有相同二阶导数的\emph{曲率圆（osculating circle）}的圆心 $(a,b)$ 与半径 $R$ 为
\[
 a=x_0-\frac{p(1+p^2)}q,\qquad
 b=f(x_0)+\frac{1+p^2}q,\qquad
 R=\frac{(1+p^2)^{3/2}}{|q|}.
\]
相应的\emph{曲率（curvature）}为 $\kappa=1/R=|q|/(1+p^2)^{3/2}$。
\end{proposition}
\begin{proof}
对圆方程 $(X-a)^2+(Y-b)^2=R^2$ 视 $Y$ 为 $X$ 的局部函数，连续求导得
$X-a+(Y-b)Y'=0$ 和 $1+(Y')^2+(Y-b)Y''=0$。在切点代入 $Y'=p$、$Y''=q$，先解得圆心，再用切点到圆心的距离求出 $R$。条件 $q\ne0$ 保证这里得到有限半径的圆。
\end{proof}

函数形状的研究应先确定定义域、连续点与不可导点，再用 $f'$ 的符号判定增减、用 $f''$ 的符号判定凸凹，并考察端点与无穷远的极限。\emph{拐点（inflection point）}要求图像在该点两侧的凸凹性发生改变；$f''(x_0)=0$ 既不是充分条件，也不适用于所有不可二次求导的情形。若 $f(x)\to\pm\infty$ 于有限端点，则可得到垂直渐近线；若在某一无穷远方向有
\[
 k=\lim\frac{f(x)}x,\qquad b=\lim(f(x)-kx)
\]
且两极限有限，则 $y=kx+b$ 是该方向的渐近线，其中 $k=0$ 为水平渐近线。左右无穷远方向应分别计算。

\subsection{实际问题的目标函数与变化率}

实际应用先把几何、运动或物理条件转成明确模型，再选择一个自变量，写出目标函数及可行范围，最后证明候选点确为全局最优。模型中出现的“速度”要说明是哪一个量对时间的导数；若使用物理平衡原理，也须说明所采用的物理假设，不能只凭求出一个驻点就断言平衡或最优。

\begin{example}
\label{ex:pei-diff-circumscribed-cone}
给定半径为 $r>0$ 的球，求外切于它的直圆锥的最小体积。
\end{example}
\begin{solve}
记圆锥高为 $h$、底面半径为 $R$，球心在圆锥轴线上且距底面为 $r$。由轴截面中球心到母线的距离等于 $r$，有
\[
 \frac{R(h-r)}{\sqrt{h^2+R^2}}=r,
 \qquad R^2=\frac{r^2h}{h-2r},\qquad h>2r.
\]
因此体积与导数分别为
\[
 V(h)=\frac{\pi r^2}{3}\frac{h^2}{h-2r},\qquad
 V'(h)=\frac{\pi r^2}{3}\frac{h(h-4r)}{(h-2r)^2}.
\]
在可行范围内，$V$ 先减后增，唯一全局最小值在 $h=4r$ 取得；此时 $R=\sqrt2r$，故 $V_{\min}=8\pi r^3/3$。可行条件 $h>2r$ 排除了代数求导所得的无效候选点。
\end{solve}

\begin{example}
\label{ex:pei-diff-angular-speed}
船从点 $B$ 以恒定线速度 $v>0$ 向东航行，观察者位于 $B$ 正南、相距 $s>0$ 的固定点。若用视线相对初始方向的角速度衡量视线转动的快慢，求它与初始角速度的比值。
\end{example}
\begin{solve}
令时间为 $t\ge0$，视线偏转角为 $\theta(t)\in[0,\pi/2)$，并以弧度为单位。直角三角形关系给出 $s\tan\theta=vt$，对时间求导得到
\[
 \theta'(t)=\frac vs\cos^2\theta(t),\qquad
 \frac{\theta'(t)}{\theta'(0)}=\cos^2\theta(t).
\]
所以比值为 $3/4,1/2,1/4$ 时，偏转角分别为 $30^\circ,45^\circ,60^\circ$。船的线速度始终为 $v$，发生变化的是模型所选定的视线角速度。
\end{solve}

导数用于求极限时，可以依据定理~\ref{thm:pei-diff-lhospital} 将比值问题转化为导数比，也可以依据 Taylor 公式提取首个非零项。前者要逐次检验法则条件，后者要确定足以越过相消项的展开阶数，并控制余项。若极限涉及移动的中值点，可参考命题~\ref{prop:pei-diff-mean-point-limit}；若涉及项数变化的求和，则需像例~\ref{ex:pei-diff-uniform-small-remainders} 一样给出对下标一致的误差估计。

\chapter{一元函数积分学}
\label{chap:pei-int-integral-calculus}

本章对应教材第四章，围绕积分与极限的互相转化、Riemann 可积性的判别、定积分估计与中值定理、经典积分不等式，以及反常积分的计算和收敛性展开。核心方法是先辨认可用的结构，再用积分和、振幅、单调性或尾部估计控制误差。有限区间上的积分默认指 Riemann 积分；在出现无界被积函数或无穷区间时，另按反常积分定义处理。除另有说明，区间均非退化、函数均取实值；有界性、非负性、连续性与绝对收敛性分别在所需位置注明，不能相互替代。

% 教材来源：数学分析/textbooks/数分习题（分章原文）/数分习题-4.pdf。
% 覆盖五节的主要结论、证明方法与代表性例题；按主题整合，非逐题转录。

\section{积分与极限}
\label{sec:pei-int-integral-limits}
% 教材对应：§4.1，PDF 第 1--17 页（原书第 250--266 页）。

\subsection{把和式、乘积与误差还原为积分}
\begin{theorem}
\label{thm:pei-int-riemann-sum}
设 $f$ 在 $[a,b]$ 上 Riemann 可积。对分划 $P_n:a=x_{n,0}<\cdots<x_{n,m_n}=b$，记网格宽度 $\lVert P_n\rVert=\max_i(x_{n,i}-x_{n,i-1})$。若 $\lVert P_n\rVert\to0$，则不论标记点 $\xi_{n,i}\in[x_{n,i-1},x_{n,i}]$ 如何选取，都有
\[
 \sum_{i=1}^{m_n}f(\xi_{n,i})(x_{n,i}-x_{n,i-1})\longrightarrow\int_a^b f(x)\,dx.
\]
\end{theorem}
识别积分和时，先确定区间与每项中的小区间长度，再确定函数值对应的标记点；\emph[caution]{分划不必等分，但最大的小区间长度必须趋于零}。例如 $\alpha>0$ 时，$n^{-\alpha-1}\sum_{k=1}^n k^\alpha\to1/(\alpha+1)$；若 $-1<\alpha<0$，被积函数 $x^\alpha$ 在零点无界，不能直接援引上述定理，须先截去零点附近的一段，再用单调性估计该段和式。这样仍得到同一极限。遇到乘积时先取对数，把乘积变成和式；取对数要求各因子为正，最后还要由指数函数的连续性还原极限。

\begin{example}
\label{ex:pei-int-geometric-partition}
设 $b\ge1$，求 $\displaystyle\lim_{n\to\infty}(b^{1/n}-1)\sum_{k=0}^{n-1}b^{k/n}\sin(b^{k/n})$。
\end{example}
\begin{solve}
当 $b=1$ 时原式恒为零。当 $b>1$ 时令 $x_k=b^{k/n}$，则 $x_{k+1}-x_k=(b^{1/n}-1)b^{k/n}$，而网格宽度不超过 $b(b^{1/n}-1)\to0$。原式为 $\sin x$ 在 $[1,b]$ 上的左端点积分和，故极限为 $\int_1^b\sin x\,dx=\cos1-\cos b$。
\end{solve}

\begin{proposition}
\label{prop:pei-int-quadrature-asymptotic}
设 $h=(b-a)/n$，$R_n=h\sum_{i=1}^n f(a+ih)$，$M_n=h\sum_{i=1}^n f(a+(i-\tfrac12)h)$。若 $f\in C^1[a,b]$，则
\begin{equation}
 \lim_{n\to\infty}n\left(R_n-\int_a^b f\right)=\frac{b-a}{2}\bigl(f(b)-f(a)\bigr).
 \label{eq:pei-int-right-sum-error}
\end{equation}
若 $f\in C^2[a,b]$，则
\begin{equation}
 \lim_{n\to\infty}n^2\left(\int_a^b f-M_n\right)=\frac{(b-a)^2}{24}\bigl(f'(b)-f'(a)\bigr).
 \label{eq:pei-int-midpoint-error}
\end{equation}
\end{proposition}
\begin{proof}
在每个右端点 $x_i$ 附近，由 $f'$ 的一致连续性，$f(x)=f(x_i)+f'(x_i)(x-x_i)+o(h)$，余项关于区间编号一致。逐段积分并求和，得到 $R_n-\int_a^b f=\tfrac{h^2}{2}\sum_i f'(x_i)+o(h)$，再应用积分和极限与 Newton--Leibniz 公式，即得式~\eqref{eq:pei-int-right-sum-error}。在中点 $m_i$ 展开至二阶，奇次项的积分为零，且 $\int_{x_{i-1}}^{x_i}(x-m_i)^2\,dx=h^3/12$，故 $\int_a^b f-M_n=\tfrac{h^3}{24}\sum_i f''(m_i)+o(h^2)$，由此得式~\eqref{eq:pei-int-midpoint-error}。余项的一致性保证 $n$ 个小区间的误差相加后仍具有所需阶数。
\end{proof}
例如取 $f(x)=1/(1+x)$、$[a,b]=[0,1]$，则 $n(\ln2-R_n)\to1/4$，$n^2(\ln2-M_n)\to1/32$。这类极限的目标是误差的首项，单独证明积分和收敛还不够。

\subsection{先分区间，再交换极限与积分}
若 $f_n\to f$ 在有限区间上一致成立，则 $|\int f_n-\int f|\le(b-a)\sup|f_n-f|\to0$。只有逐点收敛时，这个结论不自动成立；例如 $f_n(x)=n^2x(1-x)^n$ 在 $[0,1]$ 上逐点趋于零，但其积分为 $n^2/((n+1)(n+2))\to1$。教材中的许多例题采用\emph[tip]{小区间上控制长度，大区间上取得一致估计}，从而在 Riemann 积分框架内直接证明极限。

\begin{example}
证明 $\displaystyle\int_0^{\pi/2}\sin^n x\,dx\to0$。
\end{example}
\begin{solve}
固定 $0<\delta<\pi/2$，有
\[
 0\le\int_0^{\pi/2}\sin^n x\,dx
 \le\frac\pi2(\cos\delta)^n+\delta.
\]
先令 $n\to\infty$，得到上极限不超过 $\delta$，再令 $\delta\downarrow0$。这里先选并固定 $\delta$，再选 $n$，避免把依赖于 $n$ 的估计误当成一致估计。
\end{solve}

\begin{theorem}
\label{thm:pei-int-concentrating-kernel}
设 $K_n\ge0$ 在 $[a,b]$ 上可积，$\int_a^bK_n\to A<\infty$，且对每个 $\delta>0$，$\int_{[a,b]\setminus(x_0-\delta,x_0+\delta)}K_n\to0$。若 $f$ 在 $[a,b]$ 上有界、Riemann 可积，并在 $x_0\in[a,b]$ 连续，则
\[
 \int_a^b f(x)K_n(x)\,dx\longrightarrow Af(x_0).
\]
\end{theorem}
\begin{proof}
写成 $\int fK_n-f(x_0)\int K_n=\int(f-f(x_0))K_n$。在 $x_0$ 的 $\delta$ 邻域内利用连续性把 $|f-f(x_0)|$ 控制在任意给定的小量之下；在其余部分利用 $f$ 的有界性及核的质量趋于零。由于 $\int K_n$ 有界，两部分的误差都能趋于零；再用总质量的极限即得结论。
\end{proof}
这一结论称为集中核的逼近性质，关键在核的非负性、总质量和远离集中点后的质量估计。它统一了教材中“先减去待求常数，再估计误差”的拟合法。例如对连续的 $f$，
\begin{align*}
 \lim_{h\downarrow0}\int_0^1\frac{h f(x)}{h^2+x^2}\,dx&=\frac\pi2f(0),\\
 \lim_{n\to\infty}(n+1)\int_0^1x^nf(x)\,dx&=f(1),\\
 \lim_{n\to\infty}\frac{n+1}{2}\int_{-1}^1(1-|x|)^nf(x)\,dx&=f(0).
\end{align*}
第一式的核总质量为 $\arctan(1/h)\to\pi/2$，在 $[\delta,1]$ 上的积分不超过 $h/\delta$；第二式在 $[0,1-\delta]$ 上的质量为 $(1-\delta)^{n+1}$；第三式的总质量恒为一，远离零点后也具有几何衰减。各式因而无需预先假设 $f$ 可微。

\subsection{高次幂把积分集中到最大值附近}
\begin{theorem}
\label{thm:pei-int-power-integral}
设 $f\ge0$、$g>0$ 在 $[a,b]$ 上连续，$a<b$。则
\begin{equation}
 \lim_{n\to\infty}\left(\int_a^bf(x)^ng(x)\,dx\right)^{1/n}
 =\max_{[a,b]}f.
 \label{eq:pei-int-lp-max}
\end{equation}
若进一步设 $f$ 严格递减、$f(a)=1$、$f(b)=0$，记 $I_n=\int_a^bf(x)^n\,dx$，则 $I_{n+1}/I_n\to1$。
\end{theorem}
\begin{proof}
令 $M=\max f$。$M=0$ 时第一式成立；$M>0$ 时，上界为 $M(\int g)^{1/n}\to M$。对任意 $0<\varepsilon<M$，最大值点的某个单侧或双侧邻域 $J$ 具有正长度，且 $f>M-\varepsilon$；因 $\int_Jg>0$，下界为 $(M-\varepsilon)(\int_Jg)^{1/n}\to M-\varepsilon$。令 $\varepsilon\downarrow0$ 即得第一式。

对比值结论，固定 $a<c<d<b$，有
\[
 \frac{\int_d^bf^n}{I_n}\le\frac{b-d}{c-a}\left(\frac{f(d)}{f(c)}\right)^n\longrightarrow0.
\]
故 $\liminf I_{n+1}/I_n\ge f(d)$，而上界为一。令 $d\downarrow a$，便得到比值极限。这个证明还说明幂权重在最大值点附近集中；比值极限并非仅由 $I_n^{1/n}\to1$ 推出。
\end{proof}
若 $f\ge0$ 连续且严格递增，积分中值点 $x_n$ 由 $f(x_n)^n=(b-a)^{-1}\int_a^bf(x)^n\,dx$ 确定，则式~\eqref{eq:pei-int-lp-max} 给出 $f(x_n)\to f(b)$。严格单调性与连续性遂推出 $x_n\to b$；中值点的位置需要结合函数的单调性判断。

\subsection{快速振荡的平均作用与平移连续性}
\begin{theorem}
\label{thm:pei-int-periodic-average}
设 $f$ 在 $[a,b]$ 上 Riemann 可积，$g$ 以 $T>0$ 为周期且在 $[0,T]$ 上 Riemann 可积，记 $\bar g=T^{-1}\int_0^Tg(t)\,dt$。则
\begin{equation}
 \lim_{\lambda\to+\infty}\int_a^bf(x)g(\lambda x)\,dx
 =\bar g\int_a^bf(x)\,dx.
 \label{eq:pei-int-periodic-average}
\end{equation}
特别地，$\int_a^bf(x)\sin(\lambda x)\,dx\to0$，余弦的结论相同，这就是此处的 Riemann--Lebesgue 引理。
\end{theorem}
\begin{proof}
先设 $f$ 为一个区间的示性函数。将积分区间分成完整周期与至多两个残段，完整周期的积分等于区间长度乘以 $\bar g$，残段的误差为 $O(\lambda^{-1})$；由线性性，阶梯函数也成立。一般可积函数可用阶梯函数 $s$ 在积分意义下逼近，使 $\int_a^b|f-s|$ 任意小，而
\[
 \left|\int_a^b(f-s)g(\lambda x)\,dx\right|
 \le\sup|g|\int_a^b|f-s|\,dx.
\]
右边与 $\lambda$ 无关，故可先固定 $s$，再取极限，最后令逼近误差趋于零。$g$ 允许变号；也可将其分解为正部与负部分别处理。
\end{proof}

\begin{proposition}
\label{prop:pei-int-translation}
将 $[a,b]$ 上的 Riemann 可积函数 $f$ 在区间外延拓为零，则
\[
 \lim_{h\to0}\int_{\mathbb R}|f(x+h)-f(x)|\,dx=0.
\]
若 $f$ 在较大的区间上连续，则还有
\[
 \lim_{h\to0}\frac1h\int_a^b\bigl(f(x+h)-f(x)\bigr)\,dx=f(b)-f(a).
\]
\end{proposition}
\begin{proof}
第一式先对阶梯函数证明：平移后只在有限个分点附近发生变化，故积分为 $O(|h|)$。选阶梯函数 $s$ 使 $\int|f-s|<\varepsilon$，再由平移换元与三角不等式，原积分不超过 $2\varepsilon+\int|s(x+h)-s(x)|$，于是得证。第二式把分子改写为 $\int_b^{b+h}f-\int_a^{a+h}f$，除以 $h$ 后由端点连续性取极限即可，并不要求 $f'$ 存在。
\end{proof}

\section{定积分的可积性}
\label{sec:pei-int-integrability}
% 教材对应：§4.2，PDF 第 17--32 页（原书第 266--281 页）。

\subsection{定义中的量词与 Darboux 判别}
\begin{definition}
设 $f:[a,b]\to\mathbb R$，$a<b$。若存在实数 $I$，使对每个 $\varepsilon>0$，都存在 $\delta>0$，对任意满足 $\lVert P\rVert<\delta$ 的分划及任意标记点均有 $|\sum_i f(\xi_i)\Delta x_i-I|<\varepsilon$，则称 $f$ \emph{Riemann 可积（Riemann integrable）}，并记 $I=\int_a^bf$。这里 $\delta$ 只依赖于 $\varepsilon$ 与 $f$，不能依赖于其后任选的分划或标记点。
\end{definition}
可积函数必有界：若不然，任取一个足够细的固定分划，至少一个小区间上函数无界；固定其余标记点而改变这个小区间的标记点，就能使积分和无界，与定义矛盾。以下讨论均设 $f$ 有界。对 $P:a=x_0<\cdots<x_m=b$，定义
\[
 m_i=\inf_{[x_{i-1},x_i]}f,\quad M_i=\sup_{[x_{i-1},x_i]}f,\quad
 L(f,P)=\sum_i m_i\Delta x_i,\quad U(f,P)=\sum_i M_i\Delta x_i,
\]
并称 $\omega_i=M_i-m_i$ 为该区间上的振幅（oscillation）。加细分划使下和不减、上和不增；任意两分划的下和都不超过对方的上和，这可通过共同加细分划证明。因此下积分 $\sup_P L(f,P)$ 不超过上积分 $\inf_P U(f,P)$。

\begin{theorem}
\label{thm:pei-int-darboux}
对有界函数 $f$，下列条件等价：$f$ 可积；上下积分相等；对任意 $\varepsilon>0$ 存在一个分划 $P$ 使
\begin{equation}
 U(f,P)-L(f,P)=\sum_i\omega_i\Delta x_i<\varepsilon;
 \label{eq:pei-int-oscillation-sum}
\end{equation}
以及当 $\lVert P\rVert\to0$ 时，该振幅和趋于零。
\end{theorem}
\begin{proof}
由可积定义，分别选取逼近各小区间上确界、下确界的标记点，可使上和与下和任意接近同一积分值；反过来，积分和总夹在上下和之间。上下积分相等当且仅当存在使其差任意小的分划，这又由共同加细分划及上下积分的确界定义得到。

最后说明“存在一个分划”为何能推出“所有足够细分划”。固定 $P_0$ 使其上下和之差小于 $\varepsilon/2$，设 $P_0$ 有 $r$ 个内分点、$|f|\le M$。任意分划 $Q$ 中，内部不含 $P_0$ 分点的小区间可按 $P_0$ 的区间归组，其振幅和不超过 $P_0$ 的振幅和；其余小区间至多有 $r$ 个，总长不超过 $r\lVert Q\rVert$，贡献至多 $2Mr\lVert Q\rVert$。令网格宽度充分小，就得到 $U(f,Q)-L(f,Q)<\varepsilon$。这也说明两个表述中的量词不同，但在有界条件下确实等价。
\end{proof}

\begin{proposition}
\label{prop:pei-int-large-oscillation}
有界函数 $f$ 可积，当且仅当对每个 $\eta>0$ 和 $\rho>0$，都存在分划，使振幅不小于 $\eta$ 的小区间总长小于 $\rho$。
\end{proposition}
\begin{proof}
必要性由 $\eta\sum_{\omega_i\ge\eta}\Delta x_i\le\sum_i\omega_i\Delta x_i$ 得到。反过来，设 $|f|\le M$，把振幅和分成两类，得 $\sum_i\omega_i\Delta x_i\le\eta(b-a)+2M\rho$；先使 $\eta$ 小，再使 $\rho$ 小，即可应用定理~\ref{thm:pei-int-darboux}。
\end{proof}
连续函数用一致连续性控制每个振幅；单调函数则有 $\sum_i\omega_i\Delta x_i\le\lVert P\rVert|f(b)-f(a)|$。更一般地，若 $f$ 的全变差 $V_a^b(f)=\sup_P\sum_i|f(x_i)-f(x_{i-1})|$ 有限，则每个振幅不超过相应区间的变差，且变差可加，从而振幅和不超过 $\lVert P\rVert V_a^b(f)$。这证明了\emph{有界变差函数（function of bounded variation）}的可积性。

\subsection{把困难点围起来：可积而不连续的函数}
\begin{example}
\label{ex:pei-int-thomae}
沿用例~\ref{ex:pei-thomae} 中的 Thomae 函数 $R$（教材称 Riemann 函数），证明它在 $[0,1]$ 上可积且积分为零。
\end{example}
\begin{solve}
固定 $N$，分母 $q\le N$ 的点只有有限个，用总长小于 $\rho$ 的有限个小区间将它们覆盖，并选分划使覆盖区间的端点成为分点。覆盖区间对上和的贡献至多为 $\rho$，其余部分函数不超过 $1/N$，故 $U(R,P)\le\rho+1/N$；每个非退化区间含无理数，下和恒为零。令 $\rho\downarrow0$、$N\to\infty$，便得积分为零。其连续点的刻画已在例~\ref{ex:pei-thomae} 中证明，这里突出同一函数的可积性。
\end{solve}
同样的方法适用于有界函数的间断点能由总长任意小的有限个区间覆盖的情形：覆盖区间上用有界性控制，余下紧集上用局部连续性和有限覆盖控制振幅。有限个间断点、以及只有有限个聚点的某些间断点列都可这样处理。\emph[caution]{不能把“间断点可数”直接替换成“能被总长任意小的有限个区间覆盖”}；例如上例的有理间断点稠密，证明时须围住函数值较大的有限多个点，而非一次围住全部间断点。

\begin{proposition}
\label{prop:pei-int-composition}
若 $f,g$ 在 $[a,b]$ 上可积，则 $|f|$、$f^2$、$fg$、$\max(f,g)$、$\min(f,g)$ 可积。若 $f([a,b])\subset[A,B]$，而 $\varphi$ 在 $[A,B]$ 上连续，则 $\varphi\circ f$ 可积。
\end{proposition}
\begin{proof}
绝对值满足 $\bigl||u|-|v|\bigr|\le|u-v|$；在有界值域内平方函数满足 Lipschitz 估计。先由振幅判别得到 $|f|$ 和 $f^2$ 的可积性，再用 $fg=((f+g)^2-(f-g)^2)/4$ 与 $\max(f,g)=(f+g+|f-g|)/2$ 等恒等式。一般连续复合的情形，由 $\varphi$ 的一致连续性，对给定 $\varepsilon$ 选择 $\eta$，使 $|u-v|<\eta$ 时 $|\varphi(u)-\varphi(v)|<\varepsilon$；振幅小于 $\eta$ 的区间上复合函数的振幅小，而其余区间总长可以任意小，故由命题~\ref{prop:pei-int-large-oscillation} 得证。
\end{proof}

\begin{notealert}[type=caution,title={不连续的外函数不能直接套用复合结论}]
教材本节单元练习关于取整函数的提示声称取整不增大振幅，这一步一般不成立。取 $f=-R$，其中 $R$ 为例~\ref{ex:pei-int-thomae} 的函数，则 $f$ 可积，而 $\lfloor f(x)\rfloor$ 在有理点等于 $-1$、无理点等于零，每个区间上的振幅均为一，所以不可积。跨过整数阈值时，即使 $f$ 的振幅很小，取整后的振幅仍可为一；使用连续复合定理时必须核对外函数的连续性。
\end{notealert}

\subsection{零积分、连续点与导函数}
\begin{theorem}
\label{thm:pei-int-zero-integral}
可积函数的连续点在 $[a,b]$ 上稠密。对非负可积函数 $f$，$\int_a^bf=0$ 当且仅当 $f$ 在每个连续点均为零。因此连续非负函数的积分为零便处处为零；处处严格为正的可积函数的积分严格为正。
\end{theorem}
\begin{proof}
在任意非退化子区间内，由振幅和可以任意小，能找到一个振幅小于 $1/2$ 的小区间；在其内部取非退化闭区间 $J_1$。重复这一过程，构造 $J_{n+1}\subset\operatorname{int}J_n$，使 $|J_n|\to0$ 且 $f$ 在 $J_n$ 上的振幅小于 $2^{-n}$。区间套的公共点 $c$ 对每个 $J_n$ 都是内点，故 $f$ 在 $c$ 连续。初始子区间任意，因而连续点稠密。

若连续点 $c$ 满足 $f(c)>0$，则某个邻域内 $f\ge f(c)/2$，积分必为正。反过来，若所有连续点处 $f=0$，每个非退化小区间都有一个这样的点，所以每个 Darboux 下和为零，由可积性得积分为零。最后两条结论分别利用处处连续和至少存在一个连续点即得。
\end{proof}
若连续函数 $f\ge g$ 且二者积分相等，则对连续非负函数 $f-g$ 应用上述结论，得到 $f\equiv g$。积分估计中一串不等式首尾相等时，这一事实可逐步迫使被积函数相等，进而与可微性等额外条件矛盾。

\begin{proposition}
设 $F\in C^1[a,b]$，则 $F'$ 可积且 $\int_a^bF'=F(b)-F(a)$。更一般地，设 $F$ 在 $[a,b]$ 连续、在 $(a,b)$ 可微，并对导数作端点赋值，则 $F'$ 可积当且仅当存在可积函数 $g$ 使 $F(x)=F(a)+\int_a^xg(t)\,dt$。
\end{proposition}
\begin{proof}
第一条既可由连续可积性和微积分基本定理得到，也可逐段用微分中值定理，将 $F(b)-F(a)$ 写成 $\sum_i F'(\eta_i)\Delta x_i$，再用 $F'$ 的一致连续性与任意标记的和比较。

对第二条，必要性取 $g=F'$。充分性中，固定一个小区间 $J$，记 $m=\inf_Jg$、$M=\sup_Jg$；由积分平均值的上下界，对其内任意 $u<v$ 有 $m\le(F(v)-F(u))/(v-u)\le M$，令端点向任意内点逼近，得到 $m\le F'\le M$。故除有限分点的赋值外，$F'$ 的振幅受 $g$ 的相应振幅控制，由 Darboux 判别推出可积性；有限个点的赋值不影响可积性与积分。不能直接将积分表达式求导后声称处处 $F'=g$，因为 $g$ 未必处处连续。
\end{proof}

\subsection{积分定义可以怎样简化}
定义中的两种任意性不能同时随意删掉。对有界 $f$，以下两种检验仍足够：其一，只取等分分划，但要求\emph{所有标记点的积分和一致趋于同一极限}；其二，保留任意分划，每个小区间只选左、右端点之一，但要求对所有这些选择统一收敛。第一种检验可由等分分划的上下和证明；第二种检验中，把任意标记点所在的小区间在标记点处一分为二，再在两个新小区间都选该标记点，所得端点和恰好等于原标记和，且网格不增大，故回到原定义。若既只取等分分划又只取右端点，条件就不够：$[0,1]$ 上的有理数示性函数的所有等分右端点和恒为一，但函数不可积。等价地，可积性还可写成积分和的 Cauchy 条件，即任意两个充分细分划、任意两组标记点所给的积分和之差均任意小；其充分性利用实数完备性先取得一个极限，再与任意积分和比较。

\section{有关积分的几类典型问题}
\label{sec:pei-int-integral-problems}
% 教材对应：§4.3，PDF 第 33--82 页（原书第 282--331 页）。

\subsection{从已知性质建立积分上下界}
积分估计首先问“已知信息能给出什么逐点界”，再问是否需要换元、分部积分或分区间。若 $f$ 单调递增，等分步长为 $h=(b-a)/n$，则左、右端点和夹住积分，且二者之差恰为 $h(f(b)-f(a))$；因此指定精度时，可先选 $n$ 使这个差小于容许误差，再计算有限和。一般连续函数也可用 Darboux 上下和逼近，但仅知连续性时不能直接套用这一单调误差公式。

\begin{proposition}
\label{prop:pei-int-midpoint-bound}
若 $f\in C^2[a,b]$，$c=(a+b)/2$，$|f''|\le M$，则
\[
 \left|\int_a^bf(x)\,dx-(b-a)f(c)\right|\le\frac{M(b-a)^3}{24}.
\]
还存在 $\xi\in(a,b)$，使差恰等于 $f''(\xi)(b-a)^3/24$。
\end{proposition}
\begin{proof}
在 $c$ 处应用带积分余项的 Taylor 公式。线性项在对称区间上积分为零，二阶余项的绝对值不超过 $M(x-c)^2/2$，故误差不超过 $M(b-a)^3/24$。误差还可写成 $f''$ 乘以一个非负核的积分：左半区间的核为 $(t-a)^2/2$，右半区间的核为 $(b-t)^2/2$。核积分为 $(b-a)^3/24$，且在内部为正；由加权积分中值定理得到 $\xi$。
\end{proof}

\begin{example}
证明 $\displaystyle\int_0^{\sqrt{2\pi}}\sin(x^2)\,dx>0$。
\end{example}
\begin{solve}
令 $t=x^2$，并把 $[0,2\pi]$ 分成两个半周期，在第二段再令 $t=u+\pi$，得到
\[
 \int_0^{\sqrt{2\pi}}\sin(x^2)\,dx
 =\frac12\int_0^\pi\sin t\left(\frac1{\sqrt t}-\frac1{\sqrt{t+\pi}}\right)\,dt>0.
\]
零点附近的反常积分因 $t^{-1/2}\sin t=O(t^{1/2})$ 而收敛；括号在内部严格为正。换元把非线性相位变成等长半周期，配对后利用权函数递减性，便可判断正负而不求原函数。
\end{solve}
反复分部积分的作用是把导数从难估计的函数转到可控的权函数上。例如 $f,g\in C^{2m}[a,b]$，且 $f^{(j)}$、$g^{(j)}$ 在两端均为零（$0\le j<m$），则所有边界项消失，得到 $\int_a^bf^{(2m)}g=\int_a^bfg^{(2m)}$。选择 $g(x)=(x-a)^m(b-x)^m\ge0$，其最高阶导数为常数，便能把 $\int f$ 与最高阶导数的界联系起来。使用这一方法时应逐个检查边界项，不能仅因函数值为零就删去含高阶导数的边界项。

\subsection{导数约束、最佳常数与等号检查}
\begin{proposition}
\label{prop:pei-int-endpoint-slope}
若 $f\in C^1[a,b]$、$f(a)=f(b)=0$，则
\begin{equation}
 \int_a^b|f(x)|\,dx\le\frac{(b-a)^2}{4}\max_{[a,b]}|f'|.
 \label{eq:pei-int-slope-area}
\end{equation}
若 $f$ 非零，则在上述 $C^1$ 条件下不等号严格，但常数 $1/4$ 仍不能减小。
\end{proposition}
\begin{proof}
记 $D=\max|f'|$，由两端点分别应用中值定理，$|f(x)|\le D\min(x-a,b-x)$。积分后即得式~\eqref{eq:pei-int-slope-area}。若 $D>0$ 且取等号，由连续非负函数零积分的性质，必有 $|f(x)|=D\min(x-a,b-x)$；这要求中点处出现尖角，与 $C^1$ 矛盾。另一方面，把折线 $\min(x-a,b-x)$ 在中点的任意小邻域内光滑化，可使导数上界保持在一以内、面积任意接近 $(b-a)^2/4$，因此常数不能减小。
\end{proof}
\emph[caution]{最佳常数不意味着存在取等号的光滑函数}。例如若要求 $f\in C^1[0,2]$、$f(0)=f(2)=1$、$|f'|\le1$，则必有 $f(x)\ge|x-1|$，从而 $\int_0^2f\ge1$；若还要求积分恰为一，则连续性强制 $f=|x-1|$，与可微性矛盾。积分估计不仅给数值范围，也能排除函数的存在性。

\subsection{积分不等式中的单调性、对称性与求导}
若 $f$ 在 $[0,1]$ 上单调不增，则对 $0<c<1$，前段平均值不小于后段平均值，即
\[
 \frac1c\int_0^cf\ge\frac1{1-c}\int_c^1f,
 \qquad \int_0^cf\ge c\int_0^1f.
\]
可将 $f(x)\ge f(y)$ 在 $0\le x\le c\le y\le1$ 上先后积分得到，故不要求 $f$ 非负。若 $f$ 在 $[a,b]$ 单调递增，$c=(a+b)/2$，则
\[
 \int_a^b(x-c)f(x)\,dx
 =\int_0^{(b-a)/2}t\bigl(f(c+t)-f(c-t)\bigr)\,dt\ge0.
\]
对称点配对常能把两项的差变成逐点非负表达式。类似地，若 $f$ 在 $[-a,a]$ 上连续且为偶函数，则由 $x\mapsto-x$ 和 $1/(1+e^x)+1/(1+e^{-x})=1$ 得到 $\int_{-a}^a f(x)/(1+e^x)\,dx=\int_0^af(x)\,dx$。

\begin{example}
设 $f\in C^1[0,1]$、$f(0)=0$、$0<f'(x)<1$（$0<x<1$）。证明 $\displaystyle\left(\int_0^1f\right)^2>\int_0^1f^3$。
\end{example}
\begin{solve}
令 $A(x)=\int_0^xf(t)\,dt$，$H(x)=2A(x)-f(x)^2$。有 $H(0)=0$、$H'=2f(1-f')>0$，故 $H(x)>0$。再令 $F(x)=A(x)^2-\int_0^xf(t)^3\,dt$，则 $F(0)=0$、$F'=fH>0$，从而 $F(1)>0$。把常上限临时改成变量上限，再构造两端之差，是用微分学证明积分不等式的基本办法。
\end{solve}
从已知逐点不等式出发反复积分，也能生成更高阶的估计。例如由 $\cos x\le1$（$x\ge0$），依次积分得到 $\sin x\le x$、$1-\cos x\le x^2/2$，继而得到 $x-x^3/6\le\sin x\le x-x^3/6+x^5/120$。严格性须检查初始不等式是否在具有正长度的区间上严格成立。

\subsection{正交条件、零点与线性相关}
\begin{proposition}
\label{prop:pei-int-moment-zeros}
设非零连续函数 $f$ 在 $[a,b]$ 上满足 $\int_a^bx^kf(x)\,dx=0$（$k=0,\ldots,m-1$）。则 $f$ 在 $(a,b)$ 内至少有 $m$ 个不同零点。
\end{proposition}
\begin{proof}
若零点已经有无穷多个，则结论成立；否则假设内部零点少于 $m$，从中挑出真正改变符号的点 $c_1,\ldots,c_r$，$r<m$。多项式 $p(x)=\prod_{j=1}^r(x-c_j)$（$r=0$ 时取一）与 $f$ 在每个变号点同时变号，故 $pf$ 在整个区间上保持同号，且在某个非退化小区间上严格非零。因此 $\int pf\ne0$；但 $\deg p<m$，由矩条件又有 $\int pf=0$，矛盾。
\end{proof}
“选一个在零点处同步变号的乘子”是把多个积分等式转为零点数量的通用方法。例如 $\int_0^1f=\int_0^1xf=0$ 保证至少两个内部零点；仅分别对两个积分套一次中值定理，不能保证所得零点不同。

\begin{proposition}
\label{prop:pei-int-gram}
连续函数 $f_1,\ldots,f_m$ 在 $[a,b]$ 线性相关，当且仅当 Gram 矩阵 $G=(\int_a^bf_i f_j)_{i,j=1}^m$ 的行列式为零。
\end{proposition}
\begin{proof}
对 $c\in\mathbb R^m$，$c^{\mathsf T}Gc=\int_a^b(\sum_i c_if_i)^2\ge0$，所以 $G$ 半正定。若函数线性相关，相应非零系数向量属于 $G$ 的核。反过来，$\det G=0$ 时取非零核向量 $c$，其平方积分为零；被积函数连续非负，故 $\sum_i c_if_i\equiv0$。仅有可积性时，平方积分为零不能推出每一点上都为零，连续性在这里确有作用。
\end{proof}
积分还能把覆盖重数转为长度估计：若 $[0,1]$ 被 $n$ 个区间覆盖至少 $q$ 次，区间示性函数之和处处不小于 $q$，积分后得到区间总长不小于 $q$，故至少一个区间的长度不小于 $q/n$。这说明合适的示性函数可以把几何计数问题变成积分的可加性与保序性。

\subsection{凸性与积分平均}
\begin{theorem}
\label{thm:pei-int-hadamard-jensen}
若 $f$ 在 $[a,b]$ 上连续且为凸函数（convex function），则有 Hermite--Hadamard 不等式
\[
 f\left(\frac{a+b}{2}\right)\le\frac1{b-a}\int_a^bf(x)\,dx\le\frac{f(a)+f(b)}2.
\]
若 $w\ge0$ 可积、$W=\int_a^bw>0$，$u$ 可积且取值于 $[A,B]$，$\varphi$ 在 $[A,B]$ 连续且为凸函数，则加权 Jensen 不等式为
\begin{equation}
 \varphi\left(\frac1W\int_a^bwu\right)\le\frac1W\int_a^bw\varphi(u).
 \label{eq:pei-int-jensen}
\end{equation}
\end{theorem}
\begin{proof}
对称点 $x,a+b-x$ 的凸性不等式给出左侧，连接两端点的弦位于函数图像上方，积分给出右侧。Jensen 不等式令 $c=W^{-1}\int wu$；若 $A<c<B$，在 $c$ 处取一条支撑直线 $\ell(t)=\varphi(c)+s(t-c)\le\varphi(t)$，乘以 $w$ 后积分，线性项消失，便得结论。若 $c=A$，则 $\int w(u-A)=0$。对任意 $\varepsilon>0$，连续性给出 $\delta>0$，而有界性给出常数 $K$，使处处 $|\varphi(u)-\varphi(A)|\le\varepsilon+(K/\delta)(u-A)$；乘以 $w$ 积分后误差不超过 $\varepsilon W$，令 $\varepsilon\downarrow0$ 即得等号。$c=B$ 时同理，$A=B$ 时直接成立。若 $u,w$ 均连续，且 $w>0$、$\varphi$ 严格凸，等号当且仅当 $u$ 为常数。
\end{proof}
若 $f$ 在 $[0,\infty)$ 连续且凸，$A(x)=x^{-1}\int_0^xf(t)\,dt$（$x>0$），则换元给出 $A(x)=\int_0^1f(tx)\,dt$，直接对被积函数应用凸性就得到 $A$ 也凸；不必先求二阶导数。若 $g$ 单调递增，则 $F(x)=\int_c^xg(t)\,dt$ 为凸函数，因为相邻区间上的平均值随区间向右移动而不减，即割线斜率递增。这两个结论展示了“积分保持凸性”和“单调函数积分生成凸函数”的不同方向。

\subsection{变上限积分、函数方程与积分不等式}
若 $f$ 局部 Riemann 可积，$F(x)=\int_c^xf(t)\,dt$ 总是连续；在 $f$ 连续的点才可直接断言 $F'(x)=f(x)$。若 $f$ 连续、$u,v$ 可微，则
\[
 \frac{d}{dx}\int_{u(x)}^{v(x)}f(t)\,dt=f(v(x))v'(x)-f(u(x))u'(x).
\]
该公式先用一个固定基点的原函数拆成两项，再应用链式法则，负向积分亦自动包含其中。比如对局部可积的加性函数 $f(x+y)=f(x)+f(y)$，在 $y\in[0,1]$ 上积分得
\[
 f(x)=\int_x^{x+1}f(t)\,dt-\int_0^1f(t)\,dt,
\]
右端连续，从而 $f$ 连续。再由加性在有理点上确定 $f(r)=rf(1)$，用连续性延拓到实数，得到 $f(x)=xf(1)$。这里积分首先提供正则性，然后才使用连续加性函数的结论。

\begin{proposition}
\label{prop:pei-int-gronwall}
设 $u,v\ge0$ 在 $[a,b]$ 上连续，$c\ge0$，且 $u(x)\le c+\int_a^xv(t)u(t)\,dt$，则
\[
 u(x)\le c\exp\left(\int_a^xv(t)\,dt\right).
\]
\end{proposition}
\begin{proof}
记 $H(x)=c+\int_a^xvu$，则 $H'\le vH$，故 $\bigl(e^{-\int_a^xv}H(x)\bigr)'\le0$，从而 $H(x)\le ce^{\int_a^xv}$，再用 $u\le H$。当 $c=0$ 时相同论证给出 $u\equiv0$，无需除以 $u$。
\end{proof}
这就是 Gronwall 不等式的一个基本形式。教材中含未知函数积分的不等式，常可先令该积分为辅助函数，再乘积分因子；对正函数的微分不等式，也可积分其对数导数，例如 $f'/f\le-\alpha/x$、$f>0$ 推出 $f(x)\le f(a)(a/x)^\alpha$。非负函数可能有零点时，不应未经处理就除以 $f$。

\subsection{第一积分中值定理及其内点形式}
\begin{theorem}
\label{thm:pei-int-first-mean}
设 $f$ 在 $[a,b]$ 连续，$g\ge0$ 可积，则存在 $\xi\in(a,b)$，使
\[
 \int_a^bf(x)g(x)\,dx=f(\xi)\int_a^bg(x)\,dx.
\]
$g\le0$ 时同样成立；特别地 $g\equiv1$ 时，积分平均值能在区间内部取得。
\end{theorem}
\begin{proof}
若 $\int g=0$，由 $|\int fg|\le\max|f|\int g=0$，任意内点均可。否则记加权平均值为 $r$，上下界保证 $\min f\le r\le\max f$。若内部没有 $f=r$ 的点，连续性迫使 $f-r$ 在整个 $(a,b)$ 严格同号。又因 $g$ 的积分为正，可以截去足够短的两端后仍保留正积分，而 $|f-r|$ 在这个内闭区间有正下界，故 $\int(f-r)g\ne0$，矛盾。
\end{proof}
定理要求权函数不变号；若 $g$ 变号，积分平均值可能落在 $f$ 的值域之外。对于仅有界可积而无介值性的 $f$，通常只能得到 $\int fg=\theta\int g$、$\inf f\le\theta\le\sup f$，不能把 $\theta$ 写成某个函数值。强化为内点在多次使用 Rolle 定理时有实际意义：分别在互不相交的开子区间内取得中值点，才能确保它们不同，随后逐次计数导数的零点。

\subsection{第二积分中值定理与振荡估计}
\begin{theorem}
\label{thm:pei-int-second-mean}
设 $g$ 在 $[a,b]$ 可积。若 $f$ 非负单调不增，则存在 $\xi\in[a,b]$ 使 $\int_a^bfg=f(a)\int_a^\xi g$；若 $f$ 非负单调不减，则存在 $\xi$ 使 $\int_a^bfg=f(b)\int_\xi^bg$。一般地，只要 $f$ 单调，就有某个 $\xi\in[a,b]$ 满足 Bonnet 公式
\begin{equation}
 \int_a^bfg=f(a)\int_a^\xi g+f(b)\int_\xi^bg.
 \label{eq:pei-int-bonnet}
\end{equation}
\end{theorem}
\begin{proof}
先设 $f\ge0$ 单调不增，令 $G(x)=\int_a^xg$。在分划上用右端点阶梯函数逼近 $f$，其与 $f$ 的积分误差不超过 $\sup|g|\lVert P\rVert(f(a)-f(b))$。相应近似积分经 Abel 变换为
\[
 f(x_m)G(b)+\sum_{i=1}^{m-1}\bigl(f(x_i)-f(x_{i+1})\bigr)G(x_i).
\]
所有系数非负，总和为 $f(x_1)\le f(a)$；补上 $G(a)=0$ 的系数后，总和恰为 $f(a)$。故其值夹在 $f(a)\min G$ 与 $f(a)\max G$ 之间。取极限，利用 $G$ 连续及介值性即得第一种形式；递增的情形反射区间即可。一般递减的 $f$ 对 $f-f(b)\ge0$ 应用第一种形式，移项得到式~\eqref{eq:pei-int-bonnet}；递增的情形同理。
\end{proof}
例如 $f\ge0$ 递减时，$|\int_a^bf(x)\sin x\,dx|\le2f(a)$。若相位 $\varphi\in C^1[a,b]$，$\varphi'\ge m>0$ 且 $\varphi'$ 单调，令 $t=\varphi(x)$ 后，振幅 $1/\varphi'(\varphi^{-1}(t))$ 非负单调，故
\[
 \left|\int_a^be^{\mathrm i\varphi(x)}\,dx\right|\le\frac2m.
\]
这里复指数只是同时记述正弦和余弦；同一 Abel 变换证明对复数积分给出上述模长界，因为 $|\int_\alpha^\beta e^{\mathrm it}\,dt|\le2$。若只讨论实数，也可分别得到正弦和余弦积分的同样界。第一中值定理控制不变号的权，第二中值定理利用单调振幅与原函数的有界性，两者的假设和用途应分别核对。

\section{几个著名的不等式}
\label{sec:pei-int-classical-inequalities}
% 教材对应：§4.4，PDF 第 82--102 页（原书第 331--351 页）。

\subsection{Cauchy--Schwarz 不等式与等号条件}
\begin{theorem}
\label{thm:pei-int-schwarz}
对实数列 $(a_i)$、$(b_i)$，有 $|\sum_i a_ib_i|^2\le(\sum_i a_i^2)(\sum_i b_i^2)$，等号当且仅当两个向量线性相关。若 $f,g$ 在 $[a,b]$ Riemann 可积，则积分形式为
\begin{equation}
 \left|\int_a^bfg\right|^2\le\left(\int_a^bf^2\right)\left(\int_a^bg^2\right).
 \label{eq:pei-int-schwarz}
\end{equation}
一般可积函数取等号的充要条件是存在不全为零的 $\alpha,\beta$ 使 $\int_a^b(\alpha f+\beta g)^2=0$；若二者连续，则等价于 $\alpha f+\beta g\equiv0$。
\end{theorem}
\begin{proof}
二次式 $\int(f+tg)^2=A+2Bt+Ct^2$ 对每个实数 $t$ 非负。当 $C>0$ 时，其最小值 $A-B^2/C\ge0$，即得不等式，且等号恰对应某个非零线性组合的平方积分为零。若 $C=0$，对所有 $t$ 非负迫使 $B=0$，不等式也成立。离散形式把积分换成有限求和，证明相同；连续情形的等号条件由定理~\ref{thm:pei-int-zero-integral} 得到。
\end{proof}
命题~\ref{prop:pei-int-gram} 给出了多函数推广：Gram 矩阵半正定，其行列式非负，连续函数情形下行列式为零恰对应线性相关。应用 Schwarz 不等式时，选择分拆常比公式本身更关键：把 $f=\sqrt w\,(f/\sqrt w)$ 可得 $|\int f|^2\le(\int w)(\int f^2/w)$，其中 $w$ 应严格为正，所有积分须存在。

\begin{example}
设 $g\in C^1[a,b]$、$g(a)=0$，证明
\[
 \int_a^b|gg'|\le\frac{b-a}{2}\int_a^b(g')^2,
\]
并确定等号条件。若还要求 $g(b)=0$，如何改进常数？
\end{example}
\begin{solve}
令 $H(x)=\int_a^x|g'(t)|\,dt$，则 $|g|\le H$，所以
\[
 \int_a^b|gg'|\le\int_a^bHH'=\frac12\left(\int_a^b|g'|\right)^2
 \le\frac{b-a}{2}\int_a^b(g')^2.
\]
取等号要求 $|g'|$ 为常数；$g'$ 连续，非零时不能改变符号，故 $g(x)=c(x-a)$，反过来这些函数均取等号。若两端为零，在中点处分开，分别从左右端应用刚证的单端点结论，得到 $\int|gg'|\le(b-a)\int(g')^2/4$。非零 $C^1$ 函数不能取等号，但将两段直线构成的帐篷函数在顶点附近光滑化，可说明常数 $(b-a)/4$ 不能减小。
\end{solve}
另一个常用推论是 $g(a)=0$ 时 $\int_a^bg^2\le(b-a)^2\int_a^b(g')^2/2$：由 $g(x)=\int_a^xg'$ 得到 $g(x)^2\le(x-a)\int_a^b(g')^2$，再积分即可。它提供一个方便的估计，此处并不声称常数最佳；重复应用可把低阶导数的积分控制在高阶导数的积分之下，但每一步都必须有相应的端点消失条件。

\subsection{加权平均与幂平均的积分形式}
设 $w>0$、$f>0$ 在 $[a,b]$ 上连续，定义归一化平均 $\mathcal A_w(u)=(\int wu)/(\int w)$，以及\emph{幂平均（power mean）}与\emph{几何平均（geometric mean）}
\[
 M_r(f;w)=\bigl(\mathcal A_w(f^r)\bigr)^{1/r}\quad(r\ne0),
 \qquad M_0(f;w)=\exp\bigl(\mathcal A_w(\ln f)\bigr).
\]
归一化保证 $M_r(c;w)=c$。在紧区间上 $f$ 有正下界，因此负次幂与对数均可积。

\begin{theorem}
\label{thm:pei-int-power-means}
对任意实数 $r<s$，$M_r(f;w)\le M_s(f;w)$；若 $f$ 非常数，则不等号严格。并且 $\lim_{r\to0}M_r(f;w)=M_0(f;w)$，$\lim_{r\to\infty}M_r(f;w)=\max f$。
\end{theorem}
\begin{proof}
当 $0<r<s$ 时，对 $f^r$ 和凸函数 $t^{s/r}$ 应用式~\eqref{eq:pei-int-jensen}，得到 $\mathcal A_w(f^s)\ge\mathcal A_w(f^r)^{s/r}$。当 $r<s<0$ 时，对 $1/f$ 应用刚证的正指数结论并取倒数。当指数跨过零时，先由凸函数 $e^{rt}$ 的 Jensen 不等式得到 $M_r\ge M_0$（$r>0$），并对 $1/f$ 得到 $M_r\le M_0$（$r<0$）。严格性来自相应函数的严格凸性与 $w>0$。

由于 $\ln f$ 有界，$f^r=1+r\ln f+O(r^2)$ 的余项一致，故 $r^{-1}\ln\mathcal A_w(f^r)\to\mathcal A_w(\ln f)$，得到零点处的极限。无穷大处的极限由定理~\ref{thm:pei-int-power-integral} 得到，归一化常数的 $r$ 次方根趋于一。
\end{proof}
特别地，$M_{-1}\le M_0\le M_1\le M_2$ 分别对应调和、几何、算术与平方平均的次序。离散加权平均把 $\mathcal A_w(u)$ 换成 $\sum_i p_i u_i$（$p_i>0$、$\sum_i p_i=1$），同一 Jensen 证明仍适用；非归一化权重须先除以总权重。若只讨论正指数，允许 $f$ 为零，但对数、负指数及等号条件需另外处理。

\subsection{Hölder 不等式及反向形式}
\begin{theorem}
\label{thm:pei-int-holder}
设 $p,q>1$、$1/p+1/q=1$，$f,g$ 在 $[a,b]$ 可积，则
\begin{equation}
 \int_a^b|fg|\le\left(\int_a^b|f|^p\right)^{1/p}
                     \left(\int_a^b|g|^q\right)^{1/q}.
 \label{eq:pei-int-holder}
\end{equation}
若 $f,g$ 连续且两个幂积分均非零，等号当且仅当 $|f|^p=C|g|^q$，其中 $C>0$。若 $0<p<1$、$q=p/(p-1)<0$，且 $f\ge0$、$g>0$ 连续，则反向形式为
\[
 \int_a^bfg\ge\left(\int_a^bf^p\right)^{1/p}\left(\int_a^bg^q\right)^{1/q}.
\]
在 $f$ 不恒为零时，等号仍以 $f^p$ 与 $g^q$ 成比例刻画。
\end{theorem}
\begin{proof}
先证普通形式。令 $A=(\int|f|^p)^{1/p}$、$B=(\int|g|^q)^{1/q}$；$AB=0$ 时由下段的直接估计处理。否则对 $u=|f|/A$、$v=|g|/B$ 使用数值 Young 不等式 $uv\le u^p/p+v^q/q$，积分得 $\int uv\le1$。连续情形等号要求逐点 $u^p=v^q$。反向形式则在普通 Hölder 不等式中使用共轭指数 $1/p$、$1/(1-p)$：
\[
 \int f^p=\int(fg)^p(g^q)^{1-p}
 \le(\int fg)^p(\int g^q)^{1-p},
\]
移项即得结论。由于 $q<0$，$g$ 的正下界是保证各项有意义的一个充分条件。
\end{proof}
证明中若 $AB=0$，也可完全用 Riemann 积分估计处理：例如 $\int|f|^p=0$，对任意 $\delta>0$，在 $|f|\le\delta$ 时有 $|f|\le\delta$，其余情形有 $|f|\le\delta^{1-p}|f|^p$，故处处 $|f|\le\delta+\delta^{1-p}|f|^p$。积分后令 $\delta\downarrow0$，得到 $\int|f|=0$，再用 $g$ 有界得到 $\int|fg|=0$。把积分换成求和可得离散形式；在积分中改用归一化权重，则得到加权形式。$p=q=2$ 时回到 Schwarz 不等式。

\subsection{Minkowski 不等式与参数范围}
\begin{theorem}
\label{thm:pei-int-minkowski}
对 $p\ge1$，记 $\lVert f\rVert_p=(\int_a^b|f|^p)^{1/p}$，则
\begin{equation}
 \lVert f+g\rVert_p\le\lVert f\rVert_p+\lVert g\rVert_p.
 \label{eq:pei-int-minkowski}
\end{equation}
若 $f,g\ge0$ 连续、$p>1$ 且均不恒为零，等号当且仅当二者成比例。若 $0<r<1$、$f,g\ge0$ 连续，则
\[
 \left(\int_a^b(f+g)^r\right)^{1/r}
 \ge\left(\int_a^bf^r\right)^{1/r}+\left(\int_a^bg^r\right)^{1/r}.
\]
$r<0$ 时，在 $f,g>0$ 连续的条件下也有后一反向不等式。
\end{theorem}
\begin{proof}
$p=1$ 由逐点三角不等式得到。$p>1$ 时，令 $q=p/(p-1)$，利用 $|f+g|\le|f|+|g|$ 与 Hölder 不等式，
\[
 \int|f+g|^p\le\int(|f|+|g|)|f+g|^{p-1}
 \le(\lVert f\rVert_p+\lVert g\rVert_p)\lVert f+g\rVert_p^{p-1}.
\]
总范数为零时直接成立，否则约去相应幂次。非负连续情形由 Hölder 的等号条件得到比例关系。

对反向形式，设 $A=(\int f^r)^{1/r}$、$B=(\int g^r)^{1/r}$ 均为正，令 $\lambda=A/(A+B)$、$u=f/A$、$v=g/B$。当 $0<r<1$ 时，$t^r$ 凹，故 $\int(\lambda u+(1-\lambda)v)^r\ge1$；当 $r<0$ 时，$t^r$ 凸，积分不等号相反，但取负指数 $1/r$ 又反向。两种情况都得到 $\lVert f+g\rVert_r\ge A+B$。$0<r<1$ 时某函数恒为零的情形直接成立。
\end{proof}
\emph[caution]{在 $p=1$ 时，非负函数总有 $\lVert f+g\rVert_1=\lVert f\rVert_1+\lVert g\rVert_1$，无需成比例}；只有 $p>1$ 才能这样使用严格凸性给出的等号条件。对一般 Riemann 可积函数，$\lVert f\rVert_p=0$ 不保证逐点为零，因此 $p\ge1$ 时该表达式在原函数空间上是半范数（seminorm）；限制为连续函数，或将满足 $\int|f-g|^p=0$ 的函数视为同一个对象后，才得到范数。$0<r<1$ 时的表达式不满足通常的三角不等式。多函数形式可由二函数形式归纳得到，离散形式同样由求和替代积分。

\subsection{Young 不等式与反函数的面积恒等式}
\begin{theorem}
\label{thm:pei-int-young}
设 $\phi:[0,\infty)\to[0,\infty)$ 连续、严格递增、满射，且 $\phi(0)=0$。则对 $a,b\ge0$，
\begin{equation}
 ab\le\int_0^a\phi(x)\,dx+\int_0^b\phi^{-1}(y)\,dy,
 \label{eq:pei-int-young-inverse}
\end{equation}
等号当且仅当 $b=\phi(a)$。取 $\phi(x)=x^{p-1}$（$p>1$）即得 $ab\le a^p/p+b^q/q$，其中 $q=p/(p-1)$，等号当且仅当 $a^p=b^q$。
\end{theorem}
\begin{proof}
先用分划 $0=x_0<\cdots<x_n=a$ 与像分划 $\phi(x_i)$ 的积分和，或对矩形作曲边区域分割，得到
\[
 \int_0^a\phi(x)\,dx+\int_0^{\phi(a)}\phi^{-1}(y)\,dy=a\phi(a).
\]
固定 $a$ 后，式~\eqref{eq:pei-int-young-inverse} 右端减去左端等于 $\int_{\phi(a)}^b(\phi^{-1}(y)-a)\,dy$。若 $b>\phi(a)$，被积函数在内部为正；若 $b<\phi(a)$，积分方向与被积函数的符号同时反向；故均为正，只有 $b=\phi(a)$ 时为零。这一证明不要求 $\phi$ 可微。
\end{proof}

\subsection{Chebyshev 不等式：同序与反序}
\begin{proposition}
\label{prop:pei-int-chebyshev}
若可积函数 $f,g$ 同序，即 $(f(x)-f(y))(g(x)-g(y))\ge0$ 对所有 $x,y\in[a,b]$ 成立，$w\ge0$ 可积且 $W=\int w>0$，则
\[
 \left(\int wf\right)\left(\int wg\right)\le W\int wfg.
\]
反序时不等号相反。特别地，同为递增或同为递减的函数同序，取 $w\equiv1$ 得无权形式。
\end{proposition}
\begin{proof}
展开并分别积分可得
\[
 W\int wfg-\int wf\int wg
 =\frac12\int_a^b\int_a^b w(x)w(y)(f(x)-f(y))(g(x)-g(y))\,dx\,dy.
\]
由同序条件，右端非负。离散形式的差同样等于 $\tfrac12\sum_{i,j}p_ip_j(a_i-a_j)(b_i-b_j)$，因此加权推广与无权形式具有同一证明。
\end{proof}

\section{反常积分}
\label{sec:pei-int-improper-integrals}
% 教材对应：§4.5，PDF 第 102--136 页（原书第 351--385 页）。

\subsection{先识别反常点，再作计算}
\begin{definition}
若 $f$ 在每个 $[a,R]$ 上 Riemann 可积，且有限极限 $\lim_{R\to\infty}\int_a^Rf$ 存在，称\emph{反常积分（improper integral）}$\int_a^\infty f$ 收敛。若 $f$ 在每个 $[c,b]$ 上可积，则 $\int_a^bf=\lim_{c\downarrow a}\int_c^bf$ 定义端点无界情形。存在多个反常点时，必须分开每个反常点并分别取极限；各部分都收敛，才称整个反常积分收敛。若 $\int|f|<\infty$，称绝对收敛；收敛而不绝对收敛，称条件收敛。
\end{definition}
\emph[caution]{对称截断的 Cauchy 主值（Cauchy principal value）不等于通常的反常积分}。例如 $\int_{-1}^1 dx/x$ 在零点两侧均发散，虽然对称删去 $(-\varepsilon,\varepsilon)$ 后的积分恒为零，其主值为零。变量替换、分部积分和原函数计算应先在没有反常点的截断区间上进行，再分别检查边界极限；不能依靠两个发散边界项相减得到积分值。

\begin{example}
计算 $\displaystyle I=\int_0^\infty\frac{dx}{1+x^4}$。
\end{example}
\begin{solve}
被积函数在零点附近有界，在无穷远为 $O(x^{-4})$，故绝对收敛。换元 $x=1/t$ 得 $I=\int_0^\infty x^2/(1+x^4)\,dx$，相加后再令 $u=x-1/x$，得到
\[
 2I=\int_0^\infty\frac{1+x^2}{1+x^4}\,dx
 =\int_{-\infty}^\infty\frac{du}{u^2+2}=\frac\pi{\sqrt2},
 \qquad I=\frac\pi{2\sqrt2}.
\]
倒数变换先制造互补的分子，再用双射换元消去四次式；各步都可由有限截断后取极限验证。
\end{solve}

\begin{proposition}
设 $a,b>0$，$F:[2\sqrt{ab},\infty)\to\mathbb R$ 连续。若 $F\ge0$（允许积分值为 $+\infty$），或所讨论的积分绝对收敛，则
\[
 \int_0^\infty F\left(ax+\frac bx\right)\,dx
 =\frac1a\int_0^\infty F\left(\sqrt{t^2+4ab}\right)\,dt.
\]
因此，以已知 Gaussian 积分 $\int_0^\infty e^{-t^2}\,dt=\sqrt\pi/2$ 为前提，有
\[
 \int_0^\infty\exp\left(-a^2x^2-\frac{b^2}{x^2}\right)\,dx
 =\frac{\sqrt\pi}{2a}e^{-2ab}.
\]
\end{proposition}
\begin{proof}
令 $t=ax-b/x$，它把 $(0,\infty)$ 严格递增地映到 $\mathbb R$，且
\[
 ax+\frac bx=\sqrt{t^2+4ab},\qquad
 \frac{dx}{dt}=\frac1{2a}\left(1+\frac t{\sqrt{t^2+4ab}}\right).
\]
将 $t>0$ 与 $t<0$ 两部分配对，奇函数项相消，便得第一式；对非负 $F$ 可先配对有限区间再取单调极限，绝对收敛时也允许分项。取 $F(u)=e^{-u^2}$，再利用 $(ax+b/x)^2=a^2x^2+b^2/x^2+2ab$，得到第二式。
\end{proof}

\subsection{递推、对称与自我消去}
\begin{proposition}
对 $\alpha>-1$、整数 $m\ge0$，有
\[
 \int_0^1x^\alpha(\ln x)^m\,dx=\frac{(-1)^m m!}{(\alpha+1)^{m+1}}.
\]
对非负整数 $m,n$，有 $\int_0^1x^m(1-x)^n\,dx=m!n!/(m+n+1)!$。
\end{proposition}
\begin{proof}
第一式在 $[\varepsilon,1]$ 上分部积分，边界项因 $\varepsilon^{\alpha+1}|\ln\varepsilon|^m\to0$ 而消失，得到 $I_m=-mI_{m-1}/(\alpha+1)$，从 $I_0=1/(\alpha+1)$ 递推。第二式同样分部积分，得到 $B_{m,n}=nB_{m+1,n-1}/(m+1)$，递推到 $B_{m+n,0}=1/(m+n+1)$ 即得。
\end{proof}
Wallis 积分的初值、降阶递推和双阶乘表达式见式~\eqref{eq:pei-wallis-integral}，混合正弦、余弦幂的公式见式~\eqref{eq:pei-mixed-wallis}。这些公式与上面的两个递推都依靠分部积分；遇到奇点时还须额外验证边界项的极限。

\begin{example}
计算 $\displaystyle J=\int_0^{\pi/2}\ln(\sin x)\,dx$，以及 $\displaystyle K=\int_0^\infty\frac{\ln x}{1+x^2}\,dx$。
\end{example}
\begin{solve}
因 $\sin x\sim x$，第一积分在零点附近与 $\ln x$ 相差有界量，故绝对收敛。由 $x\mapsto\pi/2-x$，$J=\int_0^{\pi/2}\ln(\cos x)\,dx$，所以
\[
 2J=\int_0^{\pi/2}\ln(\sin2x)\,dx-\frac\pi2\ln2
 =J-\frac\pi2\ln2,
 \qquad J=-\frac\pi2\ln2.
\]
第二积分在零点由 $|\ln x|$ 控制，在无穷远由 $(\ln x)/x^2$ 控制，故绝对收敛；换元 $x=1/t$ 得 $K=-K$，所以 $K=0$。这里先证收敛，才可将积分分拆、换元并移项，避免以形式对称掩盖发散。
\end{solve}

\subsection{比较判别与临界阶数}
\begin{theorem}
\label{thm:pei-int-improper-cauchy}
局部可积函数的反常积分 $\int_a^\infty f$ 收敛，当且仅当对每个 $\varepsilon>0$，存在 $A>a$，使对所有 $v>u>A$ 均有 $|\int_u^vf|<\varepsilon$。若 $f\ge0$，收敛又等价于截断积分 $\int_a^Rf$ 有上界。
\end{theorem}
\begin{proof}
第一条是实数极限存在的 Cauchy 准则用于函数 $R\mapsto\int_a^Rf$；第二条利用该函数随 $R$ 单调不减。有限反常端点的版本只需把远端尾区间改成趋向该端点的小区间。
\end{proof}
若尾部 $0\le f\le g$，则 $\int g$ 收敛推出 $\int f$ 收敛，$\int f$ 发散推出 $\int g$ 发散。若 $f/g\to c\in(0,\infty)$，二者同敛散；极限为零时只保留“小者可由大者收敛控制”的方向，极限为无穷时相反。反常点附近的等价无穷小可用于非负函数的敛散性判断，不能据此直接宣称两个收敛积分的值相同。

基本标尺是 $\int_1^\infty x^{-p}\,dx$ 当且仅当 $p>1$ 收敛，而 $\int_0^1x^{-p}\,dx$ 当且仅当 $p<1$ 收敛。临界幂次必须继续看对数：
\[
 \int_e^\infty\frac{dx}{x(\ln x)^q}<\infty\iff q>1,
 \qquad
 \int_0^{e^{-1}}\frac{dx}{x|\ln x|^q}<\infty\iff q>1.
\]
两式都由 $t=|\ln x|$ 化为无穷远的幂积分。若 $f>0$ 且 $\ln f(x)/\ln x\to-\alpha$，则 $\alpha>1$ 保证无穷远积分收敛，$\alpha<1$ 保证发散，$\alpha=1$ 不定；因为对任意小的 $\varepsilon>0$，最终有 $x^{-\alpha-\varepsilon}\le f(x)\le x^{-\alpha+\varepsilon}$。例如 $1/x$ 与 $1/(x\ln^2x)$ 都具有临界指数一，却有不同敛散性。

\begin{example}
讨论参数积分 $\displaystyle I_{\alpha,k}=\int_0^\pi\frac{(\sin x)^{\alpha-1}}{|1+k\cos x|^\alpha}\,dx$ 的收敛性，其中 $\alpha,k\in\mathbb R$。
\end{example}
\begin{solve}
若 $|k|<1$，分母处处远离零，两个端点均为 $x^{\alpha-1}$ 型，故当且仅当 $\alpha>0$ 收敛。若 $k=1$，零点附近要求 $\alpha>0$；但 $x\uparrow\pi$ 时 $1+\cos x\sim(\pi-x)^2/2$，被积函数为常数倍的 $(\pi-x)^{-\alpha-1}$，要求 $\alpha<0$，两者矛盾；$k=-1$ 由反射同样发散。若 $|k|>1$，存在唯一内部零点 $c$，且 $1+k\cos x$ 在 $c$ 处为一阶零点，故内部要求 $\alpha<1$，两端要求 $\alpha>0$。综上，收敛当且仅当 $|k|<1,\alpha>0$，或 $|k|>1,0<\alpha<1$。在各反常点分别比较，才能得到全部参数条件。
\end{solve}

\subsection{Dirichlet、Abel 判别与振荡抵消}
\begin{theorem}
\label{thm:pei-int-dirichlet-abel}
设 $f,g$ 在每个有限子区间上可积。Dirichlet 判别法：若 $F(R)=\int_a^Rf$ 有界，$g$ 单调且趋于零，则 $\int_a^\infty fg$ 收敛。Abel 判别法：若 $\int_a^\infty f$ 收敛，$g$ 单调且有界，则 $\int_a^\infty fg$ 收敛。
\end{theorem}
\begin{proof}
Dirichlet 情形可在必要时把 $g$ 换成 $-g$，设 $g\ge0$ 递减。若 $|F|\le C$，由第二积分中值定理，对 $B>A$ 有 $|\int_A^Bfg|\le2Cg(A)\to0$，从而满足 Cauchy 准则。Abel 情形中，$\int f$ 收敛保证所有足够远的 $|\int_A^tf|$ 统一小于 $\varepsilon$；对单调 $g$ 使用式~\eqref{eq:pei-int-bonnet}，得 $|\int_A^Bfg|\le2\sup|g|\varepsilon$，故也收敛。
\end{proof}
两种判别法利用振荡抵消，结论只保证收敛；要区分绝对与条件收敛，仍须另查绝对值积分。对周期振荡，$|\sin x|$ 与 $\sin^2x$ 每个周期都具有固定的正积分，故其乘以单调幂权重的积分可与相应幂积分比较；不能因 $\sin x$ 变号就把 $|\sin x|$ 的平均贡献也忽略。

\begin{example}
\label{ex:pei-int-oscillatory-power}
对实数 $\alpha,\beta$，分类 $\displaystyle\int_0^\infty x^\alpha\sin(x^\beta)\,dx$ 的敛散性。
\end{example}
\begin{solve}
若 $\beta=0$，被积函数为 $x^\alpha\sin1$，两个端点不能同时收敛。若 $\beta\ne0$，令 $t=x^\beta$，注意 $\beta<0$ 时积分方向反转，原积分化为
\[
 \frac1{|\beta|}\int_0^\infty t^{r-1}\sin t\,dt,
 \qquad r=\frac{\alpha+1}{\beta}.
\]
零点处与 $t^r$ 等价，要求 $r>-1$；无穷远处当 $r<1$ 时由 Dirichlet 判别收敛，当 $r\ge1$ 时取同号半周期，其积分绝对值不趋于零，违背 Cauchy 准则。绝对值积分在无穷远处收敛当且仅当 $r<0$：可逐周期用权重的上下界比较 $\int|\sin t|$ 与级数 $\sum n^{r-1}$。因此 $-1<r<0$ 时绝对收敛，$0\le r<1$ 时条件收敛，其余情况发散。
\end{solve}

\begin{example}
对 $p>0$，讨论 $\displaystyle\int_2^\infty\frac{\sin x}{x^p+\sin x}\,dx$。
\end{example}
\begin{solve}
分母在该区间严格为正，利用恒等式
\[
 \frac{\sin x}{x^p+\sin x}
 =\frac{\sin x}{x^p}-\frac{\sin^2x}{x^p(x^p+\sin x)}.
\]
第一项的积分对所有 $p>0$ 均由 Dirichlet 判别收敛；第二项非负，且与 $x^{-2p}\sin^2x$ 有固定比例的上下界，因此当且仅当 $2p>1$ 可积。原积分故当且仅当 $p>1/2$ 收敛。其绝对值与 $x^{-p}|\sin x|$ 可比较，所以当且仅当 $p>1$ 绝对收敛，$1/2<p\le1$ 时条件收敛。仅取分母的主项会漏掉不振荡的二阶贡献，正是本题的关键。
\end{solve}

\subsection{积分收敛与函数趋于零的关系}
函数趋于零不保证反常积分收敛，例如 $1/(1+x)$；反过来，收敛也不保证函数趋于零，甚至同时要求函数非负、连续也不够。具体地，在整数 $n\ge2$ 周围放置三角尖峰 $\psi_n(x)=\max(1-2^n|x-n|,0)$，令 $f=\sum_{n=2}^\infty\psi_n$。支集互不相交，和在每个有限区间上只有有限项，故 $f$ 连续非负；而 $\int_0^\infty f=\sum_{n=2}^\infty2^{-n}<\infty$，但 $f(n)=1$。这里尖峰越来越窄，说明必须补充控制局部振荡的条件。

\begin{theorem}
\label{thm:pei-int-integral-decay}
若 $f$ 在 $[a,\infty)$ 一致连续，且 $\int_a^\infty f$ 收敛，则 $f(x)\to0$。若 $f\ge0$ 单调不增且积分收敛，则更有 $xf(x)\to0$。
\end{theorem}
\begin{proof}
若第一条不成立，存在 $\varepsilon_0>0$ 及 $x_n\to\infty$ 使 $|f(x_n)|\ge\varepsilon_0$。一致连续性给出与 $n$ 无关的 $\delta>0$，使在 $[x_n,x_n+\delta]$ 上 $f$ 与 $f(x_n)$ 同号且绝对值至少为 $\varepsilon_0/2$，故该段积分的绝对值至少为 $\delta\varepsilon_0/2$，与积分的 Cauchy 准则矛盾。第二条由单调性得到
\[
 0\le\frac x2f(x)\le\int_{x/2}^xf(t)\,dt\longrightarrow0.
 \qedhere
\]
\end{proof}
导数有界保证一致连续，因而是第一条的一个充分条件。若 $f\in C^1[a,\infty)$，$\int f$、$\int f'$ 都收敛，则 $f(x)=f(a)+\int_a^xf'$ 先给出 $f$ 的有限极限；该极限若非零，$\int f$ 必发散，所以极限只能为零。若进一步已知 $xf(x)\ge0$ 单调不增，且 $\int f$ 收敛，则 $\int_{\sqrt x}^xf(t)\,dt\ge\tfrac12xf(x)\ln x$，故还有 $xf(x)\ln x\to0$。各个衰减结论依赖的单调对象不同，不宜混用。

\subsection{反常积分的极限：截断与尾部控制}
在无穷区间上，局部一致收敛仍不足以交换积分与极限：质量可能向无穷远移动。例如在 $[n,n+1]$ 上放一个面积为一的连续三角峰，其函数列在每个固定紧区间最终恒为零，但积分始终为一。一个直接可用的充分条件是：$f_n\to f$ 在每个 $[a,A]$ 上一致成立，且存在非负可积函数 $H$ 使 $|f_n|,|f|\le H$；将误差积分拆成 $[a,A]$ 和尾部，先使 $2\int_A^\infty H$ 小，再固定 $A$ 令 $n$ 大，即可得到 $\int f_n\to\int f$。这只是用截断估计证明的充分条件，不能省去尾部的一致控制。

\begin{proposition}
\label{prop:pei-int-abel-average}
若 $f$ 局部 Riemann 可积且 $f(x)\to L$（$x\to\infty$），则
\[
 \lim_{t\downarrow0}t\int_0^\infty e^{-tx}f(x)\,dx=L.
\]
若 $f,\phi$ 局部可积，$\phi>0$ 最终成立、$\int_a^x\phi\to\infty$，且 $f(x)/\phi(x)\to L$，则
\[
 \frac{\int_a^xf(t)\,dt}{\int_a^x\phi(t)\,dt}\longrightarrow L.
\]
\end{proposition}
\begin{proof}
第一式由 $t\int_0^\infty e^{-tx}\,dx=1$ 减去 $L$。选 $A$ 使 $x\ge A$ 时 $|f-L|<\varepsilon$，则尾部误差不超过 $\varepsilon$；固定段误差不超过 $t\int_0^A|f-L|\to0$。第二式对 $f-L\phi$ 同样分段，尾部绝对值不超过 $\varepsilon\int_A^x\phi$，固定段除以发散的分母趋于零。该证明只要求局部可积性，无需使用可能不适用的 L'Hospital 法则。
\end{proof}
取 $\phi\equiv1$ 得到积分的 Cesàro 平均 $x^{-1}\int_0^xf\to L$。对 $f\in C[0,1]$，类似的集中估计还给出
\[
 \lim_{x\downarrow0}x^\alpha\int_x^1\frac{f(t)}{t^{\alpha+1}}\,dt=\frac{f(0)}\alpha\quad(\alpha>0),
 \qquad
 \lim_{x\downarrow0}\frac{\int_x^1\frac{f(t)}t\,dt}{\ln(1/x)}=f(0).
\]
第一式减去 $f(0)x^\alpha\int_x^1t^{-\alpha-1}\,dt$，在 $[x,\delta]$ 上用连续性、在 $[\delta,1]$ 上用 $x^\alpha\to0$；第二式则把分母看成同一正权 $1/t$ 的积分。端点处的函数值决定主项，而区间远端只贡献低阶项。

\begin{proposition}
\label{prop:pei-int-gaussian-kernel}
设 $f\in C[a,b]$，$a<c<b$。则
\[
 \lim_{n\to\infty}\sqrt{\frac n\pi}\int_a^be^{-n(x-c)^2}f(x)\,dx=f(c).
\]
若 $c=a$ 或 $c=b$，极限为 $f(c)/2$。若 $f$ 在 $[a,\infty)$ 局部可积且 $\int_a^\infty|f|<\infty$，$g$ 为周期 $T$ 的 Riemann 可积函数，则
\[
 \lim_{\lambda\to\infty}\int_a^\infty f(x)g(\lambda x)\,dx
 =\frac1T\int_0^Tg\,\int_a^\infty f.
\]
\end{proposition}
\begin{proof}
Gaussian 核令 $u=\sqrt n(x-c)$ 后，积分范围向整条实轴扩张；在 $c$ 的小邻域内用 $f$ 的连续性，外部由 Gaussian 尾部控制，总质量趋于一。集中点在端点时仅保留一侧，质量为一半。第二式先在 $[a,A]$ 上应用定理~\ref{thm:pei-int-periodic-average}，尾部误差由 $2\sup|g|\int_A^\infty|f|$ 控制，先固定足够大的 $A$ 再令 $\lambda\to\infty$ 即得。
\end{proof}
例如周期函数 $\psi$ 的平均值为 $C$ 时，换元 $t=nu$ 得
\[
 n\int_n^\infty\frac{\psi(t)}{t^2}\,dt
 =\int_1^\infty\frac{\psi(nu)}{u^2}\,du\longrightarrow C.
\]
绝对可积的衰减权把无穷区间的周期振荡问题化回有限区间的平均定理。

\subsection{反常积分何时仍可用和式逼近}
\begin{proposition}
\label{prop:pei-int-singular-sums}
设 $f\ge0$ 在 $(0,1]$ 单调不增，且 $\int_0^1f<\infty$，则
\[
 \frac1n\sum_{k=1}^nf(k/n)\longrightarrow\int_0^1f(x)\,dx.
\]
对每个固定 $\theta\in(0,1)$，$n^{-1}\sum_{k=0}^{n-1}f((k+\theta)/n)$ 也有同一极限。
\end{proposition}
\begin{proof}
由单调性，$\int_{1/n}^1f\le n^{-1}\sum_{k=1}^nf(k/n)\le\int_0^1f$，两端趋于同一有限值。移位和不小于右端点和，且至多比它多 $f(\theta/n)/n$；又因 $f$ 递减，$0\le f(\theta/n)/n\le\theta^{-1}\int_0^{\theta/n}f\to0$，故移位和也收敛。
\end{proof}
若 $f$ 只在奇点附近单调，而远离奇点后正常可积，可分段应用上述结论与普通积分和极限。\emph[caution]{这不是任意反常可积函数都具有的 Riemann 和定义}；无界函数在极靠近奇点处可取任意大的标记值，任意标记和通常不能统一收敛。条件中有限积分也不可略去；和式有某种广义极限并不能反推通常的反常积分收敛。

\begin{example}
设 $0<a<d$，$A_n$、$G_n$ 分别为 $a,a+d,\ldots,a+(n-1)d$ 的算术平均与几何平均，求 $\lim G_n/A_n$。
\end{example}
\begin{solve}
记 $\theta=a/d\in(0,1)$。有 $A_n/(nd)=\theta/n+(n-1)/(2n)\to1/2$，且
\[
 \ln\frac{G_n}{nd}=\frac1n\sum_{k=0}^{n-1}\ln\frac{k+\theta}{n}
 \longrightarrow\int_0^1\ln x\,dx=-1.
\]
这里对 $-\ln x$ 应用命题~\ref{prop:pei-int-singular-sums}，不能把零点无界的对数直接当作闭区间上的 Riemann 可积函数。故 $G_n/(nd)\to e^{-1}$，最终 $G_n/A_n\to2/e$。
\end{solve}

\subsection{非负截断与几个综合推论}
若 $f\ge0$ 在 $(0,\infty)$ 连续，令 $f_M=\min(f,M)$，则 $\lim_{M\to\infty}\int_0^\infty f_M=\int_0^\infty f$，允许两边为 $+\infty$。在任意 $[\delta,R]$ 上，$M$ 足够大时 $f_M=f$，因而左端极限不小于每个截断积分；另一方面 $f_M\le f$，取所有截断区间的上确界即可证明等式。这一证明依赖非负性提供的单调结构，不能不加条件地用于变号函数或任意交换两个极限。

若连续正函数 $f$ 在 $[0,\infty)$ 上可积，Schwarz 不等式给出
\[
 R^2\le\left(\int_0^Rf\right)\left(\int_0^R\frac1f\right),
\]
所以 $\int_0^\infty1/f$ 必发散。这是把有限区间上的不等式用于无穷区间的一个例子。另一方面，仅有两个非负函数各自积分发散，并不能据比较判别推出其最小值积分也发散：例如 $f(x)=\max(\sin x,0)$、$g(x)=\max(-\sin x,0)$ 连续非负，两者的积分都发散，但 $\min(f,g)\equiv0$。比较判别所需的是逐点大小关系和正确的推理方向。

若 $f$ 在实轴上局部可积，$f(x)\to A$（$x\to+\infty$）、$f(x)\to B$（$x\to-\infty$），则对固定实数 $h$，通常的两侧反常积分满足
\[
 \int_{-\infty}^\infty\bigl(f(x+h)-f(x)\bigr)\,dx=h(A-B).
\]
在有限区间 $[u,v]$ 上，原积分等于 $\int_v^{v+h}f-\int_u^{u+h}f$；令两端分别趋于无穷即得结论，且分别固定另一端也证明两条尾积分各自收敛。因此这个恒等式不依赖对称主值，也不要求 $\int f$ 本身收敛。

% 教材来源：数学分析/textbooks/数分习题（分章原文）/数分习题-5.pdf。
% 按四节主题整合正文与代表性练习；非逐题转录。
\chapter{级数}
\label{chap:pei-series}

本章依据裴礼文《数学分析中的典型问题与方法》第五章整理。级数把有限求和推广为部分和序列的极限，其困难不仅在于计算和，还在于判断极限是否存在，以及何时可以交换求和与其他运算。数项级数提供敛散判别和尾和估计；函数项级数进一步要求控制自变量变化时的误差；幂级数在收敛区间内部具有稳定的逐项运算性质；Fourier 级数则通过正交性把函数展开问题转化为系数问题。四节的共同线索是：先确定收敛的意义和范围，再说明变形、积分或求导的依据。

% 对应所给 PDF 第 1--47 页。
\section{数项级数}
\label{sec:pei-numerical-series}

\subsection{部分和、求和方法与重排}

\begin{definition}
设 $(a_n)_{n\geq1}$ 为实数列或复数列，记部分和 $S_N=\sum_{n=1}^{N}a_n$。若 $S_N$ 趋于有限极限 $S$，则称\emph{无穷级数（infinite series）} $\sum_{n=1}^{\infty}a_n$ 收敛，$S$ 为其和；否则称其发散。收敛时的余和记为 $r_N=S-S_N$。若 $\sum_{n=1}^{\infty}|a_n|$ 收敛，则称原级数\emph{绝对收敛（absolute convergence）}；原级数收敛而绝对值级数发散时，称其\emph{条件收敛（conditional convergence）}。
\end{definition}

由 $a_n=S_n-S_{n-1}$ 可知，收敛必有 $a_n\to0$，但反过来不成立。求和首先应考察有限部分和：连锁消去、错位相减、建立部分和方程等方法，都应在有限和上完成，再检查剩余边界项的极限。绝对收敛级数可以任意重排；条件收敛级数的项次不能随意交换，即使重排后仍然收敛，其和也可能改变。

\begin{example}
\label{ex:pei-series-sums}
计算 $\sum_{n=1}^{\infty}nq^n$（$|q|<1$）及 $\sum_{n=1}^{\infty}(\sqrt n-2\sqrt{n+1}+\sqrt{n+2})$。
\end{example}
\begin{solve}
第一式令 $S_N=\sum_{n=1}^{N}nq^n$，则
\[
 (1-q)S_N=\sum_{n=1}^{N}q^n-Nq^{N+1}.
\]
因 $Nq^{N+1}\to0$，所以其和为 $q/(1-q)^2$。第二式按相邻差分消去中间项，得到
\[
 \sum_{n=1}^{N}(\sqrt n-2\sqrt{n+1}+\sqrt{n+2})
 =1-\sqrt2-\sqrt{N+1}+\sqrt{N+2}\longrightarrow1-\sqrt2.
\]
这里每个通项为负并不影响差分法；关键是保留有限部分和的两端。\qedhere
\end{solve}

求 $\sum q^n\cos n\theta$、$\sum q^n\sin n\theta$ 时，可以先对有限等比和 $\sum(qe^{i\theta})^n$ 求和，再取实部和虚部；当 $|q|<1$ 时得到
\[
 \sum_{n=1}^{\infty}q^n\cos n\theta
 =\frac{q\cos\theta-q^2}{1-2q\cos\theta+q^2},\qquad
 \sum_{n=1}^{\infty}q^n\sin n\theta
 =\frac{q\sin\theta}{1-2q\cos\theta+q^2}.
\]
若已知 $\sum a_n=S$，则部分和的算术平均也趋于 $S$，因而
\begin{equation}
 \frac1N\sum_{k=1}^{N}(N-k+1)a_k
 =\frac{S_1+\cdots+S_N}{N}\longrightarrow S.
 \label{eq:pei-cesaro-series}
\end{equation}
证明只需把平均值减去 $S$ 后，将前面固定的有限项与后面的统一小量分开；这就是 Cesàro 平均的基本性质。

\begin{example}
将交错调和级数按每次取 $p$ 个正项、再取 $q$ 个负项的规则重排，其中 $p,q$ 为固定正整数，求重排后的和。
\end{example}
\begin{solve}
记 $H_m=\sum_{k=1}^{m}1/k$。完成 $N$ 个块后的部分和为
\[
 H_{2pN}-\frac12H_{pN}-\frac12H_{qN}.
\]
利用 $H_m=\log m+\gamma+o(1)$，其极限为 $\log2+\tfrac12\log(p/q)$。同一块内尚未加完的项数至多为 $p+q$，每项趋于零，所以全部部分和也趋于此极限。这里 $\gamma$ 为 Euler 常数；当 $p\ne q$ 时，重排后的和与原和 $\log2$ 不同。
\end{solve}

\subsection{Cauchy 准则与正项级数的比较}

\begin{theorem}[级数的 Cauchy 准则]
\label{thm:pei-series-cauchy}
级数 $\sum a_n$ 收敛，当且仅当
\[
 \forall\varepsilon>0\ \exists N\ \forall m>n\geq N,
 \qquad\left|\sum_{k=n+1}^{m}a_k\right|<\varepsilon.
\]
其发散的等价表述为：存在 $\varepsilon_0>0$，对每个 $N$ 都能找到 $m>n\geq N$，使上述尾段和的绝对值至少为 $\varepsilon_0$。
\end{theorem}
\begin{proof}
因为尾段和恰为 $S_m-S_n$，结论就是实数或复数序列的 Cauchy 收敛准则。
\end{proof}

\emph[caution]{这里尾段长度 $m-n$ 可以随 $n$ 改变，不能只对每个固定长度验证极限为零。}例如调和级数满足对固定 $p$ 有 $\sum_{k=n+1}^{n+p}1/k\leq p/(n+1)\to0$，但
\[
 \sum_{k=n+1}^{2n}\frac1k\geq\frac12,
\]
故仍然发散。对非负项级数，部分和单调递增，因此收敛等价于部分和有上界；比较判别法、分组估计和积分判别法都可看作寻找这一上界的方法。

\begin{theorem}[比较、比值与根值判别]
\label{thm:pei-positive-tests}
设 $a_n,b_n\geq0$。若充分大时 $a_n\leq Cb_n$，其中 $C>0$，则 $\sum b_n$ 收敛推出 $\sum a_n$ 收敛；若 $a_n/b_n\to c\in(0,\infty)$，则两级数同敛散。对最终为正的 $a_n$，若 $\limsup a_{n+1}/a_n<1$，则 $\sum a_n$ 收敛；若最终 $a_{n+1}/a_n\geq1$，则发散。对非负的 $a_n$，令 $\rho=\limsup\sqrt[n]{a_n}$，则 $\rho<1$ 时收敛，$\rho>1$ 时发散。
\end{theorem}
\begin{proof}
比较结论直接由非负部分和的有界性得到。比值的上极限小于 $1$ 时，可取 $q<1$ 使最终 $a_{n+1}\leq qa_n$，于是尾部被等比级数控制；比值最终至少为 $1$ 时，正项不能趋于零。根值的上极限小于 $1$ 时，最终有 $a_n\leq q^n$；大于 $1$ 时，有无穷多项满足 $a_n>1$，违反收敛的必要条件。
\end{proof}

比值或根值等于 $1$ 时通常无法判断，例如 $\sum1/n$ 与 $\sum1/n^2$ 在这两种检验下给出相同的临界值。比较时也不能把“比 $1/n$ 高阶的无穷小量”作为收敛的充分条件：$1/(n\log n)=o(1/n)$，但对应级数发散。更准确的做法是找到可判定的比较对象，并证明最终不等式或正的有限比值极限。

\begin{theorem}[积分判别法与凝聚判别法]
\label{thm:pei-integral-condensation}
若 $f:[1,\infty)\to[0,\infty)$ 单调递减，则 $\sum_{n=1}^{\infty}f(n)$ 与 $\int_1^{\infty}f(x)\,dx$ 同敛散；收敛时有
\begin{equation}
 \int_{N+1}^{\infty}f(x)\,dx
 \leq\sum_{n=N+1}^{\infty}f(n)
 \leq\int_N^{\infty}f(x)\,dx.
 \label{eq:pei-integral-tail}
\end{equation}
若 $a_n\geq0$ 单调递减，则 $\sum a_n$ 与 $\sum_{j=0}^{\infty}2^ja_{2^j}$ 同敛散，此为 Cauchy 凝聚判别法（condensation test）。
\end{theorem}
\begin{proof}
在每个单位区间上以端点函数值夹住积分，有限求和后取极限，即得积分判别及余和估计。对凝聚判别，在第 $j$ 个二进块上有
\[
 2^ja_{2^{j+1}}\leq\sum_{n=2^j}^{2^{j+1}-1}a_n\leq2^ja_{2^j}.
\]
两侧分别是凝聚级数的移位和常数倍，故由比较判别得到结论。
\end{proof}

因此 $\sum1/n^p$ 恰在 $p>1$ 时收敛，而
\[
 \sum_{n=2}^{\infty}\frac1{n(\log n)^p}
 \quad\text{恰在 }p>1\text{ 时收敛}.
\]
后一结论由变量代换 $t=\log x$ 化为 $\int t^{-p}\,dt$ 得到。若 $n$ 足够大且通项为 $(1-p\log n/n)^n$，其中 $p$ 是固定实数，则
\[
 n\log\!\left(1-\frac{p\log n}{n}\right)
 =-p\log n+O\!\left(\frac{(\log n)^2}{n}\right),
\]
故通项与 $n^{-p}$ 等价，仍以 $p=1$ 为临界点。使用 Taylor 展开时，必须同时检查展开参数趋于零和余项经外部倍乘后仍趋于零。

\begin{proposition}[Kummer 判别法]
设 $u_n>0$、$c_n>0$。若存在 $\delta>0$ 使充分大时
\[
 c_n\frac{u_n}{u_{n+1}}-c_{n+1}\geq\delta,
\]
则 $\sum u_n$ 收敛；若 $\sum1/c_n$ 发散且充分大时左端不大于零，则 $\sum u_n$ 发散。
\end{proposition}
\begin{proof}
第一种情形乘以 $u_{n+1}$ 后求和，得到
\[
 \delta\sum_{n=N}^{M}u_{n+1}
 \leq c_Nu_N-c_{M+1}u_{M+1}\leq c_Nu_N,
\]
故部分和有界。第二种情形说明 $c_nu_n$ 最终不减，因而最终有 $u_n\geq C/c_n$，其中 $C>0$，由比较判别得发散。
\end{proof}

取 $c_n=n$ 即得 Raabe 判别法：若 $n(u_n/u_{n+1}-1)\to L$，则 $L>1$ 时收敛，$L<1$ 时发散，$L=1$ 时需进一步分析。取辅助数列的意义是把复杂通项转化为一个可以连锁消去的不等式，而非机械计算比值。

\subsection{Abel 变换与变号级数}

\begin{proposition}[Abel 变换]
\label{prop:pei-abel-transform}
给定整数 $m\geq n$，记 $A_k=\sum_{j=n}^{k}a_j$，则
\begin{equation}
 \sum_{k=n}^{m}a_kb_k
 =A_mb_m+\sum_{k=n}^{m-1}A_k(b_k-b_{k+1}).
 \label{eq:pei-abel-transform}
\end{equation}
若 $|A_k|\leq M$，则
\[
 \left|\sum_{k=n}^{m}a_kb_k\right|
 \leq M\left(|b_m|+\sum_{k=n}^{m-1}|b_k-b_{k+1}|\right).
\]
特别地，若 $b_k\geq0$ 且单调递减，右侧不超过 $Mb_n$。
\end{proposition}
\begin{proof}
令 $A_{n-1}=0$，把 $a_k=A_k-A_{k-1}$ 代入左侧并错位相减，即得恒等式；再用三角不等式及单调性得到估计。
\end{proof}

\begin{theorem}[Dirichlet、Abel 与 Leibniz 判别法]
\label{thm:pei-dirichlet-abel}
若 $\sum_{k=1}^{n}a_k$ 对 $n$ 一致有界，且实数列 $b_n$ 单调趋于零，则 $\sum a_nb_n$ 收敛，此为 Dirichlet 判别法。若 $\sum a_n$ 收敛且 $b_n$ 单调有界，则 $\sum a_nb_n$ 收敛，此为 Abel 判别法。特别地，$b_n\downarrow0$ 时，交错级数 $\sum_{n=1}^{\infty}(-1)^{n-1}b_n$ 收敛，且余和满足 $|r_N|\leq b_{N+1}$，此为 Leibniz 判别法及其余项估计。
\end{theorem}
\begin{proof}
Dirichlet 判别中，任意尾段的 $a_k$ 部分和至多为原部分和上界的两倍。由式~\eqref{eq:pei-abel-transform}，乘积级数的尾段和被常数倍的 $|b_n|$ 控制，因此满足 Cauchy 准则。Abel 判别中令 $b=\lim b_n$，分解 $a_nb_n=a_n(b_n-b)+ba_n$，第一项用 Dirichlet 判别，第二项用原级数收敛性。交错级数的收敛是 Dirichlet 判别的特例；从余和首项开始两两分组，可知余和的符号与首项相同，绝对值不超过首项。
\end{proof}

\begin{example}
\label{ex:pei-sine-harmonic}
证明当 $\theta\notin2\pi\mathbb Z$ 时，$\sum_{n=1}^{\infty}\sin(n\theta)/n$ 收敛，并说明 $\sum_{n=1}^{\infty}\sin n/n$ 是条件收敛的。
\end{example}
\begin{solve}
有限等比和给出
\[
 \sum_{k=1}^{N}\sin(k\theta)
 =\frac{\cos(\theta/2)-\cos((N+1/2)\theta)}{2\sin(\theta/2)},
\]
对固定的 $\theta\notin2\pi\mathbb Z$，其绝对值有界，故用 Dirichlet 判别得到收敛。取 $\theta=1$，又有 $|\sin n|\geq\sin^2n=(1-\cos2n)/2$，而 $\sum\cos2n/n$ 由同一判别法收敛，调和级数发散，故 $\sum|\sin n|/n$ 发散。
\end{solve}

判断条件收敛必须分别证明原级数收敛和绝对值级数发散；仅仅发现比值或根值为 $1$，不能得到条件收敛。对交错形状也要检查绝对值的单调性，例如令 $a_{2k-1}=1/\sqrt k$、$a_{2k}=-1/(2\sqrt k)$，虽然项的符号交替且 $a_n\to0$，但每对之和为 $1/(2\sqrt k)$，级数发散。

Abel 变换还能处理单调性以外的情形。若 $\sum a_n$ 收敛且 $\sum|b_{n+1}-b_n|<\infty$，则 $b_n$ 有界且 $\sum a_nb_n$ 收敛：在式~\eqref{eq:pei-abel-transform} 中，先用原级数的 Cauchy 准则使所有尾段 $A_k$ 同时小，再以 $b_n$ 的有界性和差分绝对和作统一估计。这里起作用的是数列的\emph{有界变差（bounded variation）}，即相邻差分的绝对和有限。

\subsection{收敛性的应用与反例构造}

\begin{proposition}
\label{prop:pei-monotone-na}
若 $a_n\geq0$ 单调递减且 $\sum a_n$ 收敛，则 $na_n\to0$。若仅假设 $\sum a_n$ 收敛，则仍有
\begin{equation}
 \frac1N\sum_{k=1}^{N}ka_k\longrightarrow0.
 \label{eq:pei-weighted-small}
\end{equation}
第二个结论不要求非负性或单调性。
\end{proposition}
\begin{proof}
令 $m=\lfloor n/2\rfloor$，则
\[
 0\leq(n-m)a_n\leq\sum_{k=m+1}^{n}a_k\longrightarrow0,
\]
又 $n/(n-m)\leq2$，故第一式成立。对第二式，令 $S_k=\sum_{j=1}^{k}a_j\to S$，Abel 变换给出
\[
 \frac1N\sum_{k=1}^{N}ka_k
 =S_N-\frac1N\sum_{k=1}^{N-1}S_k\longrightarrow S-S=0.
\]
\end{proof}

\emph[caution]{单调性不能从命题~\ref{prop:pei-monotone-na} 的第一个结论中删去。}令 $a_n=1/n$ 当 $n$ 为完全平方数，其他情形令 $a_n=1/n^2$，则 $\sum a_n\leq\sum1/n^2+\sum1/k^2<\infty$，但 $na_n$ 在平方数指标处恒为 $1$。反过来，即使 $na_n\to0$ 也不能保证收敛，$a_n=1/(n\log n)$（$n\geq2$）就是反例。

条件收敛的反例常按有限块构造。令
\[
 a_{3k-2}=k^{-1/3},\qquad
 a_{3k-1}=-2k^{-1/3},\qquad
 a_{3k}=k^{-1/3}.
\]
每块之和为零，块内偏离零的部分和也趋于零，所以 $\sum a_n$ 收敛；但每块立方之和为 $-6/k$，故 $\sum a_n^3$ 发散。这说明对条件收敛级数，不能只根据 $|a_n^3|\leq|a_n|$ 就断言其立方级数收敛，因为作为比较对象的绝对值级数本身未必收敛。

\begin{example}
构造在 $[0,1]$ 上单调递增、在每个有理点间断而在每个无理点连续的函数。
\end{example}
\begin{solve}
把 $[0,1]$ 中的有理数无重复地排为 $(q_k)_{k\geq1}$，定义
\[
 f(x)=\sum_{k=1}^{\infty}2^{-k}h_k(x),\qquad
 h_k(x)=
 \begin{cases}
  \mathbf1_{\{x>0\}},&q_k=0,\\
  \mathbf1_{\{x\geq q_k\}},&q_k>0.
 \end{cases}
\]
其中 $\mathbf1_E$ 为集合 $E$ 的示性函数。各项不减，故 $f$ 不减；在 $q_k\in(0,1]$ 处，其相对于左侧的跃度至少为 $2^{-k}$，在 $q_k=0$ 处相对于右侧的跃度至少为 $2^{-k}$，所以各有理点均不连续。对无理点 $x_0$，给定 $\varepsilon>0$，先取 $N$ 使 $\sum_{k>N}2^{-k}<\varepsilon$，再取 $x_0$ 的充分小邻域避开 $q_1,\ldots,q_N$，则邻域内前 $N$ 项不变，而 $|f(x)-f(x_0)|<\varepsilon$，故在 $x_0$ 连续。
\end{solve}

\subsection{部分和归一化与反常积分}

\begin{proposition}
设 $a_n>0$，$S_n=\sum_{k=1}^{n}a_k$，$\lambda>0$。若 $\lambda>1$，则 $\sum a_n/S_n^\lambda$ 总收敛；若 $0<\lambda\leq1$，则它与 $\sum a_n$ 同敛散。
\end{proposition}
\begin{proof}
若 $S_n\to S<\infty$，则 $S_1\leq S_n\leq S$，归一化后的项与 $a_n$ 相差有界的正因子，故收敛。若 $S_n\to\infty$ 且 $\lambda>1$，对 $n\geq2$ 有
\[
 \frac{a_n}{S_n^\lambda}\leq
 \int_{S_{n-1}}^{S_n}t^{-\lambda}\,dt,
\]
求和后用 $\int_{S_1}^{\infty}t^{-\lambda}\,dt<\infty$。当 $\lambda=1$，对固定 $n$，
\[
 \sum_{k=n+1}^{m}\frac{a_k}{S_k}
 \geq\frac{S_m-S_n}{S_m}\longrightarrow1\qquad(m\to\infty),
\]
故违反 Cauchy 准则；当 $0<\lambda<1$，最终 $S_n\geq1$，有 $a_n/S_n^\lambda\geq a_n/S_n$，因此也发散。
\end{proof}

级数与反常积分之间的联系不仅用于判断敛散，还能用于余和的渐近估计和无穷区间的 Riemann 和。由式~\eqref{eq:pei-integral-tail}，若 $p>1$，则
\[
 \sum_{k=N+1}^{\infty}\frac1{k^p}
 \sim\frac{N^{1-p}}{p-1}.
\]
若 $f:[0,\infty)\to[0,\infty)$ 单调递减且可积，则对 $h>0$，单位小区间的积分估计给出
\[
 \int_h^{\infty}f(x)\,dx
 \leq h\sum_{n=1}^{\infty}f(nh)
 \leq\int_0^{\infty}f(x)\,dx.
\]
令 $h\downarrow0$ 即得
\begin{equation}
 h\sum_{n=1}^{\infty}f(nh)\longrightarrow\int_0^{\infty}f(x)\,dx.
 \label{eq:pei-infinite-riemann}
\end{equation}
例如取 $f(x)=1/(1+e^x)$，右端为 $\log2$。这种无穷区间上的求和极限不能仅凭“形式上像 Riemann 和”交换极限与无穷求和，单调可积性在此提供了对全部尾部的统一控制。

% 对应所给 PDF 第 48--95 页。
\section{函数项级数}
\label{sec:pei-function-series}

\emph{函数项级数（series of functions）}将每个数项替换为同一集合上定义的函数。固定自变量后可使用数项级数的判别方法；要研究和函数的性质，则需要进一步控制误差对自变量的依赖。

\subsection{逐点收敛、一致收敛与否定方法}

\begin{definition}
设 $E\subset\mathbb R$，函数列 $f_n:E\to\mathbb R$ 逐点收敛于 $f$，是指对每个固定 $x\in E$ 都有 $f_n(x)\to f(x)$；若
\[
 \forall\varepsilon>0\ \exists N\ \forall n\geq N\ \forall x\in E,
 \qquad |f_n(x)-f(x)|<\varepsilon,
\]
则称 $f_n$ 在 $E$ 上\emph{一致收敛（uniform convergence）}于 $f$，记为 $f_n\rightrightarrows f$。等价地，$\sup_{x\in E}|f_n(x)-f(x)|\to0$。函数项级数 $\sum u_n(x)$ 的一致收敛是指其部分和列一致收敛；在开区间每个内闭区间上一致收敛，称为\emph{内闭一致收敛或局部一致收敛（locally uniform convergence）}。
\end{definition}

逐点收敛中的指标可以依赖于 $x$，一致收敛要求所有点共用同一指标。证明一致收敛时可直接估计余项的上确界；证明非一致收敛时，可寻找随 $n$ 移动的点 $x_n$，使误差保持远离零。上确界不必在 $E$ 内取到，例如 $x^n$ 在 $[0,1)$ 上逐点趋于零，但 $\sup_{0\leq x<1}x^n=1$，故不一致收敛。

\begin{theorem}[一致 Cauchy 准则]
\label{thm:pei-uniform-cauchy}
级数 $\sum u_n(x)$ 在 $E$ 上一致收敛，当且仅当
\[
 \forall\varepsilon>0\ \exists N\ \forall m>n\geq N,
 \qquad\sup_{x\in E}\left|\sum_{k=n+1}^{m}u_k(x)\right|<\varepsilon.
\]
特别地，一致收敛必有 $\sup_{x\in E}|u_n(x)|\to0$，但此必要条件不充分。
\end{theorem}
\begin{proof}
必要性由一致收敛和三角不等式得到。充分性中，对每个固定 $x$，部分和为 Cauchy 列，故存在和 $S(x)$；将一致 Cauchy 不等式中的 $m\to\infty$，可得 $|S(x)-S_n(x)|\leq\varepsilon$ 对全部 $x$ 同时成立，因此部分和一致收敛。
\end{proof}

例~\ref{ex:pei-sine-harmonic} 中的正弦级数虽然在每点收敛，却在 $0$ 的任何邻域内都不一致收敛。事实上令 $x_n=\pi/(4n)$，则对 $n<k\leq2n$ 有 $kx_n\in[\pi/4,\pi/2]$，从而
\[
 \sum_{k=n+1}^{2n}\frac{\sin(kx_n)}k
 \geq\frac{\sqrt2}{2}\sum_{k=n+1}^{2n}\frac1k
 \geq\frac{\sqrt2}{4}.
\]
这里通项 $\sin(nx)/n$ 确实一致趋于零，仍不足以保证整个级数一致收敛；尾段长度和检验点必须允许同时变化。

\subsection{优级数与一致 Abel、Dirichlet 判别法}

\begin{theorem}[Weierstrass 的 $M$ 判别法]
\label{thm:pei-m-test}
若 $|u_n(x)|\leq M_n$ 对所有 $x\in E$ 成立，且非负项级数 $\sum M_n$ 收敛，则 $\sum u_n(x)$ 和 $\sum|u_n(x)|$ 都在 $E$ 上一致收敛。
\end{theorem}
\begin{proof}
任意尾段和的绝对值均不超过 $\sum_{k=n+1}^{m}M_k$，对 $x$ 无关；由数项级数及一致 Cauchy 准则得到结论。
\end{proof}

寻找优级数时，通常先求 $\sup_E|u_n|$，或用均值不等式、Taylor 余项和微分中值定理得到更简洁的上界。例如
\[
 \max_{0\leq x\leq1}x^n(1-x)^2
 =\left(\frac n{n+2}\right)^n\frac4{(n+2)^2}
 \leq\frac4{(n+2)^2},
\]
故 $\sum x^n(1-x)^2$ 在 $[0,1]$ 上一致收敛。$M$ 判别法是充分条件，不能因某个粗糙上界组成发散级数就断言原级数非一致收敛。

\begin{theorem}[一致 Dirichlet 与 Abel 判别法]
\label{thm:pei-uniform-dirichlet}
若 $|\sum_{k=1}^{N}a_k(x)|\leq M$ 对全部 $x\in E$ 和 $N$ 成立，$b_n(x)$ 对每个固定 $x$ 关于 $n$ 单调，且 $b_n\rightrightarrows0$，则 $\sum a_n(x)b_n(x)$ 在 $E$ 上一致收敛。若改为 $\sum a_n(x)$ 在 $E$ 上一致收敛，$b_n(x)$ 对 $n$ 单调且对 $n,x$ 一致有界，也可得到相同结论。
\end{theorem}
\begin{proof}
第一种情形，尾段的 $a_k(x)$ 部分和不超过 $2M$。对每个 $x$ 应用 Abel 变换，利用单调性将尾段和控制为常数倍的 $\sup_E|b_{n+1}|$，于是统一趋于零。第二种情形，先由 $\sum a_n(x)$ 的一致 Cauchy 性使尾段部分和不超过任意小的 $\varepsilon$；若 $|b_k(x)|\leq B$，单调性给出
\[
 |b_m(x)|+\sum_{k=n+1}^{m-1}|b_k(x)-b_{k+1}(x)|\leq3B.
\]
故乘积尾段和不超过 $3B\varepsilon$。
\end{proof}

在任意 $[\delta,2\pi-\delta]$（$0<\delta<\pi$）上，有限正弦和被 $1/\sin(\delta/2)$ 控制，因此 $\sum\sin(nx)/n$ 在这些区间上一致收敛；区间接近 $0$ 或 $2\pi$ 时，此上界不再统一，正是前述非一致收敛现象出现的位置。类似地，若数项级数 $\sum a_n$ 收敛，则 $n^{-x}$ 对 $n$ 单调且在 $x\geq0$ 上介于 $0$ 与 $1$ 之间，故 $\sum a_n/n^x$ 在 $[0,\infty)$ 上一致收敛。

\begin{example}
说明一致收敛、逐点绝对收敛和绝对值级数的一致收敛之间的区别。
\end{example}
\begin{solve}
常函数项级数 $\sum(-1)^{n-1}/n$ 在任意集合上一致收敛，但在每点都不绝对收敛。再考虑
\[
 \sum_{n=0}^{\infty}(-1)^nx^n(1-x),\qquad 0\leq x\leq1.
\]
它在每点绝对收敛，由交错余项估计有 $|r_N(x)|\leq(1-x)x^{N+1}$，右端上确界趋于零，故原级数一致收敛。然而绝对值级数的和在 $x<1$ 时为 $1$，在 $x=1$ 时为 $0$，不连续；由一致收敛保连续性定理，绝对值级数不一致收敛。因而“一致收敛且每点绝对收敛”仍不等于绝对值级数一致收敛。
\end{solve}

\subsection{Dini 定理、等度连续与紧性}

\begin{theorem}[Dini 定理]
\label{thm:pei-dini}
设连续函数列 $f_n$ 在闭区间 $[a,b]$ 上逐点收敛于连续函数 $f$，且对每个固定 $x$，$f_n(x)$ 关于 $n$ 单调，则 $f_n\rightrightarrows f$。
\end{theorem}
\begin{proof}
记 $g_n(x)=|f_n(x)-f(x)|$，则 $g_n$ 连续且随 $n$ 单调降至零。对给定 $\varepsilon>0$，集合 $U_n=\{x:g_n(x)<\varepsilon\}$ 是递增的相对开集并覆盖 $[a,b]$。由有限覆盖性质，存在 $N$ 使整个区间包含于 $U_N$，因而 $n\geq N$ 时 $g_n(x)<\varepsilon$ 对全部 $x$ 成立。
\end{proof}

Dini 定理中的单调性是对指标 $n$ 而言，不能误解为各 $f_n$ 对 $x$ 单调。若 $u_n\geq0$ 连续且 $\sum u_n$ 的和在闭区间上连续，则其部分和满足 Dini 定理，级数一致收敛。闭区间和极限连续的条件都不可任意去掉：$x^n$ 在 $[0,1)$ 上单调趋于连续的零函数，却不一致收敛；在 $[0,1]$ 上则极限不连续。

\begin{definition}
函数族 $\mathcal F$ 在 $E$ 上\emph{等度连续（equicontinuous）}，是指对每个 $\varepsilon>0$，存在与 $g\in\mathcal F$ 无关的 $\delta>0$，使 $x,y\in E$ 且 $|x-y|<\delta$ 时都有 $|g(x)-g(y)|<\varepsilon$。在闭区间上，这等价于通常逐点定义的等度连续性；这里采用统一的表述。
\end{definition}

\begin{proposition}
\label{prop:pei-equicontinuity}
若 $f_n$ 在 $[a,b]$ 上等度连续且逐点收敛于 $f$，则 $f$ 连续且 $f_n\rightrightarrows f$。若函数列一致有界且等度连续，则它存在一致收敛的子列，此为 Arzelà--Ascoli 定理在闭区间上的形式。
\end{proposition}
\begin{proof}
在等度连续不等式中令 $n\to\infty$，可知 $f$ 具有相同的连续模。给定 $\varepsilon>0$，用有限个网格点使每个 $x$ 都与某网格点相距小于相应的 $\delta$；网格点有限，可取公共指标使其上的 $|f_n-f|<\varepsilon/3$。再分别控制 $f_n$、$f$ 从网格点到 $x$ 的变化，即得全区间误差小于 $\varepsilon$。

第二个结论中，先枚举闭区间中的一个可数稠密集，在其第一点取收敛子列，再在第二点取子列，依次进行并取对角子列，使它在全部稠密点收敛。等度连续性和有限网格将这些点上的 Cauchy 性提升为全区间上的一致 Cauchy 性，故该子列一致收敛。
\end{proof}

一个常用的充分条件是导数一致有界：若 $|f_n'(x)|\leq M$，则微分中值定理给出 $|f_n(x)-f_n(y)|\leq M|x-y|$。这能把部分和导数的统一估计转化为等度连续，再结合逐点收敛得到一致收敛；仅有每个函数分别连续或分别有界，并不足以替代上述统一条件。

\subsection{交换极限、积分与微分的条件}

\begin{theorem}[一致收敛的基本运算性质]
\label{thm:pei-uniform-operations}
若 $f_n\rightrightarrows f$ 于 $E$，且各 $f_n$ 在 $x_0\in E$ 连续，则 $f$ 在 $x_0$ 连续。若各 $f_n$ 在有限闭区间 $[a,b]$ 上 Riemann 可积且一致收敛于 $f$，则 $f$ 也可积，且
\[
 \int_a^bf(x)\,dx=\lim_{n\to\infty}\int_a^bf_n(x)\,dx.
\]
若 $f_n$ 在区间 $I$ 的聚点 $x_0$ 处有极限 $c_n$，且在该点的某个去心邻域与 $I$ 的交集上一致收敛于 $f$，则 $c_n$ 收敛，并有 $\lim_{x\to x_0}f(x)=\lim_{n\to\infty}c_n$。
\end{theorem}
\begin{proof}
连续性与交换极限的结论都由三项分解得到：先用一致收敛控制 $f-f_N$，将 $N$ 固定，再利用 $f_N$ 的连续性或极限控制中间项。可积性可由 Darboux 上下和之差验证：若 $\|f-f_N\|_\infty<\eta$，则任一分割下 $f$ 的上下和之差不超过 $f_N$ 的相应差加 $2\eta(b-a)$；先取 $N$，再选取 $f_N$ 的合适分割即可。最后
\[
 \left|\int_a^b(f-f_n)\,dx\right|\leq(b-a)\|f-f_n\|_\infty\longrightarrow0.
\]
\end{proof}

\begin{theorem}[逐项求导]
\label{thm:pei-termwise-derivative}
设 $u_n\in C^1([a,b])$，$\sum u_n(x_0)$ 在某个 $x_0\in[a,b]$ 收敛，且 $\sum u_n'$ 在 $[a,b]$ 上一致收敛于 $g$。则 $\sum u_n$ 在 $[a,b]$ 上一致收敛于连续函数 $S$，并且在内部可导，$S'=g=\sum u_n'$；端点取相应的单侧导数。
\end{theorem}
\begin{proof}
对部分和 $S_N$，有 $S_N(x)=S_N(x_0)+\int_{x_0}^{x}S_N'(t)\,dt$。在导数级数中逐项积分，由积分误差不超过 $(b-a)\|S_N'-g\|_\infty$，可得 $S_N$ 一致趋于
\[
 S(x)=\sum_{n=1}^{\infty}u_n(x_0)+\int_{x_0}^{x}g(t)\,dt.
\]
$g$ 是连续函数的一致极限，因而连续，再由微积分基本定理得 $S'=g$。
\end{proof}

上述结论可在每个内闭区间上使用，因此局部一致收敛足以保证开区间内部的连续性；逐项求导则应检查\emph[caution]{导数级数}的局部一致收敛和原级数的一个收敛点。例如
\[
 F(x)=\sum_{n=1}^{\infty}ne^{-nx},\qquad x>0,
\]
在任意 $[\delta,A]$（$\delta>0$）上，其 $r$ 次导数级数的绝对值被 $\sum n^{r+1}e^{-n\delta}$ 控制，所以 $F\in C^\infty((0,\infty))$。但在整个 $(0,\infty)$ 上，通项的上确界为 $n$，原级数不一致收敛。由等比级数还可求得 $F(x)=e^{-x}/(1-e^{-x})^2$，与上述局部结论一致。

一致收敛本身不保证导数收敛，例如 $f_n(x)=\sin(nx)/n$ 在实轴上一致趋于零，而 $f_n'(0)=1$。逐项积分的充分条件则可以弱于一致收敛：若和函数及部分和可积，所需的实际条件只是余和积分趋于零。

\begin{example}
\label{ex:pei-nonuniform-integration}
令 $u_n(x)=x^n\log x$（$n\geq1$，$0<x<1$），并在 $x=0,1$ 处取连续延拓值 $0$。说明该级数不一致收敛，却可逐项积分。
\end{example}
\begin{solve}
在 $0<x<1$ 时，$S(x)=x\log x/(1-x)$，而 $S(1)=0$；因 $\lim_{x\uparrow1}S(x)=-1$，连续性定理排除 $[0,1]$ 上的一致收敛。在开区间 $(0,1)$ 上，其余和为
\[
 r_N(x)=\frac{x^{N+1}\log x}{1-x}.
\]
由 $-\log x=\int_x^1dt/t\leq(1-x)/x$，得 $|r_N(x)|\leq x^N$，从而
\[
 \int_0^1|r_N(x)|\,dx\leq\frac1{N+1}\longrightarrow0.
\]
因此 $\int_0^1S=\sum_{n=1}^{\infty}\int_0^1u_n=-\sum_{n=1}^{\infty}1/(n+1)^2$。此外固定 $N$ 时 $r_N(x)\to-1$（$x\uparrow1$），故开区间上也不一致收敛。
\end{solve}

对极限集中在端点附近的积分，可采用先控制小邻域、再控制其余区间的方法。若 $p_n\geq0$ 可积，$\int_a^bp_n=1$，且每个 $\delta>0$ 都满足 $\int_{a+\delta}^{b}p_n\to0$，则对任意连续的 $g$ 有 $\int_a^bgp_n\to g(a)$：把 $g-g(a)$ 在 $[a,a+\delta]$ 上用连续性控制，在剩余区间用有界性及质量趋零控制。非负性在此用于把局部积分估计为总质量，不能在变号函数列上直接套用。

\subsection{一致逼近及其适用区间}

\begin{theorem}[Weierstrass 多项式逼近定理]
\label{thm:pei-weierstrass-approximation}
若 $f$ 在有限闭区间 $[a,b]$ 上连续，则存在多项式列 $P_n$ 在该区间上一致收敛于 $f$。
\end{theorem}
\begin{proof}
线性换元后只需考虑 $[0,1]$。取 Bernstein 多项式
\[
 B_n(f;x)=\sum_{k=0}^{n}f(k/n)\binom nkx^k(1-x)^{n-k}.
\]
记权重 $w_{n,k}(x)=\binom nkx^k(1-x)^{n-k}\geq0$，二项式恒等式给出
\[
 \sum_{k=0}^{n}w_{n,k}=1,\qquad
 \sum_{k=0}^{n}(k/n-x)^2w_{n,k}=\frac{x(1-x)}n\leq\frac1{4n}.
\]
设 $|f|\leq M$，由一致连续性选取 $\delta>0$，使 $|s-t|<\delta$ 时 $|f(s)-f(t)|<\varepsilon$。将和分为 $|k/n-x|<\delta$ 与其余两部分，得到
\[
 |B_n(f;x)-f(x)|
 \leq\varepsilon+2M\sum_{|k/n-x|\geq\delta}w_{n,k}
 \leq\varepsilon+\frac{M}{2n\delta^2}.
\]
这一估计对全部 $x\in[0,1]$ 成立，故一致逼近成立。
\end{proof}

作为应用，若连续函数 $f$ 满足 $\int_a^bx^kf(x)\,dx=0$ 对全部非负整数 $k$ 成立，则 $f\equiv0$。事实上，对逼近 $f$ 的多项式 $P_n$ 有 $\int fP_n=0$，而
\[
 \left|\int_a^bf(f-P_n)\,dx\right|
 \leq(b-a)\|f\|_\infty\|f-P_n\|_\infty\longrightarrow0,
\]
故 $\int f^2=0$，再由连续性得到逐点为零。这里把对单项式的条件推广到一般连续函数，依靠的是一致逼近及积分误差控制。

逼近区间无界时则有本质限制。若多项式列 $P_n$ 在一个无界区间上一致收敛，则其尾部满足 $\sup|P_n-P_m|<1$；非恒定多项式在无界区间上不可能有界，所以所有充分后的 $P_n$ 只相差常数。固定其中一个多项式后，这些常数收敛，因而极限仍为多项式。故有限闭区间上的逼近定理不能直接推广为整个实轴上的多项式一致逼近。

\subsection{连续而处处不可导的级数构造}

令 $\phi(t)=\operatorname{dist}(t,\mathbb Z)$，即 $t$ 到最近整数的距离。它是周期为 $1$、值域为 $[0,1/2]$ 的连续三角波，在相邻的整数与半整数之间斜率为 $1$ 或 $-1$。考虑
\begin{equation}
 W(x)=\sum_{k=0}^{\infty}4^{-k}\phi(4^kx),\qquad x\in\mathbb R.
 \label{eq:pei-nowhere-differentiable}
\end{equation}
其通项绝对值不超过 $4^{-k}/2$，由 $M$ 判别法知 $W$ 在实轴上连续。下面核查它在每点都不存在有限导数，说明振幅趋零虽然能保证连续，却不能抵消频率不断增大对差商的影响。

固定 $x$，对每个 $n\geq1$ 选择 $h_n=4^{-n}$ 或 $h_n=-4^{-n}$，使 $x$ 与 $x+h_n$ 位于 $\phi(4^{n-1}t)$ 的同一条直线段内。这样的方向总能选到，因为该直线段长为 $2\cdot4^{-n}$。较低频率的折点都属于这一层的折点网格，所以前 $n$ 项在所选小段上均为线性；当 $k\geq n$ 时，$4^kh_n$ 是整数，相应差分由周期性全部消失。因此
\[
 \frac{W(x+h_n)-W(x)}{h_n}
 =\sum_{k=0}^{n-1}\sigma_{k,n},\qquad \sigma_{k,n}\in\{-1,1\}.
\]
若 $x$ 不属于任何折点网格，则每个较低频率在 $x$ 的斜率已经确定，$\sigma_{k,n}=\sigma_k$ 与 $n>k$ 无关；右端成为增量恒为 $1$ 或 $-1$ 的部分和列，不可能收敛。若 $x$ 属于某个折点网格，则从某一层起 $4^kx$ 都是整数，所有高频项在所选方向上的斜率相同，差商绝对值至少为 $n-C$，其中 $C$ 与 $n$ 无关，也不可能有有限极限。由此 $W$ 处处不可导；此构造同时核实了连续性和不可导性，而非仅凭图形的粗糙程度判断。

% 对应所给 PDF 第 96--131 页。
\section{幂级数}
\label{sec:pei-power-series}

\emph{幂级数（power series）}是通项为 $a_n(x-c)^n$ 的函数项级数，其中系数 $a_n$ 和展开中心 $c$ 固定。它的收敛集合具有由中心向外扩展的区间结构，这使局部一致收敛和逐项运算能够得到统一处理。

\subsection{收敛半径、缺项与端点检验}

\begin{theorem}[Cauchy--Hadamard 公式]
\label{thm:pei-radius}
幂级数 $\sum_{n=0}^{\infty}a_n(x-c)^n$ 的\emph{收敛半径（radius of convergence）}为
\[
 R=\frac1{\limsup_{n\to\infty}\sqrt[n]{|a_n|}},
\]
约定 $1/0=\infty$、$1/\infty=0$。在 $|x-c|<R$ 时绝对收敛，在 $|x-c|>R$ 时发散；当 $0<R<\infty$ 时，两个端点 $x=c\pm R$ 必须分别检验。若 $|a_n/a_{n+1}|$ 的极限存在，且系数最终非零，则也可用该极限计算 $R$。
\end{theorem}
\begin{proof}
对固定 $x$，通项绝对值的 $n$ 次根之上极限为 $|x-c|\limsup\sqrt[n]{|a_n|}$；应用根值判别即可。比值公式由比值与根值的极限关系得到。
\end{proof}

对缺项级数必须按实际幂次处理，不能把“第几项”误当作“$x$ 的几次幂”。例如
\[
 \sum_{n=1}^{\infty}nx^{n^2}
 \quad\text{的收敛范围为 }(-1,1),\qquad
 \sum_{n=1}^{\infty}\frac{x^{n^2}}{2^n}
 \quad\text{的收敛范围为 }[-1,1].
\]
两者补入零系数后的收敛半径都是 $1$：非零系数的幂次根分别为 $n^{1/n^2}\to1$、$2^{-1/n}\to1$。第一个在两端点通项不趋零，第二个在两端点都绝对收敛。对 $\sum a_n\psi(x)^n$，则先在新变量 $t=\psi(x)$ 中判断收敛范围，再解相应的不等式及端点方程。

系数逐项相加时，两个幂级数之和的半径至少为两半径的较小者；两半径不同时恰等于较小者，否则在某个原来发散的点上用相减会导致矛盾。半径相同时可能因系数抵消而增大，不能总写成等号。对系数逐项相乘的级数 $\sum a_nb_nx^n$，当两原半径为有限正数时，有 $R\geq R_aR_b$，由根值上极限的不等式直接得到；这与通常 Cauchy 乘积的系数卷积不是同一种运算。

\subsection{内部一致收敛与端点 Abel 定理}

\begin{proposition}
\label{prop:pei-power-interior}
幂级数在收敛区间内部局部一致绝对收敛，可以任意次逐项微分和积分；所得级数的收敛半径保持不变。其和函数 $f$ 满足 $a_n=f^{(n)}(c)/n!$，故幂级数展开唯一。
\end{proposition}
\begin{proof}
给定 $0<r<\rho<R$，$\sum|a_n|\rho^n$ 收敛，而 $|a_n(x-c)^n|\leq|a_n|\rho^n$ 对 $|x-c|\leq r$ 成立。导数级数中的附加因子为 $n$ 的固定次幂，由 $n^k(r/\rho)^n$ 有界，可继续用优级数控制。逐项积分亦同理；系数乘以 $n$ 或除以 $n+1$ 不改变根值上极限，所以半径不变。最后令 $x=c$ 即得到系数公式。
\end{proof}

\begin{theorem}[Abel 端点定理]
\label{thm:pei-abel-endpoint}
若 $\sum_{n=0}^{\infty}a_nR^n$ 收敛，其中 $R>0$，则 $\sum a_nx^n$ 在 $[0,R]$ 上一致收敛，且
\[
 \lim_{x\uparrow R}\sum_{n=0}^{\infty}a_nx^n
 =\sum_{n=0}^{\infty}a_nR^n.
\]
左端点的结论可将 $x$ 换为 $-x$ 后得到。一般展开中心可通过平移处理。
\end{theorem}
\begin{proof}
写成 $\sum (a_nR^n)(x/R)^n$。第一个因子组成收敛的数项级数，第二个因子对指标单调并在 $[0,R]$ 上一致有界，由一致 Abel 判别法得一致收敛；各项连续，再用一致收敛保连续性即可。
\end{proof}

\emph[caution]{端点处的级数收敛是此定理的假设，不能先取内部和函数的端点极限，再据此断言端点级数收敛。}例如 $\sum_{n=0}^{\infty}(-1)^nx^n=1/(1+x)$ 在 $x\uparrow1$ 时有极限 $1/2$，但端点级数 $1-1+1-1+\cdots$ 不收敛。另一方面，若非负系数幂级数在 $[0,1)$ 上有统一有限上界，则对每个有限 $N$ 令 $x\uparrow1$，可得 $\sum_{n=0}^{N}a_n$ 有共同上界，所以端点级数必收敛；非负性使这个逆向推理成立。

\begin{proposition}[一个 Tauber 型逆命题]
若 $f(x)=\sum_{n=1}^{\infty}a_nx^n$ 在 $|x|<1$ 收敛，$na_n\to0$，且 $\lim_{x\uparrow1}f(x)=S$ 存在，则 $\sum a_n=S$。
\end{proposition}
\begin{proof}
取 $x_N=1-1/N$（$N\geq2$），将部分和与 $f(x_N)$ 相减。因 $0\leq1-x_N^k\leq k/N$，有限部分满足
\[
 \left|\sum_{k=1}^{N}a_k(1-x_N^k)\right|
 \leq\frac1N\sum_{k=1}^{N}k|a_k|\longrightarrow0.
\]
记 $\eta_N=\sup_{k>N}k|a_k|\to0$，无穷尾部满足
\[
 \sum_{k>N}|a_k|x_N^k
 \leq\frac{\eta_N}{N}\sum_{k>N}x_N^k
 =\eta_Nx_N^{N+1}\leq\eta_N.
\]
故 $\sum_{k=1}^{N}a_k-f(x_N)\to0$，结合 $f(x_N)\to S$ 得证。
\end{proof}

\subsection{函数展开、系数递推与解析性}

常用的基本展开为
\begin{align*}
 e^x&=\sum_{n=0}^{\infty}\frac{x^n}{n!},&
 \sin x&=\sum_{n=0}^{\infty}\frac{(-1)^nx^{2n+1}}{(2n+1)!},&
 \cos x&=\sum_{n=0}^{\infty}\frac{(-1)^nx^{2n}}{(2n)!},\\
 \frac1{1-x}&=\sum_{n=0}^{\infty}x^n,&
 \log(1+x)&=\sum_{n=1}^{\infty}\frac{(-1)^{n-1}x^n}{n},&
 \arctan x&=\sum_{n=0}^{\infty}\frac{(-1)^nx^{2n+1}}{2n+1}.
\end{align*}
前三式在实轴上成立，几何级数在 $|x|<1$ 成立，对数展开在 $-1<x\leq1$ 成立，反正切展开在 $|x|\leq1$ 成立。几何级数在内闭区间上逐项积分得到后两式，再分别检验端点并用 Abel 定理确定端点和。对任意实数 $\alpha$，二项式展开为
\[
 (1+x)^\alpha=\sum_{n=0}^{\infty}\binom\alpha n x^n,\qquad |x|<1,
 \qquad
 \binom\alpha n=\frac{\alpha(\alpha-1)\cdots(\alpha-n+1)}{n!},
\]
其中 $\binom\alpha0=1$；非负整数 $\alpha$ 时该式成为有限多项式，其他情形的端点仍需另行讨论。

\begin{example}
求 $\log^2(1-x)$ 在原点的展开，并推导对 $x>0$ 成立的对数展开式。
\end{example}
\begin{solve}
在 $|x|<1$ 内绝对收敛，故可作 Cauchy 乘积。记 $H_n=\sum_{k=1}^{n}1/k$，则
\[
 \log^2(1-x)
 =\sum_{n=2}^{\infty}\left(\sum_{k=1}^{n-1}\frac1{k(n-k)}\right)x^n
 =\sum_{n=2}^{\infty}\frac{2H_{n-1}}n x^n.
\]
这里用 $1/[k(n-k)]=(1/k+1/(n-k))/n$ 合并系数。再令 $t=(x-1)/(x+1)$，对 $x>0$ 有 $|t|<1$、$x=(1+t)/(1-t)$，所以
\[
 \log x=\log(1+t)-\log(1-t)
 =2\sum_{n=0}^{\infty}\frac{t^{2n+1}}{2n+1}.
\]
当 $x$ 接近 $1$ 时，$|t|$ 小，后一式比直接在较大的变量上展开更适于近似计算。
\end{solve}

待定系数法须同时解决存在性与收敛范围。若令
\[
 \frac{x\sin\theta}{1-2x\cos\theta+x^2}=\sum_{n=0}^{\infty}c_nx^n,
\]
比较系数得到 $c_0=0$、$c_1=\sin\theta$、$c_{n+1}=2\cos\theta\,c_n-c_{n-1}$，由三角恒等式得 $c_n=\sin(n\theta)$。候选级数在 $|x|<1$ 绝对收敛，其与分母相乘确实等于分子，因而展开在该区间成立。一般地，可以先由微分方程得到系数递推，证明候选级数收敛并满足同一初值问题，再用解的唯一性确认其和；仅列出形式系数尚不是函数等于级数的证明。

\begin{definition}
函数 $f$ 在 $c$ 的某邻域内等于一个收敛幂级数，称其在 $c$ 处\emph{实解析（real analytic）}。由命题~\ref{prop:pei-power-interior}，该级数必为 $f$ 的 Taylor 级数。
\end{definition}

无穷次可微并不蕴含实解析。函数 $f(x)=e^{-1/x^2}$（$x\ne0$）、$f(0)=0$ 的每阶导数在 $x\ne0$ 时都是 $e^{-1/x^2}$ 与 $1/x$ 的某个多项式之积；任意固定幂次的 $1/|x|$ 都被此指数衰减压倒，逐阶取极限可得 $f^{(n)}(0)=0$。所以其 Taylor 级数恒为零，却不能表示邻域内的 $f$。反之，如果在某邻域有 $|f^{(n)}(x)|\leq MA^n$ 对全部 $n$ 成立，Taylor 余项被 $MA^{n+1}|x-c|^{n+1}/(n+1)!$ 控制，便能证明展开确实等于函数。

解析函数还具有局部唯一性：若一个正收敛半径的幂级数之和在趋向中心的互异点列上为零，则所有系数均为零。否则取最小的非零系数 $a_m$，可写为 $(x-c)^m(a_m+o(1))$，在充分小的去心邻域内不会为零，与点列矛盾。这种结论不能直接用于一般的光滑函数。

\subsection{借助参数求和与计算积分}

将数项级数嵌入幂级数后，分母中的 $n$ 常可通过积分处理，分子中的 $n$ 常可通过微分处理。运算在收敛区间内部完成，若最终需要端点值，再独立检查端点收敛并使用 Abel 定理。例如
\[
 \sum_{n=1}^{\infty}\frac{x^n}{n(n+1)}
 =1+\frac{1-x}{x}\log(1-x),\qquad |x|<1,
\]
其中 $x=0$ 处取连续延拓值 $0$。该式由 $1/[n(n+1)]=1/n-1/(n+1)$ 及对数展开相减得到，故端点 $x=1$ 的级数之和为 $1$。

\begin{example}
计算 $\sum_{k=1}^{\infty}(-1)^{k-1}/(3k-1)$。
\end{example}
\begin{solve}
原级数由 Leibniz 判别收敛。令
\[
 F(x)=\sum_{k=1}^{\infty}\frac{(-1)^{k-1}x^{3k-1}}{3k-1},\qquad |x|<1.
\]
逐项求导得 $F'(x)=x/(1+x^3)$，又 $F(0)=0$，所以由 Abel 端点定理，所求和为
\[
 \int_0^1\frac{x}{1+x^3}\,dx
 =\left[-\frac13\log(1+x)+\frac16\log(x^2-x+1)
 +\frac1{\sqrt3}\arctan\frac{2x-1}{\sqrt3}\right]_0^1
 =\frac{\pi}{3\sqrt3}-\frac13\log2.
 \qedhere
\]
\end{solve}

对带阶乘的和，算子 $D=x\,d/dx$ 特别方便：$D(x^n)=nx^n$，故
\[
 \sum_{n=0}^{\infty}\frac{n^kx^n}{n!}=D^k(e^x)=e^xP_k(x),
\]
其中 $k\geq1$，$P_0(x)=1$、$P_{k+1}(x)=x(P_k(x)+P_k'(x))$。递推可知 $P_k$ 为整数系数多项式，因此 $\sum_{n=0}^{\infty}n^k/n!=eP_k(1)$ 是 $e$ 的整数倍。这里原幂级数收敛半径为无穷大，逐项微分和取 $x=1$ 均有依据。

把被积函数展开为幂级数后，在内闭区间逐项积分通常没有困难，涉及奇端点时还需估计积分尾部。例如
\[
 \int_0^1\frac{\log x}{1-x}\,dx=-\sum_{n=1}^{\infty}\frac1{n^2},\qquad
 \int_0^1\frac{\log x}{1+x}\,dx=-\sum_{n=1}^{\infty}\frac{(-1)^{n-1}}{n^2}.
\]
第一式可分离 $\log x$ 后应用例~\ref{ex:pei-nonuniform-integration}；第二式有限几何展开后的余项积分绝对值不超过 $\int_0^1x^{N+1}|\log x|\,dx=1/(N+2)^2$。端点 $0$ 处以 $|\log x|$ 控制，第一式在 $1$ 处的被积函数有有限极限，故两个反常积分均存在。其具体值将在 Fourier 级数的求和例题中得到。

\subsection{不等式、误差界与端点导数}

非负系数展开常可逐项比较。例如对任意实数 $x$，
\[
 \cosh x=\sum_{n=0}^{\infty}\frac{x^{2n}}{(2n)!}
 \leq\sum_{n=0}^{\infty}\frac{x^{2n}}{2^nn!}=e^{x^2/2},
\]
因为 $(2n)!\geq2^nn!$。两边都绝对收敛且 $x^{2n}\geq0$，所以比较成立；等号仅在 $x=0$ 成立，因为 $n=2$ 项在 $x\ne0$ 时已严格不等。近似计算则应先指定余项控制方式：对 $0<t\leq1$，
\[
 \left|\arctan t-\sum_{n=0}^{N-1}\frac{(-1)^nt^{2n+1}}{2n+1}\right|
 \leq\frac{t^{2N+1}}{2N+1}.
\]
因此取项数须由目标误差反推，保留到某次幂不自动意味着对应的小数精度。

\begin{example}
\label{ex:pei-power-endpoint-derivative}
讨论 $f(x)=\sum_{n=1}^{\infty}x^n/[n^2\log(n+1)]$ 在 $[-1,1]$ 上的连续性与端点可导性。
\end{example}
\begin{solve}
原级数由 $\sum1/n^2$ 控制，在 $[-1,1]$ 上一致收敛，故 $f$ 连续。在内部
\[
 f'(x)=\sum_{n=1}^{\infty}\frac{x^{n-1}}{n\log(n+1)}.
\]
在 $[-1,0]$ 上，该导数级数由一致 Abel 判别法一致收敛，因为 $x=-1$ 处为收敛的交错级数；由逐项求导定理，$f$ 在 $-1$ 处存在有限右导数。在 $0<x<1$ 上，导数级数各项非负，且其在 $x=1$ 的部分和无界：$\sum1/[n\log(n+1)]$ 由积分比较发散。给定 $M>0$，先取有限部分和在 $1$ 处大于 $2M$，再令 $x$ 足够接近 $1$，即可得 $f'(x)>M$，故 $f'(x)\to+\infty$。对 $x<1$ 用微分中值定理，有
\[
 \frac{f(1)-f(x)}{1-x}=f'(\xi_x),\qquad x<\xi_x<1,
\]
因此左侧趋于 $+\infty$，$f$ 在 $1$ 处没有有限左导数。连续延伸到端点与导数能否延伸，是需要分别检查的问题。
\end{solve}

% 对应所给 PDF 第 131--162 页。
\section{Fourier 级数}
\label{sec:pei-fourier-series}

\subsection{正交系与 Fourier 系数}

在函数空间中，以积分 $\langle f,g\rangle=\int_a^bf(x)g(x)\,dx$ 代替向量的内积，便可讨论投影和正交展开。函数系 $(\phi_n)$ 为\emph{正交系（orthogonal system）}，是指每项平方积分非零，且 $m\ne n$ 时 $\langle\phi_m,\phi_n\rangle=0$。在 $[-\pi,\pi]$ 上，$1,\cos x,\sin x,\cos2x,\sin2x,\ldots$ 两两正交，常函数 $1$ 的平方积分为 $2\pi$，各正弦、余弦函数的平方积分为 $\pi$；这些结论可由积化和差公式和有限积分直接验证。

\begin{definition}
设实函数 $f$ 在 $[-\pi,\pi]$ 上 Riemann 可积。其\emph{Fourier 系数（Fourier coefficients）}定义为
\[
 a_n=\frac1\pi\int_{-\pi}^{\pi}f(x)\cos(nx)\,dx\quad(n\geq0),\qquad
 b_n=\frac1\pi\int_{-\pi}^{\pi}f(x)\sin(nx)\,dx\quad(n\geq1).
\]
对应的 Fourier 级数写为
\[
 f(x)\sim\frac{a_0}{2}+\sum_{n=1}^{\infty}(a_n\cos nx+b_n\sin nx).
\]
符号 $\sim$ 在此只表示系数由上述积分确定，不预先断言级数收敛，更不预先断言其和等于 $f(x)$。部分和记为 $S_Nf$。
\end{definition}

在 $[-L,L]$ 上，相应频率改为 $n\pi/L$、系数前的归一化因子改为 $1/L$，常数项仍为 $a_0/2$。一般区间可平移到这种形式；无论选取哪一段完整周期，系数计算均须使用同一周期和同一自变量。偶函数的 $b_n$ 全部为零，奇函数的 $a_n$ 全部为零；若再满足 $f(x+\pi)=-f(x)$，则常数项和全部偶数频率系数消失。这些对称性应在计算积分前检查。

\begin{example}
说明 Legendre 多项式及微分方程如何产生正交系。
\end{example}
\begin{solve}
Rodrigues 公式给出
\[
 P_n(x)=\frac1{2^nn!}\frac{d^n}{dx^n}(x^2-1)^n.
\]
对任意次数小于 $n$ 的多项式 $q$，分部积分 $n$ 次后，边界项为零，因为 $(x^2-1)^n$ 及其前 $n-1$ 阶导数在 $\pm1$ 都为零，而 $q^{(n)}=0$。故 $\int_{-1}^{1}P_nq=0$，从而 $P_n$ 两两正交。其首项系数为 $(2n)!/[2^n(n!)^2]$，再次分部积分可得
\[
 \int_{-1}^{1}P_n(x)^2\,dx
 =\frac{(2n)!}{2^{2n}(n!)^2}\int_{-1}^{1}(1-x^2)^n\,dx
 =\frac2{2n+1}.
\]
末一步中的积分由分部积分递推为 $2^{2n+1}(n!)^2/(2n+1)!$。更一般地，若 $y_m,y_n$ 满足
\[
 -(p y_j')'+q y_j=\lambda_jw y_j,
\]
其中 $p\in C^1([a,b])$ 且为正，$q,w$ 连续且 $w>0$，各解有使运算成立的二阶连续性，则将两式交叉相乘后相减、积分，得到
\[
 (\lambda_m-\lambda_n)\int_a^bw y_my_n\,dx
 =\bigl[p(y_my_n'-y_m'y_n)\bigr]_a^b.
\]
若共同边界条件使右侧为零，例如两端均为齐次 Dirichlet 条件，且 $\lambda_m\ne\lambda_n$，则两解关于权 $w$ 正交。正交性来自共同边界条件和不同参数值，不能只由“满足同类微分方程”推出。
\end{solve}

\subsection{系数衰减与点态收敛}

\begin{lemma}[Riemann--Lebesgue 引理]
若 $f$ 在有限区间 $[a,b]$ 上 Riemann 可积，则当 $|\lambda|\to\infty$ 时，
\[
 \int_a^bf(x)\cos(\lambda x)\,dx\longrightarrow0,\qquad
 \int_a^bf(x)\sin(\lambda x)\,dx\longrightarrow0.
\]
特别地，Fourier 系数 $a_n,b_n\to0$。
\end{lemma}
\begin{proof}
由 Riemann 可积性，可取阶梯函数 $g$ 使 $\int_a^b|f-g|<\varepsilon$。$f-g$ 的振荡积分绝对值不超过 $\varepsilon$；$g$ 只有有限个常值区间，每个区间上正弦或余弦的积分均为 $O(1/|\lambda|)$，所以其振荡积分趋于零。先固定 $g$ 再令频率增大，最后令 $\varepsilon\downarrow0$ 即得结论。
\end{proof}

更多正则性通常给出更快的系数衰减。例如周期函数满足 Hölder 条件 $|f(x+h)-f(x)|\leq C|h|^\alpha$（$0<\alpha\leq1$）时，有 $a_n,b_n=O(n^{-\alpha})$。令复系数 $\widehat f(n)=(2\pi)^{-1}\int_{-\pi}^{\pi}f(x)e^{-inx}\,dx$，利用平移 $h=\pi/n$ 使指数变号，可得
\[
 2|\widehat f(n)|\leq\frac1{2\pi}\int_{-\pi}^{\pi}|f(x+h)-f(x)|\,dx
 \leq C(\pi/n)^\alpha,
\]
再由 $a_n=2\operatorname{Re}\widehat f(n)$、$b_n=-2\operatorname{Im}\widehat f(n)$ 得到结论。系数趋于零只是必要条件，尚不能保证 Fourier 级数收敛。

\begin{theorem}[分段光滑函数的 Fourier 收敛定理]
\label{thm:pei-fourier-dirichlet}
设 $f$ 以 $2\pi$ 为周期，在一周期上分段 $C^1$，即有有限分割，使其在各开子区间上有连续导数，且函数及导数在分割点有有限单侧极限。则其 Fourier 级数在每点 $x$ 收敛于
\[
 \frac{f(x-0)+f(x+0)}2.
\]
在连续点上和为 $f(x)$；在周期接缝处须使用周期延拓后的左右极限。
\end{theorem}
\begin{proof}
记 Dirichlet 核
\[
 D_N(t)=\frac12+\sum_{n=1}^{N}\cos nt
 =\frac{\sin((N+1/2)t)}{2\sin(t/2)}.
\]
代入系数公式，有限求和与积分交换，得到
\[
 S_Nf(x)=\frac1\pi\int_{-\pi}^{\pi}f(x+t)D_N(t)\,dt,
 \qquad \frac1\pi\int_{-\pi}^{\pi}D_N(t)\,dt=1.
\]
令 $\ell=[f(x-0)+f(x+0)]/2$，利用核的偶性可写为
\[
 S_Nf(x)-\ell
 =\frac1\pi\int_0^\pi
 \frac{f(x+t)+f(x-t)-2\ell}{2\sin(t/2)}
 \sin((N+1/2)t)\,dt.
\]
分段光滑性使分子在 $t\downarrow0$ 时为 $O(t)$，故前面的商在零点附近有界，补充端点值后是分段连续的可积函数。由 Riemann--Lebesgue 引理，积分趋于零。
\end{proof}

在 $[0,L]$ 上求正弦展开，应先作奇延拓到 $[-L,L]$；求余弦展开则先作偶延拓。两种延拓产生不同的周期函数，所以同一个原区间上的函数可以有不同的三角展开；固定周期及原函数后，其 Fourier 系数则由积分唯一确定。端点的展开和由延拓后的单侧极限决定，不能只把原函数在孤立端点上的指定值代入。

\subsection{Fejér 平均、最小二乘与 Parseval 等式}

\begin{theorem}[Fejér 一致逼近]
\label{thm:pei-fejer}
若 $f$ 为连续的 $2\pi$ 周期函数，则其 Fourier 部分和的平均
\[
 \sigma_Nf=\frac{S_0f+S_1f+\cdots+S_Nf}{N+1}
\]
在实轴上一致收敛于 $f$。因此连续周期函数可由三角多项式一致逼近，即使原来的 Fourier 部分和不一定一致收敛。
\end{theorem}
\begin{proof}
有限几何和及取平均给出
\[
 \sigma_Nf(x)=\frac1{2\pi}\int_{-\pi}^{\pi}f(x-t)K_N(t)\,dt,
 \qquad
 K_N(t)=\frac1{N+1}\left(\frac{\sin((N+1)t/2)}{\sin(t/2)}\right)^2.
\]
核非负，其积分为 $2\pi$；在 $\delta\leq|t|\leq\pi$ 上，有 $K_N(t)\leq[(N+1)\sin^2(\delta/2)]^{-1}$。将误差积分分为 $|t|<\delta$ 和其余部分，第一部分由周期函数的一致连续性控制为任意小的 $\varepsilon$，第二部分不超过 $2\|f\|_\infty/[(N+1)\sin^2(\delta/2)]$。先取 $\delta$，再令 $N\to\infty$，即得到统一于 $x$ 的误差趋零。
\end{proof}

\begin{theorem}[最佳均方逼近与 Bessel 不等式]
\label{thm:pei-bessel}
设 $f$ 在 $[-\pi,\pi]$ 上 Riemann 可积，$a_n,b_n$ 为其 Fourier 系数。在所有次数不超过 $N$ 的三角多项式
\[
 T_N(x)=\frac{\alpha_0}{2}+\sum_{n=1}^{N}(\alpha_n\cos nx+\beta_n\sin nx)
\]
中，使\emph{均方误差（mean square error）} $\int_{-\pi}^{\pi}|f-T_N|^2$ 最小的唯一选择是 $T_N=S_Nf$。并有 Bessel 不等式
\begin{equation}
 \frac{a_0^2}{2}+\sum_{n=1}^{N}(a_n^2+b_n^2)
 \leq\frac1\pi\int_{-\pi}^{\pi}|f(x)|^2\,dx.
 \label{eq:pei-bessel}
\end{equation}
\end{theorem}
\begin{proof}
由于 $f-S_Nf$ 与所有次数不超过 $N$ 的三角多项式正交，平方展开后交叉项为零，得到
\begin{align*}
 \int_{-\pi}^{\pi}|f-T_N|^2
 &=\int_{-\pi}^{\pi}|f-S_Nf|^2
 +\frac\pi2(\alpha_0-a_0)^2\\
 &\quad+\pi\sum_{n=1}^{N}\bigl[(\alpha_n-a_n)^2+(\beta_n-b_n)^2\bigr].
\end{align*}
这证明最优性和唯一性。再展开第一项，得
\[
 \int_{-\pi}^{\pi}|f-S_Nf|^2
 =\int_{-\pi}^{\pi}|f|^2-\pi\left[\frac{a_0^2}{2}+\sum_{n=1}^{N}(a_n^2+b_n^2)\right]\geq0,
\]
即为所需不等式。
\end{proof}

\begin{theorem}[Parseval 等式]
\label{thm:pei-parseval}
对上述 Riemann 可积函数 $f$，有
\begin{equation}
 \frac1\pi\int_{-\pi}^{\pi}|f(x)|^2\,dx
 =\frac{a_0^2}{2}+\sum_{n=1}^{\infty}(a_n^2+b_n^2).
 \label{eq:pei-parseval}
\end{equation}
等价地，$\int_{-\pi}^{\pi}|f-S_Nf|^2\to0$，即 Fourier 部分和\emph{均方收敛（mean square convergence）}于 $f$。若 $g$ 的 Fourier 系数为 $\alpha_n,\beta_n$，则
\[
 \frac1\pi\int_{-\pi}^{\pi}f(x)g(x)\,dx
 =\frac{a_0\alpha_0}{2}+\sum_{n=1}^{\infty}(a_n\alpha_n+b_n\beta_n).
\]
\end{theorem}
\begin{proof}
先说明三角多项式在所需的均方意义下稠密。Riemann 可积函数有界，可用阶梯函数使平方积分误差任意小；再只在阶梯函数各跳点及周期接缝的充分小邻域内作线性连接，得到连续周期函数，并使新增平方积分误差任意小。由定理~\ref{thm:pei-fejer}，该连续函数可由三角多项式一致逼近，因而也可均方逼近。于是对每个 $\varepsilon>0$，存在三角多项式 $T$ 满足 $\int|f-T|^2<\varepsilon$。由最佳逼近性，当 $N$ 至少为 $T$ 的次数时，$\int|f-S_Nf|^2\leq\int|f-T|^2<\varepsilon$，结合误差恒等式即得 Parseval 等式。分别对 $f+g$ 与 $f-g$ 使用此式并相减，得到乘积公式；右侧由 Cauchy--Schwarz 不等式绝对收敛。
\end{proof}

均方收敛不等于逐点收敛或一致收敛，三者不可互相替换。若两个 Riemann 可积函数具有相同的 Fourier 系数，则对差函数使用式~\eqref{eq:pei-parseval} 可得其平方积分为零，等价于 $\int|f-g|=0$；如果两函数连续，便能进一步推出处处相等。没有连续性时，孤立点改值不影响系数，故不能据此确定每个点值。同样，系数趋于零不足以保证它们来自一个 Riemann 可积函数：三角级数 $\sum_{n=2}^{\infty}\sin(nx)/\log n$ 的系数虽然趋于零，但 $\sum1/(\log n)^2$ 发散，违反 Bessel 不等式。

\subsection{典型展开与级数求和}

\begin{example}
\label{ex:pei-fourier-sums}
求 $x$、$x^2$ 在 $[-\pi,\pi]$ 上的 Fourier 展开，并计算 $\sum1/n^2$、$\sum1/n^4$。
\end{example}
\begin{solve}
$x$ 为奇函数，分部积分得 $b_n=2(-1)^{n+1}/n$。周期延拓在 $\pm\pi$ 处有跳跃，故
\[
 x=2\sum_{n=1}^{\infty}\frac{(-1)^{n+1}}n\sin nx,\qquad -\pi<x<\pi,
\]
而在 $x=\pm\pi$ 处级数和为零。$x^2$ 为偶函数，系数为 $a_0=2\pi^2/3$、$a_n=4(-1)^n/n^2$，其周期延拓连续，因而
\begin{equation}
 x^2=\frac{\pi^2}{3}+4\sum_{n=1}^{\infty}\frac{(-1)^n}{n^2}\cos nx,
 \qquad -\pi\leq x\leq\pi.
 \label{eq:pei-x-square-fourier}
\end{equation}
取 $x=\pi$，得 $\sum_{n=1}^{\infty}1/n^2=\pi^2/6$；取 $x=0$，得 $\sum_{n=1}^{\infty}(-1)^{n-1}/n^2=\pi^2/12$，再扣除偶数项可得 $\sum_{k=0}^{\infty}1/(2k+1)^2=\pi^2/8$。对 $x^2$ 使用 Parseval 等式，得到
\[
 \frac1\pi\int_{-\pi}^{\pi}x^4\,dx
 =\frac{(2\pi^2/3)^2}{2}+16\sum_{n=1}^{\infty}\frac1{n^4},
\]
故 $\sum1/n^4=\pi^4/90$。这也确定了前节两个对数积分分别为 $-\pi^2/6$ 与 $-\pi^2/12$。
\end{solve}

同样由奇延拓或平移可得
\begin{equation}
 \sum_{n=1}^{\infty}\frac{\sin nx}{n}=\frac{\pi-x}{2},\qquad0<x<2\pi,
 \label{eq:pei-sawtooth-fourier}
\end{equation}
在周期端点级数和为零。在 $[0,\pi]$ 上把常函数 $1$ 作奇延拓，则
\[
 1=\frac4\pi\sum_{k=0}^{\infty}\frac{\sin((2k+1)x)}{2k+1},\qquad0<x<\pi,
\]
两端点的和仍为零。这些级数在避开跳点的内闭区间上一致收敛，却不能在整个开区间上一致收敛：若在开区间上一致 Cauchy，因各部分和在闭区间连续，其尾段估计可以取极限延至端点，从而会得到闭区间一致收敛，与和函数在端点的不连续矛盾。

对于 $|q|<1$，还可利用等比级数和对数的幂级数直接计算
\begin{align*}
 \frac{1-q^2}{1-2q\cos x+q^2}
 &=1+2\sum_{n=1}^{\infty}q^n\cos nx,\\
 \log(1-2q\cos x+q^2)
 &=-2\sum_{n=1}^{\infty}\frac{q^n}{n}\cos nx.
\end{align*}
两式对 $x\in\mathbb R$ 均绝对一致收敛。第二式可在复数 $qe^{ix}$ 的对数幂级数中取实部得到，因为 $|qe^{ix}|<1$ 且 $|1-qe^{ix}|^2=1-2q\cos x+q^2$。其常数项为零，故
\[
 \int_0^{2\pi}\log(1-2q\cos x+q^2)\,dx=0.
\]
若参数绝对值大于 $1$，先提出参数的平方再应用上述公式，积分变为 $4\pi\log|q|$；临界参数 $|q|=1$ 有对数奇点，需单独说明反常积分的存在及取极限依据。

\subsection{一致收敛与逐项运算的边界}

\begin{proposition}
\label{prop:pei-fourier-uniform}
若 $f$ 是连续的 $2\pi$ 周期函数，并在一周期上分段 $C^1$，则其 Fourier 级数在实轴上绝对一致收敛于 $f$。
\end{proposition}
\begin{proof}
分段积分并利用 $f$ 在分割点及周期接缝处连续，可得 $f'$ 的 Fourier 系数为
\[
 a_0'=0,\qquad a_n'=nb_n,\qquad b_n'=-na_n.
\]
这里带撇的 $a_n',b_n'$ 表示 $f'$ 的系数，并非把数列对指标求导。由 $f'$ 的 Bessel 不等式，$\sum n^2(a_n^2+b_n^2)<\infty$。Cauchy--Schwarz 不等式给出
\[
 \sum_{n=1}^{\infty}(|a_n|+|b_n|)
 \leq\left(\sum_{n=1}^{\infty}n^2(a_n^2+b_n^2)\right)^{1/2}
 \left(2\sum_{n=1}^{\infty}\frac1{n^2}\right)^{1/2}<\infty.
\]
故由 $M$ 判别法绝对一致收敛，再由点态收敛定理确定和为 $f$。
\end{proof}

Fourier 级数的积分比微分稳定。对任意 Riemann 可积的 $f$，即使尚未证明其 Fourier 级数逐点收敛，也有
\begin{equation}
 \int_0^xf(t)\,dt
 =\frac{a_0x}{2}+\sum_{n=1}^{\infty}
 \left[\frac{a_n}{n}\sin nx+\frac{b_n}{n}(1-\cos nx)\right],
 \quad -\pi\leq x\leq\pi.
 \label{eq:pei-fourier-integration}
\end{equation}
事实上，有限部分和积分后的误差满足
\[
 \sup_{-\pi\leq x\leq\pi}
 \left|\int_0^x(f-S_Nf)(t)\,dt\right|
 \leq\sqrt\pi\left(\int_{-\pi}^{\pi}|f-S_Nf|^2\right)^{1/2}\longrightarrow0,
\]
而右侧级数由 Bessel 不等式和 $\sum1/n^2<\infty$ 绝对一致收敛。相反，形式求导会给系数增加因子 $n$，可能彻底破坏收敛性；式~\eqref{eq:pei-sawtooth-fourier} 的和在内部可导，但形式导数为 $\sum\cos nx$，其通项不趋于零，不能逐项求导。

若 $f$ 连续周期且分段 $C^1$，其形式导数级数确实是 $f'$ 的 Fourier 级数，但还须判断后者是否收敛、收敛到哪个值。若进一步 $f'$ 也是连续周期且分段 $C^1$，命题~\ref{prop:pei-fourier-uniform} 保证导数级数一致收敛，再由定理~\ref{thm:pei-termwise-derivative} 才能把形式求导提升为实际的逐项求导。若 $f$ 本身在周期接缝处有跳跃，则分部积分会产生额外边界项，连系数关系也要相应修正。

\subsection{平移、相关积分与均方不等式的应用}

平移后的 Fourier 系数不必重新逐项积分。若 $f(x+t)$ 的系数记为 $A_n,B_n$，利用周期性换元与三角恒等式可得
\[
 A_0=a_0,\qquad
 A_n=a_n\cos nt+b_n\sin nt,\qquad
 B_n=b_n\cos nt-a_n\sin nt.
\]
令 $f$ 为连续周期函数，由 Parseval 的乘积公式得到其相关积分
\begin{equation}
 H(x)=\frac1\pi\int_{-\pi}^{\pi}f(t)f(t+x)\,dt
 =\frac{a_0^2}{2}+\sum_{n=1}^{\infty}(a_n^2+b_n^2)\cos nx.
 \label{eq:pei-fourier-correlation}
\end{equation}
右侧系数绝对可和，故在实轴上一致收敛；特别地 $H$ 为连续偶函数。取 $x=0$ 即回到 Parseval 等式。这个例子表明，积分中的平移可以转化为系数上的旋转，相关运算则把系数变成平方和。

\begin{proposition}[周期零均值函数的 Wirtinger 不等式]
设 $f\in C^1([0,L])$，$L>0$，$f(0)=f(L)$，且 $\int_0^Lf=0$。则
\[
 \int_0^L|f'(x)|^2\,dx
 \geq\left(\frac{2\pi}{L}\right)^2\int_0^L|f(x)|^2\,dx.
\]
等号成立当且仅当 $f(x)=A\cos(2\pi x/L)+B\sin(2\pi x/L)$，其中 $A,B$ 为实常数。
\end{proposition}
\begin{proof}
对 $L$ 周期展开，零均值使常数系数为零，分部积分使导数的系数增加频率因子 $2\pi n/L$。分别应用 Parseval 等式得到
\[
 \int_0^L|f|^2=\frac L2\sum_{n=1}^{\infty}(a_n^2+b_n^2),\qquad
 \int_0^L|f'|^2=\frac L2\sum_{n=1}^{\infty}
 \left(\frac{2\pi n}{L}\right)^2(a_n^2+b_n^2).
\]
逐项使用 $n^2\geq1$ 即得不等式。等号恰在全部 $n\geq2$ 的系数为零时成立，再由连续函数的系数唯一性得到所列函数形式。
\end{proof}

最后，振荡因子的平均也可用 Fourier 展开处理。由直接计算系数得到绝对一致收敛展开
\[
 |\sin x|=\frac2\pi-\frac4\pi\sum_{k=1}^{\infty}\frac{\cos(2kx)}{4k^2-1}.
\]
若 $f$ 在有限区间 $[a,b]$ 上 Riemann 可积，则
\[
 \lim_{n\to\infty}\int_a^bf(x)|\sin nx|\,dx
 =\frac2\pi\int_a^bf(x)\,dx.
\]
证明时先用展开式代入并逐项积分；第 $k$ 项积分绝对值被 $\int_a^b|f|/(4k^2-1)$ 控制，对 $n$ 无关且关于 $k$ 可和，因此可以逐项令 $n\to\infty$，再由 Riemann--Lebesgue 引理消去所有非常数频率。这里起决定作用的仍是统一可和的控制，而不是单独知道每个振荡积分趋于零。


\chapter{多元函数微分学}
\label{chap:pei-multivariable-differentiation}
% 来源：数学分析/textbooks/数分习题（分章原文）/数分习题-6.pdf，PDF 第 1--107 页。按主题整理学习笔记，并非逐题抄录。

本章研究从欧氏空间到实数或另一欧氏空间的映射。与一元情形相比，自变量可以沿不同方向趋近同一点，因而沿坐标轴得到的信息未必能控制整个邻域；微分也应当理解为一个同时适用于所有小增量的线性近似。本章沿教材第六章的次序，从极限与连续出发，讨论偏导数、全微分、Taylor 公式与极值、隐函数及函数相关，最后归纳方向导数与梯度，并把计算中需要的条件与容易混淆的逆命题一并说明。

\section{欧氏空间、多元函数的极限与连续}
\label{sec:pei-euclidean-limits-continuity}

\subsection{范数、距离与集合的基本性质}

在 $\mathbb R^m$ 中，对 $x=(x_1,\ldots,x_m)$ 定义欧氏范数（Euclidean norm）$\|x\|=(\sum_{i=1}^m x_i^2)^{1/2}$，距离为 $d(x,y)=\|x-y\|$。逐项估计与 Cauchy--Schwarz 不等式给出
\begin{equation}
\label{eq:pei-coordinate-norm-bounds}
\max_{1\le i\le m}|x_i|\le\|x\|\le\sum_{i=1}^m|x_i|\le\sqrt m\,\|x\|.
\end{equation}
因此向量收敛等价于各坐标分别收敛，使用球形邻域或各坐标同时受控的矩形邻域定义极限，得到的概念相同。

\begin{definition}
设 $E\subset\mathbb R^m$。若某个开球 $B(a,r)=\{x:\|x-a\|<r\}$ 包含于 $E$，称 $a$ 为 $E$ 的内点；若每个 $B(a,r)$ 都与 $E$ 和 $\mathbb R^m\setminus E$ 相交，称 $a$ 为边界点。若每个去心球 $B(a,r)\setminus\{a\}$ 都与 $E$ 相交，称 $a$ 为聚点。$E$ 的闭包记为 $\overline E$，边界记为 $\partial E$；开集是每一点均为内点的集合，闭集是包含全部聚点的集合。
\end{definition}

从定义可得 $\partial E=\overline E\cap\overline{\mathbb R^m\setminus E}$，故边界总是闭集。有限维欧氏空间中的有界闭集是紧集（compact set）：任意序列有收敛到该集内一点的子列，等价地，任意开覆盖有有限子覆盖。后面由局部控制获得统一控制，或由连续性获得最大、最小值时，真正使用的是紧性。

有限维空间中还可换用其他范数。若 $N$ 是满足正定性、绝对齐次性与三角不等式的范数，沿坐标分解可得 $N(x)\le\sum_i|x_i|N(e_i)\le C\|x\|$，从而 $|N(x)-N(y)|\le C\|x-y\|$，所以 $N$ 关于欧氏距离连续。它在单位球面上的最小值 $c$ 严格为正，齐次性遂给出 $c\|x\|\le N(x)\le C\|x\|$。因此有限维空间中的任意两种范数给出相同的收敛、连续和一致连续概念。

\begin{proposition}
\label{prop:pei-distance-function}
对非空集合 $E\subset\mathbb R^m$，定义点到集合的距离 $d(x,E)=\inf_{z\in E}\|x-z\|$。则
\[
|d(x,E)-d(y,E)|\le\|x-y\|,\qquad d(x,E)=0\ \Longleftrightarrow\ x\in\overline E.
\]
特别地，距离函数在整个空间上满足 Lipschitz 条件，因而一致连续；若 $E$ 闭，则 $d(x,E)=0$ 等价于 $x\in E$。
\end{proposition}
\begin{proof}
对每个 $z\in E$，由三角不等式有 $d(x,E)\le\|x-y\|+\|y-z\|$，对 $z$ 取下确界，再交换 $x,y$ 即得第一式。若 $d(x,E)=0$，每个正整数 $n$ 都可选 $z_n\in E$ 使 $\|x-z_n\|<1/n$，故 $x\in\overline E$；反向结论直接由闭包的邻域定义得到。
\end{proof}

\begin{example}
设非空紧集 $K$ 包含于开集 $U\subset\mathbb R^m$，且 $U\ne\mathbb R^m$。构造连续函数 $\varphi:\mathbb R^m\to[0,1]$，使其在 $K$ 上等于 $1$、在 $\mathbb R^m\setminus U$ 上等于 $0$。
\end{example}
\begin{solve}
令
\[
\varphi(x)=\frac{d(x,\mathbb R^m\setminus U)}{d(x,K)+d(x,\mathbb R^m\setminus U)}.
\]
两个闭集 $K$ 与 $\mathbb R^m\setminus U$ 不相交，由命题~\ref{prop:pei-distance-function}，分母在每一点都为正，所以该式定义连续函数并具有所需取值。若还要在 $K$ 与 $U$ 的外部之间留出统一宽度，可将连续正函数 $d(x,\mathbb R^m\setminus U)$ 限制在紧集 $K$ 上取最小值，得到 $d(K,\mathbb R^m\setminus U)>0$。
\end{solve}

\subsection{重极限、路径与极坐标的一致估计}

\begin{definition}
设 $a$ 是 $D\subset\mathbb R^m$ 的聚点，$f:D\to\mathbb R^p$。若对任意 $\varepsilon>0$，存在 $\delta>0$，使得对所有 $x\in D$，$0<\|x-a\|<\delta$ 蕴含 $\|f(x)-A\|<\varepsilon$，则称 $\lim_{x\to a,\,x\in D}f(x)=A$。当 $m\ge2$ 时，这也称为重极限或全面极限（joint limit）。
\end{definition}

该定义等价于序列判别法：每个满足 $x_n\in D\setminus\{a\}$、$x_n\to a$ 的序列都应满足 $f(x_n)\to A$。必要性是把极限定义应用于序列；若充分性不成立，可对某个固定的 $\varepsilon_0>0$，逐次在 $1/n$ 邻域中选出 $\|f(x_n)-A\|\ge\varepsilon_0$ 的点，得到矛盾。因此，找出两条路径或两个序列给出不同极限，就能否定重极限；验证若干条路径则不足以证明它存在。

\begin{example}
\label{ex:pei-radial-limit-failure}
证明函数 $f(x,y)=x^2y/(x^4+y^2)$ 在原点的每条固定射线上均趋于 $0$，但重极限不存在。
\end{example}
\begin{solve}
在直线 $y=kx$ 上，$k\ne0$ 时有 $f(x,kx)=kx/(x^2+k^2)\to0$；在两条坐标轴上函数恒为 $0$。但沿 $y=x^2$ 有 $f(x,x^2)=1/2$。极坐标把函数写为 $F(r,\theta)=f(r\cos\theta,r\sin\theta)$，每条射线上的收敛只表示对每个固定 $\theta$ 有 $F(r,\theta)\to0$；证明重极限需要一个同时控制所有角度的半径，即
\[
\sup_{0\le\theta<2\pi}|F(r,\theta)|\longrightarrow0\qquad(r\downarrow0).
\]
这里的非径向路径说明该一致收敛条件失败。
\end{solve}

\emph[tip]{存在性通常由统一上界证明，不存在性通常由路径反例证明。}例如
\[
0\le\frac{x^2y^2}{x^2+y^2}\le\frac{x^2+y^2}{4}\longrightarrow0.
\]
更一般地，若 $|f(r\cos\theta,r\sin\theta)-A|\le Cr^\alpha$ 对所有角度成立，其中 $\alpha>0$、$C$ 与角度无关，就得到重极限 $A$。若函数是正次数齐次函数，即 $f(tx)=t^\alpha f(x)$（$t>0$、$\alpha>0$），且在单位球面上有界，也可写成 $f(x)=\|x\|^\alpha f(x/\|x\|)$，从而得出 $f(x)\to0$。

\begin{proposition}
设 $F(r,\theta)$ 对每个固定角度满足 $F(r,\theta)\to A$（$r\downarrow0$），并且这一族角度函数在圆周上等度连续（equicontinuous）：对任意 $\varepsilon>0$，存在与 $r$ 无关的 $\eta>0$，使圆周角距离小于 $\eta$ 时 $|F(r,\theta)-F(r,\phi)|<\varepsilon$，其中 $0<r<r_0$。则 $F(r,\theta)\to A$ 关于角度一致，因而相应的二元函数在原点具有重极限 $A$。
\end{proposition}
\begin{proof}
利用圆周紧性，选有限个角度 $\theta_1,\ldots,\theta_N$，使每个角度距某个 $\theta_j$ 小于 $\eta$。由这有限条射线上的收敛，可取同一个 $\delta>0$ 使 $0<r<\delta$ 时所有 $|F(r,\theta_j)-A|<\varepsilon$。任意角度的误差均不超过 $2\varepsilon$，这就给出所需的一致估计。
\end{proof}

\subsection{累次极限与交换次序的条件}

累次极限（iterated limit）$\lim_{x\to a}\lim_{y\to b}f(x,y)$ 要先在固定 $x$ 时求内层极限，其存在性与重极限的存在性互不等价。例如 $f(x,y)=xy/(x^2+y^2)$ 的两个累次极限都是 $0$，重极限却不存在；而在 $xy\ne0$ 时定义
\[
f(x,y)=x\sin\frac1y+y\sin\frac1x,
\]
并在坐标轴上补为 $0$，则 $|f(x,y)|\le|x|+|y|$ 保证重极限为 $0$，但固定非零 $x$ 或非零 $y$ 时，相应的内层极限不存在。若重极限为 $A$，且 $g(x)=\lim_{y\to b}f(x,y)$ 在 $a$ 的某个去心邻域中存在，则将重极限估计中的 $y$ 趋于 $b$ 可得 $|g(x)-A|\le\varepsilon$，故相应的累次极限存在并等于 $A$。

\begin{theorem}
\label{thm:pei-interchange-limits}
设 $f$ 在 $(a,b)$ 附近的去心乘积邻域内定义，$g(x)=\lim_{y\to b}f(x,y)$ 对每个 $x$ 存在，且 $f(x,y)\to h(y)$（$x\to a$）关于该邻域内的 $y$ 一致。则 $\lim_{x\to a}g(x)$ 与 $\lim_{y\to b}h(y)$ 都存在且相等。
\end{theorem}
\begin{proof}
由一致收敛，对任意 $\varepsilon>0$，当 $x,x'$ 均充分接近 $a$ 时，对所有允许的 $y$ 有
\[
|f(x,y)-f(x',y)|\le|f(x,y)-h(y)|+|h(y)-f(x',y)|<2\varepsilon.
\]
令 $y\to b$，得到 $|g(x)-g(x')|\le2\varepsilon$。由实数的 Cauchy 极限准则，$g(x)\to A$。再固定一个充分接近 $a$ 的 $x_*$，使 $|g(x_*)-A|<\varepsilon$ 且 $|f(x_*,y)-h(y)|<\varepsilon$ 对所有 $y$ 成立。最后令 $y$ 足够接近 $b$，用
\[
|h(y)-A|\le|h(y)-f(x_*,y)|+|f(x_*,y)-g(x_*)|+|g(x_*)-A|
\]
得到误差小于 $3\varepsilon$，故 $h(y)\to A$。
\end{proof}

\subsection{连续、分别连续与紧集上的统一控制}

$f:D\to\mathbb R^p$ 在 $a\in D$ 连续，指 $\lim_{x\to a,\,x\in D}f(x)=f(a)$；在孤立点该条件自动成立。连续也等价于保持序列极限，或任意开集的逆像都是 $D$ 中的相对开集。后一描述中的“相对”不可省略：当定义域不是开集时，逆像未必是整个 $\mathbb R^m$ 的开集。向量值映射连续等价于各分量连续。

二元函数联合连续必然分别对每个变量连续，反向不成立；将 $xy/(x^2+y^2)$ 在原点补为 $0$ 就给出反例。若在原点附近有统一的 Lipschitz 估计 $|f(x,y)-f(x,y')|\le L|y-y'|$，且 $x\mapsto f(x,0)$ 在 $0$ 连续，则
\[
|f(x,y)-f(0,0)|\le L|y|+|f(x,0)-f(0,0)|\longrightarrow0.
\]
另一充分条件是分别连续，且对 $y$ 单调，单调方向在所讨论矩形内相同：先以 $y_0\pm\eta$ 处的函数值夹住 $f(x,y)$，再利用这两条水平线上的 $x$ 连续性和 $x_0$ 处的 $y$ 连续性，即可证明在 $(x_0,y_0)$ 联合连续。

\begin{proposition}
\label{prop:pei-maximum-continuity}
若 $K\subset\mathbb R^m$、$L\subset\mathbb R^n$ 为非空紧集，$f$ 在 $K\times L$ 上连续，则 $M(y)=\max_{x\in K}f(x,y)$ 与 $m(y)=\min_{x\in K}f(x,y)$ 都在 $L$ 上连续。
\end{proposition}
\begin{proof}
极值存在由紧性保证。对任意 $y,y'\in L$，逐点比较并取最大值得到
\[
|M(y)-M(y')|\le\max_{x\in K}|f(x,y)-f(x,y')|.
\]
$f$ 在紧集 $K\times L$ 上一致连续，所以右端随 $\|y-y'\|\to0$ 而趋于 $0$。最小值的结论应用于 $-f$ 即得。
\end{proof}

确界与下确界的交换一般只有不等式。对任意实值有界函数 $f$，由 $\inf_y f(x,y)\le f(x,y)\le\sup_x f(x,y)$，先后取确界可得
\[
\sup_x\inf_y f(x,y)\le\inf_y\sup_x f(x,y).
\]
即使 $f$ 连续也未必取等号，例如 $f(x,y)=(x-y)^2$、$x,y\in[0,1]$ 时两边分别为 $0$ 与 $1/4$。若存在鞍点 $(x_*,y_*)$ 满足 $f(x,y_*)\le f(x_*,y_*)\le f(x_*,y)$，则两边均等于该鞍点值；若 $f$ 对 $x$ 在整个区间上一致单调，也可把 $\sup_x$ 固定为同一个端点取值，直接验证等号。

\begin{remark}
若实值函数 $f$ 在 $a$ 的某邻域有界，令 $M(r)=\sup_{x\in B(a,r)}f(x)$、$m(r)=\inf_{x\in B(a,r)}f(x)$。随 $r\downarrow0$，非负差值 $M(r)-m(r)$ 单调不增，故有有限极限 $\omega_f(a)$，称为 $f$ 在 $a$ 处的振幅（oscillation）。$f$ 在 $a$ 连续当且仅当 $\omega_f(a)=0$：连续性给出邻域内所有值都接近 $f(a)$；反向利用 $m(r)\le f(a),f(x)\le M(r)$，把函数值之差控制在振幅内。
\end{remark}

\begin{theorem}
\label{thm:pei-uniform-extension}
若 $f:E\to\mathbb R^p$ 一致连续，则它能唯一地一致连续延拓到 $\overline E$。特别地，若 $D$ 是有界开集，$f:D\to\mathbb R^p$ 连续，则 $f$ 在 $D$ 上一致连续，当且仅当对每个 $a\in\partial D$，沿 $D$ 趋于 $a$ 的有限极限存在。
\end{theorem}
\begin{proof}
对 $a\in\overline E$ 取 $x_n\in E$、$x_n\to a$。一致连续性使 $f(x_n)$ 为 Cauchy 列；比较两种选点序列可知其极限与选择无关，于是定义延拓值。对 $a,b\in\overline E$，用 $E$ 内趋于它们的序列传递原来的一致连续估计，即得延拓仍一致连续，唯一性由 $E$ 在闭包中稠密得到。

对后一结论，必要性来自上述延拓。充分性需检查边界值彼此相容：按给定边界极限补定义 $\overline f$。若 $a_n\in\overline D$、$a_n\to a$，可在 $D$ 中选 $z_n$ 使 $\|z_n-a_n\|<1/n$ 且 $\|f(z_n)-\overline f(a_n)\|<1/n$，故 $z_n\to a$，由内部连续性或 $a$ 处的边界极限可知 $\overline f(a_n)\to\overline f(a)$。所以延拓在紧集 $\overline D$ 上连续，从而一致连续。
\end{proof}

\section{多元函数的偏导数与可微性}
\label{sec:pei-partials-differentiability}

\subsection{偏导数与全微分分别控制什么}

对定义在 $a\in\mathbb R^n$ 邻域内的实值函数，偏导数（partial derivative）定义为
\[
f_{x_i}(a)=\lim_{t\to0}\frac{f(a+te_i)-f(a)}{t},
\]
其中 $e_i$ 是第 $i$ 个坐标单位向量，其余坐标保持不变。分段定义函数在分界点的偏导必须回到这个定义；不能把只在去心邻域成立的求导公式直接代入分界点。下文约定 $f_{xy}=\partial_y(\partial_x f)$，先求 $x$ 偏导，再求 $y$ 偏导；在未核实条件前，不交换求导次序。

\begin{definition}
设 $f:U\subset\mathbb R^n\to\mathbb R^m$，$U$ 开，$a\in U$。若存在线性映射 $A:\mathbb R^n\to\mathbb R^m$，使
\begin{equation}
\label{eq:pei-frechet-differentiability}
f(a+h)-f(a)=Ah+o(\|h\|)\qquad(h\to0),
\end{equation}
则称 $f$ 在 $a$ 可微（differentiable），$A$ 称为导数或微分，记为 $Df(a)$；增量 $h$ 的全微分（total differential）为 $df(a)[h]=Df(a)h$。对实值函数，梯度（gradient）记为 $\nabla f(a)=(f_{x_1}(a),\ldots,f_{x_n}(a))$，因而 $Df(a)h=\nabla f(a)\cdot h=\sum_i f_{x_i}(a)h_i$；对向量值函数，其矩阵是 Jacobi 矩阵（Jacobian matrix）$J_f(a)=(\partial f_i/\partial x_j)_{m\times n}$。
\end{definition}

线性近似若存在则唯一：若 $Ah$ 与 $Bh$ 均满足式~\eqref{eq:pei-frechet-differentiability}，令 $h=tv$ 并除以 $t$，可得 $(A-B)v=0$ 对每个 $v$ 成立。可微蕴含连续和全部偏导数存在，但偏导数存在仅检查有限个方向，不能推出可微。若偏导数已求出，检验可微的直接标准是
\begin{equation}
\label{eq:pei-differentiability-test}
\frac{f(a+h)-f(a)-\sum_i f_{x_i}(a)h_i}{\|h\|}\longrightarrow0.
\end{equation}

\begin{theorem}
\label{thm:pei-continuous-partials}
若 $f$ 的一阶偏导数在 $a$ 的某邻域存在，且都在 $a$ 连续，则 $f$ 在 $a$ 可微。二元情形还可减弱为：$f_x$ 在邻域内存在且在 $a=(a_1,a_2)$ 连续，而 $f_y(a)$ 存在。
\end{theorem}
\begin{proof}
证明二元的较弱结论。将增量拆为两条坐标线段，并在第一条线上用一元中值定理，得到
\begin{align*}
f(a_1+h,a_2+k)-f(a_1,a_2)
&=f_x(a_1+\theta h,a_2+k)h\\
&\quad+f_y(a_1,a_2)k+o(|k|),\qquad 0<\theta<1.
\end{align*}
$f_x$ 在 $a$ 连续，故把第一项替换为 $f_x(a)h$ 的误差是 $o(\sqrt{h^2+k^2})$；第二个余项也具有该性质。高维结论沿各坐标逐段拆分，用同样估计即可。
\end{proof}

\begin{example}
\label{ex:pei-bounded-partials-not-differentiable}
令 $f(x,y)=xy/\sqrt{x^2+y^2}$（$(x,y)\ne(0,0)$），$f(0,0)=0$。证明它连续，两个偏导数存在且有界，但在原点不可微。
\end{example}
\begin{solve}
$|f(x,y)|\le\sqrt{x^2+y^2}/2$ 保证原点连续，沿两轴的定义给出 $f_x(0,0)=f_y(0,0)=0$。在非原点处，
\[
f_x=\frac{y^3}{(x^2+y^2)^{3/2}},\qquad f_y=\frac{x^3}{(x^2+y^2)^{3/2}},
\]
所以绝对值均不超过 $1$。但沿 $x=y=t\ne0$，$f(t,t)/\sqrt{2t^2}=1/2$，不满足式~\eqref{eq:pei-differentiability-test}。这同时说明偏导有界所带来的局部 Lipschitz 控制仍不足以保证某一点可微。
\end{solve}

\begin{example}
在原点补为 $0$ 的函数
\[
f_\alpha(x,y)=(x^2+y^2)^\alpha\sin\frac1{\sqrt{x^2+y^2}},\qquad (x,y)\ne(0,0),\quad\alpha\in\mathbb R,
\]
在原点可微的充要条件是什么？
\end{example}
\begin{solve}
沿 $y=0$，偏导定义中的差商为 $|t|^{2\alpha}\sin(1/|t|)/t$，其有限极限存在当且仅当 $2\alpha-1>0$，此时为 $0$；必要性可分别沿正弦等于 $1$ 和 $-1$ 的正数序列验证。若 $\alpha>1/2$，两个原点偏导均为 $0$，且以 $r=\sqrt{x^2+y^2}$ 记，有 $|f_\alpha|/r\le r^{2\alpha-1}\to0$，所以该条件也充分。这里先确定线性项，再比较余项与距离的阶数，避免只凭连续性判断可微。
\end{solve}

\subsection{复合函数、隐式求导与齐次函数}

\begin{theorem}
\label{thm:pei-chain-rule}
若 $g:\mathbb R^k\to\mathbb R^n$ 在 $a$ 可微，$f:\mathbb R^n\to\mathbb R^m$ 在 $g(a)$ 可微，且复合映射在相应邻域内有定义，则 $f\circ g$ 在 $a$ 可微，且
\[
D(f\circ g)(a)=Df(g(a))Dg(a).
\]
\end{theorem}
\begin{proof}
写 $g(a+h)-g(a)=Bh+o(\|h\|)=v$，则 $\|v\|=O(\|h\|)$。代入 $f(g(a)+v)-f(g(a))=Av+o(\|v\|)$，得到 $ABh+o(\|h\|)$。这也解释了链式法则需要对整个增量的控制，不能仅凭各偏导存在就无条件套用。
\end{proof}

实值函数 $z=f(x,y)$、$x=x(s,t)$、$y=y(s,t)$ 的链式法则为 $z_s=f_xx_s+f_yy_s$、$z_t=f_xx_t+f_yy_t$，其中 $f_x,f_y$ 在 $(x(s,t),y(s,t))$ 处取值。二次求导时，这些偏导仍是复合函数。例如 $F(t)=f(t,t^2)$ 且 $f\in C^2$ 时，
\[
F'(t)=f_x+2tf_y,\qquad F''(t)=f_{xx}+4tf_{xy}+4t^2f_{yy}+2f_y,
\]
右端均在 $(t,t^2)$ 处取值。末项来自内层函数 $t^2$ 的二阶导数，不能遗漏。对恒等式 $G(x,y,z(x,y))=0$，若 $G\in C^2$ 且 $G_z\ne0$，则
\[
z_x=-\frac{G_x}{G_z},\qquad z_{xx}=-\frac{G_{xx}+2G_{xz}z_x+G_{zz}z_x^2}{G_z}.
\]
这些公式在原隐式关系成立的点上使用；局部存在与唯一性由第~\ref{sec:pei-implicit-dependence} 节的定理保证。

\begin{proposition}
\label{prop:pei-euler-homogeneous}
设 $C\subset\mathbb R^n$ 为开锥域，即 $x\in C$、$t>0$ 蕴含 $tx\in C$，且 $f\in C^1(C)$。则 $f$ 是 $q$ 次齐次函数（homogeneous function），即 $f(tx)=t^qf(x)$，当且仅当
\begin{equation}
\label{eq:pei-euler-identity}
\sum_{i=1}^n x_i f_{x_i}(x)=qf(x)\qquad(x\in C).
\end{equation}
\end{proposition}
\begin{proof}
必要性是对 $t\mapsto f(tx)$ 求导后令 $t=1$。反过来，固定 $x\in C$，对 $t^{-q}f(tx)$ 求导，利用在 $tx$ 处成立的式~\eqref{eq:pei-euler-identity} 得导数为零，所以它恒等于 $f(x)$。
\end{proof}

若齐次函数 $f\in C^k(C)$，则对 $t\mapsto f(tx)=t^qf(x)$ 求 $k$ 次导数，并令 $t=1$，得到
\[
D^kf(x)[x,\ldots,x]=q(q-1)\cdots(q-k+1)f(x).
\]
二元形式的左端为 $\sum_{j=0}^k\binom{k}{j}x^jy^{k-j}\partial_x^j\partial_y^{k-j}f(x,y)$。此处 $x,y$ 是方向向量的固定分量；若把它简写为 $\bigl(x\partial_x+y\partial_y\bigr)^kf$，必须声明系数不再参与求导，否则所得算符幂会产生额外项。

\subsection{混合偏导数与矩形差商}

\begin{theorem}
\label{thm:pei-mixed-partials}
若 $f_x,f_y,f_{xy},f_{yx}$ 在 $(a,b)$ 的某邻域存在，且其中一个混合偏导数在 $(a,b)$ 连续，则 $f_{xy}(a,b)=f_{yx}(a,b)$。
\end{theorem}
\begin{proof}
设 $f_{xy}$ 连续。对 $h,k\ne0$ 引入矩形差商
\[
Q(h,k)=\frac{f(a+h,b+k)-f(a+h,b)-f(a,b+k)+f(a,b)}{hk}.
\]
先对 $x$ 再对 $y$ 使用一元中值定理，得 $Q(h,k)=f_{xy}(a+\theta h,b+\eta k)$，所以其重极限为 $f_{xy}(a,b)$。另一方面，先令 $k\to0$ 再令 $h\to0$，按定义得到 $f_{yx}(a,b)$。由重极限与已存在的累次极限的关系，两者相等。
\end{proof}

\begin{example}
令 $f(x,y)=xy(x^2-y^2)/(x^2+y^2)$（非原点），$f(0,0)=0$。说明混合偏导存在仍不保证可以交换次序。
\end{example}
\begin{solve}
按定义有 $f_x(0,y)=-y$、$f_y(x,0)=x$，包括原点处两者均为 $0$。所以 $f_{xy}(0,0)=-1$、$f_{yx}(0,0)=1$。差异不违反定理~\ref{thm:pei-mixed-partials}，因为连续性条件不成立。矩形差商的两个累次极限不同，正是次序不能交换的具体表现。
\end{solve}

\subsection{微分方程中的变量替换}

对偏微分方程换元时，应先确定哪些是新自变量、哪些是因变量，检查变换的局部可逆性，再逐阶建立导数关系。以 $x=r\cos\theta$、$y=r\sin\theta$、$r>0$ 为例，若 $U(r,\theta)=u(x,y)$，则
\[
\partial_x=\cos\theta\,\partial_r-\frac{\sin\theta}{r}\partial_\theta,\qquad
\partial_y=\sin\theta\,\partial_r+\frac{\cos\theta}{r}\partial_\theta.
\]
这些算符的系数随 $r,\theta$ 变化，二次作用时也必须求导，不能只把右端当作常系数一次式平方。由链式法则可得
\begin{align*}
U_{rr}&=u_{xx}\cos^2\theta+2u_{xy}\sin\theta\cos\theta+u_{yy}\sin^2\theta,\\
U_{\theta\theta}&=r^2\bigl(u_{xx}\sin^2\theta-2u_{xy}\sin\theta\cos\theta+u_{yy}\cos^2\theta\bigr)-rU_r.
\end{align*}
相加消去混合项，得到平面 Laplace 算子（Laplacian）的极坐标表达式
\begin{equation}
\label{eq:pei-polar-laplacian}
u_{xx}+u_{yy}=U_{rr}+\frac1rU_r+\frac1{r^2}U_{\theta\theta},\qquad r>0.
\end{equation}

\begin{example}
\label{ex:pei-quadratic-conformal-change}
设 $u\in C^2$，取 $\xi=x^2-y^2$、$\eta=2xy$，写 $u(x,y)=v(\xi,\eta)$。把方程 $u_{xx}+u_{yy}=0$ 化为新变量方程。
\end{example}
\begin{solve}
一次求导得 $u_x=2xv_\xi+2yv_\eta$、$u_y=-2yv_\xi+2xv_\eta$；再次求导并相加，得到
\[
u_{xx}+u_{yy}=4(x^2+y^2)(v_{\xi\xi}+v_{\eta\eta}).
\]
所以在非原点的局部坐标片中，新方程为 $v_{\xi\xi}+v_{\eta\eta}=0$。换元的 Jacobi 行列式为 $4(x^2+y^2)$，在原点为零；此外 $(x,y)$ 与 $(-x,-y)$ 具有相同的像，故局部可逆也不能理解为整个平面上的一一对应。
\end{solve}

对线性换元 $x=Q\xi$，若 $Q$ 是正交矩阵，则梯度长度与 Laplace 算子保持不变。事实上 $\nabla_\xi(u\circ Q)=Q^{\mathsf T}\nabla_xu$，二阶导数矩阵按 $H_\xi=Q^{\mathsf T}H_xQ$ 变换，利用 $QQ^{\mathsf T}=I$ 取迹即可。若因变量也发生变化，如 $u=e^\psi v$，还应在链式法则中同时使用乘积法则；只有把全部导数化为新变量表达式后，才代回原方程化简。

\section{多元 Taylor 公式、凸函数、几何应用与极值}
\label{sec:pei-taylor-convexity-extrema}

\subsection{沿线段化为一元问题与展开式的唯一性}

设 $a$ 与 $a+h$ 之间的线段位于函数定义域中，作 $g(t)=f(a+th)$，$0\le t\le1$。其中 $D^jf(x)$ 表示 $j$ 阶导数所定义的 $j$ 重线性映射。链式法则给出 $g^{(j)}(t)=D^jf(a+th)[h,\ldots,h]$，从而可将一元 Taylor 公式应用于 $g$。例如 $f\in C^2$ 时，在 $a=(a_1,a_2)$ 附近有
\begin{equation}
\label{eq:pei-second-order-taylor}
f(a+(h,k))=f(a)+f_x(a)h+f_y(a)k+\frac12\bigl(f_{xx}(a)h^2+2f_{xy}(a)hk+f_{yy}(a)k^2\bigr)+o(h^2+k^2).
\end{equation}
一般地，若 $f\in C^N$ 于 $a$ 的邻域，则
\begin{equation}
\label{eq:pei-multivariable-taylor}
f(a+h)=\sum_{|\alpha|\le N}\frac{\partial^\alpha f(a)}{\alpha!}h^\alpha+o(\|h\|^N).
\end{equation}
这里多重指标（multi-index）$\alpha=(\alpha_1,\ldots,\alpha_n)$ 的分量为非负整数，$|\alpha|=\sum_i\alpha_i$，$\alpha!=\prod_i\alpha_i!$，$h^\alpha=\prod_i h_i^{\alpha_i}$。余项是关于全部增量的估计，不能只在每个固定方向上分别成立。

\begin{proof}
对 $N\ge1$，沿线段应用积分余项公式，并把 $N$ 阶导数在 $a$ 的值从积分中分离，余项为
\[
\frac1{(N-1)!}\int_0^1(1-t)^{N-1}\bigl(D^Nf(a+th)-D^Nf(a)\bigr)[h,\ldots,h]\,dt.
\]
由 $D^Nf$ 在 $a$ 连续，该式绝对值不超过 $\omega(\|h\|)\|h\|^N/N!$，其中 $\omega(r)\to0$。这给出式~\eqref{eq:pei-multivariable-taylor} 所需的统一余项估计。
\end{proof}

这种多项式近似唯一。若两个次数不超过 $N$ 的多项式都与 $f(a+h)$ 相差 $o(\|h\|^N)$，把它们之差按齐次次数分组，并取最低的非零齐次项 $P_j(h)$。令 $h=tv$，除以 $t^j$ 再令 $t\downarrow0$，得到 $P_j(v)=0$ 对每个 $v$ 成立，矛盾。因此，只要通过初等展开、乘法或代换获得了正确总次数和余项，就无需逐个高阶求导来确认 Taylor 系数。

\begin{example}
把 $e^x\cos y$ 在原点展开到总次数四。
\end{example}
\begin{solve}
将 $e^x$ 与 $\cos y$ 的一元展开相乘，仅保留总次数不超过四的项，得到
\[
e^x\cos y=1+x+\frac{x^2-y^2}{2}+\frac{x^3-3xy^2}{6}+\frac{x^4-6x^2y^2+y^4}{24}+o(r^4),\quad r=\sqrt{x^2+y^2}.
\]
被舍弃的乘积项为 $O(r^5)$，一元余项也可由 $|x|,|y|\le r$ 控制，所以这是二元意义下的展开。按上述唯一性，它与直接计算偏导得到的 Taylor 多项式相同。
\end{solve}

同样的线段方法给出多元中值估计。若 $D$ 为凸开集，$f\in C^1(D)$，且 $\|\nabla f(x)\|\le M$ 对所有 $x\in D$ 成立，则
\begin{equation}
\label{eq:pei-gradient-lipschitz}
|f(y)-f(x)|=\left|\int_0^1\nabla f(x+t(y-x))\cdot(y-x)\,dt\right|\le M\|y-x\|.
\end{equation}
凸性保证积分所经线段仍在定义域内。若函数在紧凸集的某开邻域内为 $C^1$，梯度连续性自动给出该紧集上的有限上界。对向量值映射可用导数的算子范数得到相同类型的估计，但不能无条件写成 $F(y)-F(x)=DF(\xi)(y-x)$ 并要求所有分量共用同一点 $\xi$。

\subsection{凸函数的一阶与二阶刻画}

\begin{definition}
设 $D\subset\mathbb R^n$ 为凸集。若对任意 $x,y\in D$、$t\in[0,1]$ 有
\[
f((1-t)x+ty)\le(1-t)f(x)+tf(y),
\]
则称 $f:D\to\mathbb R$ 为凸函数（convex function）；若 $x\ne y$、$0<t<1$ 时均严格小于，则称为严格凸函数。
\end{definition}

\begin{theorem}
\label{thm:pei-convexity-criteria}
设 $D$ 为凸开集。可微函数 $f$ 为凸函数，当且仅当对任意 $x,y\in D$，
\begin{equation}
\label{eq:pei-convex-support}
f(y)\ge f(x)+\nabla f(x)\cdot(y-x).
\end{equation}
若进一步 $f\in C^2(D)$，则凸性等价于每一点的 Hesse 矩阵（Hessian matrix）$H_f(x)=(\partial_{x_j}\partial_{x_i}f(x))$ 半正定。
\end{theorem}
\begin{proof}
若 $f$ 凸，将定义式移项除以 $t>0$，再令 $t\downarrow0$ 即得式~\eqref{eq:pei-convex-support}。反向令 $z=(1-t)x+ty$，分别在 $z$ 处对 $x,y$ 应用该式，按权重 $1-t,t$ 相加，梯度项抵消，即得凸性。二阶结论把函数限制到任一线段上：$g(s)=f(x+sv)$ 满足 $g''(s)=v^{\mathsf T}H_f(x+sv)v$，一元凸函数的二阶判别遂给出充要条件。
\end{proof}

由式~\eqref{eq:pei-convex-support}，凸函数的每个驻点都是全局最小值点；严格凸时，最小值点若存在则至多一个。处处正定的 Hesse 矩阵保证严格凸，但不是必要条件，如 $x^4+y^4$ 严格凸，原点 Hesse 矩阵却为零。还需区分局部与全局：某一点 Hesse 矩阵正定只控制该点附近，不能直接推出整个定义域上的凸性。

\subsection{曲线与曲面的切空间}

若 $F\in C^1$，$S=\{x:F(x)=0\}$，且 $\nabla F(a)\ne0$，则 $S$ 在 $a$ 附近为正则曲面，其切空间为
\[
T_aS=\{v:\nabla F(a)\cdot v=0\}.
\]
在 $\mathbb R^3$ 中，相应切平面是 $\nabla F(a)\cdot(x-a)=0$，法线是 $x=a+t\nabla F(a)$。必要的正交关系来自对曲面内曲线 $F(\gamma(t))=0$ 求导；切空间维数及每个切向量都能由曲线实现，则由隐函数定理保证。若 $\nabla F(a)=0$，上述方程可能退化为 $0=0$，不能据此判断切平面存在或不存在。

对参数曲线 $r(t)$，若 $r'(t_0)\ne0$，切线方向为 $r'(t_0)$，法平面为 $r'(t_0)\cdot(x-r(t_0))=0$。对参数曲面 $r(s,t)$，若 $r_s\times r_t\ne0$，其法向量可取该叉积。若空间曲线由 $F=G=0$ 给出，且两梯度线性无关，则切向量可取 $\nabla F\times\nabla G$。两个正则曲面的锐夹角可通过两法向量计算：$\cos\theta=|n_1\cdot n_2|/(\|n_1\|\|n_2\|)$。

\begin{example}
求圆锥面 $r(\rho,\theta)=(\rho\cos\theta,\rho\sin\theta,\rho\cot\alpha)$ 在 $a=r(\rho_0,\theta_0)$ 处的切平面，其中 $0<\alpha<\pi/2$、$\rho_0>0$。
\end{example}
\begin{solve}
$r_\rho=(\cos\theta,\sin\theta,\cot\alpha)$，$r_\theta=(-\rho\sin\theta,\rho\cos\theta,0)$，故可取法向量 $n=(\cos\theta_0\cot\alpha,\sin\theta_0\cot\alpha,-1)$。代入 $n\cdot(X-a)=0$，利用 $n\cdot a=0$，得到
\[
X\cos\theta_0+Y\sin\theta_0-Z\tan\alpha=0.
\]
每个正则点的切平面都经过锥顶；顶点 $\rho_0=0$ 不满足参数秩条件，须另行讨论。
\end{solve}

\subsection{自由极值与退化情形}

\begin{theorem}
\label{thm:pei-unconstrained-extrema}
设 $a$ 为开集内点。若可微函数 $f$ 在 $a$ 取局部极值，则 $\nabla f(a)=0$。若 $f\in C^2$，且 $a$ 为驻点，则 $H_f(a)$ 正定时 $a$ 为严格局部极小值点，负定时为严格局部极大值点，不定时为鞍点；半正定或半负定时，一般不能据此判定。
\end{theorem}
\begin{proof}
极值的必要条件来自把 $f$ 限制在每条坐标线上，再用一元 Fermat 定理。驻点处 Taylor 公式为 $f(a+h)-f(a)=\tfrac12h^{\mathsf T}H_f(a)h+o(\|h\|^2)$。若矩阵正定，其二次型在单位球面上有正的最小值 $c$，于是主项至少为 $c\|h\|^2/2$，并能压住余项；负定同理。不定时存在正、负两种方向，各自沿射线取足够小的增量，就得到两种符号的函数增量。
\end{proof}

二元情形记 $A=f_{xx}(a)$、$B=f_{xy}(a)$、$C=f_{yy}(a)$、$\Delta=AC-B^2$。$\Delta>0$ 时按 $A$ 的正负区分极小与极大，$\Delta<0$ 时无局部极值，$\Delta=0$ 时需研究高阶项或定义。必要条件还包括：$C^2$ 函数在内点取极小值时，Hesse 矩阵必半正定，因为任意方向的一元二阶导数都非负。函数不光滑时还须检查不可微点，不能只解梯度为零的方程。

\begin{example}
\label{ex:pei-linewise-minimum-failure}
证明 $f(x,y)=(y-x^2)(y-2x^2)$ 在过原点的每条直线上都以原点为严格局部极小值点，但原点不是二元函数的局部极小值点。
\end{example}
\begin{solve}
在非水平直线 $y=kx$（$k\ne0$）上，$f(x,kx)=x^2(k-x)(k-2x)$，当非零 $x$ 足够小时为正；在水平轴上 $f(x,0)=2x^4>0$，在竖直轴上 $f(0,y)=y^2>0$。然而沿曲线 $y=3x^2/2$ 有 $f(x,3x^2/2)=-x^4/4<0$。各直线所用的邻域半径依赖于方向，不能得到一个同时控制所有方向的邻域。这与例~\ref{ex:pei-radial-limit-failure} 中的一致性障碍相同。
\end{solve}

\subsection{条件极值与 Lagrange 乘数}

\begin{theorem}
\label{thm:pei-lagrange-multipliers}
设 $f,g_1,\ldots,g_m\in C^1$，$m<n$，$a$ 满足 $g_i(a)=0$，且约束梯度 $\nabla g_1(a),\ldots,\nabla g_m(a)$ 线性无关。若 $f$ 在约束集上于 $a$ 取局部极值，则存在 $\lambda_1,\ldots,\lambda_m$ 使
\[
\nabla f(a)+\sum_{i=1}^m\lambda_i\nabla g_i(a)=0.
\]
\end{theorem}
\begin{proof}
由约束的秩条件及隐函数定理，约束集在 $a$ 附近是光滑曲面，其切空间为 $T=\{v:Dg(a)v=0\}$，其中 $g=(g_1,\ldots,g_m)$。对约束集内经过 $a$ 的曲线求导，极值性给出 $\nabla f(a)\cdot v=0$ 对所有 $v\in T$ 成立。因此 $\nabla f(a)$ 属于 $T^\perp$，而 $T^\perp$ 恰为约束梯度张成的空间，即得结论。
\end{proof}

令 $L(x,\lambda)=f(x)+\sum_i\lambda_i g_i(x)$。乘数方程只产生正则约束点中的候选点，仍须判别；若约束梯度失去秩，相应点应单独检查。若函数及约束为 $C^2$，在乘数候选点 $a$ 处，应把二次型 $v^{\mathsf T}H_xL(a,\lambda)v$ 限制在非零切向量 $v\in T$ 上：正定保证严格条件极小值，负定保证严格条件极大值，出现两种符号则无条件极值，半定仍需进一步分析。这里固定乘数，只对 $x$ 求 Hesse 矩阵；无需它在整个空间上正定。该判别由对局部参数化 $x=\psi(u)$ 的目标函数求二阶导得到：利用 $g\circ\psi=0$ 与驻点方程，恰好消去参数化二阶导项，留下 $D\psi^{\mathsf T}H_xL\,D\psi$。

\begin{example}
在约束 $xyz=1$ 下，求 $f(x,y,z)=x^4+y^4+z^4$ 的最小值，并检查乘数判别。
\end{example}
\begin{solve}
取 $L=f+\lambda(xyz-1)$。驻点方程乘以对应变量后给出 $4x^4+\lambda=4y^4+\lambda=4z^4+\lambda=0$。结合 $xyz=1$ 得 $|x|=|y|=|z|=1$，共四个符号乘积为正的点。由算术几何平均不等式，$x^4+y^4+z^4\ge3|xyz|^{4/3}=3$，所以它们都是全局最小值点。

在 $(1,1,1)$ 处 $\lambda=-4$，切向量满足 $h+k+l=0$，且
\[
d_x^2L=12(h^2+k^2+l^2)-8(hk+hl+kl)=16(h^2+k^2+l^2)>0
\]
对非零切向量成立，故也是严格条件局部极小值。其余三点由改变坐标符号得到同样结论。该约束集不紧，但下界与取等点已直接给出存在性；沿 $(x,y,z)=(t,1,t^{-1})$ 可见目标无上界。
\end{solve}

\subsection{闭区域最值、二次型与不等式}

连续函数在非空紧集上的最大、最小值必存在，求解时分别检查内部驻点、边界上的条件驻点以及边角和不光滑点，最后比较全部候选值。边界由若干段组成时需逐段处理，不能只求内部偏导。若定义域开或无界，必须另外控制边界趋近或无穷远处的行为；“只有一个驻点”本身不能保证它是全局最值点。

\begin{example}
设 $A=A^{\mathsf T}\in\mathbb R^{n\times n}$，求 $q(x)=x^{\mathsf T}Ax$ 在单位球面上的最值；若 $A$ 正定，再求椭球面 $x^{\mathsf T}Ax=1$ 的半轴长度。
\end{example}
\begin{solve}
球面紧，极值存在。乘数方程为 $Ax=\lambda x$、$\|x\|=1$，候选值 $q(x)=\lambda$，故最小、最大值分别为 $A$ 的最小、最大特征值，并在相应单位特征向量处取得。用正交矩阵对角化 $A$ 后，椭球方程成为 $\sum_i\lambda_i\xi_i^2=1$。正定性保证 $\lambda_i>0$，因此各主轴方向为特征向量，半轴长度为 $1/\sqrt{\lambda_i}$。仅有行列式非零不足以保证二次曲面是椭球。
\end{solve}

更一般地，对二次目标 $f(x)=\tfrac12x^{\mathsf T}Ax+b^{\mathsf T}x$ 与线性约束 $Cx=d$，假设 $A$ 对称且 $C\in\mathbb R^{m\times n}$ 满行秩，乘数条件化为
\[
\begin{pmatrix}A&C^{\mathsf T}\\C&0\end{pmatrix}\begin{pmatrix}x\\\lambda\end{pmatrix}=\begin{pmatrix}-b\\d\end{pmatrix}.
\]
若 $A$ 在 $\ker C$ 上正定，则该方程唯一可解，所得点为唯一全局最小值点：对任意可行增量 $v\in\ker C$，一次项由乘数条件消失，目标增量恰为 $\tfrac12v^{\mathsf T}Av$。块矩阵的可逆性也可由齐次方程推出：先左乘 $v^{\mathsf T}$ 得 $v=0$，再由 $C$ 满行秩得到乘数为零。

\begin{example}
\label{ex:pei-weighted-product-bound}
设 $a,b,c>0$。证明 $ab^2c^3\le108\bigl((a+b+c)/6\bigr)^6$，并给出等号条件。
\end{example}
\begin{solve}
固定 $a+b+c=s>0$，在紧单纯形 $a,b,c\ge0$ 上最大化连续函数 $ab^2c^3$。它在边界为零、内部为正，故最大值在内部取得。等价地，在内部最大化 $\log a+2\log b+3\log c$，乘数方程给出 $1/a=2/b=3/c$，从而 $a=s/6$、$b=s/3$、$c=s/2$，目标值为 $s^6/432=108(s/6)^6$。对数目标严格凹，且候选点唯一，因此这是最大值，等号当且仅当 $a:b:c=1:2:3$。推广时固定正权重 $w_i$，同样得到固定总和下 $\prod_i a_i^{w_i}$ 的最优比例 $a_i\propto w_i$。
\end{solve}

\subsection{建模应用：选变量、说明存在性并解释结果}

\begin{example}
把长度为 $1$ 的铁丝分为三段，分别围成圆、正方形与等边三角形，求面积之和最小时三段长度的比例。
\end{example}
\begin{solve}
以三段长度 $\ell_1,\ell_2,\ell_3\ge0$ 为变量，比以半径、边长为变量更直接地表达约束 $\ell_1+\ell_2+\ell_3=1$。面积和为
\[
S=\frac{\ell_1^2}{4\pi}+\frac{\ell_2^2}{16}+\frac{\ell_3^2}{12\sqrt3}.
\]
它在闭单纯形上连续，最小值存在且由严格凸性至多一个。乘数方程给出 $\ell_1/(2\pi)=\ell_2/8=\ell_3/(6\sqrt3)$，所以最优比例为 $\pi:4:3\sqrt3$；归一化后即可得到各段长度，最小面积为 $1/[4(\pi+4+3\sqrt3)]$。所得三段均为正且满足乘数方程，由凸性的一阶条件可知它优于所有边界分配。
\end{solve}

\begin{example}
设两个均匀介质以平面为界，传播速度为 $v_1,v_2>0$。点 $A,B$ 分居界面两侧，与界面的距离分别为 $a,b>0$，其投影的距离为 $\ell\ge0$。在分段直线传播的模型下，由最短时间原则导出折射关系。
\end{example}
\begin{solve}
先在包含两点及界面法向的平面内，以折射点相对 $A$ 的投影位置 $x\in\mathbb R$ 为变量，总时间为
\[
T(x)=\frac{\sqrt{a^2+x^2}}{v_1}+\frac{\sqrt{b^2+(\ell-x)^2}}{v_2}.
\]
$T(x)\to+\infty$ 当 $|x|\to\infty$，且 $T''(x)=a^2/[v_1(a^2+x^2)^{3/2}]+b^2/[v_2(b^2+(\ell-x)^2)^{3/2}]>0$，故存在唯一最小值。$\ell>0$ 时，$T'(0)<0<T'(\ell)$，最优点在 $(0,\ell)$；驻点方程给出
\[
\frac{x}{v_1\sqrt{a^2+x^2}}=\frac{\ell-x}{v_2\sqrt{b^2+(\ell-x)^2}},\qquad\frac{\sin\alpha}{v_1}=\frac{\sin\beta}{v_2},
\]
其中 $\alpha,\beta$ 与界面法线所成角分别为入射角和折射角。允许折射点离开上述平面时，额外的横向距离会同时增大两段长度，故最优点仍在该平面内。$\ell=0$ 时唯一最优点为 $x=0$，两角均为零；两种介质速度相同并把终点反射到同侧时，同类计算给出镜面反射的等角关系。
\end{solve}

\section{隐函数存在定理及函数相关}
\label{sec:pei-implicit-dependence}

\subsection{一个方程：先证明存在，再使用求导公式}

\begin{theorem}
\label{thm:pei-scalar-implicit}
设 $F(x,y)$ 在 $(a,b)\in\mathbb R^n\times\mathbb R$ 的邻域内连续，$F(a,b)=0$，$F_y$ 在该邻域内存在并连续，且 $F_y(a,b)\ne0$。则存在 $a$ 的邻域 $U$ 与 $b$ 的开区间 $I$，使每个 $x\in U$ 恰有一个 $y=\varphi(x)\in I$ 满足 $F(x,y)=0$，且 $\varphi$ 连续、$\varphi(a)=b$。若进一步 $F\in C^k$（$k\ge1$），则 $\varphi\in C^k$，且
\begin{equation}
\label{eq:pei-implicit-derivative}
\frac{\partial\varphi}{\partial x_i}(x)=-\frac{F_{x_i}(x,\varphi(x))}{F_y(x,\varphi(x))}.
\end{equation}
\end{theorem}
\begin{proof}
不妨设 $F_y(a,b)>0$。缩小邻域使 $F_y\ge c>0$，则对固定 $x$，函数 $y\mapsto F(x,y)$ 严格递增。先取足够小的 $\eta>0$，有 $F(a,b-\eta)<0<F(a,b+\eta)$；由连续性，这两个严格不等式在 $x$ 的某邻域仍成立。介值定理与严格单调性给出 $I=(b-\eta,b+\eta)$ 中唯一的零点。

对固定 $x_0\in U$，在 $\varphi(x_0)$ 两侧任取足够近的 $y_-,y_+$，连续性保证 $x$ 接近 $x_0$ 时仍有 $F(x,y_-)<0<F(x,y_+)$，所以 $y_-<\varphi(x)<y_+$，得到连续性。若 $F\in C^1$，以 $u=\varphi(x+h)-\varphi(x)$ 记，由 $F_y\ge c$ 的中值估计及 $F$ 对 $x$ 的局部有界导数可知 $|u|\le C\|h\|$。再将 $F(x+h,\varphi(x)+u)-F(x,\varphi(x))=0$ 作一阶展开，得
\[
0=D_xF(x,\varphi(x))h+F_y(x,\varphi(x))u+o(\|h\|).
\]
解出 $u$ 即得可微性与式~\eqref{eq:pei-implicit-derivative}，其右端连续。对更高阶导数递推使用该式，得到 $C^k$ 正则性。
\end{proof}

这是\emph[caution]{局部的充分条件}，其唯一性只限于所选区间 $I$，不排除远处存在其他解。$F_y(a,b)=0$ 时应另行分析：$y^3-x=0$ 仍唯一确定连续函数 $y=x^{1/3}$，但在原点不可微；$y^3-x^3=0$ 则唯一确定光滑函数 $y=x$；$y^2-x=0$ 在原点附近既有分支又有无实根的一侧。这些情形不能仅凭定理失效就作同一个判断。

\begin{example}
\label{ex:pei-implicit-non-invertibility}
对 $F(x,y)=x^2+y+\sin(xy)$，讨论原点附近的隐函数 $y=\varphi(x)$，并判断能否反过来唯一确定 $x$ 为 $y$ 的函数。
\end{example}
\begin{solve}
$F(0,0)=0$、$F_y(0,0)=1$，故存在唯一光滑局部分支 $\varphi(0)=0$，且
\[
\varphi'(x)=-\frac{2x+\varphi(x)\cos(x\varphi(x))}{1+x\cos(x\varphi(x))}.
\]
由 $\varphi'(0)=0$ 得 $\varphi(x)=o(x)$，所以 $\varphi'(x)=-2x+o(x)$；它在足够小的负半轴严格递增、正半轴严格递减，原点为严格极大值点。进一步 $\varphi''(0)=-2$，故 $\varphi(x)=-x^2+o(x^2)$。在两侧分支的共同值域内，每个充分接近零的负数对应两个不同的 $x$，因此无法在原点的双侧邻域中把全部解唯一表示为 $x=x(y)$。
\end{solve}

\subsection{方程组、局部反函数与二阶微分}

\begin{theorem}
\label{thm:pei-system-implicit}
设 $F:U\subset\mathbb R^n\times\mathbb R^m\to\mathbb R^m$ 为 $C^1$ 映射，$F(a,b)=0$，且 $m\times m$ 矩阵 $D_yF(a,b)$ 可逆。则在 $(a,b)$ 的适当乘积邻域内，方程组 $F(x,y)=0$ 唯一确定 $C^1$ 映射 $y=\varphi(x)$，满足
\begin{equation}
\label{eq:pei-implicit-jacobian}
D\varphi(x)=-[D_yF(x,\varphi(x))]^{-1}D_xF(x,\varphi(x)).
\end{equation}
若 $F\in C^k$，则 $\varphi\in C^k$。
\end{theorem}
\begin{proof}
记 $B=[D_yF(a,b)]^{-1}$，对固定 $x$ 考虑 $T_x(y)=y-BF(x,y)$。在中心点有 $D_yT_a(b)=0$。由导数连续性，可在一个闭球 $\overline B(b,r)$ 及 $a$ 的足够小邻域上统一保证 $\|D_yT_x\|\le q<1$，并使 $\|T_x(b)-b\|\le(1-q)r/2$。线段积分估计说明 $T_x$ 把该闭球映入自身且为压缩映射。由下一段给出的迭代论证，存在唯一不动点，即唯一零点 $\varphi(x)$。

比较两个参数对应的不动点方程可得 $\|\varphi(x)-\varphi(x')\|\le C\|x-x'\|/(1-q)$。因此隐函数增量为 $O(\|x-x'\|)$，将它代入 $F(x,\varphi(x))=0$ 的一阶展开并解线性方程，即得式~\eqref{eq:pei-implicit-jacobian} 与连续可微性。更高阶正则性由该式递推得到。
\end{proof}

\begin{proposition}
\label{prop:pei-parameter-contraction}
设 $K$ 是欧氏空间中的非空闭集，$T_x:K\to K$ 关于 $y$ 满足统一的压缩条件 $\|T_x(y)-T_x(z)\|\le q\|y-z\|$，其中 $0\le q<1$。从 $y_0\in K$ 出发作 $y_{j+1}=T_x(y_j)$，则它收敛到 $T_x$ 的唯一不动点 $\varphi(x)$，且
\[
\|y_j-\varphi(x)\|\le\frac{q^j}{1-q}\|y_1-y_0\|.
\]
若 $T_x(y)$ 关于参数 $x$ 连续，则不动点也关于 $x$ 连续。
\end{proposition}
\begin{proof}
逐次应用压缩条件有 $\|y_{j+1}-y_j\|\le q^j\|y_1-y_0\|$，求几何级数可知序列为 Cauchy 列，因 $K$ 闭而在 $K$ 中收敛；由 $T_x$ 的连续性，极限为不动点，尾和估计给出误差界。两个不动点的距离至多为自身的 $q$ 倍，故只能相同。最后固定 $x_0$，比较不动点方程得
\[
\|\varphi(x)-\varphi(x_0)\|\le\frac{\|T_x(\varphi(x_0))-T_{x_0}(\varphi(x_0))\|}{1-q}\longrightarrow0.
\]
\end{proof}

例如方程 $y=f(x)+\psi(y)$ 中，若 $f$ 连续、$\psi$ 在 $b$ 附近满足常数 $q<1$ 的 Lipschitz 条件，且 $b=f(a)+\psi(b)$，可选闭区间使该映射在 $x$ 接近 $a$ 时映入自身，再应用命题~\ref{prop:pei-parameter-contraction}。这里压缩条件可以替代可微性条件；只验证 Lipschitz 常数小于 $1$ 而不验证映入自身，仍不足以完成存在性论证。

\begin{corollary}
\label{cor:pei-inverse-function}
若 $G:U\subset\mathbb R^n\to\mathbb R^n$ 为 $C^1$ 映射，且 $\det DG(a)\ne0$，则 $G$ 在 $a$ 的某开邻域与 $G(a)$ 的某开邻域之间有唯一 $C^1$ 反函数，且 $D(G^{-1})(G(x))=[DG(x)]^{-1}$。
\end{corollary}
\begin{proof}
对方程 $G(x)-z=0$，把 $z$ 作为自变量、$x$ 作为未知量，应用定理~\ref{thm:pei-system-implicit}。同时缩小邻域，使邻域内所有解均受其局部唯一性控制，就得到互逆关系和导数公式。
\end{proof}

求隐函数的一、二阶导数时，可以保持方程组形式，先求导再在指定点代值。对 $F(x,\varphi(x))=0$，给定方向 $h$，记 $w=D\varphi(x)h$，则
\[
D_yF\,w=-D_xF\,h.
\]
若 $F\in C^2$，再求一次同方向导数，有
\begin{equation}
\label{eq:pei-implicit-second-differential}
D_yF\,D^2\varphi(x)[h,h]=-D^2F(x,\varphi(x))[(h,w),(h,w)].
\end{equation}
两次都解同一个系数矩阵的线性系统即可。式中 $h$ 固定，不对增量本身求导；若把自变量另设为非线性参数曲线，还须加上该曲线的二阶导项。

\subsection{局部非退化与全局单射的区别}

Jacobi 行列式处处非零只保证各点附近可逆，即使定义域凸，也不能单独推出全局单射。例如
\begin{equation}
\label{eq:pei-local-not-global-inverse}
G(x,y)=(e^x\cos y,e^x\sin y),\qquad(x,y)\in\mathbb R^2,
\end{equation}
有 $\det DG=e^{2x}>0$，却满足 $G(x,y+2\pi)=G(x,y)$。不能把一个二次型 $v^{\mathsf T}DG(\xi)v=0$ 误判为 $DG(\xi)v=0$，更不能由此断言行列式为零；可逆的反对称矩阵就能使前一个二次型对全部 $v$ 为零。

\begin{proposition}
\label{prop:pei-monotone-injectivity}
设 $D\subset\mathbb R^n$ 为凸开集，$G\in C^1(D,\mathbb R^n)$。若对所有 $x\in D$、$v\ne0$ 都有 $v^{\mathsf T}DG(x)v>0$，即导数矩阵的对称部分正定，则 $G$ 在 $D$ 上单射。
\end{proposition}
\begin{proof}
对不同的 $x,y\in D$，令 $v=y-x\ne0$。沿线段积分得
\[
(G(y)-G(x))\cdot(y-x)=\int_0^1v^{\mathsf T}DG(x+tv)v\,dt>0.
\]
所以 $G(x)\ne G(y)$。本结论只需 $C^1$，且不要求 $DG$ 本身对称；若对称，则它简化为通常的正定矩阵条件。
\end{proof}

另一方向，若存在 $c>0$ 使 $\|G(x)-G(y)\|\ge c\|x-y\|$ 对所有 $x,y\in D$ 成立，则固定 $x$、令 $y=x+tv$，除以 $|t|$ 并取极限，得 $\|DG(x)v\|\ge c\|v\|$。因此 $DG(x)$ 单射，在方阵情形行列式非零。这一论证直接从差商出发，无需为向量值函数假设一个所有分量共用的中值点。

\subsection{局部最小值分支、最优值导数与 Legendre 变换}

\begin{proposition}
\label{prop:pei-local-envelope}
设 $f\in C^2$ 于 $(a,b)\in\mathbb R^2$ 的邻域，$f_x(a,b)=0$、$f_{xx}(a,b)>0$。则在 $b$ 的某邻域存在唯一的近邻驻点分支 $x=\xi(y)$，每个 $\xi(y)$ 都是 $x\mapsto f(x,y)$ 的严格局部极小值点。相应局部极小值函数 $v(y)=f(\xi(y),y)$ 满足
\[
v'(y)=f_y(\xi(y),y),\qquad
v''(y)=f_{yy}-\frac{f_{xy}^2}{f_{xx}},
\]
第二式各项均在 $(\xi(y),y)$ 处取值。
\end{proposition}
\begin{proof}
对方程 $f_x(x,y)=0$ 应用隐函数定理，得 $\xi'=-f_{xy}/f_{xx}$，并由连续性缩小邻域使 $f_{xx}>0$。链式法则给出 $v'=f_x\xi'+f_y=f_y$；再求导即得二阶式。此处构造的是局部极小值分支，未断言它对全部 $x$ 给出全局最小值。
\end{proof}

教材中的另一类换元可写为局部 Legendre 变换（Legendre transform）。设 $f\in C^2$，$u=\nabla f(x)$，且 $\det H_f(a)\ne0$。由反函数定理，在 $a$ 附近可反解 $x=x(u)$，定义
\[
g(u)=u\cdot x(u)-f(x(u)).
\]
直接求微分，$dg=x\cdot du+(u-\nabla f(x))\cdot dx=x\cdot du$，因此 $\nabla g(u)=x(u)$、$H_g(u)=[H_f(x(u))]^{-1}$。若 $f$ 为二次齐次函数，由 Euler 公式 $x\cdot\nabla f(x)=2f(x)$，还有 $g(u)=f(x(u))$。这些是局部等式；若要把 $g$ 定义为全空间上确界 $\sup_x(u\cdot x-f(x))$，则还需相应的凸性及最优解存在条件。

\subsection{函数相关与 Jacobi 矩阵的秩}

设 $Y=(y_1,\ldots,y_m):U\subset\mathbb R^n\to\mathbb R^m$ 为 $C^1$ 映射，$a\in U$、$b=Y(a)$。若存在 $b$ 的邻域上的非平凡关系 $\Phi$，使 $\Phi(Y(x))=0$ 在 $a$ 附近恒成立，就得到函数间的局部依赖。使用微分判别时，应明确关系的正则性：本段称 $C^1$ 关系满足 $\nabla\Phi(b)\ne0$ 时为\emph{正则局部函数关系}。它说明像点受一个真正的局部方程约束，不只是选择一个在像集上偶然消失的函数。

\begin{theorem}
\label{thm:pei-functional-rank}
若 $\operatorname{rank}DY(a)=m$，则 $Y$ 的局部像包含 $b$ 的一个邻域，因而不存在在 $b$ 的每个邻域都不恒为零的局部函数关系。若在 $a$ 的某邻域内 $DY$ 的秩恒为 $r<m$，则重新排列输入及输出坐标后，$y_1,\ldots,y_r$ 局部无关，其余输出可写为
\[
y_j=\psi_j(y_1,\ldots,y_r),\qquad j=r+1,\ldots,m,
\]
其中 $\psi_j$ 为 $C^1$ 函数；这些方程给出正则局部函数关系。
\end{theorem}
\begin{proof}
满行秩时，选出可逆的 $m$ 阶子式，固定其余输入坐标，便得到一个可应用反函数定理的方形映射，所以像包含 $b$ 的邻域。

对常秩 $r>0$ 的情形，选定非零的 $r$ 阶子式，取新输入坐标 $(y_1,\ldots,y_r,x_{r+1},\ldots,x_n)$。其变换导数可逆，故局部可作坐标。常秩条件说明其余 $y_j$ 在这些新坐标中对后 $n-r$ 个坐标的偏导数全为零；在足够小的矩形邻域内，它们只依赖于前 $r$ 个坐标，得到所需 $\psi_j$。$r=0$ 时全部导数为零，缩小到连通邻域后各分量恒为常数，结论相应成立。
\end{proof}

仅在单独一点秩不足，不能反推局部函数相关；例如 $Y(x)=x^3$ 在原点导数为零，但局部像仍包含原点的整个邻域。对齐次线性函数组 $Y(x)=Ax$，则函数相关等价于各行线性相关：若行相关，存在非零常向量 $c$ 使 $c^{\mathsf T}Y(x)=0$；若各行独立，则满行秩条件排除任何非平凡局部关系。线性相关总给出函数相关，但非线性关系可能在线性无关的函数之间出现，如 $y_1=x$、$y_2=x^2$ 满足 $y_2-y_1^2=0$。

线性相关还可用导数构造必要条件。若一元函数 $u_1,\ldots,u_m$ 有到 $m-1$ 阶的导数，且存在不全为零的常数 $c_j$ 使 $\sum_jc_ju_j=0$，将恒等式逐次求导便知 Wronski 行列式（Wronskian）$W(x)=\det(u_j^{(i-1)}(x))_{i,j=1}^m$ 恒为零。反向结论对一般光滑度条件不成立，例如 $u_1(x)=x^2$、$u_2(x)=x|x|$ 在 $\mathbb R$ 上均为 $C^1$，Wronski 行列式恒零，但分别在正、负半轴检验可知它们线性无关；若要逆推，需另加微分方程解空间等结构条件。

当 $F$ 由函数级数给出时，应用隐函数定理之前须先证明和函数具有所需连续性和偏导数。可在考察点的闭矩形邻域上用与变量无关的可求和上界，分别控制原级数及一阶偏导级数，再按逐项求导定理核实 $F\in C^1$。仅在某点逐项算出了形式上的 $F_y$，不能代替邻域内的正则性检查。

\section{方向导数与梯度}
\label{sec:pei-directions-gradients}

\subsection{定义中的方向、参数与可微性条件}

前面已经利用梯度计算切空间，本节进一步定义方向导数，说明它与可微性的关系以及梯度的计算方法。本文对单位向量 $e$ 采用双侧极限
\[
\partial_ef(a)=\lim_{t\to0}\frac{f(a+te)-f(a)}t.
\]
有些教材采用 $t\downarrow0$ 的单侧方向导数，使用时应核对约定。给定非零向量 $v$ 时，若问题要求“沿 $v$ 的方向的单位距离变化率”，应先归一化为 $e=v/\|v\|$；直接求 $t\mapsto f(a+tv)$ 的导数得到的是沿参数 $t$ 的变化率，可微时等于 $\nabla f(a)\cdot v=\|v\|\partial_ef(a)$。

若全部方向导数存在，也不保证它们关于方向向量组成线性映射。例如将 $f(x,y)=x^3/(x^2+y^2)$ 在原点补为 $0$，沿单位向量 $(p,q)$ 的方向导数为 $p^3$，若存在表示微分的线性映射，由两轴方向导数只能得到 $A=(1,0)$，但一般的 $p^3$ 不等于 $A\cdot(p,q)=p$，故原点不可微。

可微时，方向导数关于方向向量具有线性结构，$\partial_ef=\nabla f\cdot e$；梯度不是若干方向导数的标量拼合规则，而是唯一表示微分的向量。偏导数存在、全部方向导数存在、方向导数由某线性映射表示、函数可微，这几项也应分清：最后一项要求余项对所有方向统一为距离的高阶无穷小，前三项本身均未提供这种一致性。

\begin{example}
\label{ex:pei-zero-directions-discontinuous}
将 $f(x,y)=x^6y/(x^{12}+y^2)$ 在原点补为 $0$。证明它在原点的全部方向导数均为 $0$，但不连续。
\end{example}
\begin{solve}
对单位向量 $e=(p,q)$，若 $q=0$，沿该方向函数恒为零；若 $q\ne0$，则
\[
\frac{f(tp,tq)}t=\frac{t^4p^6q}{t^{10}p^{12}+q^2}\longrightarrow0.
\]
然而沿 $y=x^6$ 有 $f(x,x^6)=1/2$，故不连续，也不可微。这比“全部方向导数存在但不组成线性映射”的反例更强：即使方向导数全为零，仍需检查邻域内的统一余项。
\end{solve}

\begin{theorem}
\label{thm:pei-lipschitz-directional-differentiability}
设 $f$ 在 $a$ 的某邻域满足 Lipschitz 条件 $|f(x)-f(y)|\le L\|x-y\|$，并且对所有单位向量 $e$，方向导数存在且 $\partial_ef(a)=A\cdot e$，其中 $A\in\mathbb R^n$ 固定。则 $f$ 在 $a$ 可微，$Df(a)h=A\cdot h$。
\end{theorem}
\begin{proof}
令 $R(r,e)=[f(a+re)-f(a)]/r-A\cdot e$，其中 $r>0$。对固定 $e$，方向导数条件保证 $R(r,e)\to0$；而
\[
|R(r,e)-R(r,e')|\le(L+\|A\|)\|e-e'\|,
\]
右端与 $r$ 无关。用单位球面的有限覆盖，先选有限个方向控制全部角度，再对这有限个方向选同一个充分小的半径，得到 $\sup_{\|e\|=1}|R(r,e)|\to0$。把任意非零 $h$ 写成 $r e$，即得可微性定义。
\end{proof}

该定理补上了例~\ref{ex:pei-zero-directions-discontinuous} 缺少的局部 Lipschitz 条件；例~\ref{ex:pei-bounded-partials-not-differentiable} 则虽然局部 Lipschitz，但在两个分量均非零的方向上连双侧方向导数也不存在。两者说明这些假设分别承担不同作用，不能只保留其中一个。

\subsection{从给定方向恢复梯度及其约束}

若 $f$ 在 $a$ 可微，在 $n$ 个线性无关的单位向量 $e_1,\ldots,e_n$ 上已知 $\partial_{e_i}f(a)=c_i$，则由线性系统 $e_i\cdot\nabla f(a)=c_i$ 可唯一恢复梯度。这些方向未必正交，但若不线性无关就未必能唯一恢复。特别地，若在连通开集内上述 $n$ 个固定独立方向的导数处处为零，则梯度处处为零，函数为常数：在每个包含于定义域的球内由线段中值定理得局部常值，再由连通性推出整体常值。

平面内若 $\nabla f(a)\ne0$，某个给定方向导数值 $c$ 对应的单位向量满足直线方程 $\nabla f(a)\cdot e=c$；直线与单位圆至多有两个交点。因此，若已知至少一个方向导数非零，却有三个不同单位方向的方向导数相同，则 $f$ 在该点不可微。又若 $N\ge2$ 个单位向量在圆周上等角排列，则 $\sum_{j=1}^N e_j=0$，由微分的线性性得 $\sum_j\partial_{e_j}f(a)=0$。这些结论都以可微性为前提，不能反过来从少量方向数据断言可微。

对正交单位基 $e_1,\ldots,e_n$，Parseval 恒等式给出
\[
\sum_{i=1}^n(\partial_{e_i}f(a))^2=\|\nabla f(a)\|^2.
\]
若 $f\in C^2$，二阶方向导数为 $\partial_e^2f(a)=e^{\mathsf T}H_f(a)e$，所以
\[
\sum_{i=1}^n\partial_{e_i}^2f(a)=\operatorname{tr}H_f(a)=\sum_{j=1}^n f_{x_jx_j}(a).
\]
这从方向导数角度解释了梯度长度和 Laplace 算子在正交坐标变换下的不变性。

\subsection{法向导数与曲线坐标中的梯度}

在正则曲面上选定单位法向量 $\nu$ 后，法向导数（normal derivative）为 $\partial_\nu f=\nabla f\cdot\nu$；若曲面是区域 $\{F<0\}$ 的边界且 $\nabla F\ne0$，其外单位法向量为 $\nu=\nabla F/\|\nabla F\|$。必须先区分内外方向，再作归一化。若 $f$ 在光滑曲面上恒为常数，则任一切向量 $v$ 均满足 $\nabla f\cdot v=0$；这一结论也可直接对曲面内的曲线求导得到。

\begin{example}
设椭球面 $x^2/a^2+y^2/b^2+z^2/c^2=1$，其中 $a\ge b\ge c>0$。求 $u(x,y,z)=x^2+y^2+z^2$ 的外法向导数，并求其最大、最小值。
\end{example}
\begin{solve}
外单位法向量为
\[
\nu=\frac{(x/a^2,y/b^2,z/c^2)}{\sqrt{x^2/a^4+y^2/b^4+z^2/c^4}},
\]
故在椭球面上有 $\partial_\nu u=2/\sqrt{x^2/a^4+y^2/b^4+z^2/c^4}$。令 $s_1=x^2/a^2$、$s_2=y^2/b^2$、$s_3=z^2/c^2$，则 $s_i\ge0$、$\sum_i s_i=1$，分母平方是 $s_1/a^2+s_2/b^2+s_3/c^2$，介于 $1/a^2$ 与 $1/c^2$ 之间。因此最小值为 $2c$，最大值为 $2a$，分别在最短与最长半轴的端点取得；半轴相等时可有更多取等点。
\end{solve}

在极坐标中，$e_r=(\cos\theta,\sin\theta)$、$e_\theta=(-\sin\theta,\cos\theta)$ 是局部正交单位基。若 $U(r,\theta)=u(r\cos\theta,r\sin\theta)$，链式法则给出 $U_r=\nabla u\cdot e_r$、$U_\theta=r\nabla u\cdot e_\theta$，故
\[
\nabla u=U_r e_r+\frac1rU_\theta e_\theta,\qquad r>0.
\]
这里 $\partial_\theta$ 对应弧长为 $r\,d\theta$ 的移动，所以角向分量带 $1/r$，不能直接把 $(U_r,U_\theta)$ 当成正交单位基下的梯度分量。

更一般地，设平面局部坐标 $r=r(s,t)$ 为 $C^1$，$r_s,r_t$ 非零且正交，记尺度因子 $H_s=\|r_s\|$、$H_t=\|r_t\|$。若 $\widetilde f(s,t)=f(r(s,t))$，则
\[
\nabla f=\frac{\widetilde f_s}{H_s^2}r_s+\frac{\widetilde f_t}{H_t^2}r_t
=\frac{\widetilde f_s}{H_s}\frac{r_s}{H_s}+\frac{\widetilde f_t}{H_t}\frac{r_t}{H_t}.
\]
因为右端与 $r_s,r_t$ 的内积分别为 $\widetilde f_s,\widetilde f_t$，与链式法则完全相同，而这两个向量构成一组基，所以该表达式唯一成立。若坐标向量不正交，则必须解包含其两两内积的线性系统，不能继续套用此式。

\subsection{方向导数的中心差分与误差界}

\begin{proposition}
设 $f\in C^3$ 于线段 $\{a+te:|t|\le h\}$ 的某邻域，$e$ 为单位向量，$h>0$。若线段上 $|D^3f(x)[e,e,e]|\le M_e$，则
\[
\left|\frac{f(a+he)-f(a-he)}{2h}-\partial_ef(a)\right|\le\frac{M_eh^2}{6}.
\]
若各坐标三阶偏导数在该线段上绝对值均不超过 $M$，可取 $M_e=Mn^{3/2}$。
\end{proposition}
\begin{proof}
对 $g(t)=f(a+te)$ 分别在 $h$ 与 $-h$ 处写到二阶的 Taylor 公式，相减时常数项与二阶项抵消；两个三阶余项的绝对值之和至多为 $M_eh^3/3$，除以 $2h$ 即得误差界。最后
\[
|D^3f(x)[e,e,e]|\le M\sum_{i,j,k}|e_ie_je_k|=M\Bigl(\sum_i|e_i|\Bigr)^3\le Mn^{3/2}.
\]
这里必须先控制沿 $e$ 的三阶导数，不能把坐标偏导的界直接当成方向导数的同一个界；二维情形得到 $\sqrt2 Mh^2/3$ 的误差上界。
\end{proof}

\chapter{多元积分学}
\label{chap:pei-multivariate-integrals}

% 教材依据：数分习题-7.pdf，第 1--203 页（书内第 655--857 页）。
% 本章按原书 §§7.1--7.5 整理理论、方法与代表性问题，不逐题转录单元练习。
本章讨论的积分有两条相互联系的主线：一条把积分看作参变量的函数，研究何时能够交换极限、微分与积分；另一条把积分放到平面区域、空间区域、曲线和曲面上，研究坐标变换以及区域内部与边界之间的关系。前一条主线的关键是控制余项，尤其是含参变量反常积分的一致收敛性；后一条主线的关键是同时辨认积分对象、体积或面积元素以及定向。以下依次整理含参变量积分、重积分、曲线积分、曲面积分与场论，计算结果均应连同其成立条件一起使用。

\section{含参变量积分}
\label{sec:pei-parameter-integrals}
% 对应原书 §7.1：PDF 第 1--58 页。

\subsection{正常积分的连续性与积分号下运算}

积分变量与参变量承担不同角色。例如
\[
F(t)=\int_a^b f(x,t)\,dx
\]
中，$x$ 在积分完成后消失，$t$ 决定函数 $F$ 的取值；把被积函数关于 $t$ 的性质传递给 $F$，需要对整个积分区间统一控制误差。积分区间有限且被积函数无不可积奇点时，称为\emph{正常积分（proper integral）}。

\begin{theorem}
\label{thm:pei-proper-parameter}
设 $a<b$，$I$ 为实区间，$f$ 在 $[a,b]\times I$ 上连续，则 $F(t)=\int_a^b f(x,t)\,dx$ 在 $I$ 上连续。若 $f_t=\partial f/\partial t$ 也连续，则在 $I$ 的内部有
\begin{equation}
\label{eq:pei-proper-leibniz}
F'(t)=\int_a^b f_t(x,t)\,dx.
\end{equation}
若 $[c,d]\subset I$，则
\begin{equation}
\label{eq:pei-proper-interchange}
\int_c^d\!\left(\int_a^b f(x,t)\,dx\right)dt
=\int_a^b\!\left(\int_c^d f(x,t)\,dt\right)dx.
\end{equation}
\end{theorem}

\begin{proof}
固定 $t_0\in I$，在其附近取一个包含于 $I$ 的紧区间 $J$；端点处取单侧区间。由 $f$ 在 $[a,b]\times J$ 上一致连续，
\[
|F(t)-F(t_0)|\le (b-a)\sup_{x\in[a,b]}|f(x,t)-f(x,t_0)|\longrightarrow0.
\]
若 $f_t$ 连续，则由微积分基本定理，差商与 $f_t(x,t)$ 的差等于
\[
\frac1h\int_t^{t+h}\bigl(f_t(x,s)-f_t(x,t)\bigr)\,ds,
\]
它关于 $x\in[a,b]$ 一致趋于零，故可对差商取积分，得到式~\eqref{eq:pei-proper-leibniz}。式~\eqref{eq:pei-proper-interchange} 可由矩形分划的有限双重求和交换次序并取极限得到；连续性保证相应的 Riemann 和均收敛。
\end{proof}

若上下限也随 $t$ 变化，设 $\alpha,\beta\in C^1(I)$，$f,f_t$ 在所有相关点的一个邻域内连续，则 \emph{Leibniz 公式（Leibniz integral rule）}为
\begin{equation}
\label{eq:pei-moving-leibniz}
\frac{d}{dt}\int_{\alpha(t)}^{\beta(t)}f(x,t)\,dx
=f(\beta(t),t)\beta'(t)-f(\alpha(t),t)\alpha'(t)
+\int_{\alpha(t)}^{\beta(t)}f_t(x,t)\,dx.
\end{equation}
证明时先把积分写成两个从同一固定点出发的变上限积分，再分别用链式法则。两个边界项反映积分区间的变化，不能因被积函数含有参数就遗漏。

\begin{example}
\label{ex:pei-absolute-kernel}
设 $f\in C([a,b])$，$F(t)=\int_a^b|x-t|f(x)\,dx$。求 $F$ 在 $(a,b)$ 内的二阶导数。
\end{example}
\begin{solve}
绝对值在 $x=t$ 处不可直接按光滑核处理，应先分段：
\[
F(t)=\int_a^t(t-x)f(x)\,dx+\int_t^b(x-t)f(x)\,dx.
\]
利用式~\eqref{eq:pei-moving-leibniz}，两个活动端点的被积函数均为零，故
\[
F'(t)=\int_a^t f(x)\,dx-\int_t^b f(x)\,dx,
\qquad F''(t)=2f(t)\quad(a<t<b).
\]
在 $t<a$ 或 $t>b$ 时，$F$ 是 $t$ 的一次函数，二阶导数为零；因此不能把内部的二阶导数公式未经检验地延伸到端点。
\end{solve}

\subsection{反常积分的一致收敛与一致 Cauchy 准则}

对无穷区间上的积分，有限区间上的连续性只能控制积分的前段；还必须控制被截去的尾部。设对每个 $t\in I$，$f(\cdot,t)$ 在任意 $[a,A]$ 上可积，且
\[
F(t)=\int_a^{+\infty}f(x,t)\,dx
=\lim_{A\to+\infty}\int_a^A f(x,t)\,dx
\]
存在。这里的收敛最初只是对每个固定 $t$ 成立。

\begin{definition}
\label{def:pei-uniform-improper}
若对任意 $\varepsilon>0$，存在与 $t$ 无关的 $A_0>a$，使得 $A\ge A_0$ 时，对所有 $t\in I$ 均有
\[
\left|\int_A^{+\infty}f(x,t)\,dx\right|<\varepsilon,
\]
则称该反常积分关于 $t\in I$ \emph{一致收敛（uniform convergence）}。若它在 $I$ 的每个紧子区间上一致收敛，则称为\emph{局部一致收敛（locally uniform convergence）}，亦称内闭一致收敛。
\end{definition}

\begin{theorem}
\label{thm:pei-improper-uniform-cauchy}
上述反常积分在 $I$ 上一致收敛，当且仅当对任意 $\varepsilon>0$，存在 $A_0>a$，使得对所有 $B>A\ge A_0$ 及所有 $t\in I$，均有
\begin{equation}
\label{eq:pei-uniform-cauchy}
\left|\int_A^B f(x,t)\,dx\right|<\varepsilon.
\end{equation}
此条件本身也保证每个固定 $t$ 对应的反常积分收敛。
\end{theorem}
\begin{proof}
一致收敛时，将 $[A,B]$ 上的积分写成两个尾积分之差，并分别用 $\varepsilon/2$ 控制。反之，先固定 $t$，由实数完备性及通常的 Cauchy 准则得到 $F(t)$ 存在；再将式~\eqref{eq:pei-uniform-cauchy} 中的 $B$ 趋于无穷，即得到与 $t$ 无关的尾积分估计。
\end{proof}

在有限端点 $b$ 有奇点时，把尾积分换成 $\int_{b-\delta}^b f(x,t)\,dx$，并令 $\delta\downarrow0$；若有多个固定奇点，则应分别控制各奇点附近的积分。有活动奇点 $x=t$ 时，应明确截去的是哪一个邻域，并分别估计左右两侧，不能只检查无穷远处。证明非一致收敛时，常选取 $A_n\to\infty$、$B_n>A_n$ 与参数 $t_n$，使 $|\int_{A_n}^{B_n}f(x,t_n)\,dx|$ 保持大于某个固定正常数；有限区间上的这一估计通常比直接估计无限尾积分容易。

\begin{example}
\label{ex:pei-nonuniform-exponential}
证明 $\int_0^{+\infty}xe^{-tx}\,dx$ 在每个 $[t_0,+\infty)$（$t_0>0$）上一致收敛，但在 $(0,+\infty)$ 上非一致收敛。
\end{example}
\begin{solve}
直接计算尾积分得到
\[
\int_A^{+\infty}xe^{-tx}\,dx
=e^{-tA}\left(\frac At+\frac1{t^2}\right).
\]
当 $t\ge t_0$ 时，它不超过 $e^{-t_0A}(A/t_0+1/t_0^2)$，后者趋于零，故一致收敛。另一方面，对任意固定 $A>0$，当 $t\downarrow0$ 时，该尾积分趋于无穷，不可能有适用于所有 $t>0$ 的统一尾部估计。这里每个积分都存在且等于 $1/t^2$，仍不足以推出整个参数区间上的一致收敛。
\end{solve}

\subsection{控制尾部的三个判别法}

\begin{theorem}
\label{thm:pei-improper-tests}
设下述函数对积分变量局部可积，所出现的乘积也局部可积。

若存在非负可积函数 $M$，使 $|f(x,t)|\le M(x)$ 对所有 $x\ge a,t\in I$ 成立，且 $\int_a^{+\infty}M(x)\,dx<\infty$，则 $\int_a^{+\infty}f(x,t)\,dx$ 一致收敛，$\int_a^{+\infty}|f(x,t)|\,dx$ 也一致收敛。这是 \emph{Weierstrass 优函数判别法（Weierstrass $M$-test）}。

设 $G(X,t)=\int_a^X g(x,t)\,dx$ 对 $X\ge a,t\in I$ 一致有界。若每个固定 $t$ 对应的 $h(x,t)$ 关于 $x$ 单调，且 $\sup_{t\in I}|h(X,t)|\to0$，则 $\int_a^{+\infty}g(x,t)h(x,t)\,dx$ 一致收敛。这是 \emph{Dirichlet 判别法（Dirichlet test）}。

若 $\int_a^{+\infty}g(x,t)\,dx$ 已在 $I$ 上一致收敛，而 $h(x,t)$ 对每个固定 $t$ 关于 $x$ 单调，且对 $x,t$ 一致有界，则乘积积分 $\int_a^{+\infty}g(x,t)h(x,t)\,dx$ 一致收敛。这是 \emph{Abel 判别法（Abel test）}。
\end{theorem}
\begin{proof}
优函数判别法直接来自 $\int_A^B|f|\le\int_A^{+\infty}M$。为说明另外两个判别法，先取光滑单调的 $h$，置 $G_A(x,t)=\int_A^xg(s,t)\,ds$，分部积分得
\[
\int_A^B gh\,dx=G_A(B,t)h(B,t)-\int_A^B G_A(x,t)h_x(x,t)\,dx.
\]
因此其绝对值不超过
\[
\sup_{A\le x\le B}|G_A(x,t)|
\bigl(|h(B,t)|+|h(B,t)-h(A,t)|\bigr).
\]
Dirichlet 条件使第一因子一致有界、第二因子在 $A\to\infty$ 时一致趋零；Abel 条件使第一因子由一致 Cauchy 准则一致趋零、第二因子一致有界。一般单调 $h$ 可用积分第二中值定理得到同样的估计，不要求其可微。最后应用定理~\ref{thm:pei-improper-uniform-cauchy}。
\end{proof}

优函数适合控制绝对值；有振荡时应保留抵消结构，优先检查 Dirichlet 或 Abel 条件。例如，对 $p\in[p_0,+\infty)$、$p_0>0$，$\int_1^{+\infty}x^{-p}\sin x\,dx$ 由 Dirichlet 判别法一致收敛，而当 $0<p\le1$ 时它不绝对收敛。\emph[caution]{参数相关的单调性不等于关于参数的一致有界或一致趋零}，使用这些判别法时，三个量词都要逐项核验；优函数判别法只是充分条件，找不到优函数不能反推出非一致收敛。

\begin{example}
\label{ex:pei-moving-singularity}
设 $0<q<1$，$f$ 是 $[0,1]$ 上有界的 Riemann 可积函数。说明活动奇点积分
\[
F(t)=\int_0^1\frac{f(x)}{|x-t|^q}\,dx,\qquad 0\le t\le1,
\]
的奇点附近可被一致控制。
\end{example}
\begin{solve}
设 $|f|\le M$。在奇点左右分别取不超过 $\delta$ 的区间，便有
\[
\int_{[0,1]\cap\{|x-t|<\delta\}}
\frac{|f(x)|}{|x-t|^q}\,dx
\le\frac{2M}{1-q}\delta^{1-q},
\]
右端与 $t$ 无关且趋于零。离开奇点后，被积函数是 Riemann 可积函数与连续有界函数的乘积，故正常可积；上述估计由此保证全部 $F(t)$ 都绝对收敛，并给出关于活动奇点的统一余项控制。若还要证明 $F$ 连续，可把积分分为靠近参数极限点的小邻域与其补集，前者用此估计，后者用核函数的一致连续性。
\end{solve}

\subsection{反常积分的极限、连续性、求导与换序}

\begin{theorem}
\label{thm:pei-improper-limit}
设 $f_n$ 在 $[a,+\infty)$ 上局部可积，在每个 $[a,A]$ 上一致收敛于 $f$，且积分 $\int_a^{+\infty}f_n(x)\,dx$ 关于 $n\in\mathbb N$ 一致收敛。则 $f$ 的反常积分存在，且
\begin{equation}
\label{eq:pei-improper-limit}
\lim_{n\to\infty}\int_a^{+\infty}f_n(x)\,dx
=\int_a^{+\infty}f(x)\,dx.
\end{equation}
相应地，若 $f(x,t)$ 连续且其反常积分关于 $t$ 局部一致收敛，则该积分所定义的函数在参数区间内连续。
\end{theorem}
\begin{proof}
由一致 Cauchy 准则，充分大的 $B>A$ 满足 $|\int_A^B f_n|<\varepsilon$，估计与 $n$ 无关。先在固定的 $[A,B]$ 上令 $n\to\infty$，得到 $f$ 的 Cauchy 估计，故其反常积分存在。再分解
\[
\int_a^{+\infty}(f_n-f)
=\int_a^A(f_n-f)+\int_A^{+\infty}f_n-\int_A^{+\infty}f.
\]
先取 $A$ 同时控制后两项，再取 $n$ 控制第一项，即证式~\eqref{eq:pei-improper-limit}。参数连续性的证明采用同样的分段估计，并在有限前段利用定理~\ref{thm:pei-proper-parameter}。
\end{proof}

\begin{theorem}
\label{thm:pei-improper-differentiation}
设 $I$ 是开区间，$f,f_t$ 在 $[a,+\infty)\times I$ 上连续；存在 $t_*\in I$ 使 $\int_a^{+\infty}f(x,t_*)\,dx$ 收敛，且 $\int_a^{+\infty}f_t(x,t)\,dx$ 关于 $t$ 局部一致收敛。则原积分对每个 $t\in I$ 收敛，并定义一个 $C^1$ 函数，满足
\begin{equation}
\label{eq:pei-improper-differentiation}
\frac{d}{dt}\int_a^{+\infty}f(x,t)\,dx
=\int_a^{+\infty}f_t(x,t)\,dx.
\end{equation}
\end{theorem}
\begin{proof}
在包含 $t_*,t$ 的任意紧子区间 $J\subset I$ 上，对有限截断积分应用正常积分的公式，得到
\[
F_A(t)=F_A(t_*)+\int_{t_*}^t\left(\int_a^A f_t(x,s)\,dx\right)ds.
\]
右端内层积分关于 $s\in J$ 一致收敛，因此可令 $A\to\infty$，得到
\[
F(t)=F(t_*)+\int_{t_*}^t H(s)\,ds,
\qquad H(s)=\int_a^{+\infty}f_t(x,s)\,dx.
\]
由定理~\ref{thm:pei-improper-limit}，$H$ 连续，再用微积分基本定理即得结论。
\end{proof}

求导定理控制的是\emph[caution]{求导后的积分}；原积分一致收敛本身不足以保证积分号下求导。对积分号下再求积分，若 $f$ 在 $[a,+\infty)\times[c,d]$ 连续且原积分关于 $t\in[c,d]$ 一致收敛，则
\begin{equation}
\label{eq:pei-improper-parameter-interchange}
\int_c^d\!\left(\int_a^{+\infty}f(x,t)\,dx\right)dt
=\int_a^{+\infty}\!\left(\int_c^d f(x,t)\,dt\right)dx.
\end{equation}
证明是先在 $[a,A]\times[c,d]$ 上换序，再用一致尾部估计令 $A\to\infty$。若两个积分区间都无界，不能直接套用此式；应另外建立非负性或绝对可积性等换序依据，参见第~\ref{sec:pei-multiple-integrals} 节。

\subsection{引入参数、建立微分方程与收敛因子}

计算含对数、反三角函数或振荡因子的积分时，直接求原函数未必有效。\emph[tip]{把待求积分嵌入一个易于求导的函数族}，可先确定导数或微分方程，再用某个参数值或边界极限确定常数。引入参数后应先核验定理~\ref{thm:pei-improper-differentiation} 的条件；若条件失效，可考虑加入衰减因子，最后再证明去掉该因子的极限合法。

\begin{example}
\label{ex:pei-frullani}
设 $f$ 在 $[0,+\infty)$ 上连续，$\lim_{x\to+\infty}f(x)=L$，$a,b>0$。证明 \emph{Frullani 公式（Frullani integral）}
\begin{equation}
\label{eq:pei-frullani}
\int_0^{+\infty}\frac{f(ax)-f(bx)}x\,dx
=(f(0)-L)\log\frac ba.
\end{equation}
\end{example}
\begin{solve}
先设 $a<b$，对 $0<\delta<R$ 换元并消去公共区间，得到
\[
\int_\delta^R\frac{f(ax)-f(bx)}x\,dx
=\int_{a\delta}^{b\delta}\frac{f(u)}u\,du
-\int_{aR}^{bR}\frac{f(u)}u\,du.
\]
第一项与 $f(0)\log(b/a)$ 的差不超过
$\sup_{a\delta\le u\le b\delta}|f(u)-f(0)|\log(b/a)$，第二项与 $L\log(b/a)$ 的差也有类似估计。因此两个端点的极限可分别控制，反常积分存在且满足式~\eqref{eq:pei-frullani}。$a=b$ 时两端均为零，$a>b$ 时交换 $a,b$ 即可。例如取 $f(x)=e^{-x}$，得到 $\int_0^{+\infty}(e^{-ax}-e^{-bx})\,dx/x=\log(b/a)$。
\end{solve}

\begin{example}
\label{ex:pei-dirichlet-integral}
利用收敛因子计算 $\int_0^{+\infty}\sin(bx)\,dx/x$，其中 $b\in\mathbb R$。
\end{example}
\begin{solve}
原积分在零点补充连续定义后正常可积，在无穷远由 Dirichlet 判别法收敛，但直接对 $b$ 求导会产生通常不收敛的 $\int_0^{+\infty}\cos(bx)\,dx$。先取 $s>0$，令
\[
H(s,b)=\int_0^{+\infty}e^{-sx}\frac{\sin(bx)}x\,dx.
\]
对固定 $s$，求导后的被积函数被可积函数 $e^{-sx}$ 控制，故
\[
H_b(s,b)=\int_0^{+\infty}e^{-sx}\cos(bx)\,dx
=\frac{s}{s^2+b^2},
\qquad H(s,0)=0.
\]
积分得 $H(s,b)=\arctan(b/s)$。对固定 $b\ne0$，在 $[A,+\infty)$ 上把 $\sin(bx)$ 与单调因子 $e^{-sx}/x$ 配对，由 Dirichlet 估计得到
\[
\left|\int_A^{+\infty}e^{-sx}\frac{\sin(bx)}x\,dx\right|
\le\frac{C}{|b|A}\qquad(s\ge0),
\]
其中 $C$ 是绝对常数。有限前段则可直接令 $s\downarrow0$，故去掉收敛因子是合法的，最终
\begin{equation}
\label{eq:pei-dirichlet-value}
\int_0^{+\infty}\frac{\sin(bx)}x\,dx
=\begin{cases}\pi/2,&b>0,\\0,&b=0,\\-\pi/2,&b<0.\end{cases}
\end{equation}
这个结果在 $b=0$ 处不连续，说明关于 $b$ 的一致收敛性不能跨过零点。
\end{solve}

\begin{example}
\label{ex:pei-gaussian-transform}
已知 $\int_0^{+\infty}e^{-x^2}\,dx=\sqrt\pi/2$，求
$G(t)=\int_0^{+\infty}e^{-x^2}\cos(2tx)\,dx$。
\end{example}
\begin{solve}
求导后的绝对值不超过 $2xe^{-x^2}$，故可在积分号下求导。利用 $(e^{-x^2})'=-2xe^{-x^2}$ 分部积分，端点项均为零，得到
\[
G'(t)=-2\int_0^{+\infty}xe^{-x^2}\sin(2tx)\,dx=-2tG(t).
\]
由初值 $G(0)=\sqrt\pi/2$，得 $G(t)=(\sqrt\pi/2)e^{-t^2}$。高斯积分的值将在例~\ref{ex:pei-gaussian-double} 中用二重积分独立求出，因此本例同时展示了先建立微分方程、再确定初值的计算路线。
\end{solve}

\subsection{Euler 积分及其转换}

\begin{definition}
\label{def:pei-euler-integrals}
对 $s>0$ 定义 \emph{Gamma 函数（Gamma function）}
\begin{equation}
\label{eq:pei-gamma-definition}
\Gamma(s)=\int_0^{+\infty}x^{s-1}e^{-x}\,dx;
\end{equation}
对 $p,q>0$ 定义 \emph{Beta 函数（Beta function）}
\begin{equation}
\label{eq:pei-beta-definition}
B(p,q)=\int_0^1x^{p-1}(1-x)^{q-1}\,dx.
\end{equation}
它们分别称为第二型与第一型 Euler 积分。
\end{definition}

$\Gamma$ 在零点附近与 $x^{s-1}$ 同阶，在无穷远由指数衰减控制；$B$ 的两个端点分别要求 $p,q>0$。在任意紧参数集 $s\in[\alpha,\beta]\subset(0,+\infty)$ 上，$\Gamma$ 的 $k$ 阶形式导数可分别由零点附近的 $x^{\alpha-1}|\log x|^k$ 和无穷远的 $x^{\beta-1}e^{-x}|\log x|^k$ 控制，因此
\[
\Gamma^{(k)}(s)=\int_0^{+\infty}x^{s-1}e^{-x}(\log x)^k\,dx
\quad(k=0,1,\ldots).
\]
分部积分及 $x=u^2$ 的代换给出
\[
\Gamma(s+1)=s\Gamma(s),\qquad
\Gamma(1)=1,\qquad \Gamma(n+1)=n!,\qquad
\Gamma\!\left(\frac12\right)=\sqrt\pi.
\]
这里的积分定义只适用于 $s>0$；若使用递推关系延拓到负数，应另行说明定义域，并排除 $0,-1,-2,\ldots$。

\begin{proposition}
\label{prop:pei-beta-gamma}
对 $p,q>0$，有
\begin{align}
B(p,q)&=B(q,p)=\frac{\Gamma(p)\Gamma(q)}{\Gamma(p+q)},
\label{eq:pei-beta-gamma}\\
B(p,q)&=2\int_0^{\pi/2}\sin^{2p-1}\theta\cos^{2q-1}\theta\,d\theta
=\int_0^{+\infty}\frac{u^{p-1}}{(1+u)^{p+q}}\,du.
\label{eq:pei-beta-variants}
\end{align}
\end{proposition}
\begin{proof}
对称性来自 $x=1-u$，式~\eqref{eq:pei-beta-variants} 分别来自 $x=\sin^2\theta$ 与 $x=u/(1+u)$。为证明式~\eqref{eq:pei-beta-gamma}，将两个非负积分的乘积写为第一象限上的二重积分，再作 $x=rv,y=r(1-v)$（$r>0,0<v<1$），Jacobian 的绝对值为 $r$，于是
\[
\Gamma(p)\Gamma(q)
=\int_0^{+\infty}e^{-r}r^{p+q-1}\,dr
\int_0^1v^{p-1}(1-v)^{q-1}\,dv.
\]
非负性保证截断后取极限与换序的合法性，故得到所求关系。
\end{proof}

由积分定义与式~\eqref{eq:pei-beta-gamma} 可得
\[
B(p+1,q)=\frac{p}{p+q}B(p,q),\qquad
B(p,q+1)=\frac{q}{p+q}B(p,q),\qquad
B(m,n)=\frac{(m-1)!(n-1)!}{(m+n-1)!}
\]
（$m,n$ 为正整数）。教材还给出了 \emph{Euler 反射公式（Euler reflection formula）}与 \emph{Legendre 倍增公式（Legendre duplication formula）}：
\begin{equation}
\label{eq:pei-gamma-reflection-duplication}
\Gamma(s)\Gamma(1-s)=\frac{\pi}{\sin\pi s}\quad(0<s<1),\qquad
\Gamma(s)\Gamma\!\left(s+\frac12\right)
=2^{1-2s}\sqrt\pi\,\Gamma(2s)\quad(s>0).
\end{equation}
反射公式在此引用教材结论，不展开证明。倍增公式则可由 $x=(1+u)/2$ 与 $v=u^2$ 得到
$B(s,s)=2^{1-2s}B(1/2,s)$，再代入式~\eqref{eq:pei-beta-gamma} 得证。运用这些公式时，先检查参数范围，再进行形式化简。

\begin{example}
\label{ex:pei-euler-model-integrals}
将幂函数与指数、三角函数或有理函数组合的积分化为 Euler 积分。
\end{example}
\begin{solve}
对 $\alpha,\beta,c>0$，令 $u=cx^\beta$，得
\[
\int_0^{+\infty}x^{\alpha-1}e^{-cx^\beta}\,dx
=\frac1\beta c^{-\alpha/\beta}\Gamma\!\left(\frac\alpha\beta\right).
\]
对 $m,n>-1$，由式~\eqref{eq:pei-beta-variants}，
\[
\int_0^{\pi/2}\sin^m x\cos^n x\,dx
=\frac12B\!\left(\frac{m+1}2,\frac{n+1}2\right).
\]
对 $0<a<b$，令 $u=x^b$，再用反射公式，得到
\[
\int_0^{+\infty}\frac{x^{a-1}}{1+x^b}\,dx
=\frac1b B\!\left(\frac ab,1-\frac ab\right)
=\frac{\pi}{b\sin(\pi a/b)}.
\]
最后一个积分在零点要求 $a>0$，在无穷远要求 $a<b$，这些条件与 Beta 参数的正性完全一致。
\end{solve}

\section{重积分}
\label{sec:pei-multiple-integrals}
% 对应原书 §7.2：PDF 第 58--119 页。

\subsection{积分定义、可积性与积分和的极限}

设 $D\subset\mathbb R^2$ 是有界 Jordan 可测区域，即其边界可由总面积任意小的有限个矩形覆盖；$f$ 在 $D$ 上有界。把 $D$ 分成有限个可求面积的小区域 $D_j$，在每块取代表点 $\xi_j$。若当分划最大直径趋于零时，积分和
\[
\sum_j f(\xi_j)|D_j|
\]
趋向一个与分划及代表点无关的数，就称 $f$ 在 $D$ 上 \emph{Riemann 可积（Riemann integrable）}，极限记为 $\iint_Df\,dA$，其中 $dA=dx\,dy$。三重积分与 $n$ 重积分的定义相同，只需把面积替换为体积或 $n$ 维体积。

\begin{proposition}
\label{prop:pei-darboux-integrable}
记 $\omega_j=\sup_{D_j}f-\inf_{D_j}f$。有界函数 $f$ 在 $D$ 上可积，当且仅当对任意 $\varepsilon>0$，存在分划使
\[
\sum_j\omega_j|D_j|<\varepsilon.
\]
特别地，连续函数在紧的 Jordan 可测区域上可积；在有限条分段光滑曲线上有间断、其余处连续的有界函数也可积。后一个结论允许在这些曲线上任意改变有界函数值而不改变积分。
\end{proposition}
\begin{proof}
上述和正是上和与下和之差，而积分和始终位于上下和之间，故结论由 Darboux 判别法得到。连续函数的一致连续性可使每块的振幅同时变小。对所述间断集，先用总面积很小的有限个矩形邻域盖住它，其上贡献由全局有界性控制；在补集上再用一致连续性控制振幅，两个部分的贡献都可任意小。
\end{proof}

\begin{example}
\label{ex:pei-monotone-double-integrable}
设 $f$ 在矩形 $[a,b]\times[c,d]$ 上有界，且固定任一坐标后关于另一坐标均单调递增。证明 $f$ 可积。
\end{example}
\begin{solve}
把两个坐标区间均分为 $n$ 段，网点记为 $x_i,y_j$。每块上的振幅不超过
\[
f(x_i,y_j)-f(x_{i-1},y_{j-1})
=\bigl[f(x_i,y_j)-f(x_{i-1},y_j)\bigr]
+\bigl[f(x_{i-1},y_j)-f(x_{i-1},y_{j-1})\bigr].
\]
设 $|f|\le M$，对所有 $i,j$ 求和时，每个中括号沿一个方向望远镜相消，其总和不超过 $4Mn$。乘每块面积 $(b-a)(d-c)/n^2$ 后，上下和之差不超过 $4M(b-a)(d-c)/n$，趋于零，故可积。
\end{solve}

反过来，积分定义也是计算极限的工具。若 $f$ 在 $[0,1]^2$ 连续，则
\[
\lim_{n\to\infty}\frac1{n^2}\sum_{i,j=1}^n
f\!\left(\frac in,\frac jn\right)=\int_0^1\!\int_0^1f(x,y)\,dy\,dx.
\]
使用这种识别时，应明确网格面积与采样点；若求和含指数、乘积或对数，还应先估计展开余项的总和，不能只对每一项作点态近似后直接相加。

\subsection{累次积分、换序、分区与对称性}

\begin{theorem}
\label{thm:pei-proper-fubini}
设 $D=\{(x,y):a\le x\le b,\ \varphi_1(x)\le y\le\varphi_2(x)\}$，其中 $\varphi_1,\varphi_2$ 连续，$\varphi_1\le\varphi_2$。若 $f$ 在 $D$ 上连续，则
\begin{equation}
\label{eq:pei-double-iterated}
\iint_Df(x,y)\,dA
=\int_a^b\left(\int_{\varphi_1(x)}^{\varphi_2(x)}f(x,y)\,dy\right)dx.
\end{equation}
若 $D$ 也能沿水平方向表示为若干简单区域之并，则对这些区域分别积分，可以得到相应的先对 $x$、再对 $y$ 的表达式。
\end{theorem}
\begin{proof}
先对矩形用积分和分组，将有限双重和写成迭代求和；一致连续性保证内层近似的误差可统一控制，故可依次取极限。一般区域可放入一个矩形，把 $f$ 在区域外延为零；新增间断只可能位于面积为零的边界上，仍可积，再按竖直截线分组，即得式~\eqref{eq:pei-double-iterated}。
\end{proof}

\emph[tip]{改变积分次序时，先把原有的上下限还原成区域不等式，再沿新的方向求截线}。例如，$a<b$ 时
\begin{equation}
\label{eq:pei-triangle-interchange}
\int_a^b\!\int_a^x f(x,y)\,dy\,dx
=\int_a^b\!\int_y^b f(x,y)\,dx\,dy.
\end{equation}
两端共同描述 $a\le y\le x\le b$。区域的边界表达式发生变化，或被积函数因绝对值、符号函数、取整函数而分段时，应先分区；换序不是把两个积分符号和上下限原样互换。

对三重积分，若 $V=\{(x,y,z):(x,y)\in D,\ z_1(x,y)\le z\le z_2(x,y)\}$ 且 $f$ 连续，则
\[
\iiint_Vf\,dV=\iint_D\!\left(\int_{z_1(x,y)}^{z_2(x,y)}f(x,y,z)\,dz\right)dA.
\]
这称为投影法；若先固定 $z$，以 $D_z$ 表示水平截面，则也可写成
$\int_{z_{\min}}^{z_{\max}}(\iint_{D_z}f\,dA)\,dz$，称为截面法。选法取决于哪一种截面的边界和被积函数更简单。

对称性需要同时检查区域与被积函数。例如 $D$ 在反射 $(x,y)\mapsto(-x,y)$ 下不变，且 $f(-x,y)=-f(x,y)$，则积分为零；若变换后函数不变，可在半个区域上计算后乘二。更一般地，对可积函数，保持区域的正交变换 $T$ 满足
\[
\int_D f(x)\,dx=\int_D f(Tx)\,dx,
\]
因为 $|\det T|=1$。区域对称本身并不意味着任意函数的积分都能折半。

\subsection{坐标变换、Jacobian 与典型区域}

\begin{theorem}
\label{thm:pei-change-variables}
设 $T:U\to V$ 是开集之间的 $C^1$ 微分同胚，即一一对应且 $T,T^{-1}$ 均为 $C^1$；$E\subset U$ 是紧的 Jordan 可测区域，$f$ 在 $T(E)$ 上连续。则
\begin{equation}
\label{eq:pei-change-variables}
\int_{T(E)}f(x)\,dx
=\int_E f(T(u))|\det DT(u)|\,du.
\end{equation}
$\det DT$ 称为 \emph{Jacobian 行列式（Jacobian determinant）}。极坐标等变换若只在边界或可忽略的接缝处退化，可先避开这些位置，再取极限使用公式。
\end{theorem}
\begin{proof}[证明思路]
在足够小的分块上，$T$ 由线性变换 $DT(u)$ 作一阶近似；线性变换把体积乘以 $|\det DT(u)|$。把各块上的函数变化与体积近似误差一致控制，再对 Riemann 和取极限，即得公式。一一对应保证内部各点不会被重复计数；绝对值保证体积元素非负。
\end{proof}

平面极坐标 $x=r\cos\theta,y=r\sin\theta$ 给出 $dA=r\,dr\,d\theta$；广义极坐标 $x=ar\cos\theta,y=br\sin\theta$（$a,b>0$）给出 $dA=ab\,r\,dr\,d\theta$。三维柱坐标与球坐标分别给出
\begin{align*}
&x=r\cos\theta,\ y=r\sin\theta,\ z=z,
&&dV=r\,dr\,d\theta\,dz,\\
&x=\rho\sin\varphi\cos\theta,\ y=\rho\sin\varphi\sin\theta,
\ z=\rho\cos\varphi,
&&dV=\rho^2\sin\varphi\,d\rho\,d\varphi\,d\theta,
\end{align*}
其中 $r,\rho\ge0$，球坐标的极角为 $0\le\varphi\le\pi$，方位角可取 $0\le\theta<2\pi$。选取新坐标时，应兼顾被积函数与区域；出现 $xy$ 和 $y/x$ 可考虑把它们作为新变量，但要先检查所在象限及一一对应性。

\begin{example}
\label{ex:pei-gaussian-double}
计算高斯积分 $I=\int_0^{+\infty}e^{-x^2}\,dx$。
\end{example}
\begin{solve}
$e^{-x^2}\le e^{-x}$ 在 $x\ge1$ 成立，故 $I$ 收敛。对非负函数可以按递增区域取极限；在第一象限改用极坐标，得到
\[
I^2=\iint_{x,y\ge0}e^{-(x^2+y^2)}\,dx\,dy
=\int_0^{\pi/2}\!\int_0^{+\infty}e^{-r^2}r\,dr\,d\theta
=\frac\pi4.
\]
因 $I>0$，故 $I=\sqrt\pi/2$。这一步利用的是二重积分的非负性与坐标变换，并未假定可以在两个反常积分之间任意交换极限。
\end{solve}

\begin{example}
\label{ex:pei-ellipsoid-moment}
设 $V=\{(x,y,z):x^2/a^2+y^2/b^2+z^2/c^2\le1\}$，其中 $a,b,c>0$，求其体积及 $\iiint_Vx^2\,dV$。
\end{example}
\begin{solve}
令 $x=au,y=bv,z=cw$，体积元素变为 $abc\,du\,dv\,dw$，区域变为单位球 $B$。因此
\[
|V|=\frac{4\pi}3abc,
\qquad
\iiint_Vx^2\,dV=a^3bc\iiint_Bu^2\,du\,dv\,dw.
\]
由球的轮换对称性，三个坐标平方的积分相等，其总和为
$\iiint_B(u^2+v^2+w^2)\,dV=4\pi\int_0^1\rho^4\,d\rho=4\pi/5$，故
\[
\iiint_Vx^2\,dV=\frac{4\pi}{15}a^3bc.
\]
同理可求另外两个坐标的二阶矩。对非负密度 $\mu$，总质量为 $M=\iiint_V\mu\,dV$；若 $M>0$，质心为 $M^{-1}\iiint_V(x,y,z)\mu\,dV$，绕 $z$ 轴的转动惯量为 $\iiint_V(x^2+y^2)\mu\,dV$。
\end{solve}

\subsection{反常重积分：截断方式、绝对收敛与奇点阶数}

无界区域或无界被积函数需要先截断再积分。\emph[caution]{一般反常重积分的定义要求有限极限不依赖容许的截断方式}：若只沿特定的圆、球或矩形截断得到极限，尚不能据此断言一般反常积分收敛。对非负的局部可积函数，截断区域递增时积分单调增加，故只需检查它们是否有界；任何真正穷尽区域的递增截断都会给出相同的有限值或 $+\infty$。对变号函数，绝对可积性保证这种独立性。

在允许任意容许截断的通常定义下，局部连续函数的反常重积分收敛等价于绝对收敛。其 Cauchy 条件要求：在一个足够大的紧的正常积分区域之外，任取有限个正常小区域，积分总和的绝对值都小。若正部与负部的剩余积分不能同时变小，就可优先截取其中一种符号的区域，使此条件失效。因而这一结论不能照搬到只用固定端点截断的一元条件收敛积分，也不能把特定对称截断所定义的\emph{主值（principal value）}与一般反常重积分混同。

若 $f\ge0$，可用非负积分的换序性质，即 \emph{Tonelli 定理（Tonelli's theorem）}，把重积分化为非负累次积分，结果允许为 $+\infty$。若 $\int_D|f|<\infty$，则由 \emph{Fubini 定理（Fubini's theorem）}可换序计算变号函数。这里截线积分一般只保证对几乎所有外层变量存在；在初等反常积分计算中，应核对实际出现的截线积分，并对个别异常参数作相应解释。

\begin{example}
\label{ex:pei-fubini-failure}
说明两个累次积分各自存在仍不足以保证可以换序。取
\[
f(x,y)=\frac{x^2-y^2}{(x^2+y^2)^2},\qquad 0<x,y<1.
\]
\end{example}
\begin{solve}
对 $x>0$，$f=\partial_y\bigl(y/(x^2+y^2)\bigr)$，故
\[
\int_0^1\!\left(\int_0^1f(x,y)\,dy\right)dx
=\int_0^1\frac{dx}{1+x^2}=\frac\pi4.
\]
同理 $f=-\partial_x\bigl(x/(x^2+y^2)\bigr)$，反向换序所得值为 $-\pi/4$。原点附近用极坐标，$|f|\,dA=|\cos2\theta|\,dr\,d\theta/r$，其积分发散，因此绝对可积条件确实失效。对称地挖去原点小圆时，由交换 $x,y$ 的反对称性还会得到零；这只是特定截断的值，不是一般反常二重积分。
\end{solve}

\begin{proposition}
\label{prop:pei-radial-singularity}
在 $\mathbb R^n$ 中，径向核 $|x|^{-p}$ 在原点附近可积当且仅当 $p<n$，在无穷远可积当且仅当 $p>n$。临界情形 $p=n$ 均为对数发散。
\end{proposition}
\begin{proof}
设单位球面面积为 $\sigma_{n-1}>0$，球坐标将两个积分分别化为
\[
\sigma_{n-1}\int_0^1r^{n-1-p}\,dr,
\qquad
\sigma_{n-1}\int_1^{+\infty}r^{n-1-p}\,dr,
\]
按幂函数积分的敛散条件即得结论。
\end{proof}

与此核比较时，上界给绝对收敛的充分条件；下界给发散结论则要求函数非负，或应用于绝对值，并且该下界要在具有正角度的锥形区域内成立。若积分区域向奇点收缩得很窄，结论会变化。例如在 $D=\{0<x<1,\ 0<y<x^2\}$ 上，
\[
\iint_D\frac{dA}{x^2+y^2}
=\int_0^1\frac{\arctan x}{x}\,dx<\infty,
\]
因为被积的一元函数在零点有有限极限。又若奇点位于光滑曲面上，且核与到该曲面的距离的 $-p$ 次幂同阶，局部收敛阈值通常是 $p<1$，应沿法向坐标判断，不能套用孤立点的三维阈值 $p<3$。

\subsection{高维积分、缩放与降维}

$n$ 重积分沿用相同的换元原则。相似缩放 $x=au$（$a>0$）使体积乘以 $a^n$；若 $A$ 是实对称正定的 $n\times n$ 矩阵，变换 $u=A^{1/2}x$ 把椭球 $x^{\mathsf T}Ax\le1$ 化为单位球，体积元素乘以 $(\det A)^{-1/2}$。这类问题的困难常来自二次型，先作正交对角化比逐层直接定限更有效。

\begin{proposition}
\label{prop:pei-ball-simplex}
$n$ 维半径为 $R$ 的球体体积及球面面积分别为
\begin{equation}
\label{eq:pei-ball-volume}
V_n(R)=\frac{\pi^{n/2}}{\Gamma(n/2+1)}R^n,
\qquad
\sigma_{n-1}(R)=\frac{2\pi^{n/2}}{\Gamma(n/2)}R^{n-1}.
\end{equation}
截距均为 $a>0$ 的单纯形
$\{x_i\ge0,\ \sum_{i=1}^nx_i\le a\}$ 的体积为 $a^n/n!$。
\end{proposition}
\begin{proof}
记单位球体积为 $v_n$。由缩放关系，径向积分等于 $nv_n\int_0^{+\infty}h(r)r^{n-1}\,dr$。另一方面，由高斯积分的乘积，
\[
\pi^{n/2}=\int_{\mathbb R^n}e^{-|x|^2}\,dx
=nv_n\int_0^{+\infty}e^{-r^2}r^{n-1}\,dr
=\frac{nv_n}{2}\Gamma\!\left(\frac n2\right),
\]
从而得到球体积公式，球面面积为其关于 $R$ 的导数。单纯形先缩放到 $a=1$；固定 $x_n=t$ 时，截面是截距为 $1-t$ 的 $(n-1)$ 维单纯形，故其体积满足
\[
s_n(1)=s_{n-1}(1)\int_0^1(1-t)^{n-1}\,dt
=\frac1n s_{n-1}(1).
\]
由 $s_1(1)=1$ 递推得 $s_n(1)=1/n!$，再乘 $a^n$ 即得结论。
\end{proof}

进一步，设 $\alpha_1,\ldots,\alpha_{n+1}>0$，逐次使用 Beta 积分可得 \emph{Dirichlet 单纯形积分（Dirichlet integral）}
\begin{equation}
\label{eq:pei-dirichlet-simplex}
\int_{\substack{x_i\ge0\\\sum_{i=1}^nx_i\le1}}
\left(\prod_{i=1}^nx_i^{\alpha_i-1}\right)
\left(1-\sum_{i=1}^nx_i\right)^{\alpha_{n+1}-1}\,dx
=\frac{\prod_{i=1}^{n+1}\Gamma(\alpha_i)}
{\Gamma(\sum_{i=1}^{n+1}\alpha_i)}.
\end{equation}
证明时固定前 $n-1$ 个坐标，在最内层令 $x_n=(1-\sum_{i<n}x_i)t$，先取得一个 Beta 因子，再降一维递推；非负性保证这些操作可由截断取极限完成。各参数正性正是相应边界奇点可积的条件。所有参数取一，就回到单位单纯形体积公式。

反复求原函数也能降维。设 $f\in C([a,b])$，$n\ge1$，则
\begin{equation}
\label{eq:pei-repeated-integration}
\int_a^xdt_1\int_a^{t_1}dt_2\cdots
\int_a^{t_{n-1}}f(t_n)\,dt_n
=\frac1{(n-1)!}\int_a^x(x-t)^{n-1}f(t)\,dt
\quad(a\le x\le b).
\end{equation}
固定最内层 $t_n=t$ 后，其余 $n-1$ 个变量构成一个体积为 $(x-t)^{n-1}/(n-1)!$ 的有序单纯形，故由换序立即得到公式。这说明某些高重累次积分的本质是一维积分乘以截面体积。

\section{曲线积分与 Green 公式}
\label{sec:pei-line-green}
% 对应原书 §7.3：PDF 第 119--150 页。

\subsection{两类曲线积分及其参数化}

设 $C$ 是分段 $C^1$ 的正则曲线，一组按指定方向遍历的参数方程为 $\boldsymbol r(t)=(x(t),y(t),z(t))$，$\alpha\le t\le\beta$，每段上 $\boldsymbol r'(t)\ne0$。\emph{第一型曲线积分（scalar line integral）}以弧长为积分元素：
\begin{equation}
\label{eq:pei-scalar-line}
\int_Cf\,ds=\int_\alpha^\beta f(\boldsymbol r(t))\|\boldsymbol r'(t)\|\,dt.
\end{equation}
\emph{第二型曲线积分（line integral of a vector field）}则对向量场 $\boldsymbol F=(P,Q,R)$ 定义为
\begin{equation}
\label{eq:pei-vector-line}
\int_C\boldsymbol F\cdot d\boldsymbol r
=\int_C P\,dx+Q\,dy+R\,dz
=\int_\alpha^\beta\boldsymbol F(\boldsymbol r(t))\cdot\boldsymbol r'(t)\,dt.
\end{equation}
两式均对各光滑段求和。第一型积分与方向无关，第二型积分反向后变号；参数化应按所需次数遍历曲线，重复遍历会按重数计入积分。

令 $\boldsymbol\tau$ 是沿指定方向的单位切向量，则 $d\boldsymbol r=\boldsymbol\tau\,ds$，从而
\[
\int_C\boldsymbol F\cdot d\boldsymbol r
=\int_C(\boldsymbol F\cdot\boldsymbol\tau)\,ds,
\qquad
\left|\int_C\boldsymbol F\cdot d\boldsymbol r\right|
\le\int_C\|\boldsymbol F\|\,ds
\le L\max_C\|\boldsymbol F\|,
\]
其中 $L$ 为曲线长度。第一型积分的非负性、单调性和连续函数的积分中值定理来自 $ds\ge0$；第二型积分的 $dx,dy,dz$ 带符号，不能逐分量照搬这些性质。特别地，不能一般地写成 $\int_Cf\,dx=f(\xi)\int_Cdx$。

\begin{example}
\label{ex:pei-line-orientation}
设 $C$ 为单位圆，取逆时针方向。计算 $\int_Cx^2\,ds$ 与 $\int_Cy\,dx$。
\end{example}
\begin{solve}
令 $x=\cos t,y=\sin t$，$0\le t\le2\pi$，则 $ds=dt$、$dx=-\sin t\,dt$，故
\[
\int_Cx^2\,ds=\int_0^{2\pi}\cos^2t\,dt=\pi,
\qquad
\int_Cy\,dx=-\int_0^{2\pi}\sin^2t\,dt=-\pi.
\]
反向后第一式仍为 $\pi$，第二式变为 $\pi$。又因 $\int_Cdx=0$，第二式同时反驳了上述未经条件限制的积分中值表达式。
\end{solve}

对称性也应通过参数或积分元素检查：反射可能保持 $ds$，却改变某个坐标微分和曲线方向。计算空间交线上的积分时，可先消去一个变量，把交线投影到易于参数化的平面；若曲线由球面与平面相交而成，旋转坐标使该平面水平通常更方便。

作为第一型积分的几何应用，若平面曲线 $C$ 绕同平面直线 $\ell:ax+by+c=0$ 旋转，局部弧段 $ds$ 产生的面积为 $2\pi\,\operatorname{dist}((x,y),\ell)\,ds$，因此
\[
S=\frac{2\pi}{\sqrt{a^2+b^2}}\int_C|ax+by+c|\,ds.
\]
这一式子按参数化的重数计算旋转面积；若曲线不同部分旋转后覆盖同一片曲面，要先选取只生成一次的弧段，否则会重复计数。

\subsection{Green 公式、补线与挖去奇点}

\begin{theorem}
\label{thm:pei-green}
设 $D\subset\mathbb R^2$ 为有界区域，边界由有限条互不相交的分段光滑简单闭曲线组成，$P,Q$ 在 $\overline D$ 的一个开邻域内有连续的一阶偏导数。边界正向规定为沿行进方向看区域始终在左侧，即外边界逆时针、洞的边界顺时针。则 \emph{Green 公式（Green's theorem）}为
\begin{equation}
\label{eq:pei-green}
\oint_{\partial D}P\,dx+Q\,dy
=\iint_D(Q_x-P_y)\,dx\,dy.
\end{equation}
\end{theorem}
\begin{proof}[证明思路]
若 $D$ 是竖直简单区域，沿上下两条边界曲线积分并利用微积分基本定理，可得 $\oint_{\partial D}P\,dx=-\iint_DP_y\,dA$；在水平简单区域上同理得到 $\oint_{\partial D}Q\,dy=\iint_DQ_x\,dA$。对可分成有限个简单区域的区域求和时，公共边被相反方向遍历而抵消，留下外边界及洞壁边界。一般所述分段光滑区域可通过局部化与区域逼近得到同一公式。
\end{proof}

若 $\boldsymbol n$ 为平面边界的单位外法向量，边界取正向时有 $\boldsymbol n\,ds=(dy,-dx)$，因此 Green 公式还有平面通量形式
\begin{equation}
\label{eq:pei-planar-flux}
\oint_{\partial D}(P,Q)\cdot\boldsymbol n\,ds
=\oint_{\partial D}P\,dy-Q\,dx
=\iint_D(P_x+Q_y)\,dA.
\end{equation}
取适当的 $P,Q$，可得面积公式
\begin{equation}
\label{eq:pei-green-area}
|D|=\oint_{\partial D}x\,dy
=-\oint_{\partial D}y\,dx
=\frac12\oint_{\partial D}(x\,dy-y\,dx).
\end{equation}
如果原曲线是从 $A$ 到 $B$ 的开曲线，可补一条从 $B$ 回到 $A$ 的方便曲线 $L$，先用 Green 公式计算闭路 $C+L$，再减去 $L$ 上的积分。

\begin{example}
\label{ex:pei-punctured-circulation}
设 $C$ 为不过原点的逆时针分段光滑简单闭曲线，计算
\[
\oint_C\frac{-y\,dx+x\,dy}{x^2+y^2}.
\]
\end{example}
\begin{solve}
置 $P=-y/(x^2+y^2)$，$Q=x/(x^2+y^2)$。在原点以外，直接求导得 $Q_x-P_y=0$。若 $C$ 不包围原点，Green 公式给出积分为零。若 $C$ 包围原点，则不能直接在整个内部使用公式；应挖去一个足够小的圆盘，内圆 $C_\varepsilon$ 取逆时针方向时，剩余区域的正向边界为 $C-C_\varepsilon$，故
\[
\oint_C P\,dx+Q\,dy=\oint_{C_\varepsilon}P\,dx+Q\,dy
=\int_0^{2\pi}dt=2\pi.
\]
结果只由是否包围奇点决定。一般闭路绕原点 $k$ 圈时，积分为 $2\pi k$，其中 $k$ 为带方向的绕数；这也说明局部导数条件并未包含全部的全局信息。
\end{solve}

\begin{example}
\label{ex:pei-linear-ellipse-line}
设 $a d-b c\ne0$，$C$ 是椭圆 $(ax+by)^2+(cx+dy)^2=1$ 的逆时针边界，求
\[
\oint_C\frac{x\,dy-y\,dx}{(ax+by)^2+(cx+dy)^2}.
\]
\end{example}
\begin{solve}
分母在曲线上恒为一，故由式~\eqref{eq:pei-green-area}，所求值是椭圆面积的两倍。线性变换 $u=ax+by,v=cx+dy$ 把椭圆域映为单位圆盘，面积为 $\pi/|ad-bc|$，因此积分为 $2\pi/|ad-bc|$。若改用参数 $u=\cos t,v=\sin t$，还须检查逆变换是否翻转方向；用面积公式已经把这一方向问题统一处理。
\end{solve}

\subsection{路径无关、恰当微分与势函数}

\begin{theorem}
\label{thm:pei-path-independent}
设 $D\subset\mathbb R^2$ 为连通开集，$P,Q\in C^1(D)$。以下三个条件等价：从任意 $A$ 到任意 $B$ 的分段光滑曲线积分 $\int(P\,dx+Q\,dy)$ 与路径无关；每条分段光滑闭路的积分为零；存在 $u\in C^2(D)$ 使
\begin{equation}
\label{eq:pei-exact-differential}
du=P\,dx+Q\,dy,\qquad u_x=P,\quad u_y=Q.
\end{equation}
此时称 $P\,dx+Q\,dy$ 为\emph{恰当微分（exact differential）}，$u$ 为势函数或原函数，在 $D$ 上相差常数。上述条件均推出 $Q_x=P_y$；若 $D$ 还\emph{单连通（simply connected）}，即每条闭路都能在 $D$ 内连续收缩为一点，则 $Q_x=P_y$ 也是充分条件。
\end{theorem}
\begin{proof}
两条同端点路径一正一反组成闭路，因此路径无关与闭路积分为零等价。固定基点 $A$，定义 $u(X)=\int_A^X(P\,dx+Q\,dy)$；在 $X$ 附近用短水平或竖直线段计算增量，得到 $u_x=P,u_y=Q$。反之，若有势函数，链式法则给出任意曲线上的积分为 $u(B)-u(A)$。由混合偏导相等可得 $P_y=Q_x$。最后，在单连通域内，利用闭路的收缩及局部 Green 公式，闭路积分保持不变而最终为零，故导数条件也充分。
\end{proof}

例~\ref{ex:pei-punctured-circulation} 的定义域 $\mathbb R^2\setminus\{0\}$ 不是单连通的，虽然 $Q_x=P_y$ 处处成立，但绕原点的闭路积分不为零，故不存在全域单值势函数。在不跨越所选射线的角度分支内，它可以写成 $d\theta$；角度的局部表达式不能直接充当穿孔平面上的全局原函数。若平面区域只有有限个洞，且 $Q_x=P_y$，则只需再检查绕各个洞的一组生成闭路积分均为零，就能保证全域路径无关。

求势函数时可先由 $u_x=P$ 积分，写成 $u(x,y)=\int P(x,y)\,dx+C(y)$，再用 $u_y=Q$ 决定 $C$；也可选取完全落在定义域内的方便折线计算 $\int_A^X$。例如 $P=2xy+e^x,Q=x^2+2y$ 定义在全平面，$P_y=Q_x=2x$，积分后得到 $u=x^2y+e^x+y^2+C$。最后应对两个偏导数分别验算；只在某一分支或局部区域成立的表达式不能越出其定义域使用。

\subsection{由边界积分恢复微分性质与调和函数}

Green 公式还可把积分条件还原为点态条件。若连续函数 $g$ 在每个充分小的圆盘上积分均为零，则 $g\equiv0$：若某点 $g>0$ 或 $g<0$，由连续性可取一个保号小圆盘，其积分不可能为零。例如 $P,Q\in C^1(D)$，若每个闭圆盘 $B\subset D$ 都满足 $\oint_{\partial B}(P,Q)\cdot\boldsymbol n\,ds=0$，由式~\eqref{eq:pei-planar-flux} 便得到 $P_x+Q_y=0$。

设 $u\in C^2(D)$，记 $\nabla u=(u_x,u_y)$，$\Delta u=u_{xx}+u_{yy}$；满足 $\Delta u=0$ 的函数称为\emph{调和函数（harmonic function）}。对边界分段光滑的有界区域，若 $u$ 在闭包的邻域内为 $C^2$，将式~\eqref{eq:pei-planar-flux} 用于 $u\nabla u$，得到能量恒等式
\begin{equation}
\label{eq:pei-planar-green-energy}
\iint_D|\nabla u|^2\,dA
=\oint_{\partial D}u\frac{\partial u}{\partial n}\,ds
-\iint_Du\Delta u\,dA,
\qquad \frac{\partial u}{\partial n}=\nabla u\cdot\boldsymbol n.
\end{equation}
其中散度的展开为 $\operatorname{div}(u\nabla u)=|\nabla u|^2+u\Delta u$。

\begin{proposition}
\label{prop:pei-harmonic-uniqueness}
在上述正则性条件下，有界连通区域内的调和函数由其边界值唯一确定。
\end{proposition}
\begin{proof}
若 $u_1,u_2$ 调和且边界值相同，令 $w=u_1-u_2$，则 $\Delta w=0$、$w|_{\partial D}=0$。由式~\eqref{eq:pei-planar-green-energy}，$\iint_D|\nabla w|^2\,dA=0$。被积函数连续非负，故每点都为零；$D$ 连通保证 $w$ 为常数，再由零边界值知 $w=0$。
\end{proof}

调和函数还满足圆周平均值性质。若闭圆盘 $\overline{B_r(X_0)}\subset D$，置
$M(r)=(2\pi)^{-1}\int_0^{2\pi}u(X_0+r(\cos\theta,\sin\theta))\,d\theta$，则
\[
M'(r)=\frac1{2\pi r}\oint_{\partial B_r(X_0)}\frac{\partial u}{\partial n}\,ds
=\frac1{2\pi r}\iint_{B_r(X_0)}\Delta u\,dA=0.
\]
由 $M(r)\to u(X_0)$（$r\downarrow0$），得到 $u(X_0)$ 等于任一这样的圆周上的平均值。这里圆周必须完全位于调和区域内，不能跨过奇点。

\section{曲面积分、Gauss 公式与 Stokes 公式}
\label{sec:pei-surface-gauss-stokes}
% 对应原书 §7.4：PDF 第 150--189 页。

\subsection{第一型曲面积分与面积元素}

设曲面片由正则参数方程 $\boldsymbol r(u,v)=(x(u,v),y(u,v),z(u,v))$ 给出，$(u,v)\in E$，$\boldsymbol r\in C^1$ 且 $\boldsymbol r_u\times\boldsymbol r_v\ne0$；除边界接缝外不重复覆盖。\emph{第一型曲面积分（scalar surface integral）}定义为对面积的积分，其计算公式为
\begin{equation}
\label{eq:pei-scalar-surface}
\iint_Sf\,dS
=\iint_E f(\boldsymbol r(u,v))\|\boldsymbol r_u\times\boldsymbol r_v\|\,du\,dv.
\end{equation}
记 $E_1=\boldsymbol r_u\cdot\boldsymbol r_u$、$F_1=\boldsymbol r_u\cdot\boldsymbol r_v$、$G_1=\boldsymbol r_v\cdot\boldsymbol r_v$，则 $dS=\sqrt{E_1G_1-F_1^2}\,du\,dv$。面积积分不依赖法向的选择，分片曲面应分别积分再相加。

若曲面是图像 $z=g(x,y)$，则 $\boldsymbol r=(x,y,g)$，故
\begin{equation}
\label{eq:pei-graph-area}
dS=\sqrt{1+g_x^2+g_y^2}\,dx\,dy.
\end{equation}
若选投影后曲面有多层，应逐层表示，不能把整个曲面只算成一张图像。球面 $x=R\sin\varphi\cos\theta,y=R\sin\varphi\sin\theta,z=R\cos\varphi$ 的面积元素为 $R^2\sin\varphi\,d\varphi\,d\theta$；不能混用球面面积元素与球坐标体积元素 $\rho^2\sin\varphi\,d\rho\,d\varphi\,d\theta$。

\begin{example}
\label{ex:pei-spherical-average}
设 $\boldsymbol a\in\mathbb R^3$，$f$ 在 $[-\|\boldsymbol a\|,\|\boldsymbol a\|]$ 上连续。证明单位球面上的公式
\begin{equation}
\label{eq:pei-spherical-average}
\iint_{S^2}f(\boldsymbol a\cdot\boldsymbol x)\,dS
=2\pi\int_{-1}^1f(\|\boldsymbol a\|t)\,dt.
\end{equation}
\end{example}
\begin{solve}
$\boldsymbol a\ne0$ 时旋转坐标，使新的 $z$ 轴沿 $\boldsymbol a$；球面与面积元素在正交变换下不变，内积变为 $\|\boldsymbol a\|\cos\varphi$。于是积分为
\[
\int_0^{2\pi}\!\int_0^\pi f(\|\boldsymbol a\|\cos\varphi)\sin\varphi\,d\varphi\,d\theta,
\]
令 $t=\cos\varphi$ 即得结论。$\boldsymbol a=0$ 时，两端均为 $4\pi f(0)$。特别地，取 $\boldsymbol a=(0,0,1)$、$f(t)=t^2$，得 $\iint_{S^2}z^2\,dS=4\pi/3$，与三个坐标平方的积分相等且总和为球面积 $4\pi$ 相一致。
\end{solve}

\subsection{第二型曲面积分与定向}

若曲面能够连续地选取单位法向量 $\boldsymbol n$，称其\emph{可定向（orientable）}。选定 $\boldsymbol n$ 后，向量场 $\boldsymbol F=(P,Q,R)$ 穿过曲面的\emph{通量（flux）}为
\begin{equation}
\label{eq:pei-flux-definition}
\iint_S\boldsymbol F\cdot\boldsymbol n\,dS
=\iint_SP\,dy\,dz+Q\,dz\,dx+R\,dx\,dy,
\end{equation}
右端称为\emph{第二型曲面积分（flux integral）}。这里 $dy\,dz,dz\,dx,dx\,dy$ 表示带方向的投影面积元素，满足
\[
(dy\,dz,dz\,dx,dx\,dy)=\boldsymbol n\,dS.
\]
因此反向后通量变号，判断对称性时必须同时比较函数值与定向面积元素。

若 $\boldsymbol r_u\times\boldsymbol r_v$ 与指定法向同向，则
\begin{equation}
\label{eq:pei-parametric-flux}
\iint_S\boldsymbol F\cdot\boldsymbol n\,dS
=\iint_E\boldsymbol F(\boldsymbol r(u,v))\cdot
(\boldsymbol r_u\times\boldsymbol r_v)\,du\,dv;
\end{equation}
反向则加负号。对图像 $z=g(x,y)$ 的上侧，有 $\boldsymbol n\,dS=(-g_x,-g_y,1)\,dx\,dy$，故
\begin{equation}
\label{eq:pei-graph-flux}
\iint_S\boldsymbol F\cdot\boldsymbol n\,dS
=\iint_D\bigl[-P(x,y,g)g_x-Q(x,y,g)g_y+R(x,y,g)\bigr]\,dx\,dy.
\end{equation}
下侧将整式变号。这里有向投影已经吸收了面积因子，不能再乘一次 $\sqrt{1+g_x^2+g_y^2}$。

\subsection{Gauss 公式与补面计算}

\begin{theorem}
\label{thm:pei-gauss}
设 $V\subset\mathbb R^3$ 是有界区域，边界分片光滑，$\boldsymbol F=(P,Q,R)$ 在 $\overline V$ 的一个开邻域内为 $C^1$。令 $\boldsymbol n$ 为 $\partial V$ 的外单位法向量，则 \emph{Gauss 公式（divergence theorem）}为
\begin{equation}
\label{eq:pei-gauss}
\iint_{\partial V}\boldsymbol F\cdot\boldsymbol n\,dS
=\iiint_V\operatorname{div}\boldsymbol F\,dV,
\qquad
\operatorname{div}\boldsymbol F=P_x+Q_y+R_z.
\end{equation}
$\operatorname{div}\boldsymbol F$ 称为\emph{散度（divergence）}。若区域有空腔，洞壁上的外法向指向空腔内部。
\end{theorem}
\begin{proof}[证明思路]
对沿 $x$ 方向简单的区域，利用一元微积分基本定理，$\iiint_VP_x\,dV$ 等于左右边界上 $P$ 的有向通量；对 $Q_y,R_z$ 同理。三式相加得到简单区域的公式。分块求和时，内部公共面法向相反，通量抵消，只剩区域的外边界；一般分片光滑区域可经局部化与逼近得到结论。
\end{proof}

取 $\boldsymbol F=(x,0,0)$ 或 $(x,y,z)/3$，有体积公式
\[
|V|=\iint_{\partial V}x\,dy\,dz
=\frac13\iint_{\partial V}\boldsymbol x\cdot\boldsymbol n\,dS.
\]
若原曲面 $S$ 不封闭，可以补上易算的面 $S_0$，使 $S\cup S_0=\partial V$。当原定向与外侧一致时，
\[
\iint_S\boldsymbol F\cdot\boldsymbol n\,dS
=\iiint_V\operatorname{div}\boldsymbol F\,dV
-\iint_{S_0}\boldsymbol F\cdot\boldsymbol n\,dS.
\]
先在封闭曲面上统一指定外侧，再决定补面符号；补面的积分只有经核验才可以省略。

\begin{example}
\label{ex:pei-paraboloid-flux}
计算 $\boldsymbol F=(x,y,z)$ 穿过抛物面 $S:z=1-x^2-y^2$、$z\ge0$ 上侧的通量。
\end{example}
\begin{solve}
补上 $z=0$ 内的单位圆盘 $S_0$，其外法向为向下。因 $z=0$，补面通量为零；$\operatorname{div}\boldsymbol F=3$，故
\[
\iint_S\boldsymbol F\cdot\boldsymbol n\,dS
=3\iiint_{\substack{x^2+y^2\le1\\0\le z\le1-x^2-y^2}}dV
=3\int_0^{2\pi}\!\int_0^1(1-r^2)r\,dr\,d\theta
=\frac{3\pi}2.
\]
也可用式~\eqref{eq:pei-graph-flux}，得到单位圆盘上 $1+x^2+y^2$ 的二重积分，结果相同。
\end{solve}

\subsection{奇点附近的通量与局部散度}

\begin{example}
\label{ex:pei-inverse-square-flux}
设 $S=\partial V$ 为有界区域的分片光滑边界，原点不在 $S$ 上，取外侧。计算
\[
\iint_S\frac{\boldsymbol x}{|\boldsymbol x|^3}\cdot\boldsymbol n\,dS.
\]
\end{example}
\begin{solve}
在 $\boldsymbol x\ne0$ 处，直接求导得 $\operatorname{div}(\boldsymbol x/|\boldsymbol x|^3)=0$。若 $0\notin\overline V$，直接应用 Gauss 公式得零。若 $0\in V$，挖去完全位于 $V$ 内的半径 $\varepsilon$ 小球；对剩余区域，小球面的外法向是指向球心的，故
\[
\iint_S\frac{\boldsymbol x}{|\boldsymbol x|^3}\cdot\boldsymbol n\,dS
=\iint_{|\boldsymbol x|=\varepsilon}
\frac{\boldsymbol x}{\varepsilon^3}\cdot\frac{\boldsymbol x}{\varepsilon}\,dS
=\frac1{\varepsilon^2}\,4\pi\varepsilon^2=4\pi.
\]
因此奇点外散度为零，并不允许忽略被曲面包围的奇点；此处无需把原点处的场强或散度强行定义为某个有限值。
\end{solve}

散度也能由局部通量恢复。若 $\boldsymbol F\in C^1(D)$，$X_0\in D$，取充分小的球 $B_r(X_0)\Subset D$，由 Gauss 公式及散度连续性，
\begin{equation}
\label{eq:pei-local-divergence}
\operatorname{div}\boldsymbol F(X_0)
=\lim_{r\downarrow0}\frac1{|B_r|}
\iint_{\partial B_r(X_0)}\boldsymbol F\cdot\boldsymbol n\,dS.
\end{equation}
证明只需把右端写成散度在球内的平均值，其与中心值的差不超过球内振幅。因此每个充分小球的通量为零，当且仅当散度处处为零；若要推到域内任意闭曲面的通量，还须保证该曲面连同所围区域都落在场的定义域内。

\subsection{Stokes 公式与空间曲线积分}

\begin{theorem}
\label{thm:pei-stokes}
设 $S$ 是紧的、可定向的分片光滑曲面（各片可取正则 $C^2$ 参数化），边界为有限条分段光滑闭曲线；$\boldsymbol F=(P,Q,R)$ 在 $S$ 的一个开邻域内为 $C^1$。选定 $S$ 的单位法向量 $\boldsymbol n$，边界按右手法则取诱导正向：从法向指向的一侧看边界，其正向使曲面局部位于左侧。则 \emph{Stokes 公式（Stokes' theorem）}为
\begin{equation}
\label{eq:pei-stokes}
\oint_{\partial S}\boldsymbol F\cdot d\boldsymbol r
=\iint_S\operatorname{curl}\boldsymbol F\cdot\boldsymbol n\,dS,
\qquad
\operatorname{curl}\boldsymbol F=(R_y-Q_z,\ P_z-R_x,\ Q_x-P_y).
\end{equation}
$\operatorname{curl}\boldsymbol F$ 称为\emph{旋度（curl）}，亦记作 $\operatorname{rot}\boldsymbol F$。
\end{theorem}
\begin{proof}
对一张参数曲面 $\boldsymbol r(u,v)$，曲线积分变为参数域边界上
\[
\oint_{\partial E}(\boldsymbol F(\boldsymbol r)\cdot\boldsymbol r_u)\,du
+(\boldsymbol F(\boldsymbol r)\cdot\boldsymbol r_v)\,dv.
\]
在参数域应用 Green 公式。利用链式法则并消去 $\boldsymbol r_{uv}=\boldsymbol r_{vu}$ 的项，所得被积函数为
\[
\partial_u(\boldsymbol F(\boldsymbol r)\cdot\boldsymbol r_v)
-\partial_v(\boldsymbol F(\boldsymbol r)\cdot\boldsymbol r_u)
=(\operatorname{curl}\boldsymbol F)(\boldsymbol r)\cdot
(\boldsymbol r_u\times\boldsymbol r_v).
\]
这就是有向曲面积分公式。多张曲面片相加时，内部公共边按相反方向抵消，留下 $\partial S$。
\end{proof}

计算空间闭曲线积分时，先求旋度，再选边界为该曲线且积分方便的曲面；若边界位于平面上，通常取平面区域最简便。只要两张候选曲面都满足 Stokes 公式的正则性与定义域条件，所得旋度通量便相同；不能把曲面移过场的奇点后仍默认公式适用。

\begin{example}
\label{ex:pei-stokes-triangle}
设 $a>0$，$C$ 为按 $A(a,0,0)\to B(0,a,0)\to C_0(0,0,a)\to A$ 顺序遍历的三角形边界，求 $\oint_C z\,dx+x\,dy+y\,dz$。
\end{example}
\begin{solve}
取平面三角形 $S:x+y+z=a$。其向上法向与所给边界方向匹配，向量场 $\boldsymbol F=(z,x,y)$ 的旋度为 $(1,1,1)$。将 $S$ 投影到 $xy$ 平面，有
\[
\boldsymbol n\,dS=(1,1,1)\,dx\,dy,
\qquad D=\{x,y\ge0,\ x+y\le a\}.
\]
因此由 Stokes 公式，原积分为 $\iint_D3\,dx\,dy=3a^2/2$。若边界反向则结果变号；这里向量面积元素已经包含曲面倾斜的影响。
\end{solve}

\section{场论}
\label{sec:pei-vector-fields}
% 对应原书 §7.5：PDF 第 189--203 页。

\subsection{梯度、散度、旋度及其局部含义}

设 $D\subset\mathbb R^3$ 为开集，$u$ 为标量场，$\boldsymbol F$ 为向量场。\emph{Hamilton 算符（nabla operator）}记为
\[
\nabla=\left(\frac\partial{\partial x},\frac\partial{\partial y},\frac\partial{\partial z}\right),
\]
它是微分算符的记号，必须结合所作用的函数及运算类型解释。\emph{梯度（gradient）}为 $\operatorname{grad}u=\nabla u=(u_x,u_y,u_z)$；散度与旋度分别记为 $\operatorname{div}\boldsymbol F=\nabla\cdot\boldsymbol F$ 与 $\operatorname{curl}\boldsymbol F=\nabla\times\boldsymbol F$，坐标表达式已在式~\eqref{eq:pei-gauss} 与式~\eqref{eq:pei-stokes} 中给出。\emph{Laplace 算符（Laplacian）}为
\[
\Delta=\nabla\cdot\nabla=\partial_x^2+\partial_y^2+\partial_z^2;
\]
作用于向量场时约定逐分量计算，即 $\Delta\boldsymbol F=(\Delta F_1,\Delta F_2,\Delta F_3)$。

对单位方向 $\boldsymbol e$，方向导数为 $D_{\boldsymbol e}u=\nabla u\cdot\boldsymbol e$。由 Cauchy--Schwarz 不等式，若 $\nabla u\ne0$，最大方向导数为 $\|\nabla u\|$，在梯度方向达到；梯度同时垂直于正则等值面 $u=\text{常数}$。散度是单位体积的净外流率，精确表达为式~\eqref{eq:pei-local-divergence}。旋度的方向分量则是单位面积的局部环量：若 $D_r$ 是以 $X_0$ 为中心、单位法向为 $\boldsymbol e$ 的半径 $r$ 平面圆盘，边界取诱导正向，则
\[
\operatorname{curl}\boldsymbol F(X_0)\cdot\boldsymbol e
=\lim_{r\downarrow0}\frac1{\pi r^2}
\oint_{\partial D_r}\boldsymbol F\cdot d\boldsymbol r.
\]
这由 Stokes 公式及旋度连续性得到。三个算符中，梯度把标量变成向量，散度把向量变成标量，旋度把向量变成向量，计算前先核对输入与输出类型可避免许多形式错误。

\subsection{乘积法则与混合微分恒等式}

以下各式在相关一阶导数连续时成立；涉及二阶导数的恒等式要求相应场为 $C^2$。对标量场 $u,v$、向量场 $\boldsymbol A,\boldsymbol B$ 及可微一元函数 $\phi$，有
\begin{align}
\nabla(uv)&=u\nabla v+v\nabla u,
&\nabla\phi(u)&=\phi'(u)\nabla u,
\label{eq:pei-gradient-rules}\\
\nabla\cdot(u\boldsymbol A)&=\nabla u\cdot\boldsymbol A+u\nabla\cdot\boldsymbol A,
&\nabla\times(u\boldsymbol A)&=\nabla u\times\boldsymbol A+u\nabla\times\boldsymbol A.
\label{eq:pei-field-product-rules}
\end{align}
当 $v\ne0$ 时，商法则为 $\nabla(u/v)=(v\nabla u-u\nabla v)/v^2$。更复杂的混合运算为
\begin{align}
\nabla\cdot(\boldsymbol A\times\boldsymbol B)
&=\boldsymbol B\cdot(\nabla\times\boldsymbol A)
-\boldsymbol A\cdot(\nabla\times\boldsymbol B),
\label{eq:pei-div-cross}\\
\nabla\times(\boldsymbol A\times\boldsymbol B)
&=\boldsymbol A(\nabla\cdot\boldsymbol B)
-\boldsymbol B(\nabla\cdot\boldsymbol A)
+(\boldsymbol B\cdot\nabla)\boldsymbol A
-(\boldsymbol A\cdot\nabla)\boldsymbol B,
\label{eq:pei-curl-cross}\\
\nabla(\boldsymbol A\cdot\boldsymbol B)
&=(\boldsymbol A\cdot\nabla)\boldsymbol B
+(\boldsymbol B\cdot\nabla)\boldsymbol A
+\boldsymbol A\times(\nabla\times\boldsymbol B)
+\boldsymbol B\times(\nabla\times\boldsymbol A).
\label{eq:pei-grad-dot}
\end{align}
这里 $\boldsymbol A\cdot\nabla=A_1\partial_x+A_2\partial_y+A_3\partial_z$ 是作用在后面函数上的方向微分算符；不能把 $\nabla$ 当作一个普通常向量任意移动，否则会遗漏对系数求导的项。

\begin{proposition}
\label{prop:pei-differential-identities}
若 $u\in C^2(D)$、$\boldsymbol A\in C^2(D;\mathbb R^3)$，则
\begin{equation}
\label{eq:pei-curl-grad-div-curl}
\nabla\times\nabla u=\boldsymbol0,
\qquad\nabla\cdot(\nabla\times\boldsymbol A)=0,
\qquad\nabla\times(\nabla\times\boldsymbol A)
=\nabla(\nabla\cdot\boldsymbol A)-\Delta\boldsymbol A.
\end{equation}
\end{proposition}
\begin{proof}
前两式把坐标表达式展开后，分别由成对的混合偏导相等而抵消。第三式的第一分量为
\[
\partial_y(\partial_x A_2-\partial_y A_1)
-\partial_z(\partial_z A_1-\partial_x A_3)
=\partial_x(\partial_x A_1+\partial_y A_2+\partial_z A_3)-\Delta A_1.
\]
其余两分量循环换标即可。乘积恒等式也可按同样的方法逐分量核验，例如 $\nabla\cdot(u\boldsymbol A)=\sum_i\partial_i(uA_i)$，逐项用普通乘积法则即得式~\eqref{eq:pei-field-product-rules} 的第一式。
\end{proof}

\begin{example}
\label{ex:pei-rigid-rotation}
设刚体以常角速度向量 $\boldsymbol\omega$ 绕过原点的轴旋转，求速度场与加速度场的散度和旋度。
\end{example}
\begin{solve}
记位置向量为 $\boldsymbol x$，则
\[
\boldsymbol v=\boldsymbol\omega\times\boldsymbol x,
\qquad
\boldsymbol a=\boldsymbol\omega\times(\boldsymbol\omega\times\boldsymbol x)
=(\boldsymbol\omega\cdot\boldsymbol x)\boldsymbol\omega
-\|\boldsymbol\omega\|^2\boldsymbol x.
\]
旋转坐标使 $\boldsymbol\omega=(0,0,\omega)$，有 $\boldsymbol v=(-\omega y,\omega x,0)$、$\boldsymbol a=(-\omega^2x,-\omega^2y,0)$，逐项求导得
\[
\nabla\cdot\boldsymbol v=0,\quad
\nabla\times\boldsymbol v=2\boldsymbol\omega,\quad
\nabla\cdot\boldsymbol a=-2\|\boldsymbol\omega\|^2,\quad
\nabla\times\boldsymbol a=\boldsymbol0.
\]
这里旋度为角速度的两倍。另由式~\eqref{eq:pei-grad-dot} 令 $\boldsymbol A=\boldsymbol B=\boldsymbol v$，可得
$(\boldsymbol v\cdot\nabla)\boldsymbol v
=\frac12\nabla\|\boldsymbol v\|^2-\boldsymbol v\times(\nabla\times\boldsymbol v)$，这给出同一加速度的另一种计算方式。
\end{solve}

\subsection{Green 恒等式与边值唯一性}

设 $V$ 为边界分片光滑的有界区域，$u,v$ 在 $\overline V$ 的邻域内有连续二阶偏导数，$\boldsymbol n$ 为外单位法向量。把 Gauss 公式用于 $u\nabla v$，得到 \emph{Green 第一恒等式（Green's first identity）}
\begin{equation}
\label{eq:pei-green-first-identity}
\iiint_Vu\Delta v\,dV
=\iint_{\partial V}u\frac{\partial v}{\partial n}\,dS
-\iiint_V\nabla u\cdot\nabla v\,dV.
\end{equation}
交换 $u,v$ 后相减，得到 \emph{Green 第二恒等式（Green's second identity）}
\begin{equation}
\label{eq:pei-green-second-identity}
\iiint_V(u\Delta v-v\Delta u)\,dV
=\iint_{\partial V}\left(u\frac{\partial v}{\partial n}
-v\frac{\partial u}{\partial n}\right)dS.
\end{equation}
两式的证明只用乘积法则与 Gauss 公式；二维版本则来自平面通量公式，与式~\eqref{eq:pei-planar-green-energy} 一致。

若 $V$ 连通，两个调和函数具有相同的边界值，其差 $w$ 满足 $\Delta w=0$、$w|_{\partial V}=0$，由第一恒等式得到 $\iiint_V|\nabla w|^2\,dV=0$，故 $w=0$。若改为法向导数相同，则同一推导给出 $\nabla w=0$，只能推出两函数相差常数。对于方程 $\Delta u=f$ 与边界条件 $\partial_nu=g$，Gauss 公式还给出存在解的必要兼容条件
\[
\iiint_V f\,dV=\iint_{\partial V}g\,dS.
\]
这里的能量方法证明的是给定正则性下的唯一性或必要条件，不能据此自动断言解一定存在。

\subsection{保守场、有势场、无旋场与径向势}

设 $D$ 为连通开集，$\boldsymbol F\in C^1(D;\mathbb R^3)$。若线积分只依赖端点，称为\emph{保守场（conservative field）}；若存在标量函数 $u$ 使 $\boldsymbol F=\nabla u$，称为\emph{有势场（potential field）}；若 $\nabla\times\boldsymbol F=\boldsymbol0$，称为\emph{无旋场（irrotational field）}。本文采用 $\boldsymbol F=\nabla u$ 的势函数约定；若采用力学中的 $\boldsymbol F=-\nabla V$，则 $V=-u$。

沿用定理~\ref{thm:pei-path-independent} 的证明，在任意连通开集内都有
\begin{equation}
\label{eq:pei-conservative-equivalences}
\boldsymbol F\text{ 为保守场}
\quad\Longleftrightarrow\quad
\text{每条闭路上的环量为零}
\quad\Longleftrightarrow\quad
\boldsymbol F=\nabla u
\quad\Longrightarrow\quad
\nabla\times\boldsymbol F=\boldsymbol0.
\end{equation}
若 $D$ 单连通，最后一个条件也充分；若没有这一全局条件，只能先得到局部势函数。把例~\ref{ex:pei-punctured-circulation} 的平面场扩展为
$\boldsymbol F=(-y/(x^2+y^2),x/(x^2+y^2),0)$，其定义域为去掉 $z$ 轴的空间，便得到三维无旋但不保守的例子。

在\emph{星形区域（star-shaped domain）}中，存在一点 $X_*$，使每个 $X\in D$ 与 $X_*$ 的连线段都在 $D$ 内。此时若 $\boldsymbol F$ 无旋，可以直接构造
\begin{equation}
\label{eq:pei-star-shaped-potential}
u(X)=\int_0^1\boldsymbol F(X_*+t(X-X_*))\cdot(X-X_*)\,dt.
\end{equation}
为核验公式，可平移使 $X_*=0$。对第 $j$ 个坐标求导，并用无旋条件 $\partial_jF_i=\partial_iF_j$，得到
\[
\partial_j u(X)
=\int_0^1\left[F_j(tX)+t\sum_iX_i\partial_iF_j(tX)\right]dt
=\int_0^1\frac d{dt}\bigl(tF_j(tX)\bigr)dt=F_j(X).
\]
因而这不仅证明存在性，也给出了不需要猜测的势函数计算方法。

\begin{example}
\label{ex:pei-radial-potential}
设 $r=|\boldsymbol x|>0$，$\boldsymbol F(\boldsymbol x)=f(r)\boldsymbol x$，其中 $f:(0,+\infty)\to\mathbb R$ 连续。求势函数，并在 $f\in C^1((0,+\infty))$ 时判定何时散度为零。
\end{example}
\begin{solve}
因为 $\nabla r=\boldsymbol x/r$，任取 $r_0>0$，令
\[
u(\boldsymbol x)=\int_{r_0}^{r}s f(s)\,ds,
\]
便有 $\nabla u=rf(r)\nabla r=f(r)\boldsymbol x$。这在穿孔空间上仍给出全局单值势函数，无须依赖单连通性判别。若 $f\in C^1$，则
\[
\nabla\cdot\boldsymbol F=3f(r)+rf'(r).
\]
因此无源的径向场满足 $(r^3f(r))'=0$，即 $\boldsymbol F=c\boldsymbol x/r^3$，其势为 $u=-c/r+C$。等价地，三维径向调和函数 $u=h(r)$ 满足
\[
\Delta u=h''(r)+\frac2r h'(r)=0,
\qquad h(r)=C_1+\frac{C_2}r.
\]
这些公式仅在 $r>0$ 成立；若要求函数在原点也为 $C^2$，则 $C_2$ 必须为零。
\end{solve}

例如，位于不同点 $a_i$ 的质量 $m_i>0$ 所产生的引力场，在所有质量点以外可写为
\[
\boldsymbol F(\boldsymbol x)
=-G\sum_i m_i\frac{\boldsymbol x-a_i}{|\boldsymbol x-a_i|^3}
=\nabla\left(G\sum_i\frac{m_i}{|\boldsymbol x-a_i|}\right),
\]
其中 $G>0$ 是引力常数。它在每个非奇点处既无旋又无源，但闭曲面若包围质量点，其通量一般不为零；可由例~\ref{ex:pei-inverse-square-flux} 平移并相加计算。

\subsection{无源场、向量位与流管守恒}

满足 $\nabla\cdot\boldsymbol F=0$ 的向量场称为\emph{无源场或管量场（solenoidal field）}。对任何连同所围区域一起位于定义域内的闭曲面，Gauss 公式保证其通量为零；反过来，若每个充分小球面的通量都为零，式~\eqref{eq:pei-local-divergence} 保证散度为零。\emph[caution]{“曲面位于定义域内”与“曲面所围区域也位于定义域内”是不同条件}，例~\ref{ex:pei-inverse-square-flux} 正是两者不能混同的反例。

若存在 $\boldsymbol B$ 使 $\boldsymbol F=\nabla\times\boldsymbol B$，称 $\boldsymbol B$ 为 $\boldsymbol F$ 的\emph{向量位（vector potential）}。对 $C^2$ 向量位，由式~\eqref{eq:pei-curl-grad-div-curl}，$\boldsymbol F$ 必无源；而且任一位于其定义域内的闭定向曲面上的通量均为零，因为 Stokes 公式在无边界曲面上的右端对应零环量。无源条件本身一般不足以保证全局向量位存在：在 $\mathbb R^3\setminus\{0\}$ 上，$\boldsymbol x/r^3$ 散度为零，但单位球面通量为 $4\pi$，所以它不可能在整个穿孔空间上写成一个光滑场的旋度。这个区域虽单连通，却仍有不能忽略的二维拓扑障碍。

若 $D$ 为星形区域，则无源场存在全局向量位。对 $\boldsymbol F\in C^2(D)$，把星形中心平移到原点后，一个构造为
\begin{equation}
\label{eq:pei-star-shaped-vector-potential}
\boldsymbol B(X)=\int_0^1t\bigl(\boldsymbol F(tX)\times X\bigr)\,dt.
\end{equation}
利用式~\eqref{eq:pei-curl-cross} 及 $\nabla\cdot\boldsymbol F=0$，有
\[
\nabla\times\bigl(\boldsymbol F(tX)\times X\bigr)
=2\boldsymbol F(tX)+t(X\cdot\nabla)\boldsymbol F(tX).
\]
再乘 $t$ 并积分，恰为 $\int_0^1\frac d{dt}(t^2\boldsymbol F(tX))\,dt=\boldsymbol F(X)$，从而验证了构造。若 $\boldsymbol B$ 为一个向量位，则 $\boldsymbol B+\nabla\phi$ 仍为向量位；反过来，两个向量位之差无旋，只有在相应定义域满足全局势函数条件时，才能断言其差必为某个全局梯度。

最后，若一段流管的侧壁处处与 $\boldsymbol F$ 相切，且管段闭包位于无源场的定义域内，则侧壁通量为零。把管段的两个横截面均按同一流向定向，Gauss 公式给出
\[
\iint_{S_2}\boldsymbol F\cdot\boldsymbol n_2\,dS
-\iint_{S_1}\boldsymbol F\cdot\boldsymbol n_1\,dS=0.
\]
负号来自入口横截面的指定方向与管段的外法向相反。由此同一流管不同横截面的通量相等；这一守恒关系同时说明了“管量场”一名的几何含义。

\end{document}
